\newcommand{\DoNotLoadEpstopdf}{}
\documentclass[11pt,reqno]{amsart}
\usepackage[T1]{fontenc}
\usepackage{lmodern,amsmath,amssymb,amsthm,mathtools,microtype,array}
\usepackage[margin=1.1in]{geometry}
\usepackage[hidelinks]{hyperref}
\hypersetup{pdftitle={Collision records, arithmetic, and raw local inversion in a triangular Lorentz gas},pdfauthor={Qian Qi}}
\setlength{\emergencystretch}{2em}
\allowdisplaybreaks[2]
\numberwithin{equation}{section}
\newtheorem{theorem}{Theorem}[section]
\newtheorem{maintheorem}{Theorem}
\renewcommand{\themaintheorem}{\Alph{maintheorem}}
\newtheorem{lemma}[theorem]{Lemma}
\newtheorem{proposition}[theorem]{Proposition}
\newtheorem{corollary}[theorem]{Corollary}
\theoremstyle{definition}\newtheorem{definition}[theorem]{Definition}
\theoremstyle{remark}\newtheorem{remark}[theorem]{Remark}
\newcommand{\R}{\mathbb R}\newcommand{\Z}{\mathbb Z}\newcommand{\T}{\mathbb T}
\newcommand{\dd}{\,\mathrm d}\newcommand{\one}{\mathbf1}
\newcommand{\supp}{\operatorname{supp}}\newcommand{\TV}{\operatorname{TV}}
\title[Collision records and raw local inversion]{Collision records, arithmetic, and raw local inversion in a triangular Lorentz gas}
\author{Qian Qi}
\date{October 6, 2026. A2-DYN, revision 16}
\subjclass[2020]{37D50, 37A50, 60F05, 42A38}
\keywords{Sinai billiards, collision records, local limits, inducing, Fourier inversion, geometric response}
\begin{document}
\raggedbottom
\begin{abstract}
We study the joint displacement, collision count and flight time of
actual first returns in a triangular periodic Lorentz gas, uniformly
for disk radii in $[0.45,0.47]$. A bounded collision compensation yields
an absolutely convergent, parameter-continuous Green--Kubo matrix.
Smoothing on a shrinking scale and stopping at the genuine return
times give a uniform Gaussian limit for the unsmoothed record, with a
quantitative central-frequency error. A fourth-moment argument yields
uniform functional limits for the return process and for unconditioned
physical displacement and collision fluctuations. The kernel of the
return Gaussian covariance is characterized by an actual $L^2$
coboundary. An explicit smoothing and stopping balance also gives
an integrated Gaussian comparison on a polynomially expanding rescaled
Fourier band, including signed initial insertions. A two-sided compensation at a
marked return extends that comparison to terminal and intermediate
collision-state insertions, uniformly in the position of the mark.
The estimate separates the supremum and variation norms, allowing
controlled growth of the latter. Gap-dependent smoothing further gives fourth moments for unsmoothed
collision sums and a quantitative second-moment identification of the
actual return covariance, also with a mark at any return. The collision-count
variance is positive by an integer-valued cohomology argument. A measurable
phase argument using symplectic area, triangular rotations and finite-cover
ergodicity proves uniform positive definiteness of the entire joint
covariance. It treats the discontinuous section compensation without
assigning values to a measurable transfer function at periodic points.
Complete physical phase arithmetic is obtained from countable phase
ranges and finite covers. On fixed compact nonzero collision-frequency
bands we prove logarithmic defect bounds for bounded-variation phases,
and logarithmic-square bounds for normalized complex functions without
a lower-modulus assumption. At fixed regularity the separation is uniform
in the radius. These results supply the nondegenerate Gaussian, central
integrals and a quantitative phase-reconstruction step for the exact raw
inversion problem. A diffusive middle-block argument also yields uniform
quadratic phase coercivity at frequencies tending to zero. An exact
tower lift and a quantitative finite-record variation estimate transfer
it to the actual unbounded return record, with an explicit loss for
complex functions and growing regularity budgets.

The same physical family admits explicit periodic records, quantitative
phase separation, exact critical-edge subtraction, and a raw local
inversion criterion. We establish cumulative return tails, common
renewal identities, moving-domain continuity, and conditioned control
of unfinished returns. The Gaussian theorems require neither a spectral
realization of the unbounded induced twists nor a presumed induced
correlation sum. The raw mixed-density local limit remains the organizing
problem: its complementary-frequency estimates and complete
critical/singular residual sum remain separate from the Gaussian,
moment and nondegeneracy theorems.
\end{abstract}
\maketitle

\paragraph{The physical record.}
Let $T_R$ be the collision map, $Y_R^*$ the section constructed below,
and $F_R^*$ its actual first-return map. Write
$J_{n,R}=(K_{n,R},N_{n,R},T_{n,R})$ for the four-coordinate record of
$n$ returns, with the two coordinates of $K_{n,R}$ displayed separately.
Its law is taken under the normalized collision measure $\nu_R^*$.
The following synopsis records the unconditional probabilistic
conclusions. The geometry, normalization, and limiting matrices are
specified intrinsically in the cited theorems.

\begin{maintheorem}[Gaussian laws of the physical record]
\label{thm:intro-gaussian}
For $R\in I=[0.45,0.47]$ set
\[
 c_* = \frac{91}{10000\pi},\qquad
 \bar\tau_R=\frac{\sqrt3/2-\pi R^2}{2R},\qquad
 \bar G_R=(0,0,c_*^{-1},\bar\tau_R/c_*).
\]
There is a continuous positive semidefinite matrix $D_R$ such that,
for every $K<\infty$ and $n\ge2$,
\[
 \sup_{R\in I,\ |t|\le K}
 \left|\int e^{it\cdot(J_{n,R}-n\bar G_R)/\sqrt n}\dd\nu_R^*
              -e^{-t^{\mathsf T}D_Rt/2}\right|
 \le C_Kn^{-1/26}\sqrt{\log(2+n)}.
\]
The polygonal centered return processes converge, uniformly in $R$,
to Brownian motion with covariance $D_R$. The stationary physical
processes of lattice displacement and centered collision count also
satisfy a parameter-uniform functional Gaussian limit.

More precisely, $D_R=c_*^{-1}\Gamma_R$, where $\Gamma_R$ is the
absolutely convergent collision Green--Kubo matrix of
$h_R=f_R-\bar G_R\one_{Y_R^*}$. Its kernel consists exactly of directions
$v$ for which $v\cdot(G_R-\bar G_R)$ is an $L^2$ coboundary for
$F_R^*$. No nondegeneracy assumption is used in these conclusions.
\end{maintheorem}
\begin{proof}
Theorems~\ref{thm:collision-covariance},
\ref{thm:actual-record-gaussian}, \ref{thm:gaussian-kernel-coboundary},
and \ref{thm:actual-return-functional}, together with
Corollary~\ref{cor:physical-functional-gaussian}, prove the assertions.
\end{proof}

\begin{maintheorem}[An integrated central band]
\label{thm:intro-major-arc}
For the same actual return record and the matrix $D_R$ in Theorem A,
\[
 \sup_{R\in I}\int_{|v|\le2n^{1/200}}
 \left|\int e^{iv\cdot(J_{n,R}-n\bar G_R)/\sqrt n}\dd\nu_R^*
                       -e^{-v^{\mathsf T}D_Rv/2}\right|\dd v
       \le Cn^{-3/280}\sqrt{\log(2+n)}.
\]
The variable $v$ is rescaled: the corresponding physical frequencies
have size at most $2n^{-1/2+1/200}$. The same estimate holds for complex
initial insertions whose zero extensions to the collision cylinder
have uniformly bounded supremum and variation norms, with the Gaussian
multiplied by the integral of the insertion.
\end{maintheorem}
\begin{proof}
Apply Theorem~\ref{thm:growing-major-arc}; the original probability is
the admissible insertion $a_R=s_R$.
\end{proof}

\paragraph{Complex insertion convention.}
In Theorem B the insertion is a density with respect to the collision
probability $\nu$. Its amplitude $\alpha_R=\int a_R\dd\nu$ may be
complex; the associated pushforward is a finite complex measure, not
in general a probability. The original section probability corresponds
to $a_R=s_R=c_*^{-1}\one_{Y_R^*}$.

\begin{maintheorem}[A mark at any actual return]
\label{thm:intro-marked-return}
Let $a$ be a complex function on the collision cylinder, zero outside
$Y_R^*$, and put $M_a=\|a\|_\infty$, $V_a=\|a\|_{\mathrm{BV}}$,
and $\alpha_a=\int a\dd\nu$. For $0\le k\le n$ set
\[
 C^{[k],a}_{n,R}(v)=\int_{Y_R^*}a((F_R^*)^k x)
       e^{iv\cdot(J_{n,R}(x)-n\bar G_R)/\sqrt n}\dd\nu(x).
\]
There is a constant $C$, independent of $R,n,k,a$, such that, for $n\ge2$,
\begin{align*}
 \int_{|v|\le2n^{1/200}}
 \left|C^{[k],a}_{n,R}(v)-\alpha_a e^{-v^{\mathsf T}D_Rv/2}\right|\dd v
 \le C\left(M_an^{-3/280}\sqrt{\log(2+n)}+V_an^{-9/175}\right).
\end{align*}
Thus the same central band accommodates initial, terminal, and one
intermediate return-state insertion, including a mark whose index
varies with $n$. The displayed transform is integrated against the
unnormalized restricted collision measure $\nu|_{Y_R^*}$; division by
$c_*$ gives the corresponding transform under the section probability
$\nu_R^*$. No positivity of $a$ or nonsingularity of $D_R$ is assumed.
\end{maintheorem}
\begin{proof}
This is Theorem~\ref{thm:marked-return-band}.
\end{proof}


\begin{maintheorem}[Actual second moments and a marked quadratic law]
\label{thm:intro-actual-moments}
Put $U_{n,R}=J_{n,R}-n\bar G_R$. Uniformly for $R\in I$,
\[
 \int |U_{n,R}|^4\dd\nu_R^*\le Cn^2,\qquad
 \left\|\frac1n\int U_{n,R}\otimes U_{n,R}\dd\nu_R^*-D_R\right\|
       \le Cn^{-1/6}.
\]
The count component satisfies $e_3^{\mathsf T}D_Re_3\ge d_N>0$
uniformly in $R$, where $e_3=(0,0,1,0)$.
For the complex insertion and normalization of Theorem C,
\begin{align*}
 &\left\|\int_{Y_R^*}a((F_R^*)^kx)
       \frac{U_{n,R}(x)\otimes U_{n,R}(x)}n\dd\nu(x)-\alpha_aD_R\right\|\\
 &\hspace{20mm}\le C\left(M_an^{-1/6}+(M_a+V_a)n^{-1}\right),
       \qquad 0\le k\le n.
\end{align*}
The matrix $D_R$ is also the uniform limit of the induced Green--Kubo
formula with Cesaro weights. No induced mixing or absolute convergence
of its correlation series is assumed.
\end{maintheorem}
\begin{proof}
Apply Theorems~\ref{thm:marked-L2-stopping} and
\ref{thm:marked-quadratic-moments}, and
Corollaries~\ref{cor:induced-cesaro-covariance} and
\ref{cor:positive-count-variance}.
\end{proof}

\begin{maintheorem}[Uniform nondegeneracy of the joint record]
\label{thm:intro-joint-nondegeneracy}
The covariance of the same actual four-coordinate return record satisfies
\[
 d_0 I_4\le D_R\le D_0 I_4\qquad(R\in I)
\]
for constants $0<d_0\le D_0<\infty$. Thus every nonzero real joint
direction has positive Gaussian variance. For all sufficiently large
$n$, the normalized actual covariance in Theorem D is at least
$(d_0/2)I_4$, uniformly in $R$. The three-dimensional physical-time
fluctuation covariance is uniformly positive definite as well.
\end{maintheorem}
\begin{proof}
Apply Theorem~\ref{thm:joint-uniform-nondegeneracy} and
Corollary~\ref{cor:actual-uniform-ellipticity}.
\end{proof}
In every subsequent use of the matrix from Theorems A--D, the stronger
uniform positive definiteness of Theorem E is available. The positive
semidefinite formulation of the earlier statements records their weaker
logical input, not a remaining covariance hypothesis.


\begin{maintheorem}[Physical phase arithmetic and quantitative defects]
\label{thm:intro-quantitative-phases}
Let $\omega=(u,s,t)\in\T^2\times\T\times\R$ and
$\phi_{R,\omega}=u\cdot\kappa_R+s+t\tau_R$ on the collision cylinder.
A measurable circle-valued solution of
$q\circ T_R=e^{i\phi_{R,\omega}}q$ exists exactly when $\omega=0$;
in that case it is constant almost everywhere.

For each fixed $R$ and compact nonzero frequency set $\mathcal K$,
there is $c_{R,\mathcal K}>0$ such that, for $H\ge2$,
\[
 \inf_{\omega\in\mathcal K}
 \|q\circ T_R-e^{i\phi_{R,\omega}}q\|_1
                  \ge\frac{c_{R,\mathcal K}}{1+\log H}
 \quad (|q|=1,\ \|q\|_{\mathrm{BV}}\le H).
\]
For complex $f$ with $\|f\|_2=1$ and
$\|f\|_\infty+\|f\|_{\mathrm{BV}}\le H$, the corresponding bound is
\[
 \inf_{\omega\in\mathcal K}
 \|f\circ T_R-e^{i\phi_{R,\omega}}f\|_2
                  \ge\frac{c_{R,\mathcal K}}{(1+\log H)^2}.
\]
No lower bound on $|f|$ is assumed. For fixed $H$ and $\mathcal K$,
the second infimum is also positive uniformly over $R\in I$.
The latter assertion does not specify a uniform logarithmic rate in $H$.
\end{maintheorem}
\begin{proof}
Theorem~\ref{thm:complete-physical-phase} proves the exact assertion.
Theorems~\ref{thm:compact-log-phase-defect},
\ref{thm:compact-vector-defect} and
\ref{thm:uniform-compact-bv-separation} prove the defect bounds.
\end{proof}

\begin{maintheorem}[Near-origin defects for the actual return record]
\label{thm:intro-actual-return-defects}
Set $g_R=G_R-\bar G_R$ and
\[
 \Lambda=1+\log(H/r),\qquad K=\Lambda\log(2+\Lambda),
                 \qquad H\ge2,\quad 0<r<1/2.
\]
There are $a,c,r_*>0$, independent of $R$, such that $r\le r_*$
and $rK\le a$ imply the following conclusions for $|z|=r$.
Every unit-modulus section function $q$ whose zero extension has
variation at most $H$ satisfies
\[
 \|q\circ F_R^*-e^{iz\cdot g_R}q\|_{L^1(\nu_R^*)}\ge cr^2.
\]
Every complex section function $f$ with
$\|f\|_{L^2(\nu_R^*)}=1$ and
$\|f\|_\infty+\|f^0\|_{\mathrm{BV}}\le H$ satisfies
\[
 \|f\circ F_R^*-e^{iz\cdot g_R}f\|_{L^2(\nu_R^*)}
                                    \ge cr^2/K.
\]
The function may vanish. With $H\le H_0n^\zeta$, these bounds apply
uniformly on the physical annulus
$2n^{-99/200}\le|z|\le n^{-2/5}$ for all sufficiently large $n$.
The corresponding rescaled radii are $2n^{1/200}$ and $n^{1/10}$.
\end{maintheorem}
\begin{proof}
Apply Theorems~\ref{thm:actual-return-circle-defect} and
\ref{thm:actual-return-vector-defect}, and
Corollary~\ref{cor:actual-annulus-defect}.
\end{proof}

\paragraph{From Gaussian laws to raw local inversion.}
The local problem concerns densities with respect to counting measure
in the lattice coordinates and Lebesgue measure in flight time, not
convolution against a fixed test function. Theorem~\ref{thm:intro-gaussian}
therefore does not replace the complementary-frequency and raw residual
estimates in Theorem~\ref{thm:LLT}. The explicit periodic geometry
and edge terms are retained as inputs to that problem. Theorem E proves
nondegeneracy by a different route: the section indicator is removed
before measurable phase rigidity is applied on positive-measure
stable--unstable product sets. No periodic representative is assumed.
The article develops the Gaussian limit, joint nondegeneracy, geometric
estimates and exact local inversion requirements for the same record.

\paragraph{Structure of the proof.}
The joint periodic arithmetic and quantitative coercivity in
Sections~\ref{sec:joint-arithmetic} and~\ref{sec:periodic-coercivity}
concern actual physical orbits. The collision compensation, Gaussian
limits, and functional time changes lead to Theorem A.
Section~\ref{sec:growing-major-arc} then tracks every frequency factor
in the smoothing and stopping errors to prove Theorem B. Its corollary
inserts that proved central band into the exact local decomposition:
positive definiteness is supplied by Theorem E, whereas the
complementary-frequency integral and local edge errors remain distinct
analytical requirements. In particular, an outer-tail
estimate beginning at $n^{-2/5}$ would not cover the intervening annulus.
Appendices~\ref{app:bv-details} and~\ref{app:functional-details} expose
the single-collision variation, smooth-density embeddings,
chronological products, and interpolation estimates. Section~\ref{sec:marked-return-band} recenters the physical orbit at a
chosen return and applies the collision splitting on both sides of a
single smooth multiplier. This proves Theorem C without assuming
mixing of the induced map, and without replacing a randomly evaluated
weight by its deterministic-clock value. Section~\ref{sec:unsmoothed-moments} then proves a finite-product
decoupling estimate by smoothing at the chosen collision gap. This
yields unsmoothed fourth moments, an $L^2$ comparison with the actual
two-sided stopping, and Theorem D.
Sections~\ref{sec:measurable-phase-rigidity} and
\ref{sec:joint-nondegeneracy} use the exact collision action one-form
and conditional pairs to exclude nonzero roof phases. Rotation products
exclude a nonintegral constant phase, and ergodicity of a finite physical
cover excludes a real spatial coboundary. This proves Theorem E for
every joint direction. Sections~\ref{sec:complete-phase-arithmetic} and
\ref{sec:quantitative-phase-defects} then complete the physical circle
arithmetic, quantify the conditional-pair argument, and reconstruct a
bounded-variation phase from an approximate complex vector. This proves
Theorem F on its stated collision-function class.
Sections~\ref{sec:near-origin-defects} and~\ref{sec:actual-return-defects}
then use diffusive damping, untwisted endpoint blocks and an explicitly
regularized actual tower to prove Theorem G. The annulus conclusion is
a function-defect bound; converting it to a complementary transform
estimate requires the stated operator comparison. The raw inversion
routes retain the same mixed-density endpoint.

\paragraph{Relation to the limit theory of billiards.}
The existence of Gaussian and functional laws for finite-horizon
billiards is not in itself new: tower and invariance-principle methods
already give such laws for collision maps and flows
\cite{Young,MN}. Perturbative collision-space methods give stability of
spectral data under geometric changes \cite{DZ}. Local limit theorems
for cell displacement, including randomly deformed tables, provide a
further comparison \cite{SV,DPZ}. Our problem instead keeps the lattice
displacement, collision count, and flight time of a genuine moving-section
first return in one four-coordinate record. The compensation identity
allows quantitative, radius-uniform central integrals for this
unbounded induced record using bounded collision observables. Its role
here is to connect that central analysis to the arithmetic and exact
edge subtraction of the raw mixed-density problem. None of the cited
cell-displacement theorems is invoked as a theorem for this joint raw
record, and the central-band result is not identified with a complete
local limit theorem.

\section{The mechanical question and its interfaces}
Let $a=(1,0)$, $b=(1/2,\sqrt3/2)$, $\Lambda=\Z a+\Z b$, and
\[
 I=[9/20,47/100],\qquad Q_R=(\R^2/\Lambda)\setminus\overline B(0,R).
\]
The speed is one and every reflection is specular. Three indices will be kept distinct: a physical collision count $N$, an induced return count $n$, and the number of experimental bits. An induced record has values in $\Z^2\times\Z\times\R$, with counting measure in the first three coordinates and Lebesgue measure in the last. A four-dimensional Lebesgue density would be the wrong object.

The geometric inverse of the preceding article identifies the table from two fixed collision fields \cite{A2G}. It supplies neither independent flights nor an induced transfer-operator estimate. Conversely, the stationary law of the equilibrium billiard is not the unknown preparation law of that inverse experiment. The distinction permits geometric estimates to be used for equilibrium statistics without first estimating the entire preparation law.

The joint arithmetic result is Theorem~\ref{thm:joint-periodic}. It concerns
the complete periodic phase of an explicitly enlarged physical first-return
section. Lemma~\ref{lem:excursion-family} constructs its infinite period
family, and Theorem~\ref{thm:period-excess} proves the strict excess
inequalities. These results go beyond the zero-roof determinant by ruling
out every nonzero roof frequency, without a numerical irrationality claim.
The real-rank and compact-frequency consequences are stated separately
from measurable cohomology regularity and operator contraction.

The original physical record, winding geometry, critical-edge and
conditional inversion arguments remain in the article. Theorem
\ref{thm:local-edge-LLT} refines the inversion requirement by using local
edge errors instead of small total edge mass. Propositions
\ref{prop:exact-induced-means} and \ref{prop:induced-suspension} give exact
normalizations and stationary endpoint control for the enlarged section.
The full raw local limit is not inferred from the arithmetic theorem:
its spectral, residual and weighted-insertion requirements are listed in
the final section. All probability laws here refer to actual specular
orbits, not independent flight surrogates.

The relevant prior local-limit theory includes the Lorentz-process theorem of Sz\'asz--Varj\'u \cite{SV} and the suspension framework of Dolgopyat--N\'andori \cite{DN}. Definition 2.2 and Proposition 6.3 of the latter concern a mixing local limit tested against fixed compactly supported functions under minimality and group assumptions. They do not state absolute integrability of every raw finite-time characteristic function, or the particular geometrically uniform density theorem sought here. Demers--Zhang \cite{DZ} give a perturbation framework for Lorentz-gas maps; its spectral continuity is not automatically twice differentiability in a moving geometric parameter. We use the corrected version of the flow-mixing work of Baladi--Demers--Liverani \cite{BDL}. No novelty is claimed for Fourier inversion, the Morse lemma, or the suspension identity as general methods.


Theorem~\ref{thm:return-stability} compares the actual moving return
events in density $L^1$ and first moment, at every fixed return count.
Corollary~\ref{cor:uniform-stationary-endpoints} then supplies uniform
stationary endpoint control over the whole radius interval. The proof uses
finite-itinerary stability, coarea and exact Kac means, not a hypothetical
common spectral space. These finite-record moduli do not purport to supply
the quantitative long-time estimates in the local-limit interface.

Theorem~\ref{thm:quantitative-excess} strengthens strict increase to
uniform two-sided exponential bounds, obtained from the exact nonlinear
half action. Theorem~\ref{thm:log-period-separation} uses their upper and
lower bounds together: the former selects periods of logarithmic length,
and the latter supplies a polynomial discrepancy rather than mere
nonvanishing. Corollary~\ref{cor:approximate-phase-bound} states precisely
the continuous-phase consequence. It does not assume the measurable
regularity needed for an operator application.

At the level of the clock, Theorems~\ref{thm:uniform-return-fluid} and
\ref{thm:uniform-physical-clock} turn fixed-record continuity and exact
means into a uniform functional first-order law. This includes the maximum
unfinished return over the observation window at scale $t$, whereas the
retained stationary endpoint result is at scale $\sqrt t$ for finitely
many deterministic endpoints. The distinction of scales is essential.
Corollary~\ref{cor:physical-rate-error} gives a geometric error term for
the physical rate without assuming independence of the radius estimate
and the observed orbit. No originality is claimed for the general ergodic,
subadditive or renewal inversion arguments; the point is their verified
application to the actual moving return family.

\section{Physical records and an exact marked identity}
Write an outgoing section point as
\[
 q=R n_\alpha,\qquad v=\sqrt{1-p^2}\,n_\alpha+p\,t_\alpha,
 \qquad (\alpha,p)\in\T\times(-1,1).
\]
The collision map is $T_R$, its next physical flight is $\tau_R$, and
\begin{equation}\label{eq:flux}
 \dd\nu=\frac{\dd\alpha\dd p}{4\pi},\qquad
 \bar\tau_R=\frac{\sqrt3/2-\pi R^2}{2R},\qquad
 \dd\mu_R=\bar\tau_R^{-1}\dd\nu\dd s\quad(0\le s<\tau_R).
\end{equation}
These formulas follow by integrating the boundary flux $\cos\phi\dd s_{\partial Q}\dd\phi$: total section flux is $4\pi R$ and phase volume is $2\pi(\sqrt3/2-\pi R^2)$. Reflection preserves this flux. Thus $T_R$ preserves $\nu$.

Each next disk has a lattice label $d\ne0$. With $D=d-q$ and $L=D\cdot v$, its candidate entry time is
\begin{equation}\label{eq:entry}
 \tau_d=L-\sqrt{R^2-|D|^2+L^2}.
\end{equation}
The physical branch is the positive entry root with the smallest time, not an arbitrarily selected root. Its endpoint normal is $n'=(q+\tau_dv-d)/R$, and its outgoing velocity is $v'=v-2(v\cdot n')n'$. These expressions define the label, point, both velocities, and time on every regular branch. On a compact finite itinerary separated from tangencies and competing roots they depend analytically on $R$ and the initial coordinates, by the implicit function theorem. This is a finite-itinerary assertion, not a common norm for all iterates.

Distinct disks have distance at least $g_R=1-2R\ge3/50$. A straight outgoing flight cannot return to its initial convex disk. In particular there is no finite-time collision accumulation. This family has uniform finite horizon: a rational lattice direction has row spacing at most $\sqrt3/2$, so every line in that direction meets an open disk when $R>\sqrt3/4$; an irrational line is dense on the torus. Compactness then excludes arbitrarily long free segments at the smallest radius. Increasing the radius preserves this upper bound. Tangencies and their finite preimages have section measure zero.

Put $S_j\tau(y)=\sum_{i=0}^{j-1}\tau(T_R^iy)$. For $y$ interpreted as the first future outgoing collision and $u$ as its arrival time, the first $k$ future collisions have times $u+S_j\tau(y)$, $0\le j<k$. Their points, lattice labels, and incoming and outgoing velocities are obtained by iterating \eqref{eq:entry}. The final incomplete flight has duration $t-u-S_{k-1}\tau(y)$.

\begin{proposition}[Marked stationary record]\label{prop:marked}
For every bounded measurable functional $\Psi$ of this record and its terminal flight,
\begin{align}\label{eq:marked}
 \mathbb E_{\mu_R}[\Psi;N_t=k]
 =\frac1{\bar\tau_R}\int\dd\nu(y)\int_0^{\tau(T_R^{-1}y)}
 &\one_{\{u+S_{k-1}\tau(y)\le t<u+S_k\tau(y)\}}\nonumber\\[-2mm]
 &\hspace{-8mm}\cdot\Psi(\mathcal R_{k,t}(y,u))\dd u,\qquad k\ge1.
\end{align}
The initial phase point is $\Phi_R^{-u}y$, so the entire initial residual flight is retained as well. For $t=11/100$ there are at most two future collisions, and \eqref{eq:marked} with $k=1,2$, together with the zero-collision case, gives the full marked law without an independence assumption.
\end{proposition}
\begin{proof}
Under \eqref{eq:flux}, set $y=T_Rx$ and $u=\tau(x)-s$. Invariance of $\nu$ gives the displayed domain for $u$ with Jacobian one. The $k$th and $(k+1)$st future collision times are exactly the two endpoints in the indicator. This proves the identity for arbitrary bounded measurable $\Psi$, first for indicators and then by integration. The event of equality at an endpoint has zero suspension measure. Three collisions require at least two complete intercollision flights after the first event; their total duration is at least $6/50>11/100$.
\end{proof}


\section{Explicit periodic geometry and a genuine inducing set}
Two elementary cycles have zero winding. The normal nearest-neighbor orbit has two flights, total length
\begin{equation}\label{eq:two-three}
 L_2=2(1-2R),\qquad
 L_3=3(1-\sqrt3 R)
\end{equation}
is the length of the inner equilateral three-cycle on centers $0,a,b$. Its contact angles are $\pi/6,5\pi/6,3\pi/2$ and its outgoing $p$ coordinates are all $-1/2$. The reflection cosines are respectively $1$ and $\sqrt3/2$. The equilateral segments lie inside their elementary triangle and outside the three vertex disks; any other lattice center is at distance at least $\sqrt3/2$ from that triangle. Hence all other disks are strictly avoided.

\begin{lemma}[A winding four-cycle]\label{lem:winding}
For every $R\in I$ there is an analytic angle $\theta_R\in(0.92,0.967)$ satisfying
\begin{equation}\label{eq:theta}
 f_R(\theta)=2R\sin\theta-\sin\theta\cos\theta
                       +\frac{\sqrt3}{2}\cos(2\theta)=0.
\end{equation}
The itinerary with centers $0,a,a-b,2a-b,a$ and normals at angles
\[
 -\theta_R,\quad\pi+\theta_R,\quad\theta_R,\quad\pi-\theta_R
\]
is a physical four-collision periodic orbit modulo $\Lambda$, with winding $a$. All its incidence cosines exceed $1/2$ and its clearance from every third disk exceeds $1/200$. Rotation through $\pi/3$ gives another cycle with winding $b$ and the same length.
\end{lemma}
\begin{proof}
Put $c=\cos\theta$, $s=\sin\theta$. Its first four contact points and the translated fifth point are
\begin{align*}
 q_0&=(Rc,-Rs),&q_1&=(1-Rc,-Rs),\\
 q_2&=(1/2+Rc,-\sqrt3/2+Rs),&q_3&=(3/2-Rc,-\sqrt3/2+Rs),\\
 q_4&=q_0+a.
\end{align*}
Two segments are horizontal. Write $d=2Rc-1/2$, $h=\sqrt3/2-2Rs$. The other two are $(d,-h)$ and $(d,h)$. The scalar equation $d\sin2\theta+h\cos2\theta=0$ is precisely \eqref{eq:theta}. Since $d,h>0$, it makes the diagonal direction the reflection of the horizontal direction at the intervening normal. The remaining reflection equations follow by symmetry, with cosine $c$ at every contact. The total length is
\begin{equation}\label{eq:fourlength}
 L_4=2(1-2Rc)+2\sqrt{d^2+h^2}.
\end{equation}

Here are reproducible rational bounds for the inequalities and absence of occlusion. Divide $I$ into 32 closed intervals of width $1/1600$. On each, bisect the endpoint equations \eqref{eq:theta} between $9/10$ and $1$ to width at most $10^{-10}$. The interval derivative
\[
 2R\cos\theta-\cos2\theta-\sqrt3\sin2\theta
\]
is less than $-1/2$. The endpoint signs and positivity of $\partial_R f$ show that the two endpoint root brackets enclose the root for the entire parameter interval. The resulting angles lie in $(0.92,0.967)$; $d,h>0$ and $c>1/2$ on every interval.

For clarity, the arithmetic of this finite certificate uses no floating root solver. Square roots are enclosed by integer square roots at denominator $10^{22}$. For sine and cosine use the Taylor polynomials through degrees 25 and 24 at the rational midpoint, with errors at most $|x|^{26}/26!$ and $|x|^{25}/25!$, respectively; add the interval radius, since both functions are 1-Lipschitz. All interval operations are rational.

For each segment check centers $ia+jb$, $|i|,|j|\le4$, excluding its two endpoints. The distance from the rational midpoint segment is computed by clamping the rational orthogonal projection to $[0,1]$. Subtract the endpoint and center enclosure radii and the largest obstacle radius. Each lower bound exceeds $1/200$. Every contact point has norm less than 3; a center outside the index box has norm at least $5\sqrt3/2$, so it too has distance greater than the obstacle radius from every segment. The complete deterministic calculation is specified in \texttt{tools/certify\_winding.py}; it obtains lower bounds greater than $0.00925$ for clearance and $0.569$ for incidence, and a derivative upper bound below $-0.731$. The signs and strict margins prove the assertions on the full intervals, not only at sampled radii. The implicit function theorem gives analytic dependence and uniqueness in the printed root interval.
\end{proof}

Choose $\epsilon=1/200$. In $(\alpha,p)$ coordinates let $Y_R$ be the union of the open squares of half-side $\epsilon$ centered at
\begin{equation}\label{eq:Y}
 (0,0),\quad(\pi/6,-1/2),\quad
 (-\theta_R,\sin\theta_R),\quad
 (\pi/3-\theta_R,\sin\theta_R).
\end{equation}
Angles are taken modulo $2\pi$. The squares are disjoint and
$\nu(Y_R)=4\epsilon^2/\pi$. Let $r_R$ be the first return time to $Y_R$ and $F_R=T_R^{r_R}$ on its recurrent full-measure domain, with invariant probability $\nu_R^Y=\nu|_{Y_R}/\nu(Y_R)$. Define
\begin{equation}\label{eq:induced}
 \kappa_R=\sum_{j<r_R}\kappa_{\rm coll}\circ T_R^j,
 \qquad \tau_R^Y=\sum_{j<r_R}\tau_R\circ T_R^j.
\end{equation}
Poincar\'e recurrence and invariance give this induced probability; they do not give an exponential return tail, a Markov partition, or a differentiable family of spectral projectors.

\begin{proposition}[Full augmented lattice at zero roof frequency]\label{prop:lattice}
The four cycles above give fixed points of $F_R$ with data $(\kappa_1,\kappa_2,r,\text{return period})$ equal to the rows of
\begin{equation}\label{eq:matrix}
 V=\begin{pmatrix}0&0&2&1\\0&0&3&1\\1&0&4&1\\0&1&4&1\end{pmatrix},
                  \qquad |\det V|=1.
\end{equation}
Consequently a circle-valued coboundary identity
$e^{i(u\cdot\kappa+s r)}=e^{ic}g/(g\circ F_R)$ that can be evaluated at these fixed points forces $u=s=c=0$ modulo $2\pi$.
\end{proposition}
\begin{proof}
Only one contact of each cycle belongs to \eqref{eq:Y}. For the four-cycle the successive $p$ coordinates are $s,s,-s,-s$. The same is true after rotation. The two positive-$p$ unselected contacts are separated from the selected centers; the smallest relevant angular separation is $2\pi/3-2\theta_R>0.16$. All negative-$p$ contacts are separated in $p$ from the selected winding contacts, and also from $0,-1/2$. The remaining contacts of the two- and three-cycles are separated by their displayed angles. Thus the physical first return counts are $2,3,4,4$, not freely assigned labels. The four selected contacts are interior to their squares, so the same finite return branches exist in neighborhoods. Direct subtraction of the first two rows gives determinant of magnitude one. Evaluation first gives $2s=c=3s$, hence $s=c=0$; the last two rows then give $u_1=u_2=0$.
\end{proof}

The qualification about evaluation is essential. A merely measurable transfer function is not automatically defined at a periodic point. A spectral application requires its own regularity or Liv\v{s}ic argument. Moreover, adding a roof twist $b\tau_R^Y$ to this proposition leaves $bL_j$ in every periodic equation. The four rows do not eliminate it. Joint roof nonarithmeticity and returned temporal nonintegrability remain distinct requirements. As one smaller unconditional conclusion, \eqref{eq:two-three} excludes a continuous cohomology $\tau_R=c+g-g\circ T_R$: the two periodic averages are $1-2R$ and $1-\sqrt3R$, which differ by $(2-\sqrt3)R>0$.


\section{A physical period family and the complete periodic annihilator}
\label{sec:joint-arithmetic}
The four records in Proposition~\ref{prop:lattice} settle the integer
coordinates at zero roof frequency. We now add a family of actual periodic
trajectories that settles every roof frequency. A slightly larger section is
specified explicitly; the original section and every assertion about it
remain as stated in the preceding sections.

\subsection{The enlarged section}
Put
\begin{equation}\label{eq:enlarged-section}
 U=\{(\alpha,p):|\alpha-\pi/6|<3/50,\ |p|<3/20\},
 \qquad Y_R^*=Y_R\cup U,\qquad c_*:=\nu(Y_R^*)=\frac{91}{10000\pi}.
\end{equation}
The rectangle $U$ is disjoint from the four squares defining $Y_R$.
At the square with the same angular center the $p$ coordinate is near
$-1/2$; the other three squares are separated in angle or in $p$.
The area of $U$ is $4(3/50)(3/20)$, and the section density is
$1/(4\pi)$, proving the value of $c_*$. Write
$F_R^*=T_R^{r_R^*}$ for the first return to $Y_R^*$, and let
$\kappa_R^*$ and $\tau_R^*$ be the sums of the physical displacement
and flight length before that return. Its invariant probability is
$\nu_R^*=\nu|_{Y_R^*}/c_*$. An assertion about $F_R$ does not become
an assertion about $F_R^*$ merely by changing notation.

Rotate coordinates through $-\pi/6$ and put
\begin{gather*}
 D=\sqrt3,\qquad O=(0,0),\quad C=(D,0),\quad
 A=(D/2,-1/2),\quad B=(D/2,1/2),\\
 d=\frac12-R,\qquad g=D-2R.
\end{gather*}
The rotated lattice consists of $(Dk/2,l/2)$ with integers $k,l$ of
the same parity. The normal orbit between $O$ and $C$ is unobstructed:
the two intermediate centers $A,B$ have distance $1/2$ from its
supporting line, and all other centers are farther away. It has two
physical flights of total length $2g$, zero winding, and exactly one
visit to $Y_R^*$ per period. Its selected state is $(\pi/6,0)$ in the
unrotated coordinates. The four cycles of Proposition~\ref{prop:lattice}
still have exactly one visit per period. Indeed, the equilateral cycle
has $p=-1/2$, the two-cycle has angles $0,\pi$, and the winding cycles
have $|p|=\sin\theta_R>0.79$. None gains a visit to $U$.

\subsection{A convex optical action}
For $-d\le y,z\le0$ define
\begin{align}
 x_R(y)&=\sqrt{R^2-y^2},\qquad u_R(y)=R-x_R(y),\nonumber\\
 e_R(y)&=\sqrt{(D/2-x_R(y))^2+(d+y)^2},\label{eq:excursion-lengths}\\
 \ell_R(y,z)&=\sqrt{(D-x_R(y)-x_R(z))^2+(z-y)^2}.\nonumber
\end{align}
The function $e_R$ is the length from the top contact
$q_A=(D/2,-d)$ of the disk at $A$ to a contact at height $y$ on the
facing arc of the disk at $O$ or $C$. The function $\ell_R$ is a
flight length between the two facing arcs. For $m\ge1$ let
\begin{equation}\label{eq:excursion-action}
 \mathcal A_{m,R}(y_1,\ldots,y_{2m})
 =e_R(y_1)+\sum_{j=1}^{2m-1}\ell_R(y_j,y_{j+1})+e_R(y_{2m}),
 \qquad L_m(R)=\min_{[-d,0]^{2m}}\mathcal A_{m,R}.
\end{equation}

\begin{lemma}[A uniformly regular physical family]\label{lem:excursion-family}
For each $m\ge1$ and $R\in I$, the minimum in
\eqref{eq:excursion-action} is unique. Its coordinates satisfy
\begin{equation}\label{eq:excursion-height}
 -\frac{2d}{5}<y_j<0,\qquad y_j=y_{2m+1-j}.
\end{equation}
The contacts with these heights form a genuine specular periodic orbit
with center word $A,(O,C)^m,A$. Every flight avoids all third disks.
The orbit has zero winding, $2m+1$ physical collisions, and exactly
$m$ returns to $Y_R^*$. The period length $L_m(R)$ is analytic in $R$
for each fixed $m$.
\end{lemma}
\begin{proof}
We supply the variational, mechanical and return checks separately.
On the whole box, $x_R(y)>11/25$, $x_R(y)\le47/100$, and
$1.73<D<1.74$. The positive horizontal quantities
\[
 a(y)=D/2-x_R(y),\qquad h(y,z)=D-x_R(y)-x_R(z)
\]
satisfy $a>79/200$ and $h>79/100$. The second derivative of
$-x_R(y)$ is $R^2/x_R(y)^3>0$. The Hessian formula for a Euclidean
norm shows that the Hessian of $e_R$ or $\ell_R$ is the sum of a
positive semidefinite term and the Hessian of its horizontal
component multiplied by its positive horizontal direction cosine.
The latter cosine exceeds $9/10$ on this box. Thus each segment
length is strictly convex in its variables; in fact the Hessian of
$\mathcal A_{m,R}$ dominates the identity. One may check the last
numerical bound from
$(9/10)(9/20)^2/(47/100)^3>1$. Compactness and strict convexity give
a unique minimizer.

No coordinate equals $-d$. For such a coordinate $y$ and an incident
interior segment, the derivative is
\[
 \partial_y\ell_R(y,z)=\frac{h(y,z)y/x_R(y)+y-z}{\ell_R(y,z)}<0,
\]
since $z\ge y=-d$. For an end segment,
$e_R'(-d)=-a(-d)d/(x_R(-d)e_R(-d))<0$.
Increasing that coordinate therefore decreases the action.
At an end coordinate equal to zero, $e_R'(0)=d/e_R(0)>0$, and
all other incident derivatives are nonnegative. At an interior zero
adjacent to a negative coordinate, its derivative is positive.
If a block of zero coordinates had no such neighbor it would reach
an endpoint. A variation into the box excludes all these cases.
Every coordinate is consequently in $(-d,0)$. Reflection in the
line $x=D/2$ reverses the center word and preserves the action, so
uniqueness gives the symmetry in \eqref{eq:excursion-height}.

Here is a uniform improvement of the height bound. A smallest
coordinate $y=-w$ cannot be interior, because both incident
derivatives there are negative. Take it to be $y_1$, using reversal
if needed, and write $z=y_2$. Stationarity gives
\[
 d-w=\left(a+e_R(y)\frac{h}{\ell_R(y,z)}\right)\frac{w}{x_R(y)}
           +\frac{e_R(y)}{\ell_R(y,z)}(w+z).
\]
The last term is nonnegative since $z\ge-w$. Moreover,
\[
 \frac{a+e_R(y)h/\ell_R(y,z)}{x_R(y)}
 >\frac{(19/10)(79/200)}{47/100}=\frac{1501}{940}>\frac32.
\]
It follows that $w<2d/5$. This bounds every coordinate.

Interior stationarity of the optical length is precisely equality
of tangential incoming and outgoing velocities at each contact.
The outgoing directions point into the exterior normal half-plane.
For example, at a contact on $O$, its outward normal is
$(x_R(y),y)/R$. Its dot product with a direction to $C$ has numerator
at least $x_R(y)h-|y||z-y|>0$; for a direction to $q_A$ both the
horizontal contribution and $y(-d-y)$ are positive. The checks at
$C$ are reflections of these. The symmetry of the two end contacts
implies the reflection law at $q_A$, whose normal is vertical. There
$e_R(y_1)<11/25$ and $d+y_1>3d/5\ge9/500$, so the normal
incidence cosine exceeds $9/220>1/25$. The other contacts are
uniformly transverse by the preceding horizontal-component bounds. Thus the stationary polygon
is specular rather than the straight-through solution of the
stationarity equation.

The long flights lie in $-2d/5<y<0$, so their distance from the disk
at $A$ is at least $3d/5\ge9/500$, and from the disk at $B$ at
least $d\ge3/100$. End flights join $q_A$ to the appropriate facing
arc and leave both endpoint disks in their exterior supporting
half-planes. Their third disk $O$ or $C$ has horizontal separation
at least $D/2>R$, and $B$ is vertically separated. All segments
lie in the rectangle
\[
       11/25<x<D-11/25,\qquad -d\le y\le0.
\]
For the remaining lattice centers, the possibilities $k=0,2$ with
nonzero $l$ have vertical separation at least $1-d\ge0.95$;
$k=1$, $|l|\ge3$ have separation at least $3/2-d$; and
$k\notin\{0,1,2\}$ have horizontal separation exceeding $1.3$.
This enumerates the whole lattice and proves nonocclusion, not just
a check of a finite list of nearby obstacles.

At an outgoing contact on $O$, the normal angle differs from
$\pi/6$ by at most $|y|/x_R(y)<1/22<3/50$. The outgoing velocity
points toward $C$ and its angular deviation from the horizontal is
at most $|z-y|/h<2/79$. Thus
\[
 |p|\le \frac1{22}+\frac2{79}<\frac3{20}.
\]
All $m$ contacts on $O$ belong to $U$. Contacts on $C$ have angles
near $7\pi/6$, and $q_A$ has angle $2\pi/3$; none belongs to
$Y_R^*$. The return counts are $2$ between successive contacts on
$O$, with one count $3$ through $C,A,O$. Their sum is $2m+1$,
and there are $m$ induced returns. The orbit is closed in the plane,
so its winding is zero. All selected contacts and unselected contacts
have strict membership margins, giving actual regular return branches.
Finally the positive Hessian and the interior minimum allow the
analytic implicit function theorem. This proves analyticity for each
fixed $m$; no bound for high parameter derivatives uniform in $m$ is
being inferred.
\end{proof}

\subsection{Strictly increasing bounded excess lengths}
The following argument supplies an infinite sequence of period relations.
It does not infer nonarithmeticity from finitely many decimal lengths.

\begin{theorem}[Strict period excess]\label{thm:period-excess}
For the family of Lemma~\ref{lem:excursion-family}, define
$E_m(R)=L_m(R)-2m g$. Then, for every $R\in I$,
\begin{equation}\label{eq:period-excess}
 0<E_m(R)<E_{m+1}(R)\le \frac{2d^2}{g}<\frac1{100}.
\end{equation}
In particular $E_{m+1}(R)-E_m(R)>0$ and these differences tend to
zero. The convergence of $E_m$ is used pointwise in $R$; the subsequent
compact-frequency assertion supplies its own uniformity argument.
\end{theorem}
\begin{proof}
Suppress $R$ and set $K(y,z)=\ell(y,z)-g$ on $[-d,0]^2$.
The horizontal component of the flight gives
\begin{equation}\label{eq:kernel-coercivity}
 K(y,z)\ge u(y)+u(z),
\end{equation}
with equality exactly when $y=z$. Define continuous value functions
\[
 V_0(y)=u(y),\qquad
 V_{j+1}(y)=\min_{z\in[-d,0]}\{K(y,z)+V_j(z)\}.
\]
They are nonnegative and vanish at zero. For $j\ge1$, choosing
$z=0$ gives $V_j(y)\le K(y,0)=O(y^2)$ near zero; $V_0$ has the
same bound. The implied bound is independent of $j$.
Equation~\eqref{eq:kernel-coercivity} implies $V_1(y)>V_0(y)$ for
$y<0$: equality at a minimizing $z$ would require both $z=0$ and
equality in \eqref{eq:kernel-coercivity}, hence $y=0$.

For any $y<0$, zero is not a minimizer of $K(y,z)+V_j(z)$. Indeed,
\[
 \partial_zK(y,0)=\frac{-y}{\ell(y,0)}>0,
\]
whereas $0\le V_j(z)=O(z^2)$. A sufficiently small negative $z$
therefore reduces the value. If $V_j>V_{j-1}$ on $[-d,0)$, evaluate
at a minimizer $z_*<0$ for $V_{j+1}(y)$. It follows that
\[
 V_{j+1}(y)=K(y,z_*)+V_j(z_*)
   >K(y,z_*)+V_{j-1}(z_*)\ge V_j(y).
\]
Induction proves the strict inequality for every $j$ and $y<0$.

The minimizing path is symmetric, and its middle edge has length
$\ell(y_m,y_m)=g+2u(y_m)$. Splitting the action at that edge gives
the exact identity
\begin{equation}\label{eq:bellman-excess}
 E_m=2\min_{y\in[-d,0]}\{e(y)-g/2+V_{m-1}(y)\}.
\end{equation}
The outer minimizer cannot be zero, since $e'(0)=d/e(0)>0$ and
$V_{m-1}(y)=O(y^2)$. Evaluating the formula for $E_{m+1}$ at its
negative minimizer, and using strict increase of $V_j$ there, proves
$E_{m+1}>E_m$. Since $e(y)\ge D/2-x_R(y)\ge g/2$, equality
$E_m=0$ would force $y=0$ and then $d=0$, which is excluded.
Taking all heights zero gives
\[
 E_m\le2\sqrt{(g/2)^2+d^2}-g
      =\frac{2d^2}{\sqrt{(g/2)^2+d^2}+g/2}\le\frac{2d^2}{g}.
\]
Here $d\le1/20$ and $g>79/100$, proving the last bound in
\eqref{eq:period-excess}. A bounded increasing sequence has summable
positive increments, which must tend to zero.
\end{proof}

\subsection{Joint arithmetic and its exact scope}
A regular periodic return orbit has a total record $(K,N,T,h)$:
winding $K\in\Z^2$, physical collision count $N$, flight length $T$,
and induced period $h$. Define its phase at
$(u,s,\beta,c)\in\T^2\times\T\times\R\times\T$ to be
$u\cdot K+sN+\beta T-ch$, with $\T=\R/(2\pi\Z)$.

\begin{theorem}[Trivial joint periodic annihilator]\label{thm:joint-periodic}
For every $R\in I$, if
\begin{equation}\label{eq:joint-periodic-phase}
 \exp i(u\cdot K+sN+\beta T-ch)=1
\end{equation}
for every regular periodic orbit of $F_R^*$, then
$\beta=0$ and $u=s=c=0$ on their tori. It suffices to test the four
cycles of Proposition~\ref{prop:lattice}, the long normal two-cycle,
and the family of Lemma~\ref{lem:excursion-family}.
\end{theorem}
\begin{proof}
The long normal cycle gives $\exp i(2s+2\beta g-c)=1$.
The $m$th excursion gives
$\exp i((2m+1)s+\beta L_m-mc)=1$. Dividing the latter by the
$m$th power of the former yields
\[
                    \exp i(s+\beta E_m)=1.
\]
Taking consecutive quotients shows that
$\beta(E_{m+1}-E_m)\in2\pi\Z$ for all $m$. If $\beta\ne0$,
Theorem~\ref{thm:period-excess} makes these quantities nonzero and
strictly smaller than $2\pi$ in absolute value for large $m$, a
contradiction. Hence $\beta=0$. The four unchanged first-return
records in \eqref{eq:matrix} now give $u=s=c=0$.
\end{proof}

\begin{corollary}[Finite uniform separation on compact frequency sets]
\label{cor:compact-periodic}
For every $B<\infty$ and $\varepsilon>0$, there are a finite list of
these physical periodic words and $\eta>0$ such that, uniformly in
$R\in I$ and in $|\beta|\le B$ at distance at least $\varepsilon$
from the trivial phase, at least one word has
$|\exp i(u\cdot K+sN+\beta T-ch)-1|\ge\eta$.
\end{corollary}
\begin{proof}
Each word has fixed integer record and continuous length on $I$.
Theorem~\ref{thm:joint-periodic} gives a nonzero discrepancy at every
point of the compact parameter--frequency set in question. The open
sets on which a given word has positive discrepancy cover that set.
Take a finite subcover and minimize the maximum of its discrepancies.
This minimum is positive and supplies $\eta$.
\end{proof}

\begin{corollary}[Full real period rank]\label{cor:real-period-rank}
The real vectors $(K_1,K_2,N,h,T)$ of five fixed cycles span $\R^5$.
The absolute determinant is $2(\sqrt3-1)$, independently of $R$.
Consequently no nonzero real linear combination of the four induced
record coordinates is cohomologous to a constant by a transfer
function for which periodic evaluation is valid.
\end{corollary}
\begin{proof}
Append lengths $L_2,L_3,L_4,L_4$ to the four rows in
\eqref{eq:matrix}, and append $(0,0,2,1,2g)$ as the fifth row.
Subtracting the first row from the fifth leaves only the last entry
$2g-L_2=2(\sqrt3-1)$. The remaining determinant is that of $V$,
of absolute value one. Evaluation of a real cohomology on these
cycles therefore forces all coefficients, including the constant,
to vanish.
\end{proof}

The preceding theorems are statements about real billiard orbits, not
merely formal integer records. They also exclude a circle-valued
coboundary with a continuous transfer function on the regular return
branches: an almost-everywhere identity then extends to each regular
periodic point by continuity and full support of the section measure.
A transfer function assumed only measurable requires a separate
regularity theorem. Similarly, a covariance conclusion requires a
proved variance--coboundary criterion in the relevant function space.
If that criterion and a continuous covariance family are established,
Corollary~\ref{cor:real-period-rank} gives positive definiteness and,
on $I$, a uniform positive lower eigenvalue. Existence and continuity
of that covariance are not consequences of the determinant alone.
Compact periodic separation is not a quantitative operator contraction
or a returned temporal nonintegrability estimate.

\section{Quantitative separation by logarithmic-length periods}
\label{sec:quantitative-periods}
The complete periodic annihilator was determined in
Theorem~\ref{thm:joint-periodic}. We now estimate the period increments
uniformly in both the radius and the number of collisions. This gives an
explicit separation at large roof frequencies, using three actual
periodic words. Throughout this section $g=\sqrt3-2R$, $d=1/2-R$,
and $e$, $\ell$ and $u$ are the functions in+\eqref{eq:excursion-lengths}. The coordinates below are heights, not
arc coordinates.

Define the half action on $[-d,0]^m$ by
\begin{equation}\label{eq:half-action-v4}
 W_{m,R}(y)=e_R(y_1)-g/2+
       \sum_{j=1}^{m-1}\bigl(\ell_R(y_j,y_{j+1})-g\bigr)+u_R(y_m).
\end{equation}
Reflection symmetry and uniqueness of the full minimizer imply
$E_m(R)=2\min W_{m,R}$. Denote its minimizing coordinates by
$y^{(m,R)}_1,\ldots,y^{(m,R)}_m$.

\begin{lemma}[Uniform Hessian and endpoint estimates]
\label{lem:uniform-half-hessian}
On the entire height box, the Hessian $H=\nabla^2 W_{m,R}$ is
tridiagonal and satisfies
\begin{equation}\label{eq:half-hessian-bounds}
 3I_m\le H\le12I_m,\qquad H_{jj}\le10,\qquad
 \frac12\le-H_{j,j+1}\le2\quad(1\le j<m).
\end{equation}
The last condition is empty when $m=1$. For all $1\le j\le m$,
\begin{equation}\label{eq:height-two-sided}
 \frac3{440}\,20^{-(j-1)}
 \le -y^{(m,R)}_j
 \le \frac1{21}\left(\frac7{10}\right)^{j-1}.
\end{equation}
All constants are independent of $m$ and $R\in I$.
\end{lemma}
\begin{proof}
Write $x(y)=\sqrt{R^2-y^2}$, $a=D/2-x(y)$,
$h=D-x(y)-x(z)$, $v=z-y$, with $D=\sqrt3$.
The elementary bounds on the whole box are
\[
 \begin{gathered}
 x>11/25,\quad |x'|\le5/44,\quad79/200<a<43/100,\\
 79/100<h<86/100,\quad e<11/25,\quad\ell<9/10.
 \end{gathered}
\]
Also $u''=R^2/x^3\ge1/R\ge100/47$ and
$u''\le55225/21296<3$. For the norm of a vector-valued function,
the Hessian is its transverse positive semidefinite quadratic form
plus the Hessian of the vector dotted with its unit direction.
Consequently
\[
 e''\ge(a/e)u''>\frac{79}{88}\frac{100}{47}>\frac32,
 \qquad
 \nabla^2\ell\ge(h/\ell)\operatorname{diag}(u''(y),u''(z))
 >\frac32 I_2.
\]
Each variable of $W_m$ occurs in two such terms, with $u''>3/2$
at the terminal endpoint. This proves the lower bound $H\ge3I_m$,
including $m=1$.

The same norm formula gives
\[
 \ell_{yy},\ell_{zz}
 \le\frac{1+(5/44)^2}{79/100}+\frac{55225}{21296}<4,
 \qquad
 e''\le\frac{1+(5/44)^2}{79/200}+\frac{55225}{21296}<6.
\]
An explicit mixed derivative is
\begin{equation}\label{eq:mixed-length-hessian}
 -\ell_{yz}
 =\frac{(h+vy/x(y))(h-vz/x(z))}{\ell^3}.
\end{equation}
Each numerator factor lies between $78/100$ and $87/100$.
Using $79/100<\ell<9/10$ proves $1/2<-\ell_{yz}<2$.
The diagonal bounds for $W_m$ are therefore $10$ at its first
endpoint, $8$ in its interior, and $7$ at its last endpoint; when
$m=1$ the bound is $9$. Absolute row sums are at most $12$.
The symmetric matrix bound $H\le12I_m$ follows, proving
\eqref{eq:half-hessian-bounds}.

At zero the gradient is $a_0 e_1$, where
$a_0=d/e_R(0)$ satisfies $3/44\le a_0\le1/7$.
Integrate the Hessian on the segment from zero to the minimizer and
write the resulting matrix as $A$. Stationarity gives
$Ay^{(m,R)}=-a_0e_1$. The matrix $B=10I_m-A$ is entrywise
nonnegative and tridiagonal. Its eigenvalues lie in $[-2,7]$, so
$\|B/10\|_2\le7/10$. Hence
\[
 A^{-1}=\frac1{10}\sum_{k\ge0}(B/10)^k.
\]
For the $(j,1)$ entry every power before $j-1$ vanishes. The
nearest-neighbor product at power $j-1$ is at least $20^{-(j-1)}$.
Entrywise nonnegativity and the norm bound yield
\[
 \frac1{10}20^{-(j-1)}
 \le (A^{-1})_{j1}
 \le\frac1{10}\sum_{k\ge j-1}(7/10)^k
 =\frac13(7/10)^{j-1}.
\]
Multiplication by the two bounds on $a_0$ proves
\eqref{eq:height-two-sided}. The matrix is an averaged Hessian of
the actual nonlinear action; no quadratic approximation of the
minimizing orbit has been made.
\end{proof}

\begin{theorem}[Two-sided excess increments]
\label{thm:quantitative-excess}
For every $R\in I$ and $m\ge1$,
\begin{equation}\label{eq:quantitative-excess}
 \frac1{10000}\,400^{-m}
 \le E_{m+1}(R)-E_m(R)
 \le\frac1{300}\left(\frac{49}{100}\right)^{m-1}.
\end{equation}
In particular $E_m$ converges uniformly on $I$ to a continuous
function $E_\infty$, with
\begin{equation}\label{eq:uniform-excess-limit}
 0<E_\infty(R)-E_m(R)
 \le\frac1{153}\left(\frac{49}{100}\right)^{m-1}.
\end{equation}
\end{theorem}
\begin{proof}
For a minimizer $y$ of $W_m$, append the coordinate zero. With
$K(y,z)=\ell(y,z)-g$, the resulting change in half action is
\[
 K(y_m,0)-u(y_m)
 =\sqrt{(g+u(y_m))^2+y_m^2}-(g+u(y_m))
 \le\frac{y_m^2}{2g}.
\]
It follows that $E_{m+1}-E_m\le y_m^2/g$. The upper bound in
\eqref{eq:height-two-sided} and $g>79/100$ give the upper
coefficient $100/34839<1/300$ in \eqref{eq:quantitative-excess}.

Conversely, truncate a minimizer $z$ of $W_{m+1}$ after its first
$m$ coordinates. The removed contribution is
\[
 K(z_m,z_{m+1})+u(z_{m+1})-u(z_m)\ge2u(z_{m+1}),
\]
by \eqref{eq:kernel-coercivity}. Thus
\[
 E_{m+1}-E_m\ge4u(z_{m+1})\ge\frac{2z_{m+1}^2}{R}
 \ge\frac9{45496}\,400^{-m}>\frac1{10000}\,400^{-m}.
\]
Here $u(z)=z^2/(R+x(z))\ge z^2/(2R)$ and
\eqref{eq:height-two-sided} were used. Both comparisons concern
minima of the exact actions. Summing the geometric upper bound
proves \eqref{eq:uniform-excess-limit}. Each $E_m$ is continuous
by Lemma~\ref{lem:excursion-family}; the uniform limit is continuous.
\end{proof}

Let $P_0$ denote the long normal orbit and $P_m$ the $m$th excursion.
For a periodic word $P$ write
\[
 Z_P(u,s,\beta,c;R)
 =\exp i\bigl(u\cdot K_P+sN_P+\beta T_P-ch_P\bigr).
\]
Put $a=\log(100/49)$ and, for $b\ge1$, define
\begin{equation}\label{eq:frequency-period-budget}
 m(b)=1+\left\lceil\frac{\log b}{a}\right\rceil,
 \qquad \gamma=\frac{\log400}{a}-1.
\end{equation}

\begin{theorem}[Explicit large-roof phase separation]
\label{thm:log-period-separation}
For $|\beta|=b\ge1$, all $R\in I$, and arbitrary torus
frequencies $u,s,c$, the three true periods $P_0,P_{m(b)},P_{m(b)+1}$
satisfy
\begin{equation}\label{eq:log-period-separation}
 \max_P|Z_P-1|
 \ge\frac{b\,400^{-m(b)}}{15000\pi}
 \ge\frac{b^{-\gamma}}{2400000000\pi}.
\end{equation}
The exponent is $\gamma=\log400/\log(100/49)-1$.
For a whole band $1\le|\beta|\le B$, the fixed three periods
$P_0,P_{m(B)},P_{m(B)+1}$ suffice, with lower bound
$400^{-m(B)}/(15000\pi)$. Their longest physical word has
$2m(B)+3$ collisions. All assertions are uniform in the radius.
\end{theorem}
\begin{proof}
The actual counts and lengths give the exact identity
\begin{equation}\label{eq:three-period-cancellation}
 Z_{P_{m+1}}\overline{Z_{P_m}}\overline{Z_{P_0}}
       =\exp\bigl(i\beta(E_{m+1}-E_m)\bigr).
\end{equation}
The winding, physical count and induced period cancel. With $m=m(b)$,
\eqref{eq:quantitative-excess} implies
\[
 \frac{b\,400^{-m}}{10000}
 \le b(E_{m+1}-E_m)\le\frac1{300}.
\]
For $|v|\le\pi$, $|e^{iv}-1|\ge2|v|/\pi$. The distance
from one of a product of three unit complex numbers is at most
the sum of their individual distances from one. Applying these
facts to \eqref{eq:three-period-cancellation} proves the first
bound in \eqref{eq:log-period-separation}. Since
$m(b)\le2+\log b/a$, the second bound follows. For the band
assertion use $m(B)$; the upper phase bound still holds, while
$b\ge1$ supplies the printed lower bound. The number of collisions
was proved in Lemma~\ref{lem:excursion-family}.
\end{proof}

\begin{corollary}[An approximate-phase lower bound]
\label{cor:approximate-phase-bound}
Let $q:Y_R^*\to\mathbb C$ have modulus one and be continuous in
neighborhoods of all states in the three cycles used above. Suppose
$q$ is measurable elsewhere and let
\[
 D(q)=\mathop{\rm ess\,sup}_{x\in Y_R^*}
 \left|e^{i(u\cdot\kappa_R^*(x)+s r_R^*(x)+\beta\tau_R^*(x))}
                 q(F_R^*x)-e^{ic}q(x)\right|.
\]
For $b=|\beta|\ge1$ and $m=m(b)$,
\begin{equation}\label{eq:approximate-phase-bound}
 D(q)\ge\frac{b\,400^{-m}}{10000\pi(m+1)}.
\end{equation}
\end{corollary}
\begin{proof}
The return branches through these cycles are regular, and their
section measure has full support. Continuity therefore extends the
essential bound to every state of each selected cycle. Multiplication
around a cycle of induced period $h$ telescopes the unit phases and
gives $|Z_P-1|\le hD(q)$. In
\eqref{eq:three-period-cancellation} the sum of the three periods is
$1+m+(m+1)=2(m+1)$. The lower sine bound in the preceding proof
therefore gives
$2b(E_{m+1}-E_m)/\pi\le2(m+1)D(q)$. Apply the lower bound
in \eqref{eq:quantitative-excess}.
\end{proof}

The estimates are quantitative statements about the specified mechanical
return map, and in particular about continuous circle-valued approximate
phases. They do not assert regularity of a merely measurable transfer
function, a lower bound for a spectral eigenfunction on these orbits,
or a transfer-operator resolvent estimate. Those are distinct analytical
steps when using \eqref{eq:approximate-phase-bound} in a high-frequency
argument. The raw characteristic-function tail is not replaced by the
period discrepancy.

\section{Uniform phase separation and positive-measure coercivity}
\label{sec:periodic-coercivity}
We strengthen the frequency estimate by comparing consecutive exact
minimizers. The choice of a period will depend on the frequency and
possibly on the radius. This dependence is essential: a fixed triple
of periods is not asserted to separate the entire unbounded frequency
axis by a positive constant.

Retain $K_R(y,z)=\ell_R(y,z)-g$ and $W_{m,R}$ from
\eqref{eq:half-action-v4}. Write
\[
 a_m(R)=-y_m^{(m,R)}>0,\qquad
 \Delta_m(R)=E_{m+1}(R)-E_m(R)>0.
\]
The symbol $a_m$ denotes a terminal height; it is unrelated to the
logarithmic constant in \eqref{eq:frequency-period-budget}.

\subsection{Comparison of successive exact minima}
\begin{lemma}[The terminal value function]
\label{lem:terminal-response}
For $y\in[-d,0]$ put
\[
 V_R(y)=\min_{z\in[-d,0]}\{K_R(y,z)+u_R(z)\},\qquad
 \delta_R(y)=V_R(y)-u_R(y).
\]
The minimizer $f_R(y)$ is unique. It belongs to $(-d,0)$ for $y<0$
and equals zero for $y=0$. The functions are twice continuously
differentiable, with one-sided derivatives at the endpoints, and
\begin{equation}\label{eq:terminal-response}
 f_R'(y)\ge\frac1{14},\qquad
 |f_R(y)|\ge\frac{|y|}{14},\qquad
 |\delta_R''(y)|\le2,\qquad |\delta_R'(y)|\le2|y|.
\end{equation}
These estimates hold on the entire height interval, uniformly in $R$.
\end{lemma}
\begin{proof}
The second derivative in $z$ is $K_{zz}+u''>0$. At $z=-d$,
\[
 K_z(y,-d)=\frac{h(y,-d)(-d)/x_R(-d)-d-y}{\ell_R(y,-d)}<0,
 \qquad u_R'(-d)<0.
\]
At $z=0$ and $y<0$ the derivative is $-y/\ell_R(y,0)>0$.
Consequently the minimizer is interior and unique. For $y=0$,
\eqref{eq:kernel-coercivity} gives the unique minimizer zero.
The implicit function theorem applies to the stationarity equation,
including its smooth extension near the endpoints, and gives
\[
 f_R'(y)=\frac{-K_{yz}(y,f_R(y))}
                    {K_{zz}(y,f_R(y))+u_R''(f_R(y))}
 \ge\frac{1/2}{4+3}=\frac1{14}.
\]
Here the individual segment bounds established in the proof of
Lemma~\ref{lem:uniform-half-hessian} were used. Integration from $y$
to zero proves the second inequality of \eqref{eq:terminal-response}.

The envelope formula is
\[
 V_R''(y)=K_{yy}(y,f_R(y))-
       \frac{K_{yz}(y,f_R(y))^2}{K_{zz}(y,f_R(y))+u_R''(f_R(y))}.
\]
The Hessian of $K_R(y,z)+u_R(z)$ dominates $(3/2)I_2$.
Its Schur complement therefore is at least $3/2$: evaluate its
quadratic form on $(1,t)$ and minimize over $t$. Also
$V_R''\le K_{yy}\le4$. Since $100/47\le u_R''<3$, subtraction
shows $|\delta_R''|\le2$. The first derivatives of $K_R$ and $u_R$
vanish at the origin, so $\delta_R'(0)=0$. Integration proves the
last inequality. All estimates concern the nonlinear action itself.
\end{proof}

\begin{proposition}[Consecutive terminal heights]
\label{prop:terminal-comparison}
For every $m\ge1$ and $R\in I$,
\begin{equation}\label{eq:terminal-comparison}
 a_{m+1}(R)\ge\frac3{70}a_m(R).
\end{equation}
\end{proposition}
\begin{proof}
Let $y$ minimize $W_{m,R}$ and let $(X,z)$ minimize $W_{m+1,R}$,
where $X$ consists of the first $m$ coordinates. Eliminating the
last coordinate gives exactly
\[
 \min_z W_{m+1,R}(X,z)=W_{m,R}(X)+\delta_R(X_m),
 \qquad z=f_R(X_m).
\]
All minimizing coordinates are interior by
Lemma~\ref{lem:excursion-family}. Stationarity and the fundamental
theorem of calculus thus yield
\[
 A(X-y)=-\delta_R'(X_m)e_m,\qquad
 A=\int_0^1\nabla^2W_{m,R}(y+t(X-y))\dd t\ge3I_m.
\]
The segment stays in the height box. In particular
$0<(A^{-1})_{mm}\le1/3$. Taking the last coordinate and applying
Lemma~\ref{lem:terminal-response}, we obtain
\[
 |y_m|\le |X_m|+\frac13|\delta_R'(X_m)|
             \le\frac53|X_m|.
\]
Finally $|z|=|f_R(X_m)|\ge|X_m|/14$. These two inequalities give
\eqref{eq:terminal-comparison}. No monotonicity of the terminal
heights, and no quadratic replacement of $W_{m,R}$, is assumed.
\end{proof}

\begin{corollary}[A lower ratio for the excess increments]
\label{cor:increment-ratio}
For every $m\ge1$ and $R\in I$,
\begin{equation}\label{eq:increment-ratio}
 \Delta_{m+1}(R)\ge10^{-5}\Delta_m(R).
\end{equation}
\end{corollary}
\begin{proof}
The append and truncate comparisons in the proof of
Theorem~\ref{thm:quantitative-excess} give
\[
 \Delta_m\le\frac{a_m^2}{g},\qquad
 \Delta_{m+1}\ge\frac2R a_{m+2}^2.
\]
Applying \eqref{eq:terminal-comparison} twice, and using
$g>79/100$ and $R\le47/100$, yields
\[
 \Delta_{m+1}\ge\frac{2g}{R}\left(\frac3{70}\right)^4\Delta_m
 >\frac{158}{47}\left(\frac3{70}\right)^4\Delta_m
 >\frac{\Delta_m}{100000}.
\]
This is a comparison between adjacent increments, stronger for the
present purpose than separate exponential upper and lower bounds.
\end{proof}

\subsection{A frequency-dependent triple and a fixed finite menu}
Put
\begin{equation}\label{eq:uniform-phase-constants}
 \eta_0=\frac1{4000000},\qquad
 \sigma_0=\frac1{6000000\pi},\qquad
 M(b)=m(b)+1,
\end{equation}
where $m(b)$ is defined in \eqref{eq:frequency-period-budget}.

\begin{theorem}[Uniform positive phase separation]
\label{thm:uniform-phase-separation}
For every $R\in I$ and $b\ge1$ the integer
\begin{equation}\label{eq:adaptive-period}
 j(R,b)=\min\{j\ge1:b\Delta_j(R)\le1/2\}
\end{equation}
is well defined and satisfies $j(R,b)\le m(b)$. If $|\beta|=b$,
then, for arbitrary $u,s,c$,
\begin{equation}\label{eq:uniform-phase-separation}
 \max_{P\in\{P_0,P_{j(R,b)},P_{j(R,b)+1}\}}
       |Z_P(u,s,\beta,c;R)-1|\ge\sigma_0.
\end{equation}
The longest of these physical words has at most $2m(b)+3$ collisions.
For the entire band $1\le|\beta|\le B$, the fixed menu
\begin{equation}\label{eq:fixed-period-menu}
 \mathcal P_B(R)=\{P_0(R),P_1(R),\ldots,P_{m(B)+1}(R)\}
\end{equation}
satisfies the same lower bound for its maximum, uniformly in $R$.
It contains $m(B)+2$ periods.
\end{theorem}
\begin{proof}
The upper bound of \eqref{eq:quantitative-excess} gives
$b\Delta_{m(b)}\le1/300$. Thus the defining set in
\eqref{eq:adaptive-period} is nonempty. If $j(R,b)=1$, its lower
bound gives $b\Delta_1\ge\Delta_1\ge1/4000000=\eta_0$.
If $j(R,b)>1$, minimality and \eqref{eq:increment-ratio} give
\[
 b\Delta_{j(R,b)}>10^{-5}/2\ge\eta_0.
\]
In either case
\begin{equation}\label{eq:selected-phase-window}
 \eta_0\le b\Delta_{j(R,b)}\le1/2.
\end{equation}
The exact three-period identity \eqref{eq:three-period-cancellation}
therefore has distance at least $2\eta_0/\pi$ from one. The distance
of a product of three unit complex numbers from one is at most the
sum of their three distances. This proves
\eqref{eq:uniform-phase-separation}, since
$2\eta_0/(3\pi)=\sigma_0$. The physical and induced counts are
those of Lemma~\ref{lem:excursion-family}; hence the stated word
length. Finally $m(b)\le m(B)$ on the band, so the selected triple
is contained in \eqref{eq:fixed-period-menu}.

The index selection uses a first crossing, not monotonicity of
$\Delta_j$. It need not be continuous in $R$ or $b$. The menu has
fixed indices and consists of analytically varying true periods,
which is the asserted parameter-uniform conclusion.
\end{proof}

\begin{corollary}[Logarithmic uniform-norm coercivity]
\label{cor:logarithmic-coercivity}
Let $q$ be circle-valued and measurable on $Y_R^*$, and continuous
in neighborhoods of every state of the triple in
\eqref{eq:uniform-phase-separation}. With $D(q)$ as in
Corollary~\ref{cor:approximate-phase-bound}, for $b=|\beta|\ge1$,
\begin{equation}\label{eq:logarithmic-coercivity}
 D(q)\ge\frac1{4000000\pi(j(R,b)+1)}
       \ge\frac1{4000000\pi M(b)}.
\end{equation}
\end{corollary}
\begin{proof}
Continuity and full support extend the essential bound to the
selected regular periodic states. Telescoping around each period
bounds its phase discrepancy by its induced length times $D(q)$.
The three induced lengths sum to $2(j(R,b)+1)$. Apply the product
identity and \eqref{eq:selected-phase-window} to obtain
\[
 2\eta_0/\pi\le2(j(R,b)+1)D(q).
\]
This proves \eqref{eq:logarithmic-coercivity}. The modulus-one
condition is used in telescoping and is not a consequence of an
arbitrary operator normalization.
\end{proof}

\subsection{Regular tubes of the actual return map}
We use the Euclidean metric in local $(\alpha,p)$ coordinates. The
following uniformity concerns the specified periodic family, not all
branches of the induced map.

\begin{lemma}[Uniform return neighborhoods of the period family]
\label{lem:uniform-period-tubes}
There are constants $0<\rho_0<1$, $L\ge2$ and $B_0\ge1$, independent
of $R$ and $m$, with the following properties. For every selected
state $x_P$ of $P_0$ or any $P_m$, the ball $B(x_P,\rho_0)$ lies
in $Y_R^*$ and in one regular first-return branch. On this ball
$r_R^*$ and $\kappa_R^*$ are constant, $r_R^*\in\{2,3\}$, and
\[
 \operatorname{Lip}(F_R^*)\le L,\qquad
 \operatorname{Lip}(\tau_R^*)\le B_0.
\]
If $P$ has induced period $h$ and $0<\rho\le\rho_0L^{-h}$, every
$x\in B(x_P,\rho)$ follows the same $h$ return branches as $x_P$,
and, writing $F=F_R^*$,
\begin{align}\label{eq:period-tube-bounds}
 d(F^k x,F^k x_P)&\le L^k\rho &&(0\le k\le h),\\
 |T_{h,R}(x)-T_{h,R}(x_P)|&\le2B_0L^h\rho,\\
 \nu_R^*(B(x_P,\rho))&=\rho^2/(4c_*).
\end{align}
\end{lemma}
\begin{proof}
We verify that the family stays a uniform distance from every local
singularity relevant to a single return. The estimates in
Lemma~\ref{lem:excursion-family} hold for all its contact heights:
$-2d/5<y<0$, the long-flight third-disk clearances are at least
$9/500$, and the incidence cosine at $A$ exceeds $1/25$.
At the facing contacts on $O,C$ it is bounded below by a positive
constant from the horizontal-component inequalities in that proof.
End flights have a uniform positive distance from every third disk:
the opposite facing center is horizontally separated by $D/2$,
$B$ is vertically separated, and the remaining lattice centers obey
the enumerated uniform separation bounds. All flight lengths are
bounded above and below by positive constants.

Return membership has uniform margins as well. The selected contacts
on $O$ satisfy
\[
 |\alpha-\pi/6|\le1/22<3/50,\qquad
 |p|\le1/22+2/79<3/20.
\]
The unselected contacts on $C$ and $A$ have angles near $7\pi/6$
and equal to $2\pi/3$, respectively, uniformly separated from $U$
and the four squares in $Y_R$. The long normal orbit has the same
strict margins. Each selected-to-selected path consequently contains
two physical flights, or three through $C,A,O$.

For clarity, compactness is applied only to these bounded-length
paths. The parameters $R$ and their contact coordinates range in a
compact box. Every limiting path satisfies the non-grazing,
nonocclusion and return-membership margins just established. In a
neighborhood of such a path the selected entry roots are simple;
the physical collision map and flight time are analytic in the
initial state and parameter. There are only the finitely many
center itineraries just listed, up to lattice translation. A finite
cover of their compact closures therefore gives a common neighborhood
radius and uniform first-derivative bounds for the actual two- or
three-flight first returns. Decrease the radius so its coordinate
balls lie in $Y_R^*$ and contain no change of lattice labels; enlarge
the derivative bounds to $L\ge2$, $B_0\ge1$. This proves the first
part without taking a derivative of an arbitrarily long return.

Starting in a ball of radius at most $\rho_0L^{-h}$, induction gives
the first inequality in \eqref{eq:period-tube-bounds} and preserves
the same return decisions at each step. Summing the roof variations
gives $B_0\rho\sum_{k<h}L^k\le2B_0L^h\rho$. Finally
$\dd\nu_R^*=(4\pi c_*)^{-1}\dd\alpha\dd p$ on the ball, whose
Euclidean area is $\pi\rho^2$. This proves the last equality.
\end{proof}

\begin{theorem}[Positive-measure coercivity for regular phases]
\label{thm:positive-measure-coercivity}
Fix $0<\alpha\le1$, $H\ge1$ and $1\le p<\infty$. Suppose $q$ is
measurable and circle-valued on $Y_R^*$ and has an $\alpha$-H\"older
representative, of constant at most $H$, on the balls of radius
$\rho_0$ around all states of the selected triple. Define
\[
 \mathcal D_q(x)=e^{i(u\cdot\kappa_R^*(x)+s r_R^*(x)+\beta\tau_R^*(x))}
                      q(F_R^*x)-e^{ic}q(x).
\]
For $b=|\beta|\ge1$, put
\begin{equation}\label{eq:coercivity-radius}
 \rho(b,H)=L^{-M(b)}
 \min\left\{\frac{\rho_0}{2},\frac{\sigma_0}{8bB_0},
                       \left(\frac{\sigma_0}{8H}\right)^{1/\alpha}\right\}.
\end{equation}
Then the actual invariant return probability satisfies
\begin{equation}\label{eq:positive-measure-coercivity}
 \|\mathcal D_q\|_{L^p(\nu_R^*)}
 \ge\frac{\sigma_0}{2M(b)}
             \left(\frac{\rho(b,H)^2}{4c_*}\right)^{1/p}.
\end{equation}
All geometric constants are uniform over $R\in I$. No independence
of the returns is assumed.
\end{theorem}
\begin{proof}
Select from the triple a period $P$ with $|Z_P-1|\ge\sigma_0$.
Let its selected state be $x_P$ and its induced period be $h\le M(b)$.
Write $\varphi=u\cdot\kappa_R^*+s r_R^*+\beta\tau_R^*$, $F=F_R^*$
and $S_h\varphi=\sum_{k<h}\varphi\circ F^k$. The lattice and count
terms are constant along the corresponding branches in the tube.
Thus for $x\in B(x_P,\rho(b,H))$, Lemma~\ref{lem:uniform-period-tubes}
gives
\[
 |S_h\varphi(x)-S_h\varphi(x_P)|\le2bB_0L^h\rho(b,H).
\]
At $x_P$, the modulus of
$e^{iS_h\varphi(x_P)}q(F^h x_P)-e^{ich}q(x_P)$ is $|Z_P-1|$.
Since $F^h x_P=x_P$, the H\"older bounds at the two endpoints and
\eqref{eq:coercivity-radius} imply, throughout the smaller ball,
\begin{align*}
 \left|e^{iS_h\varphi(x)}q(F^h x)-e^{ich}q(x)\right|
 &\ge\sigma_0-2bB_0L^h\rho(b,H)-2H(L^h\rho(b,H))^\alpha\\
 &\ge\sigma_0/2.
\end{align*}
Choose the indicated representative on these balls; changes on null
sets do not affect the following integrals. Telescoping $h$ times
bounds the left side by $\sum_{k<h}|\mathcal D_q(F^k x)|$ almost
everywhere. Minkowski's inequality on the smaller ball and invariance
of $\nu_R^*$ give
\[
 \frac{\sigma_0}{2}\nu_R^*(B(x_P,\rho(b,H)))^{1/p}
 \le\sum_{k<h}\|\mathcal D_q\circ F^k\|_{L^p(B(x_P,\rho(b,H)),\nu_R^*)}
 \le h\|\mathcal D_q\|_{L^p(\nu_R^*)}.
\]
Use the exact ball mass and $h\le M(b)$ to conclude. The estimate
uses measure preservation, not disjointness or probabilistic
independence of the iterates of the ball.
\end{proof}

\begin{corollary}[A polynomial budget for the phase regularity]
\label{cor:polynomial-regularity-budget}
Fix $\alpha,p,H_0$ as above and $\zeta\ge0$. If the H\"older constants
satisfy $H\le H_0b^\zeta$, set
\begin{equation}\label{eq:coercivity-exponent}
 \kappa=\frac{\log L}{\log(100/49)}+
                       \max\{1,\zeta/\alpha\}.
\end{equation}
There is $C>0$, depending only on these fixed constants and the
mechanical family, such that
\begin{equation}\label{eq:polynomial-coercivity}
 \|\mathcal D_q\|_{L^p(\nu_R^*)}
             \ge C\frac{b^{-2\kappa/p}}{1+\log b}\qquad(b\ge1).
\end{equation}
\end{corollary}
\begin{proof}
Writing $a=\log(100/49)$, we have $M(b)\le3+\log b/a$ and hence
$L^{-M(b)}\ge L^{-3}b^{-\log L/a}$. The minimum in
\eqref{eq:coercivity-radius} is at least
\[
 \min\left\{\frac{\rho_0}{2},\frac{\sigma_0}{8B_0},
           \left(\frac{\sigma_0}{8H_0}\right)^{1/\alpha}\right\}
                     b^{-\max\{1,\zeta/\alpha\}}.
\]
Substitution into \eqref{eq:positive-measure-coercivity} proves the
claim with a positive constant independent of $R$ and $b$.
\end{proof}

\subsection{The operator interface}
Theorems~\ref{thm:uniform-phase-separation} and
\ref{thm:positive-measure-coercivity} have different hypotheses.
The former is an unconditional estimate for the specified mechanical
periods. The latter is an estimate for a stated regularity class of
circle-valued phases, in the actual invariant measure. In particular
it supplies a positive-measure obstruction rather than a conclusion
obtained by evaluating an unspecified measurable representative at
isolated periodic points.

To apply \eqref{eq:polynomial-coercivity} to a transfer operator,
one must construct from its approximate spectral vectors phases in
that class, bound the H\"older constant, and compare their physical
$L^p$ defect with the operator defect. Such a comparison would
contradict an error smaller than the right side of
\eqref{eq:polynomial-coercivity}. It is not proved merely by writing
an eigenvalue equation. In particular a division by the modulus of a
spectral vector requires a separate lower-modulus estimate.
General Liv\v{s}ic regularity results for Markov systems
\cite{BHN} concern precise dynamical and cohomological hypotheses;
they do not supply a uniform quantitative regularity bound for these
approximate phases without checking those hypotheses and the defect
comparison. No such comparison is assumed here for the full induced
operator. The common operator realization and the all-branch raw
residual estimates in the final section remain the requirements for
the original joint density theorem.

\section{Exact return means and stationary unfinished flights}
\label{sec:exact-return-means}
The enlarged section gives explicit normalization constants without any
local-limit assumption. We use ergodicity of the finite-horizon
periodic dispersing collision map, as supplied by the standard billiard
framework \cite{Young}. No spectral conclusion for the new first-return
map is imported from that fact.

\begin{proposition}[Exact induced means]\label{prop:exact-induced-means}
For the actual return records to $Y_R^*$,
\begin{equation}\label{eq:induced-mean-vector}
 \int(\kappa_{R,1}^*,\kappa_{R,2}^*,r_R^*,\tau_R^*)\dd\nu_R^*
 =\left(0,0,\frac1{c_*},\frac{\bar\tau_R}{c_*}\right),
 \quad
 \bar\tau_R=\frac{\sqrt3}{4R}-\frac{\pi R}{2}.
\end{equation}
The mean roof satisfies
\begin{equation}\label{eq:mean-roof-modulus}
 \left|\frac{\mathrm d}{\mathrm dR}\int\tau_R^*\dd\nu_R^*\right|
 <\frac{15}{4c_*}\qquad(R\in I).
\end{equation}
The equilibrium physical collision rate remains $1/\bar\tau_R$.
\end{proposition}
\begin{proof}
For any integrable physical collision observable $f$, the disjoint
return-tower levels partition the collision section up to null sets.
Invariance, first for nonnegative functions and then for positive and
negative parts, gives the Kac identity
\[
 \int_{Y_R^*}\sum_{j=0}^{r_R^*(x)-1} f(T_R^j x)\dd\nu(x)
                         =\int f\dd\nu.
\]
Taking $f=1$ and dividing by $c_*$ gives the mean return count.
Taking $f=\tau_R$ gives the mean induced roof. The physical
suspension normalization proved earlier gives the displayed value
of $\bar\tau_R$. The physical displacement is integrable by the
finite-horizon bound. Spatial inversion of the entire lattice maps
its value to its negative and preserves $\nu$, so its mean is zero.
This proves \eqref{eq:induced-mean-vector}.

The section volume $c_*$ is independent of $R$. Differentiating the
explicit expression for $\bar\tau_R$ yields
$-\sqrt3/(4R^2)-\pi/2$. The bounds $\sqrt3<7/4$,
$\pi<22/7$ and $R\ge9/20$ show its absolute value is less than
$15/4$. Finally the ratio of mean physical collisions per return to
mean physical time per return is $1/\bar\tau_R$.
\end{proof}

\begin{proposition}[The complete induced suspension]\label{prop:induced-suspension}
The equilibrium flow is represented by the suspension over $F_R^*$
with roof $\tau_R^*$ and probability
\begin{equation}\label{eq:induced-suspension}
 \frac{\nu_R^*(\mathrm dx)\,\mathrm ds}{\int\tau_R^*\dd\nu_R^*},
               \qquad 0\le s<\tau_R^*(x).
\end{equation}
If $A_R$ and $B_R$ denote respectively the age since the last visit
and the time to the next visit to $Y_R^*$ under this stationary law,
then, for $q\ge0$,
\begin{equation}\label{eq:residual-tails}
 \Pr(A_R>q)=\Pr(B_R>q)
 =\frac{\int(\tau_R^*-q)_+\dd\nu_R^*}
              {\int\tau_R^*\dd\nu_R^*}.
\end{equation}
At either one fixed endpoint of a time interval of length $t$,
the incomplete induced flight, its physical collision count and its
lattice displacement divided by $\sqrt t$ converge to zero in
probability as $t\to\infty$, for each fixed $R$.
\end{proposition}
\begin{proof}
In the physical suspension over $T_R$, group the consecutive
single-flight intervals before each first return to $Y_R^*$.
The tower identity and the normalization
$c_*\int\tau_R^*\dd\nu_R^*=\bar\tau_R$ convert the original
suspension measure exactly to \eqref{eq:induced-suspension}.
This rearranges the actual dependent path; it introduces no new
launch law and no independent flight variables. The marked formula
of Proposition~\ref{prop:marked} can consequently be grouped into
these complete return records while keeping both endpoint pieces.

For a given roof $\tau$, the set of $s\in[0,\tau)$ with age $s>q$
has length $(\tau-q)_+$. The set with remaining time $\tau-s>q$
has the same length. Integration proves \eqref{eq:residual-tails}.
Its right side tends to zero by the integrability in
Proposition~\ref{prop:exact-induced-means}. Each endpoint roof is
finite almost surely under the length-biased stationary law;
stationarity gives the same marginal at the other fixed endpoint.
Thus its ratio to $\sqrt t$ tends to zero in probability.

Every physical flight has length at least $1-2R\ge3/50$. The
number of collisions in an incomplete interval of length $s$ is
at most $1+s/(1-2R)$. Its displacement between fundamental cells
is bounded by a constant plus a fixed multiple of $s$, since the
speed is one and the lattice cell has bounded diameter. These
bounds prove the same endpoint statement for the two marked rewards.
\end{proof}

A fixed-endpoint assertion is not a bound on the maximum residual over
all times in $[0,t]$. Nor does pointwise integrability in $R$ imply a
uniform residual estimate over $I$. Such uniform or process-level
claims require uniform integrability or a quantitative tail for the
induced roof. The functional limit assumptions in
Proposition~\ref{prop:clock} therefore remain explicit. The exact
mean and the modulus \eqref{eq:mean-roof-modulus}, by contrast, are
already available for geometric error propagation. They do not
require recovery of the entire preparation distribution in A2-GEOM.

\section{Stability of the actual return records}
\label{sec:return-stability}
We compare the first returns themselves, including their changing domains.
The comparison is made on the common collision space
$M=\T\times(-1,1)$ with the probability $\nu$ in \eqref{eq:flux}.
It is not a comparison between independently chosen inducing schemes.
The squares defining $Y_R$ move analytically with $R$, and $U$ is fixed.
Consequently $Y_R^*$ has constant measure $c_*$, and its boundary is a
finite union of smooth arcs, continuously varying with $R$.

We use only ergodicity, not a parameter-uniform mixing rate, from the
finite-horizon Sinai theory. In the present family the obstacles are
smooth, strictly dispersing and disjoint, and the horizon is finite.
The flow-mixing theorem \cite[Corollary 1.3 and correction note]{BDL}
therefore implies ergodicity of the collision map: the suspension of a
nontrivial invariant collision set would be a nontrivial invariant flow
set. Thus every $Y_R^*$ has full saturation. This explicitly supplies the
full-tower hypothesis in the Kac identity used above.

Write $r=r_R^*$, $F=F_R^*$ and $\tau^*=\tau_R^*$ in sums indexed by
induced returns. For $n\ge1$ define the actual cumulative record
\[
 J_{n,R}=\left(\sum_{j<n}\kappa_R^*\circ F^j,
                   \sum_{j<n}r_R^*\circ F^j,
                   \sum_{j<n}\tau_R^*\circ F^j\right)
       =(K_{n,R},N_{n,R},T_{n,R})
\]
on $Y_R^*$. Its values lie in $\Z^3\times\R$; in particular the last
coordinate is one cumulative time, not $n$ independent roof coordinates.
A tilde denotes extension by zero to $M$. Let $p_{n,R}$ denote the raw
mixed density of $J_{n,R}$ under $\nu_R^*$, whose existence is included
in Theorem~\ref{thm:return-stability} below. Norms of densities use
counting measure times Lebesgue measure, denoted by $\mathfrak m$.

\begin{lemma}[Almost-everywhere stability of a finite return record]
\label{lem:finite-return-stability}
Fix $R_0\in I$ and a finite $n$. Outside a $\nu$-null subset of $M$,
$\widetilde J_{n,R}(x)\to\widetilde J_{n,R_0}(x)$ as $R\to R_0$.
For $x\in Y_{R_0}^*$ outside that null set, the physical word, every
intermediate return decision, $K_{n,R}(x)$ and $N_{n,R}(x)$ are constant
for all sufficiently close $R$. On a neighborhood of such a point the
cumulative time is the physical analytic length of that fixed word.
The neighborhood and parameter tolerance may depend on $x$ and $n$.
\end{lemma}
\begin{proof}
Remove the boundary of $Y_{R_0}^*$, all singular initial states and all
preimages under $T_{R_0}$ of that boundary. These sets are null: the
boundary has zero area, the map preserves $\nu$, and the billiard
singularities and their iterates have zero flux measure. Remove also
the null set on which a required return is infinite. For a remaining
$x\in Y_{R_0}^*$, the first $n$ returns contain finitely many regular
physical collisions. Every tested state lies strictly inside or strictly
outside the return set. The finite-itinerary analyticity following
\eqref{eq:entry}, and the continuous motion of the return boundary,
preserve all these decisions for nearby $(R,x)$. They also preserve
the selected entry roots and lattice labels. The final physical count
and displacement are therefore unchanged, while the sum of flight
lengths varies analytically. A remaining point outside $Y_{R_0}^*$
stays outside $Y_R^*$ for nearby $R$, and its zero extension stays zero.
These statements prove the asserted convergence. No common neighborhood
for infinitely many words is used.
\end{proof}

\begin{theorem}[Finite-record stability in density and first moment]
\label{thm:return-stability}
For every fixed $n\ge1$, the maps
\begin{equation}\label{eq:return-Lone}
 R\longmapsto\widetilde J_{n,R}\in L^1(\nu;\R^4),\qquad
 R\longmapsto p_{n,R}\in L^1(\mathfrak m)
\end{equation}
are continuous. Both maps are uniformly continuous on $I$.
Moreover $\{\widetilde N_{n,R}:R\in I\}$ and
$\{\widetilde T_{n,R}:R\in I\}$ are uniformly integrable.
If $\eta_{n,R}$ is the record law, then
\begin{equation}\label{eq:record-Wone}
 W_1(\eta_{n,R},\eta_{n,S})
 \le c_*^{-1}\|\widetilde J_{n,R}-\widetilde J_{n,S}\|_{L^1(\nu)},
\end{equation}
where $W_1$ uses the Euclidean distance on $\Z^3\times\R$.
In particular, for the original, unmodified characteristic functions,
\begin{align}\label{eq:raw-frequency-comparison}
 |\Phi_{n,R}(\omega)-\Phi_{n,S}(\omega)|
 &\le \|p_{n,R}-p_{n,S}\|_{L^1(\mathfrak m)},\nonumber\\
 |\Phi_{n,R}(\omega)-\Phi_{n,S}(\omega)|
 &\le c_*^{-1}|\omega|\,
       \|\widetilde J_{n,R}-\widetilde J_{n,S}\|_{L^1(\nu)}.
\end{align}
No modulus uniform in $n$, and no frequency-integrated estimate, is
asserted by these inequalities.
\end{theorem}
\begin{proof}
The Kac identity and invariance of $F_R^*$ give, exactly,
\begin{equation}\label{eq:block-Kac-mass}
 \int_M\widetilde N_{n,R}\dd\nu=n,\qquad
 \int_M\widetilde T_{n,R}\dd\nu=n\bar\tau_R.
\end{equation}
For nonnegative functions $h_j\to h$ almost everywhere with
$\int h_j\to\int h$, the identity
\[
 \|h_j-h\|_1=\int h_j+\int h-2\int\min(h_j,h)
\]
and dominated convergence for $\min(h_j,h)\le h$ prove $L^1$
convergence. Apply this to the two nonnegative coordinates, using
Lemma~\ref{lem:finite-return-stability}, \eqref{eq:block-Kac-mass}
and continuity of $\bar\tau_R$. Uniform finite horizon bounds every
single-flight lattice label by a fixed constant. Thus
$|\widetilde K_{n,R}|\le C\widetilde N_{n,R}$. The already established
$L^1$ convergence of the counts makes these bounds uniformly integrable
along each convergent parameter sequence. Almost-everywhere convergence
and truncation then give $L^1$ convergence of the displacement as well.
This proves the first continuity in \eqref{eq:return-Lone}.

A continuous image of the compact interval $I$ is compact in $L^1$.
For completeness, an $L^1$-compact nonnegative family is uniformly
integrable. Cover it by finitely many $L^1$ balls of radius $\epsilon$
with centers $h_1,\ldots,h_q$. For a member $h$ in the ball about $h_i$,
\[
 \int_{\{h>A\}}h
 \le 2\|h-h_i\|_1+
              2\int_{\{h_i>A/2\}}h_i.
\]
Choose $A$ for the finite set of centers and then let $\epsilon$
decrease. This proves the stated uniform integrability without a tail
assumption or an exponential return estimate.

We next prove existence and $L^1$ continuity of the density. Partition
$Y_R^*$, up to null sets, by total physical count and ordered lattice
word. The argument of Corollary~\ref{cor:raw-density} applies to this
actual first-return event as well: on each regular word the cumulative
length has at most one critical point by
Proposition~\ref{prop:general-edge}; elsewhere it is a submersion.
Coarea on a countable exhaustion and then summation give absolute
continuity. This argument recomputes the return event for $Y_R^*$;
it does not transfer a density normalization from $Y_R$.

Fix $R_0$ and $\epsilon>0$. After deleting the null sets just described,
choose finitely many smooth nonnegative cutoffs $\psi_i$ on $M$, with
$\sum_i\psi_i\le1$, whose supports lie compactly in regular
first-$n$-return patches for $R_0$, such that
\[
 c_*^{-1}\int\sum_i\psi_i\dd\nu>1-\epsilon.
\]
Refine the patches so that on each one a coordinate derivative of the
cumulative length is nonzero. Compact support and the finite-return
lemma preserve the same word, all return decisions and this derivative
for $R$ close to $R_0$. On each patch the map consisting of the other
initial coordinate and the cumulative length is a local diffeomorphism.
It and its inverse depend continuously in $C^1$ on $R$ on slightly
larger compact patches. Changing variables and integrating the other
coordinate therefore shows that the pushforward of
$c_*^{-1}\psi_i\nu$ has a density continuous in $L^1$ with $R$.
One can see this directly by extending the transformed smooth weight
by zero off its coordinate patch: the compactly supported cutoff
vanishes near the patch boundary, the Jacobian is bounded away from
zero, and dominated convergence applies on a fixed bounded coordinate
box. The lattice coordinates are constant on the patch.

The remainder is a positive measure. Its mass is
$1-c_*^{-1}\int\sum_i\psi_i\dd\nu<\epsilon$, for all these $R$,
because the cutoff supports stay inside $Y_R^*$ and $\nu(Y_R^*)=c_*$.
Consequently
\[
 \|p_{n,R}-p_{n,R_0}\|_1
 \le\sum_i\|p_{i,R}-p_{i,R_0}\|_1+2\epsilon.
\]
Let $R\to R_0$ and then $\epsilon\to0$. This proves the second
continuity in \eqref{eq:return-Lone}. Compactness of $I$ gives uniform
continuity of both maps.

For a 1-Lipschitz function $f$, subtract the constant $f(0)$. Then
\[
 \int f\dd\eta_{n,R}
 =c_*^{-1}\int_M f(\widetilde J_{n,R})\dd\nu.
\]
The same identity holds for $S$. Subtraction and the Lipschitz bound,
followed by the dual characterization of $W_1$, give
\eqref{eq:record-Wone}. All first moments are finite by
\eqref{eq:block-Kac-mass}. Direct integration against the densities
gives the first bound in \eqref{eq:raw-frequency-comparison}. For the
second, use the zero extensions and
$|e^{ia}-e^{ib}|\le|a-b|$; the off-section constants cancel because
both section measures equal $c_*$. These are estimates of the raw
transforms, with no convolution or change of observable.
\end{proof}

\begin{proposition}[Weighted records and a positive denominator]
\label{prop:weighted-return-stability}
Fix $n$. Suppose measurable insertions $w_R:M\to\R$ satisfy
$|w_R|\le M_0$ and, for every $R_0$, $w_R(x)\to w_{R_0}(x)$ for
$\nu$-almost every $x$ as $R\to R_0$. The $w_R$-weighted record
densities $p^w_{n,R}$ are continuous in $L^1(\mathfrak m)$.
Set
\[
 d_n(\delta)=\sup_{|R-S|\le\delta}\|p_{n,R}-p_{n,S}\|_1,
 \qquad
 d_n^w(\delta)=\sup_{|R-S|\le\delta}\|p^w_{n,R}-p^w_{n,S}\|_1.
\]
These moduli tend to zero. For any Borel observation event $B$ in the
record space, if $\eta_{n,R}(B)\ge b>0$ and
$d_n(\delta)<b$, then, whenever $|R-S|\le\delta$,
\begin{equation}\label{eq:finite-conditioning-stability}
 \left|\frac{\int_Bp^w_{n,R}\dd\mathfrak m}{\eta_{n,R}(B)}
       -\frac{\int_Bp^w_{n,S}\dd\mathfrak m}{\eta_{n,S}(B)}\right|
 \le\frac{d_n^w(\delta)+M_0d_n(\delta)}{b-d_n(\delta)}.
\end{equation}
\end{proposition}
\begin{proof}
Use the compact regular patches in the preceding proof. Replacing
$w_R$ by $w_{R_0}$ costs at most
$c_*^{-1}\int|w_R-w_{R_0}|\dd\nu\to0$, by dominated convergence
and contraction of total variation under pushforward. On each fixed
patch approximate the bounded measurable weight $w_{R_0}$ in $L^1$
by a smooth one. The approximation error in its pushforward is at
most the same $L^1$ error, uniformly in $R$. For the smooth weight,
the coordinate-change proof above applies. The leftover measure has
total variation at most $M_0\epsilon$. This proves weighted density
continuity, and compactness gives the modulus.

Write $a_R=\int_Bp^w_{n,R}\dd\mathfrak m$ and
$b_R=\eta_{n,R}(B)$. Then
$|a_R-a_S|\le d_n^w(\delta)$,
$|b_R-b_S|\le d_n(\delta)$ and $|a_R|\le M_0b_R$.
Since $b_S\ge b-d_n(\delta)>0$, subtraction of the two ratios gives
\eqref{eq:finite-conditioning-stability}.
\end{proof}

The event $B$ may fix the lattice coordinates and restrict the roof
to an interval. The denominator in
\eqref{eq:finite-conditioning-stability} is an actual event probability.
The estimate does not remove its deterioration for small events. In
particular density $L^1$ continuity alone is not the shrinking-interval
LLT of Proposition~\ref{prop:conditioning}; that proposition retains
its stronger pointwise density hypotheses. An error bound
$|\widehat R-R|\le\delta$ can be inserted into the displayed moduli,
but no numerical sampling budget follows from a noneffective modulus.

\begin{corollary}[Uniform stationary endpoint control]
\label{cor:uniform-stationary-endpoints}
Put $H_R=\widetilde T_{1,R}$ and
$b_*=\min_{R\in I}\bar\tau_R>0$. The stationary induced suspensions,
regarded as probabilities on $M\times\R_+$, have densities
\[
 h_R(x,s)=\bar\tau_R^{-1}\one_{\{0\le s<H_R(x)\}}
\]
with respect to $\nu(\dd x)\dd s$, and
\begin{equation}\label{eq:suspension-Lone}
 \|h_R-h_S\|_1\le
 b_*^{-1}\bigl(\|H_R-H_S\|_1+|\bar\tau_R-\bar\tau_S|\bigr).
\end{equation}
Furthermore
\begin{equation}\label{eq:uniform-roof-tail}
 \Psi(q):=b_*^{-1}\sup_{R\in I}
                   \int_M H_R\one_{\{H_R>q\}}\dd\nu
             \longrightarrow0\quad(q\to\infty).
\end{equation}
Under the stationary flow, at either deterministic endpoint of an
interval of length $t$, the incomplete induced time, physical collision
count and lattice displacement are $o_P(\sqrt t)$, uniformly over
$R\in I$. The same is true jointly at any fixed finite number of
deterministic endpoints.
\end{corollary}
\begin{proof}
The normalization in \eqref{eq:induced-suspension} and
$c_*\int\tau_R^*\dd\nu_R^*=\bar\tau_R$ give $h_R$ exactly.
For fixed $x$, the symmetric difference of the intervals
$[0,H_R(x))$ and $[0,H_S(x))$ has length $|H_R(x)-H_S(x)|$.
Integrating, and then comparing the two normalizing constants, proves
\eqref{eq:suspension-Lone}. The uniform integrability in
Theorem~\ref{thm:return-stability} proves
\eqref{eq:uniform-roof-tail}.

The full roof containing a stationary endpoint has tail
\[
 \Pr_R\{\text{containing roof}>q\}
 =\bar\tau_R^{-1}\int_MH_R\one_{\{H_R>q\}}\dd\nu
 \le\Psi(q).
\]
This is the length-biased law; replacing it by
$\nu_R^*(\tau_R^*>q)$ would be incorrect. Both endpoint pieces
are at most this full roof. The collision count of an incomplete
piece of length $s$ is at most $1+(50/3)s$, and its lattice
displacement is bounded by $C(1+s)$, by the physical estimates in
Proposition~\ref{prop:induced-suspension}. Hence there are fixed
constants $C_0,C_1$ such that each marked endpoint norm is bounded
by $C_0+C_1H$ when its containing roof has length $H$. Its probability of exceeding
$\epsilon\sqrt t$ is at most
$\Psi((\epsilon\sqrt t-C_0)/C_1)$ for large $t$, uniformly in $R$.
Stationarity applies at every deterministic endpoint, including an
endpoint depending deterministically on $t$. A union bound proves
the finite-endpoint assertion; independence is not used.
\end{proof}

This removes the additional uniform-integrability premise for fixed
stationary endpoints in this family. It is not a bound on a supremum
over a growing time interval or at an arbitrary stopping time. The
functional clock theorem still needs its printed process hypotheses.
Likewise \eqref{eq:raw-frequency-comparison} is a uniform comparison
of raw characteristic functions, not an integrable tail for any one
of them. The high-frequency and central all-branch remainder problem
is unchanged by an $L^1$ comparison of probability densities.

\section{Exponential tails for the genuine return record}
\label{sec:exponential-returns}
We now use the spectral theory of the collision map, in addition to the
geometric and ergodic inputs used above. The operator in this section is
not the operator of the induced map. A smooth killing weight on the
collision space gives a scalar estimate for avoiding the actual return
section. In particular no multiplication by its discontinuous indicator
on an anisotropic space will be needed.

\subsection{The collision-space input}
We use the strong and weak distribution spaces of Demers--Zhang
\cite[Section 3.3]{DZ}. Their properties needed here are the local uniform
spectral gap and power bounds for unforced dispersing collision maps,
including small changes of the scatterers and their boundary lengths
\cite[Theorems 2.1--2.6 and Remark 2.7(b)]{DZ}. The spectral gap for an
individual unforced Lorentz map is the input from \cite{DZ11} used there.
Here is the precise consequence, together with the checks for this family.

\begin{lemma}[Local uniform collision spaces]
\label{lem:collision-spectral-input}
There is a finite cover of $I$ by parameter neighborhoods $I_j$. On each
$I_j$ there is a Banach space $\mathcal B_j$ of distributions on a common
collision space, a continuous mass functional $\ell_j$, and transfer
operators $\mathcal L_R$, such that
\begin{equation}\label{eq:collision-split}
 \mathcal L_R=\Pi_R+\mathcal N_R,\quad
 \Pi_R h=\ell_j(h)\nu,\quad
 \Pi_R\mathcal N_R=\mathcal N_R\Pi_R=0,\quad
 \|\mathcal N_R^k\|_{\mathcal B_j}\le C\vartheta^k
 \quad(k\ge0),
\end{equation}
where $0<\vartheta<1$ and $C<\infty$ can be chosen for all the finitely
many neighborhoods. The norms of $\nu$, $\ell_j$ and $\mathcal L_R$ are
uniformly bounded. Multiplication by a smooth function $g$ satisfies
\begin{equation}\label{eq:smooth-multiplier}
 \|g h\|_{\mathcal B_j}\le C\|g\|_{C^2}\|h\|_{\mathcal B_j}.
\end{equation}
The transfer convention is $(\mathcal L_R h)(f)=h(f\circ T_R)$.
No assertion about the unbounded induced twists is included in this lemma.
\end{lemma}
\begin{proof}
For the collision-space construction use coordinates $(\alpha,\varphi)$,
where $p=\sin\varphi$. The invariant probability is the same for all
parameters: $\dd\nu=(4\pi)^{-1}\cos\varphi\dd\alpha\dd\varphi$.
We do not conjugate a distribution space by the inverse of
$p=\sin\varphi$ at grazing. At a fixed $R_0$, use the boundary coordinate
$r_0=R_0\alpha$. The parametrization of the nearby circle is then
$r_0\mapsto R n_{r_0/R_0}$. It is a $C^3$-bounded deformation, $C^2$-close
to the original parametrization, and its constant speed $R/R_0$ tends to
one. This is precisely the boundary-length reparametrization permitted in
\cite[Remark 2.7(b)]{DZ}.

The physical gap between different disks is at least $3/50$. The
curvatures belong to $[100/47,20/9]$, and the parametrized boundaries have
uniform $C^3$ bounds. The finite horizon is uniform on $I$: a line with
no disk intersection at the smallest radius would be an infinite free
line in that table, which has no corridor; compactness of position and
direction then bounds the free length, and increasing the radius cannot
increase it. If a rectangular fundamental domain is desired, use the
two-scatterer cover with periods $(1,0)$ and $(0,\sqrt3)$ and disk centers
$0,b$. Its deck translation commutes with the collision map. Restricting
to the deck-invariant distributions gives the original table. Thus the
local billiard estimates are applied without an affine distortion of
specular reflection.

The uniform geometric hypotheses of \cite[Theorem 2.5]{DZ} and the
perturbation hypothesis of its Theorem 2.6 are satisfied. The unforced
collision map at each $R_0$ has the simple eigenvalue one and a spectral
gap. The strong/weak perturbation theorem, specifically the complementary
power bound in \cite[Theorem 2.1(3)]{DZ}, gives \eqref{eq:collision-split}
on a neighborhood of $R_0$. The eigenprojection has the displayed form
because the invariant probability is $\nu$ and mass is preserved. A
finite cover of the compact parameter interval makes the constants
uniform. This uses local common spaces and does not assert that all
arbitrary inducing schemes share a Banach space.

For completeness, smooth multiplication follows directly from the
stable-test and unstable-comparison norms defining these spaces. On a
single admissible stable curve, multiplying its test function by $g$
changes its H\"older norm by at most $C\|g\|_{C^1}$. For two matched
curves at distance $\epsilon$, split the difference of the weighted tests
into the old test difference and the difference of the two restrictions
of $g$. After graph matching, the latter has $C^q$ norm at most
$C\|g\|_{C^2}\epsilon^{1-q}$ by interpolation between its $C^0$ and
$C^1$ bounds. The defining exponents satisfy $\beta<1-q$, so division
by $\epsilon^\beta$ is bounded for $\epsilon\le1$. The stable part of the
norm controls this second term and the unstable part controls the first.
The weak norm has the same single-curve multiplication bound. Density of
the defining smooth distributions extends these bounds to
$\mathcal B_j$, proving \eqref{eq:smooth-multiplier}. The coordinate
changes above and the finite cover have uniformly bounded derivatives.
\end{proof}
\subsection{A smooth avoidance estimate}
Choose once and for all $\chi\in C^\infty(M)$ such that
\begin{equation}\label{eq:fixed-killing-weight}
 0\le\chi\le1,\quad
 \supp\chi\Subset\{ |\alpha|<1/400,\ |p|<1/400\},\quad
 a_\chi:=\int\chi\dd\nu>0.
\end{equation}
The indicated rectangle is inside the square of $Y_R$ centered at
$(0,0)$ for every $R$; hence it is inside $Y_R^*$. Its support is away
from grazing, so it is smooth also in the collision coordinates used in
Lemma~\ref{lem:collision-spectral-input}. This is a proof weight, not an
extra field in the collision record.

\begin{proposition}[Uniform soft-killing estimate]
\label{prop:soft-killing}
There exist $z_0,c,C>0$ such that for every $R\in I$ and $k\ge0$,
\begin{equation}\label{eq:soft-killing}
 \int_M\exp\left(-z_0\sum_{j=0}^{k-1}\chi\circ T_R^j\right)\dd\nu
       \le Ce^{-ck}.
\end{equation}
\end{proposition}
\begin{proof}
On a local collision space set
$\mathcal L_{R,z}=\mathcal L_R M_{e^{-z\chi}}$ for complex $z$.
The multiplier bound proves analyticity in the strong operator norm and,
uniformly in $R$, $\|\mathcal L_{R,z}-\mathcal L_R\|\le C|z|$ near
zero. Equation~\eqref{eq:collision-split} bounds the unperturbed
resolvents uniformly on a circle around one and on a second circle
$|w|=\vartheta_1$ with $\vartheta<\vartheta_1<1$, chosen disjointly.
A uniformly convergent resolvent Neumann series therefore defines a
rank-one analytic projection and a single eigenvalue $\lambda_R(z)$
near one for all $|z|<z_*$. The complementary powers are bounded by
$C\vartheta_1^k$, after decreasing $z_*$ if necessary.

Since $\ell_j(\nu)=1$, differentiation at zero gives
\[
 \lambda_R(0)=1,\qquad
 \lambda_R'(0)=-\ell_j(\mathcal L_R(\chi\nu))=-a_\chi.
\]
The fixed complex disk and the uniform resolvent bounds give a uniform
bound on the second derivative of $\lambda_R$ on a smaller disk, by
Cauchy's formula. Thus
$|\lambda_R(z)-1+a_\chi z|\le C_2|z|^2$ uniformly in $R$.
For real small $z$ the eigenvalue is real, by conjugation and uniqueness,
and positive, by its closeness to one. Choose a fixed $z_0>0$ such that
$C_2z_0\le a_\chi/2$ and $z_0<z_*$. Then
$0<\lambda_R(z_0)\le1-a_\chi z_0/2$. The spectral decomposition gives
$\|\mathcal L_{R,z_0}^k\|\le Ce^{-ck}$. Only finitely many collision
spaces occur, so the same $z_0,c,C$ work on all of $I$.

Iteration of the transfer identity, first for densities and then as
bounded functionals, yields
\[
 \ell_j(\mathcal L_{R,z_0}^k\nu)
 =\int\prod_{j=0}^{k-1}e^{-z_0\chi(T_R^j x)}\dd\nu(x).
\]
The norms of $\ell_j$ and $\nu$ are bounded. This proves the asserted
scalar estimate. In particular the argument does not apply a local
large-deviation formula outside its stated neighborhood of the mean.
\end{proof}

\begin{theorem}[Exponential moments of the actual return block]
\label{thm:exponential-return}
Let
$Q_R=r_R^*+\tau_R^*+|\kappa_R^*|$ on the genuine first-return section.
There are constants $c_1,a_1,C_1>0$, independent of $R\in I$, such that
\begin{equation}\label{eq:exponential-return}
 \nu_R^*(r_R^*>k)\le C_1e^{-c_1k}\quad(k\ge0),\qquad
 \int e^{a_1 Q_R}\dd\nu_R^*\le C_1.
\end{equation}
For $Q_{n,R}=\sum_{j<n}Q_R\circ(F_R^*)^j$ and every $n\ge1$,
\begin{equation}\label{eq:block-exponential-moment}
 \int\exp(a_1 Q_{n,R}/n)\dd\nu_R^*\le C_1,\qquad
 \nu_R^*(N_{n,R}>L)\le C_1e^{-a_1L/n}.
\end{equation}
These are estimates for the original deterministic mechanical records.
\end{theorem}
\begin{proof}
If $x\in Y_R^*$ and $r_R^*(x)>k$, none of $T_Rx,\ldots,T_R^kx$
belongs to $Y_R^*$. Every term of $\sum_{j=1}^k\chi\circ T_R^j$ is
therefore zero. Positivity and invariance of $\nu$ imply
\[
 \nu_R^*(r_R^*>k)
 \le c_*^{-1}\int_M e^{-z_0\sum_{j=1}^k\chi\circ T_R^j}\dd\nu
 \le c_*^{-1}Ce^{-ck}.
\]
This comparison uses neither the regularity of $\one_{Y_R^*}$ as a
multiplier nor an independent-return construction.

Uniform finite horizon and the uniform bound on a single-flight lattice
step give $Q_R\le C_0 r_R^*$ for a fixed $C_0$. Summation of the
exponential tail of the positive integer $r_R^*$ gives a uniform
exponential moment of $Q_R$ for $a_1>0$ sufficiently small. Enlarge
$C_1$ to include the first bound. Finally H\"older's inequality with
$n$ equal exponents and invariance of $F_R^*$ give
\[
 \int\prod_{j<n}\exp\bigl(a_1Q_R\circ(F_R^*)^j/n\bigr)\dd\nu_R^*
 \le\prod_{j<n}\left(\int e^{a_1Q_R\circ(F_R^*)^j}\dd\nu_R^*\right)^{1/n}
 \le C_1.
\]
Since $N_{n,R}\le Q_{n,R}$, exponential Markov inequality proves the
last assertion. No independence is used in this step either.
\end{proof}

\begin{corollary}[All finite moments under parameter variation]
\label{cor:all-moment-continuity}
For fixed $n$ and $1\le p<\infty$, the map
$R\mapsto\widetilde J_{n,R}$ is continuous into $L^p(\nu;\R^4)$.
Every polynomial moment of a fixed finite list of return records is
finite and continuous in $R$. In particular every finite-lag covariance
of the induced record is continuous.
\end{corollary}
\begin{proof}
Almost-everywhere convergence is Lemma~\ref{lem:finite-return-stability}.
Equation~\eqref{eq:block-exponential-moment}, the fixed section mass and
$|J_{n,R}|\le Q_{n,R}$ make all fixed powers uniformly integrable.
Truncation followed by dominated convergence proves $L^p$ convergence
along every convergent parameter sequence. A finite list of records can
be expressed as differences of cumulative records. H\"older's inequality
and the same exponential moment bound give uniform integrability of each
polynomial in that list. Truncation proves its continuity. This proves
finite-lag, not yet infinite Green--Kubo, covariance continuity.
\end{proof}

\section{Exponential control of cumulative collision counts}
\label{sec:cumulative-returns}
The avoidance estimate controls more than a single return. It also
controls the time of the $n$th visit to the section, without separating
the successive return blocks. This observation improves the truncation
budget in \eqref{eq:physical-count-truncation} and supplies exponential
moments on a fixed complex neighborhood for the entire physical record.

Retain $z_0,c,C$ from Proposition~\ref{prop:soft-killing}. Write
$N_{n,R}=\sum_{j<n}r_R^*\circ(F_R^*)^j$, and enlarge a fixed constant
$D\ge1$, if necessary, so that
\begin{equation}\label{eq:cumulative-record-domination}
 |J_{n,R}|_1\le D N_{n,R},\qquad Q_{n,R}\le D N_{n,R}.
\end{equation}
Such a constant exists because each physical flight has uniformly
bounded length and lattice displacement. All probabilities and
expectations in this section are with respect to $\nu_R^*$.

\begin{theorem}[The time of the $n$th actual return]
\label{thm:cumulative-return-tail}
There are constants $A\ge1$, $a,c>0$, independent of $n$, $L$ and $R$,
such that
\begin{equation}\label{eq:cumulative-return-tail}
 \nu_R^*\{N_{n,R}>L\}\le A\exp(an-cL)
       \qquad(n\ge1,\ L\ge0).
\end{equation}
One may take $a=z_0$ and the exponent $c$ of the soft-killing estimate.
Put $v=a/c$. For $0<s<c$ there is a constant $B_s<\infty$, uniform in
$n,R$, for which
\begin{equation}\label{eq:cumulative-exponential-moment}
 \int e^{s(N_{n,R}-vn)_+}\dd\nu_R^*\le B_s,
 \qquad \int e^{sN_{n,R}}\dd\nu_R^*\le B_s e^{svn}.
\end{equation}
Consequently, for every $0<\lambda<c/D$,
\begin{equation}\label{eq:fixed-strip-record-moment}
 \int e^{\lambda Q_{n,R}}\dd\nu_R^*
          \le B_{\lambda D}e^{\lambda Dvn}.
\end{equation}
The constants $B_s$ are locally bounded for $0\le s<c$, with $B_0=1$.
The number $v$ is an upper envelope, not the Kac mean $c_*^{-1}$.
\end{theorem}
\begin{proof}
For an integer $L\ge0$ let
\[
 H_{L,R}(x)=\sum_{j=1}^{L}\one_{Y_R^*}(T_R^j x).
\]
Since $N_{n,R}$ is the physical collision time of the $n$th return,
$N_{n,R}>L$ if and only if $H_{L,R}<n$ on $Y_R^*$. The fixed smooth
weight satisfies $0\le\chi\le\one_{Y_R^*}$. Thus, on this event,
$\sum_{j=1}^{L}\chi(T_R^j x)\le n-1$. In particular, pointwise on
$Y_R^*$,
\[
 \one_{\{N_{n,R}>L\}}
 \le e^{z_0(n-1)}
       \exp\left(-z_0\sum_{j=1}^{L}\chi(T_R^j x)\right).
\]
Integrate on $Y_R^*$, enlarge the domain to $M$, and use invariance to
shift the sum by one collision. Equation~\eqref{eq:soft-killing} gives
\[
 \nu_R^*\{N_{n,R}>L\}\le c_*^{-1}C e^{z_0(n-1)-cL}.
\]
For real $L$ replace it by $\lfloor L\rfloor$ and enlarge the constant
by $e^c$. This proves \eqref{eq:cumulative-return-tail}. It uses a single
orbit's visit count, not a product of probabilities of return blocks.

It follows that, for $u\ge0$,
$\nu_R^*\{N_{n,R}>vn+u\}\le A e^{-cu}$. For a nonnegative random
variable $X$, Tonelli's theorem gives
\[
 \int e^{sX}\dd\nu_R^*
       =1+s\int_0^\infty e^{su}\nu_R^*\{X>u\}\dd u.
\]
Apply this with $X=(N_{n,R}-vn)_+$ to obtain
$B_s=1+As/(c-s)$. Since $N_{n,R}\le vn+X$, the second bound in
\eqref{eq:cumulative-exponential-moment} follows. Finally apply
\eqref{eq:cumulative-record-domination} with $s=\lambda D$.
\end{proof}

We now apply the tail to the actual record measure. For a measurable
insertion $w_{n,R}$ define the finite signed or complex measure
\[
 \mu_{n,R}^{w}=(J_{n,R})_*(w_{n,R}\nu_R^*),\qquad
 \mu_{n,R}^{w,L}=(J_{n,R})_*
       (w_{n,R}\one_{\{N_{n,R}\le L\}}\nu_R^*).
\]
The transform uses the convention $\widehat\mu(\omega)=
\int e^{i\omega\cdot j}\dd\mu(j)$, with
$\omega=(u_1,u_2,s,b)\in\T^3\times\R$.

\begin{proposition}[Weighted finite-band truncation]
\label{prop:linear-count-budget}
Fix a nonnegative integer $d$ and suppose
$|w_{n,R}|\le M(1+Q_{n,R})^d$. For every multi-index $\gamma$ there is a
constant $C_{d,\gamma}$, independent of $n,R,L,w$, such that
\begin{align}\label{eq:weighted-count-truncation}
 \|\mu_{n,R}^{w}-\mu_{n,R}^{w,L}\|_{\TV}
  &\le C_{d,0}M(1+L)^d e^{an-cL},\\
 \sup_{\omega\in\T^3\times\R}
 \left|\partial_\omega^\gamma
   (\widehat\mu_{n,R}^{w}-\widehat\mu_{n,R}^{w,L})(\omega)\right|
  &\le C_{d,\gamma}M(1+L)^{d+|\gamma|}e^{an-cL}.\nonumber
\end{align}
For a bounded insertion ($d=0$), a frequency region of finite volume
$V>0$, and $\varepsilon>0$, the sufficient integer cutoff for integrated
transform error at most $\varepsilon$ is
\begin{equation}\label{eq:linear-count-budget}
 L\ge\left\lceil\frac{a}{c}n+
       \frac1c\max\left\{0,\log\frac{C_{0,0}MV}{\varepsilon}\right\}
        \right\rceil.
\end{equation}
If $V_n\le V_0e^{\kappa n}$ and the required error is $e^{-\delta n}$,
then a cutoff linear in $n$ suffices. This remains true for each fixed
$d,\gamma$, with any fixed positive slack in the exponential budget.
\end{proposition}
\begin{proof}
For an integer $m\ge0$, integration of the tail gives
\begin{align*}
 &\int (1+N_{n,R})^m\one_{\{N_{n,R}>L\}}\dd\nu_R^*\\
 &\quad=(1+L)^m\nu_R^*\{N_{n,R}>L\}
    +m\int_L^\infty(1+t)^{m-1}\nu_R^*\{N_{n,R}>t\}\dd t\\
 &\quad\le C_m(1+L)^m e^{an-cL}.
\end{align*}
The integral term is zero when $m=0$. For $m\ge1$, the last inequality
follows by writing $t=L+u$ and using
$1+L+u\le(1+L)(1+u)$. Pushforward does not increase total variation.
For the derivative estimate differentiate under the integral: its
absolute integrand is at most
$M(1+Q_{n,R})^d|J_{n,R}|_1^{|\gamma|}
\one_{\{N_{n,R}>L\}}$. Equation~\eqref{eq:cumulative-record-domination}
and the preceding tail integral prove both estimates.

Multiplication by $V$ bounds the integral of the transform error. Solving
$C_{0,0}MV e^{an-cL}\le\varepsilon$ gives
\eqref{eq:linear-count-budget}. More generally choose a slope
$\ell>(a+\kappa+\delta)/c$. Then, for $L_n=\lceil\ell n\rceil$,
\[
 V_n(1+L_n)^{d+|\gamma|}e^{an-cL_n}
       \le C(1+n)^{d+|\gamma|}e^{-(\delta+\epsilon_0)n}
\]
for some $\epsilon_0>0$. A fixed additive increase of the cutoffs absorbs
the remaining constant and polynomial, proving the last assertion.
\end{proof}

\begin{remark}
This improves the sufficient cutoff of order
$n\log(V/\varepsilon)$ obtained from
\eqref{eq:physical-count-truncation} to order
$n+\log(V/\varepsilon)$. On an exponentially growing frequency band the
former budget is quadratic in $n$, whereas the new one is linear.
Neither estimate bounds the number of critical words, the variation of
an inverse coarea Jacobian, or a pointwise raw density. In particular the
all-branch derivative sum in Lemma~\ref{lem:raw-residual-sum} remains a
separate requirement. The truncation removes only actual high-count
records in an error estimate; the limiting datum is still the full
unmodified physical record.
\end{remark}

\section{A common realization and the physical renewal identity}
\label{sec:common-renewal}
The preceding moment estimate gives a common realization of the
unbounded record twists on a scale of Lebesgue spaces. It is useful to
construct this realization before seeking an anisotropic spectral
estimate: the collision and induced operators, their damping variables,
and their chronological composition can then be identified without
regularity assumptions on an indicator multiplier. The renewal algebra
is classical; compare \cite{GouezelRenewal}. The assertion here is its
exact realization for the moving physical first-return record.

\subsection{First-return operators on the common collision space}
Use $M=\T\times(-1,1)$ and $\dd\nu=(4\pi)^{-1}\dd\alpha\dd p$.
Let $P_R$ be multiplication by $\one_{Y_R^*}$ and let $E_R=I-P_R$.
Recall locally that
\[
 c_* = \nu(Y_R^*)=\frac{91}{10000\pi},\qquad
 \dd\nu_R^*=c_*^{-1}\one_{Y_R^*}\dd\nu.
\]
On $L^1(M,\nu)$, define
\[
 (\mathcal A_{R,\omega}h)(y)
   =e^{i\omega\cdot f_R(T_R^{-1}y)}h(T_R^{-1}y),\qquad
 f_R=(\kappa_{{\rm coll},1},\kappa_{{\rm coll},2},1,\tau_R),
\]
for real $\omega=(u_1,u_2,s,b)$. Collision singularities are discarded
as null sets. These are isometries. The operators $P_R,E_R$ are
contractions on every $L^p$, although no such assertion about their
action on anisotropic distribution spaces is needed or made here.

\begin{theorem}[Exact damped renewal realization]
\label{thm:common-renewal}
For $m\ge1$ set
\begin{equation}\label{eq:first-return-operator}
 \mathcal R_{m,R,\omega}
 =P_R\mathcal A_{R,\omega}
       (E_R\mathcal A_{R,\omega})^{m-1}P_R,
 \qquad
 \mathcal R_R(z,\omega)=\sum_{m\ge1}z^m\mathcal R_{m,R,\omega}.
\end{equation}
For $|z|<1$ the series converges in operator norm on $L^1(M,\nu)$ and
\begin{equation}\label{eq:induced-damped-bound}
 \|\mathcal R_R(z,\omega)h\|_1\le |z|\|P_Rh\|_1.
\end{equation}
On the closed subspace $P_RL^1$, with identity $P_R$, one has
\begin{equation}\label{eq:physical-renewal-identity}
 P_R(I-z\mathcal A_{R,\omega})^{-1}P_R
       =(I-\mathcal R_R(z,\omega))^{-1}.
\end{equation}
Moreover,
\begin{equation}\label{eq:schur-return-operator}
 \begin{split}
 \mathcal R_R(z,\omega)={}&zP_R\mathcal A P_R\\
 &+z^2P_R\mathcal A E_R
 (I-zE_R\mathcal A E_R)^{-1}
             E_R\mathcal A P_R,
 \end{split}
\end{equation}
where $\mathcal A=\mathcal A_{R,\omega}$ and the middle inverse acts
on $E_R L^1$.

For $|z|\le1$, the strongly defined boundary series represents the
actual first-return twist: for $y\in Y_R^*$ and $x=(F_R^*)^{-1}y$,
\begin{equation}\label{eq:actual-induced-realization}
 (\mathcal R_R(z,\omega)h)(y)
       =z^{r_R^*(x)}e^{i\omega\cdot G_R(x)}h(x).
\end{equation}
It is zero off $Y_R^*$. In particular, for every $n\ge1$,
\begin{equation}\label{eq:actual-record-pairing}
 c_*^{-1}\int_M\mathcal R_R(z,\omega)^n\one_{Y_R^*}\dd\nu
       =\int_{Y_R^*}z^{N_{n,R}}e^{i\omega\cdot J_{n,R}}\dd\nu_R^*.
\end{equation}
\end{theorem}
\begin{proof}
The collision map is invertible and preserves $\nu$. Poincar\'e
recurrence in both time directions makes $F_R^*$ invertible and
measure-preserving on $Y_R^*$, modulo null sets. In the product
\eqref{eq:first-return-operator}, the initial and final states lie in
$Y_R^*$ and the $m-1$ intervening states lie outside it. Thus its initial
domain is exactly $\{r_R^*=m\}$. Its weight is the exponential of the
sum of the $m$ actual collision records, hence $e^{i\omega\cdot G_R}$.
These initial domains partition $Y_R^*$ up to a null set. Their images
under the corresponding return branches also partition $Y_R^*$, by
invertibility of the induced map. Consequently
\[
 \sum_{m\ge1}|z|^m\|\mathcal R_{m,R,\omega}h\|_1
       =\int_{Y_R^*}|z|^{r_R^*}|h|\dd\nu
       \le |z|\|P_Rh\|_1 \quad (|z|\le1).
\]
Each summand has operator norm at most one, so for $|z|<1$ the
operator-norm series converges geometrically. At the boundary the last
display proves strong convergence for each $h$ and gives
\eqref{eq:actual-induced-realization}. It does not assert operator-norm
convergence of the boundary series.

Let $\mathcal U_0=P_R$ and $\mathcal U_m=P_R\mathcal A^mP_R$ for
$m\ge1$. Partition a trajectory with endpoints in $Y_R^*$ according to
its first return time. Chronological composition, with the first segment
on the right, gives
\[
 \mathcal U_m=\sum_{r=1}^m\mathcal U_{m-r}\mathcal R_{r,R,\omega}.
\]
All generating series converge in operator norm for $|z|<1$. Summing
the identity gives $\mathcal U(z)=P_R+\mathcal U(z)\mathcal R_R(z,\omega)$.
The inverse of $I-\mathcal R_R$ exists by
\eqref{eq:induced-damped-bound}, and
$\mathcal U(z)=P_R(I-z\mathcal A)^{-1}P_R$, proving
\eqref{eq:physical-renewal-identity}. Separating $m=1$ from the series
and summing the intervening outside-section iterates proves
\eqref{eq:schur-return-operator}. All these inverses are used only in
$|z|<1$.

Finally iterate \eqref{eq:actual-induced-realization} and change
variables by $(F_R^*)^n$. Its collision counts and record sums are
$N_{n,R}$ and $J_{n,R}$, respectively. Division by $c_*$ yields
\eqref{eq:actual-record-pairing}. No independence of successive
excursions enters the identity.
\end{proof}

\subsection{A fixed complex neighborhood for the unbounded twists}
To cross the damping boundary at $z=1$, use $z=e^w$. Define the operator
directly by the actual $n$th return, rather than by an unproved
continuation of a resolvent. For $y\in Y_R^*$ let $x=(F_R^*)^{-n}y$ and set
\begin{equation}\label{eq:block-scale-operator}
 (\mathcal T_{n,R}(w,\omega)h)(y)
       =e^{wN_{n,R}(x)+i\omega\cdot J_{n,R}(x)}h(x),
\end{equation}
with zero extension to $M\setminus Y_R^*$. Here $w\in\mathbb C$ and
$\omega\in\mathbb C^4$. Holomorphy is stated on this universal cover.
The formulas are invariant under real shifts by $2\pi\mathbb Z^3$
in the first three frequency coordinates. Local complex charts about
real torus frequencies are restrictions of the same formulas; no
additional global complex-torus construction is being asserted. The
symbols $J_{n,R}$ and $G_R$ denote four real coordinates, with the
displacement contributing two coordinates.

\begin{theorem}[Holomorphic realization on a common Lebesgue scale]
\label{thm:common-scale-holomorphy}
Fix $1\le q<p\le\infty$ and put $t=(q^{-1}-p^{-1})^{-1}$. Let
\[
 \sigma(w,\omega)=|\operatorname{Re}w|+D|\operatorname{Im}\omega|_\infty,
 \qquad \Omega_{p,q}=\{(w,\omega):t\sigma(w,\omega)<c\}.
\]
For every $n\ge1$, \eqref{eq:block-scale-operator} defines a
holomorphic map from $\Omega_{p,q}$ to
$\mathcal L(L^p(M,\nu),L^q(M,\nu))$. Its norm satisfies
\begin{equation}\label{eq:scale-operator-bound}
 \|\mathcal T_{n,R}(w,\omega)\|_{p\to q}
       \le c_*^{1/t}B_{t\sigma}^{1/t}e^{\sigma vn}.
\end{equation}
On every compact subset of $\Omega_{p,q}$, all fixed-order derivatives in $(w,\omega)$ have bounds $C_k e^{C_k n}$ uniform in $R$. At every fixed
$(w,\omega)\in\Omega_{p,q}$ and every $h\in L^p(M,\nu)$,
\begin{equation}\label{eq:scale-parameter-continuity}
 R\longmapsto\mathcal T_{n,R}(w,\omega)h
       \quad\hbox{is continuous in }L^q(M,\nu).
\end{equation}
For $\operatorname{Re}w<0$ and real $\omega$, this operator agrees
with $\mathcal R_R(e^w,\omega)^n$ wherever the latter is considered
as an $L^1$ operator. The resulting inserted record pairings extend+holomorphically across $w=0$ on this fixed neighborhood.
\end{theorem}
\begin{proof}
Change variables $y=(F_R^*)^n x$ in the $q$th power of the norm. The
absolute weight is at most $e^{\sigma N_{n,R}(x)}$, by
\eqref{eq:cumulative-record-domination}. H\"older's inequality with
$1/q=1/p+1/t$ and Theorem~\ref{thm:cumulative-return-tail} give
\[
 \|\mathcal T_{n,R}(w,\omega)h\|_q
 \le\left(\int_{Y_R^*}e^{t\sigma N_{n,R}}\dd\nu\right)^{1/t}
       \|h\|_p
 \le c_*^{1/t}B_{t\sigma}^{1/t}e^{\sigma vn}\|h\|_p.
\]
For $\sigma=0$ use $B_0=1$. This proves the operator bound, including
$p=\infty$.

On a compact subset of the stated domain choose a slightly larger
compact neighborhood for which $t\sigma'<c$. A derivative of order
$k$ in $(w,\omega)$ inserts a monomial in $N_{n,R}$ and the four
coordinates of $J_{n,R}$, bounded by $C_kN_{n,R}^k$. Every such
polynomial is absorbed by $e^{(\sigma'-\sigma)N_{n,R}}$. H\"older's
inequality and \eqref{eq:cumulative-exponential-moment} therefore bound
the derivative integrands and the Taylor remainders in operator norm.
The exponential power series on a sufficiently small polydisc converges
in that norm. This proves joint holomorphy, not just weak holomorphy of
individual pairings. The same estimates on the larger neighborhood,
or Cauchy's inequalities there, give the claimed derivative bounds.

We give the parameter-continuity argument because the underlying return
domains move. Fix $R_0$. Remove $\partial Y_{R_0}^*$, the forward and
backward collision singularities, and all finite collision iterates of
$\partial Y_{R_0}^*$. Each is a countable union of null curves, up to
singular endpoints. Backward recurrence holds almost everywhere. At a
remaining point of $Y_{R_0}^*$ the preceding $n$ returns use finitely many
nonsingular physical collisions; none is on a section boundary. The
implicit collision-root argument of
Theorem~\ref{thm:return-stability}, applied backwards, preserves this
finite itinerary for $R$ sufficiently near $R_0$. The backward starting
point and all flight lengths vary continuously; the lattice labels and
collision counts are fixed. At a remaining point outside $Y_{R_0}^*$
the section indicator stays zero. Thus
\eqref{eq:block-scale-operator} converges pointwise almost everywhere
when $h$ is bounded and continuous.

Choose $q'>q$ sufficiently close to $q$ so that $q'<p$ and
$t'\sigma<c$, where $t'=(1/q'-1/p)^{-1}$. The preceding bound in
$L^{q'}$ is uniform in nearby $R$. It gives uniform integrability of the
$q$th powers, so pointwise convergence implies convergence in $L^q$.
For $p<\infty$, approximate an arbitrary $h$ in $L^p$ by bounded smooth
functions and use the uniform $p\to q$ bound. For $p=\infty$, choose a
finite $p_0>q$ so large that $(1/q-1/p_0)^{-1}\sigma<c$. Approximation
in $L^{p_0}$ and its uniform $p_0\to q$ bound give the same conclusion
for $h\in L^\infty$. This avoids assuming norm density of smooth
functions in $L^\infty$.

In the damped region the two operators have the same pointwise formula,
by Theorem~\ref{thm:common-renewal}. For example, for $h\in L^p$ and
$b\in L^{q^\dagger}$, where $q^\dagger$ is the conjugate exponent of $q$,
\[
 c_*^{-1}\int b\,\mathcal T_{n,R}(w,\omega)h\dd\nu
 =\int_{Y_R^*}h(x)b((F_R^*)^nx)
       e^{wN_{n,R}(x)+i\omega\cdot J_{n,R}(x)}\dd\nu_R^*.
\]
This continuous pairing is holomorphic on $\Omega_{p,q}$. It agrees
with the damped renewal pairing there and hence is the asserted
extension across $w=0$.
\end{proof}

\subsection*{Compatibility of the realizations}
For real $\omega$ and $|z|<1$, the constructions fit the diagram
\[
 \begin{array}{ccc}
 \mathcal A_{R,\omega}&\longmapsto&\mathcal R_R(z,\omega)\\
 \Big\downarrow {\scriptstyle P_R(I-z\mathcal A)^{-1}P_R}
 &&\Big\downarrow {\scriptstyle (I-\mathcal R)^{-1}}\\
 \mathcal U_R(z,\omega)&=&\mathcal U_R(z,\omega).
 \end{array}
\]
The upper arrow is the first-return sum
\eqref{eq:schur-return-operator}; the bottom equality is
\eqref{eq:physical-renewal-identity}. On the common domain
$z=e^w$, $\operatorname{Re}w<0$, $\omega\in\mathbb R^4$, the second
compatibility is
\[
 \mathcal R_R(e^w,\omega)^n=\mathcal T_{n,R}(w,\omega),\qquad
 c_*^{-1}\langle b,\mathcal T_{n,R}h\rangle_\nu
 =\int h\,(b\circ(F_R^*)^n)e^{wN_{n,R}+i\omega\cdot J_{n,R}}
                        \dd\nu_R^*.
\]
The first identity is interpreted on $L^p\subset L^1$, with values in
$L^q$, when the scale operator is used; the second has the insertion
exponents of the theorem. Any anisotropic realization must reproduce
these pairings on its admissible test class. Such an additional
realization is not an arrow already established by this diagram.

\begin{remark}[Relation to the spectral and raw-density estimates]
The operator identities concern the genuine four-coordinate record,
and all their spaces live on the same collision probability space.
They specify a common realization with fixed exponential integrability
and strong parameter continuity. They do not assert operator-norm
continuity in $R$, quasi-compactness on $L^p$, or a holomorphic family of
endomorphisms of a single anisotropic space. At real frequencies the
induced $L^1$ operator is an isometry on its section, so its construction
alone gives no spectral gap. Likewise, maps $L^p\to L^q$ with $q<p$
cannot simply be iterated on one space outside the damped region; the
$n$-block operator above was defined directly.

The next analytical step is to establish an anisotropic realization
compatible with \eqref{eq:actual-record-pairing} and
\eqref{eq:physical-renewal-identity}, and to prove its phase
reconstruction and raw-residual bounds. The collision compensation
in Sections~\ref{sec:collision-covariance}--\ref{sec:functional-gaussian}
will establish the actual Gaussian covariance and functional limits
by a different route, without identifying an induced spectral family. The criteria in
Theorem~\ref{thm:phase-resolvent-criterion},
Proposition~\ref{prop:covariance-closure}, and
Lemma~\ref{lem:raw-residual-sum} continue to state these requirements
separately. In particular the common Lebesgue realization is not used
as a substitute for any of those estimates in the raw local limit.
\end{remark}

\section{Common charts for moving return domains}
\label{sec:common-charts}
We give a local change-of-variables estimate and then apply it to the
actual return domains. This makes explicit the domination used in
Theorem~\ref{thm:return-stability}. Critical neighborhoods are discarded
in the initial collision space, not by assuming a uniform density bound
near their images.

\begin{lemma}[A common submersion chart]
\label{lem:common-submersion-chart}
Let $U=U_1\times U_2\Subset\R^2$ be a bounded open rectangle and let
$f(R,u,v)$ be $C^2$ on a neighborhood of $J\times\overline U$, where
$J$ is a compact parameter interval. Suppose
$\partial_v f\ge\eta>0$ there. For $a\in C_c^1(U)$, let $g_R$ be the
density of the pushforward of $a(u,v)\dd u\dd v$ by $f_R$. Then
\begin{equation}\label{eq:chart-Lipschitz}
 \|g_R-g_S\|_{L^1(\R)}\le C_a|R-S|\qquad(R,S\in J).
\end{equation}
The constant depends only on $U$, a common bounded range of $f$, $\eta$,
the $C^2$ bound for $f$, and the $C^1$ bound and support of $a$.
The same assertion holds if $\partial_v f\le-\eta$.
\end{lemma}
\begin{proof}
For fixed $R,u$ the map $v\mapsto f(R,u,v)$ is one-to-one. Write its
inverse as $v_R(u,t)$. The pushforward under $(u,v)\mapsto(u,f_R(u,v))$
has density
\[
 A_R(u,t)=\frac{a(u,v_R(u,t))}{\partial_v f(R,u,v_R(u,t))}
\]
on its image, extended by zero elsewhere. All images lie in one bounded
$(u,t)$ box. Since $a$ vanishes on a fixed neighborhood of $\partial U$,
this zero extension introduces no boundary jump as $R$ varies. Wherever
$a$ or its derivatives are nonzero, the inverse is defined on a common
parameter neighborhood. If $M$ bounds the relevant derivatives of $f$,
implicit differentiation gives
\[
 |\partial_R v_R|\le M/\eta,\qquad
 |\partial_R A_R|
 \le \|a_v\|_\infty M/\eta^2
       +\|a\|_\infty(M/\eta^2+M^2/\eta^3).
\]
At points outside the inverse chart both sides are interpreted by the
zero extension, which vanishes on an open neighborhood of its moving
edge. The same bound therefore gives an $L^1$ majorant on the fixed box.
Integrating in $R$ proves
$\|A_R-A_S\|_1\le C_a|R-S|$. Since
$g_R(t)=\int A_R(u,t)\dd u$, integration in $u$ proves
\eqref{eq:chart-Lipschitz}. Reversing the $v$ coordinate handles the
negative derivative. The result also applies on a constant lattice
component of a mixed density.
\end{proof}

\begin{proposition}[Finite common patches with an explicit remainder]
\label{prop:finite-common-patches}
Fix $n\ge1$, $R_0\in I$, and $\epsilon>0$. There are finitely many
nonnegative smooth cutoffs $\psi_i$ on the common collision space, a
parameter neighborhood $J$ of $R_0$, and $C_{\epsilon,n,R_0}<\infty$
such that $\sum_i\psi_i\le1$, every support lies in one common regular
first-$n$-return patch for all $R\in J$, and
\begin{equation}\label{eq:common-patch-mass}
 c_*^{-1}\int\sum_i\psi_i\dd\nu>1-\epsilon.
\end{equation}
The lattice labels and all intermediate return decisions are fixed on
each patch. If $p^i_{n,R}$ is the pushforward density of
$c_*^{-1}\psi_i\nu$, then
\begin{align}\label{eq:common-patch-bound}
 \sum_i\|p^i_{n,R}-p^i_{n,S}\|_1
       &\le C_{\epsilon,n,R_0}|R-S|,\\
 \|p_{n,R}-p_{n,S}\|_1
       &\le2\epsilon+C_{\epsilon,n,R_0}|R-S|
               \qquad(R,S\in J).\nonumber
\end{align}
Before choosing these patches one may discard $N_{n,R_0}>L$ at cost at
most $C_1e^{-a_1L/n}$, with the constants of
Theorem~\ref{thm:exponential-return}.
\end{proposition}
\begin{proof}
The exponential moment bound proves the last assertion. Choose $L$ so
that this cost is less than $\epsilon/3$. On the remainder remove the
initial boundary of the section, all grazing and competing-root
singularities encountered before collision $L$, and the preimages of the
moving return boundary at $R_0$ through those collisions. These sets have
zero collision measure. On every remaining regular word also remove the
critical points of its cumulative flight length. Proposition
\ref{prop:general-edge} shows that there is at most one such point on a
regular optical word; the submersion complement has full measure.
There are only finitely many possible physical lattice words of length
at most $L$, since each flight label belongs to a fixed finite set.
No estimate for the number of unbounded-length words is used here.

Inner regularity of the collision measure and a countable exhaustion of
the regular submersion patches give a compact subset of mass greater
than $1-\epsilon$ in the normalized section. At every point of this
compact subset, choose a small coordinate rectangle on which one
coordinate derivative of the cumulative time is bounded away from zero.
The finite-itinerary lemma preserves the entire physical word and every
return decision on a slightly larger rectangle for $R$ close to $R_0$.
Select finitely many such rectangles and a subordinate family of smooth
cutoffs with sum at most one, equal to one on the compact subset.
Their supports have positive distance from all the excluded sets at
$R_0$. Finiteness gives a common parameter neighborhood $J$ on which
these distances, return decisions and derivative bounds persist. The
time functions and their first two derivatives are uniformly bounded
on these finitely many compact rectangles. Lemma
\ref{lem:common-submersion-chart} applies with
$a=\psi_i/(4\pi c_*)$ in $(\alpha,p)$ coordinates. Summing its bounds
proves the first line of \eqref{eq:common-patch-bound}.

For every $R\in J$ the remainder of the original probability is a
positive measure. Its mass is exactly
$1-c_*^{-1}\int\sum_i\psi_i\dd\nu<\epsilon$, because the section
measure is the fixed number $c_*$. The raw densities exist by the
submersion/coarea argument of Theorem~\ref{thm:return-stability}.
The two remainder densities have total $L^1$ norm less than
$2\epsilon$. This gives the second line. Nothing has been assumed about
the size of a density at a critical value or a sharp image boundary.
\end{proof}

The constants in \eqref{eq:common-patch-bound} expose the order of
choices: physical-count truncation, deletion of small initial-state
neighborhoods, compact common charts, and finally parameter tolerance.
They are not asserted to remain bounded as $n$ or $\epsilon^{-1}$ grows.
A quantitative global density modulus requires quantitative control of
those additional choices. Bounded measurable insertions are handled by
$L^1$ approximation on the same patches, as in
Proposition~\ref{prop:weighted-return-stability}; no artificial observation
coordinate is introduced.

\section{A uniform functional law for the mechanical clock}
\label{sec:uniform-fluid-clock}
Finite-record continuity also has a long-time consequence that does not
require a mixing rate. We prove a parameter-uniform functional law at the
first-order scale, including the maximum unfinished return over an entire
observation interval. The \(\sqrt t\)-scale process theorem in
Proposition~\ref{prop:clock} remains a separate statement.

Set
\[
 G_R=(\kappa_R^*,r_R^*,\tau_R^*),\qquad
 \bar G_R=(0,0,c_*^{-1},\bar\tau_R/c_*),\qquad J_{0,R}=0.
\]
All probabilities on the return section use $\nu_R^*$, not an
independent sequence of records. Write $\bar r=c_*^{-1}$ and
$\bar\tau_R^*=\bar\tau_R/c_*$. Both $\bar\tau_R^*$ and its reciprocal
are bounded on $I$.

\begin{lemma}[A compact subadditive argument]
\label{lem:compact-subadditive}
Let $a_n:I\to[0,\infty)$ be continuous, $a_0=0$, and
$a_{n+m}(R)\le a_n(R)+a_m(R)$. Suppose
$\inf_{n\ge1}a_n(R)/n=0$ for every $R$. Then
\[
 \lim_{n\to\infty}\sup_{R\in I}\frac{a_n(R)}n=0.
\]
More precisely, for every $\epsilon>0$ there is a finite integer $M$
such that, with $C=\sup_R a_1(R)$,
\begin{equation}\label{eq:subadditive-uniform-budget}
 \sup_R a_n(R)/n\le\epsilon+CM/n\qquad(n\ge1).
\end{equation}
\end{lemma}
\begin{proof}
For each $R$ choose a block length $m_R$ with
$a_{m_R}(R)/m_R<\epsilon$. Continuity preserves this strict inequality
on a neighborhood of $R$. Take a finite subcover and let $M$ be the
largest of its block lengths. For an arbitrary parameter choose one of
the covering neighborhoods, with block length $m\le M$, and write
$n=qm+r$, $0\le r<m$. Subadditivity and $a_r\le r a_1$ give
$a_n\le q\epsilon m+Cr\le\epsilon n+CM$. This proves the
estimate and then the limit. The neighborhoods need not have a common
block length; replacing them by an unproved common mixing time would
not be justified.
\end{proof}

\begin{theorem}[Uniform first-order return process]
\label{thm:uniform-return-fluid}
The actual induced records satisfy
\begin{equation}\label{eq:uniform-return-mean-law}
 \lim_{n\to\infty}\sup_{R\in I}
       \int\left|J_{n,R}/n-\bar G_R\right|\dd\nu_R^*=0.
\end{equation}
For every fixed $A<\infty$ and $\epsilon>0$,
\begin{equation}\label{eq:uniform-return-functional}
 \lim_{n\to\infty}\sup_{R\in I}\nu_R^*\left\{
 \sup_{0\le u\le A}
 \left|n^{-1}J_{\lfloor nu\rfloor,R}-u\bar G_R\right|>\epsilon
 \right\}=0.
\end{equation}
\end{theorem}
\begin{proof}
The ergodic collision map in Section~\ref{sec:return-stability}
induces an ergodic map on its positive-measure return section. Indeed,
the complete tower over an invariant subset of the induced map is an
invariant collision set. The two towers over that subset and its
complement partition the full collision space modulo null sets;
collision ergodicity forces one of the two to have zero measure.
Thus the integrable vector $G_R$ obeys the $L^1$ ergodic theorem for
each fixed $R$, with the exact mean in
Proposition~\ref{prop:exact-induced-means}.

Put $a_n(R)=\int|J_{n,R}-n\bar G_R|\dd\nu_R^*$.
Invariance of $F_R^*$ and the cocycle identity make $a_n$ subadditive.
It is continuous: on the common collision space it equals
\[
 c_*^{-1}\int_M
 \left|\widetilde J_{n,R}-n\bar G_R\one_{Y_R^*}\right|\dd\nu.
\]
Theorem~\ref{thm:return-stability} gives the first term's $L^1$
continuity, the section indicators are continuous in $L^1$ because
their boundaries move continuously and have zero measure, and the
mean is continuous by its explicit formula. The pointwise ergodic
theorem gives $a_n(R)/n\to0$. Lemma~\ref{lem:compact-subadditive}
proves \eqref{eq:uniform-return-mean-law}.

For the functional conclusion, first test a fixed finite mesh of
$[0,A]$. Equation~\eqref{eq:uniform-return-mean-law} and Markov's
inequality give convergence in probability at all its mesh points,
uniformly in $R$. The count and roof coordinates of $J_n$ are
nondecreasing. Between mesh points their normalized deviations are
therefore bounded by the endpoint deviations plus a uniformly bounded
mean times the mesh width, with an $O(n^{-1})$ integer-part error.
For the displacement coordinate, finite horizon gives a constant
$C_\kappa$ such that for every $l\ge j$,
\[
 |K_{l,R}-K_{j,R}|\le C_\kappa(N_{l,R}-N_{j,R}).
\]
Consequently its oscillation on a mesh interval is controlled by the
count increment on that interval. Its mean is zero. On the event
where all mesh errors are small, every coordinate has a uniformly
small error between mesh points as well. Let $n$ grow, then let the
mesh width and the mesh error tolerance decrease. This proves
\eqref{eq:uniform-return-functional} without an independence or
maximal-martingale assumption.
\end{proof}

\begin{lemma}[Sublinear maximum return blocks]
\label{lem:uniform-maximal-block}
Put $Q_R=r_R^*+\tau_R^*+|\kappa_R^*|$ and
\[
 U(L)=\sup_R\int Q_R\one_{\{Q_R>L\}}\dd\nu_R^*.
\]
Then $U(L)\to0$. For $n\ge1$, $A\ge0$, and $\epsilon>0$,
\begin{equation}\label{eq:uniform-maximal-block}
 \sup_R\nu_R^*\left\{
 \max_{0\le j\le\lfloor An\rfloor}Q_R\circ(F_R^*)^j>\epsilon n
 \right\}\le\frac{A+1}{\epsilon}U(\epsilon n).
\end{equation}
Both \eqref{eq:uniform-return-functional} and the convergence to zero
in \eqref{eq:uniform-maximal-block} remain valid when the initial
base point has the stationary length-biased law
\begin{equation}\label{eq:length-biased-base-v4}
 \dd\rho_R=\frac{\tau_R^*}{\bar\tau_R^*}\dd\nu_R^*.
\end{equation}
\end{lemma}
\begin{proof}
Finite horizon bounds $\tau_R^*$ and $|\kappa_R^*|$ by a uniform
constant times $r_R^*$. The uniform integrability in
Theorem~\ref{thm:return-stability}, and the constant section mass,
therefore give $U(L)\to0$. Invariance of $F_R^*$, the union bound
and truncated Markov inequality give
\[
 \nu_R^*\{\max_{j\le\lfloor An\rfloor}Q_R\circ F^j>\epsilon n\}
 \le (\lfloor An\rfloor+1)\nu_R^*(Q_R>\epsilon n)
 \le\frac{A+1}{\epsilon}\int Q_R\one_{Q_R>\epsilon n}\dd\nu_R^*.
\]
No independence of the events is used.

Let $d_*:=\inf_R\bar\tau_R^*>0$. For any event $B_R$ determined
by the forward base orbit and every $L>0$,
\begin{equation}\label{eq:length-bias-transfer-v4}
 \rho_R(B_R)\le d_*^{-1}\left(
 L\nu_R^*(B_R)+\int\tau_R^*\one_{\{\tau_R^*>L\}}\dd\nu_R^*\right).
\end{equation}
If $\sup_R\nu_R^*(B_R)\to0$, first hold $L$ fixed and pass to that
limit, then let $L\to\infty$. Uniform integrability of the roof
makes the right side tend to zero uniformly. Apply this argument to
the exceptional events in the two displayed convergence statements.
This transfer uses a length-biased law, not the unweighted launch law.
\end{proof}

For the stationary physical flow, let $V_R(t)$ count future visits to
$Y_R^*$, let $C_R(t)$ count future physical collisions, and let
$W_R(t)\in\Z^2$ be the net fundamental-cell displacement since time
zero. Changing the choice of a fixed fundamental cell changes the
continuous-position comparison only by a bounded term.

\begin{theorem}[Uniform stationary physical clock]
\label{thm:uniform-physical-clock}
For every $\epsilon>0$, under the actual stationary billiard laws
$\mu_R$, uniformly for $R\in I$,
\begin{equation}\label{eq:uniform-physical-clock}
 \sup_{0\le u\le1}\left|
 \frac1t\bigl(V_R(tu),C_R(tu),W_R(tu)\bigr)
 -\left(\frac{uc_*}{\bar\tau_R},
        \frac{u}{\bar\tau_R},0\right)\right|
 \longrightarrow0
 \quad\hbox{in probability as }t\to\infty.
\end{equation}
The maximum marked unfinished return encountered during $[0,t]$,
including the initial return containing time zero, is $o_P(t)$
uniformly in $R$.
\end{theorem}
\begin{proof}
Represent the starting state by $(x,s)$ in the induced suspension,
$0\le s<\tau_R^*(x)$. The marginal law of $x$ is exactly $\rho_R$
in \eqref{eq:length-biased-base-v4}. Write $T_{j,R}$ for the roof
coordinate of $J_{j,R}$. Count future events in $(0,t]$. With this
convention the following inverse identity also holds at event times:
\[
 V_R(tu)=\max\{j\ge0:T_{j,R}(x)\le s+tu\}.
\]
Choose a fixed $A$ with $A d_*>2$. The functional return law,
transferred by Lemma~\ref{lem:uniform-maximal-block}, shows that
$T_{\lfloor At\rfloor,R}>s+t$ with probability tending to one
uniformly. Here $s/t\to0$ uniformly because $s\le Q_R(x)$ and
the same lemma applies at $j=0$. Integer parts do not affect these
statements; the return-process theorem may be applied with
$n=\lceil t\rceil$ and a slightly larger fixed horizon.

On this event all inverse indices lie below $At$. Uniformly for
$0\le j\le At$ we have
$|T_{j,R}-j\bar\tau_R^*|/t=o_P(1)$ and
$\max_{j\le At}Q_R\circ F^j/t=o_P(1)$. The inequalities
\[
 T_{V_R(tu),R}\le s+tu<T_{V_R(tu)+1,R}
\]
therefore imply
\[
 \sup_{u\le1}\left|
 \bar\tau_R^* V_R(tu)/t-u\right|=o_P(1)
\]
uniformly in $R$. Division by $\bar\tau_R^*\ge d_*$ yields the
first component in \eqref{eq:uniform-physical-clock}.

The completed-return count $N_{V_R(tu),R}$ differs from the actual
physical count by the initial prefix and the final incomplete return.
Their sum is at most a fixed constant plus twice the maximal $Q_R$
over these return indices. The displacement comparison has the same
bound, up to a uniform multiplicative constant, since each physical
lattice step is bounded. These bounds hold simultaneously for all
$u\in[0,1]$. The functional return law at the random inverse indices
now gives
\[
 C_R(tu)/t=\bar r\,V_R(tu)/t+o_P(1),\qquad
 W_R(tu)/t=o_P(1)
\]
with uniform suprema. Since $\bar r/\bar\tau_R^*=1/\bar\tau_R$,
this proves the remaining components. Every unfinished piece in
$[0,t]$ belongs to one of the same return blocks and is bounded by
its marked size, proving the final assertion.
\end{proof}

\begin{corollary}[Geometric error in the physical collision rate]
\label{cor:physical-rate-error}
Let $b_*:=\min_{R\in I}\bar\tau_R$ and put
$L_*=15/(4b_*^2)$. There is a function $\omega_t(\epsilon)\to0$
for every fixed $\epsilon>0$ such that, for any $I$-valued radius
estimate $\widehat R$ on a probability space carrying the stationary
billiard at the true radius $R$, and any $\delta>0$,
\begin{align}\label{eq:physical-rate-error}
 \Pr_R\left\{\sup_{0\le u\le1}
 \left|C_R(tu)/t-u/\bar\tau_{\widehat R}\right|
                       >\epsilon+L_*\delta\right\}
 \le\omega_t(\epsilon)+\Pr_R\{|\widehat R-R|>\delta\}.
\end{align}
No independence of $\widehat R$ and the observed orbit is required.
For the visit rate, replace the deterministic error constant by $c_*L_*$.
\end{corollary}
\begin{proof}
The explicit mean satisfies $|\bar\tau_R'|<15/4$, so
$|(1/\bar\tau_R)'|\le L_*$. On the event
$|\widehat R-R|\le\delta$ the two proposed rates differ by at
most $L_*\delta$, uniformly over $u\in[0,1]$. The triangle
inequality and Theorem~\ref{thm:uniform-physical-clock} prove the
claim with $\omega_t$ equal to the supremum, over $R$, of the true
rate's exceptional probability. Multiplication by $c_*$ gives the
visit-rate statement.
\end{proof}

The supremum statement here is at scale $t$. Uniform integrability
alone does not turn it into an $o_P(\sqrt t)$ bound for a growing
window. Likewise \eqref{eq:subadditive-uniform-budget} and
\eqref{eq:physical-rate-error} do not supply a numerical convergence
rate for $\omega_t$ without quantitative finite-block estimates.
They establish the full first-order clock comparison for this moving
family. The Gaussian covariance and processes require the additional
collision estimates proved in Sections~\ref{sec:collision-covariance}--
\ref{sec:functional-gaussian}; raw density and shrinking-conditioning
conclusions require the further local estimates stated there. The
stationary law throughout is not substituted for an unknown experimental
preparation law.

\section{Bounded collision compensation and its covariance}
\label{sec:collision-covariance}
The discontinuity of the section can be treated on the collision clock
without making its indicator a multiplier on a distribution space. The
argument uses bounded variation in the initial collision coordinates,
smooth approximation, and the collision-space spectral splitting. It
will give a covariance for the actual return record by an exact stopping
identity, rather than by an assumed spectral expansion of an induced
operator.

Write $\eta_R=\one_{Y_R^*}$ and recall
\begin{equation}\label{eq:compensation-normalizations}
 c_* = \frac{91}{10000\pi},\qquad
 \bar\tau_R=\frac{\sqrt3/2-\pi R^2}{2R},\qquad
 \bar G_R=\left(0,0,c_*^{-1},\frac{\bar\tau_R}{c_*}\right).
\end{equation}
Here $\bar G_R=\int G_R\dd\nu_R^*$ is the mean of the actual induced
record, by the Kac identities. On the common collision cylinder put
\begin{equation}\label{eq:bounded-compensation}
 h_R=f_R-\bar G_R\eta_R,
 \qquad s_R=c_*^{-1}\eta_R,
 \qquad S_{m,R}u=\sum_{j=0}^{m-1}u\circ T_R^j.
\end{equation}
Thus $\int h_R\dd\nu=0$ and $s_R\nu=\nu_R^*$, with the latter
measure regarded as a probability on $M$. The four components of $h_R$
are bounded, although those of $G_R$ need not be.

\subsection{Variation of a single physical collision}
For a function on $\T\times(-1,1)$, $\mathrm{BV}$ denotes bounded
variation with respect to the flat coordinates $(\alpha,p)$: both
first distributional derivatives are finite signed measures. All
variation norms below include the $L^1(\nu)$ norm. They are not norms
of the density of a pushed-forward record.

\begin{lemma}[Uniform single-collision variation]\label{lem:physical-bv}
There is a constant $B$ such that
\[
 \|f_R\|_\infty+\|f_R\|_{\mathrm{BV}}
 +\|\eta_R\|_\infty+\|\eta_R\|_{\mathrm{BV}}
 +\|h_R\|_\infty+\|h_R\|_{\mathrm{BV}}
 +\|s_R\|_\infty+\|s_R\|_{\mathrm{BV}}\le B
 \qquad(R\in I).
\]
The norms of vector-valued functions mean the sum of the component norms.
\end{lemma}
\begin{proof}
The uniform horizon bounds all possible next-disk centers in a fixed
finite subset of the triangular lattice. Cover the angular circle by
finitely many bounded rational charts
$t=\tan((\alpha-\alpha_0)/2)$, with uniformly bounded chart derivatives
and inverse derivatives. In such a chart $n_\alpha$ and $t_\alpha$ are
rational functions of $t$. The outgoing velocity
$\sqrt{1-p^2}n_\alpha+pt_\alpha$ is semialgebraic in $(t,p)$. For each
candidate center $d$, the admissibility inequalities and the entry root
\[
 (d-Rn_\alpha)\cdot v
 -\sqrt{R^2-|d-Rn_\alpha|^2+((d-Rn_\alpha)\cdot v)^2}
\]
are semialgebraic in $(R,t,p)$. Taking the least admissible positive
entry root and its lattice label involves only finitely many comparisons.
Choose any fixed tie convention on the singular set. Thus the bounded
single-collision record, extended arbitrarily by zero where no regular
root is defined, is a bounded semialgebraic family on each chart.
Only these one-collision functions, not an arbitrary long itinerary,
are used in this assertion.

We spell out the uniformity implication. A bounded semialgebraic family
$u(\lambda,x)$, with $x$ in a bounded interval, has a uniform finite
number of monotonicity intervals in $x$. To see this, apply differentiable
cell decomposition to its graph with the parameters $\lambda$ ordered
first, and refine by the sign of its fiber derivative where defined.
Uniform finiteness bounds the number of remaining exceptional points
and fiber intervals. On each interval the function is monotone or
constant. These are the cell decomposition and monotonicity results of
\cite[Sections 2.1--2.2 and 6.2]{Coste}. In particular the number does
not depend on the parameter or on the other fixed coordinate of a slice.
If $|u|\le M_0$ and there are at most $b_0$ intervals and exceptional
points, its one-dimensional distributional variation is at most
$C b_0 M_0$, including the jumps.

Apply this statement first with $(R,p)$ as parameters and then with
$(R,t)$ as parameters. Integrating the two slice-variation bounds
against the remaining coordinate proves the bounds for the two
first distributional derivatives. Equivalently, integration by parts
on almost every slice bounds the action on every compactly supported
$C^1$ test function by its supremum norm. A fixed finite chart partition
transfers these bounds to $\alpha$; any jumps at the chart endpoints
have uniformly bounded size and length. This proves the assertion for
$f_R$, including the square-root behavior at grazing.

The section is a fixed finite union of coordinate rectangles with
moving centers. Its perimeter and the number of its sides are uniformly
bounded, so the same assertion for $\eta_R$ follows directly, without
requiring the centers to be semialgebraic in $R$. The explicit means in
\eqref{eq:compensation-normalizations} are bounded on $I$. The remaining
bounds now follow from \eqref{eq:bounded-compensation}.
\end{proof}

\begin{lemma}[Mean-preserving smoothing]\label{lem:bv-smoothing}
There are positive linear operators $\mathsf S_\delta$, $0<\delta<1/2$,
which preserve the integral with respect to $\nu$ and satisfy
\begin{equation}\label{eq:bv-smoothing-bounds}
 \begin{split}
 \|\mathsf S_\delta u\|_\infty&\le\|u\|_\infty,\qquad
 \|u-\mathsf S_\delta u\|_1\le C\delta\|u\|_{\mathrm{BV}},\\
 \|\mathsf S_\delta u\|_{\mathrm{BV}}&\le C\|u\|_{\mathrm{BV}},\qquad
 \|\mathsf S_\delta u\|_{C^2(\alpha,\varphi)}
       \le C\delta^{-2}\|u\|_\infty .
 \end{split}
\end{equation}
The last norm uses $p=\sin\varphi$, as in the collision spaces.
\end{lemma}
\begin{proof}
Reflect $u$ evenly across $p=1$ and $p=-1$ to a function periodic with
period four in $p$, retaining the angular periodicity. Convolve with a
nonnegative, even, smooth product kernel of mass one and scale $\delta$,
then restrict to $-1<p<1$. Reflection symmetry is preserved, so the
integral over this half of the extended cylinder remains the original
integral. The reflected function has variation bounded by a fixed
multiple of the original variation; its traces match at the reflection
boundaries. Translation of a bounded-variation function by a vector $a$
changes its $L^1$ norm by at most $|a|$ times its variation. This follows
first for smooth functions by the fundamental theorem of calculus,
then for distributional derivatives by convolution and passage to the
limit. Integrating this inequality against the kernel gives the $L^1$
error. Convolution contracts variation on the extended cylinder.
The $L^1$ norms of the first two derivatives of the kernel are
$O(\delta^{-1})$ and $O(\delta^{-2})$. These facts prove all the flat
coordinate estimates. Composition with $p=\sin\varphi$ preserves the
stated $C^2$ upper bound, since the first two derivatives of $\sin$
are bounded. No inverse coordinate change at grazing is used.
\end{proof}

\subsection{An exponentially summable collision covariance}
\begin{proposition}[Correlations of bounded-variation observables]
\label{prop:bv-correlation}
For every $B_0<\infty$ there are $C,\gamma>0$, uniform in $R$, such
that centered real functions $a,b$ with
$\|a\|_\infty+\|a\|_{\mathrm{BV}}+
\|b\|_\infty+\|b\|_{\mathrm{BV}}\le B_0$ satisfy
\begin{equation}\label{eq:bv-correlation}
 \left|\int a(b\circ T_R^k)\dd\nu\right|\le Ce^{-\gamma k}
 \qquad(k\ge0).
\end{equation}
If also $\|a\|_1\le C_0\delta$, then for $r=a$ and $b=a$,
\begin{equation}\label{eq:small-bv-variance}
 \int|S_{m,R}r|^2\dd\nu
       \le C m\delta(1+|\log\delta|).
\end{equation}
The second constant may depend on $C_0$ and $B_0$, but not on $m,R$.
\end{proposition}
\begin{proof}
For smooth centered $a,b$, apply Lemma~\ref{lem:collision-spectral-input}
to the distribution $a\nu$ and evaluate the result on $b$ using smooth
multiplication and the mass functional. It gives
\[
 \left|\int a(b\circ T_R^k)\dd\nu\right|
       \le C\vartheta^k\|a\|_{C^2}\|b\|_{C^2}.
\]
Smooth both bounded-variation functions at scale $\epsilon$. Invariance
of $\nu$, boundedness, and \eqref{eq:bv-smoothing-bounds} bound the
change in the correlation by $C\epsilon$. The smooth correlation is
bounded by $C\vartheta^k\epsilon^{-4}$. Choose
$\epsilon=\tfrac14\vartheta^{k/5}$ and enlarge $C$ for $k=0$.
This proves \eqref{eq:bv-correlation}. In particular the constants depend
only on the displayed variation and supremum bounds, not on a modulus
of continuity or on a chosen smoothing scale.

For the last assertion, boundedness gives
$|\int r(r\circ T_R^k)\dd\nu|\le C\delta$ at every lag, while
\eqref{eq:bv-correlation} gives the exponential bound. Hence
\[
 \sum_{k\ge0}\left|\int r(r\circ T_R^k)\dd\nu\right|
       \le C\sum_{k\ge0}\min\{\delta,e^{-\gamma k}\}
       \le C\delta(1+|\log\delta|).
\]
Expanding the square of the stationary sum proves
\eqref{eq:small-bv-variance}.
\end{proof}

Put $h_{R,\delta}=\mathsf S_\delta h_R$. Define, on the collision clock,
\begin{equation}\label{eq:collision-gk}
 \begin{split}
 \Gamma_R={}&\int h_R\otimes h_R\dd\nu\\
 &+\sum_{k=1}^{\infty}\int
 \bigl[h_R\otimes(h_R\circ T_R^k)
       +(h_R\circ T_R^k)\otimes h_R\bigr]\dd\nu .
 \end{split}
\end{equation}
Use the same formula with $h_{R,\delta}$ for $\Gamma_{R,\delta}$.

\begin{theorem}[A continuous collision Green--Kubo matrix]
\label{thm:collision-covariance}
The series \eqref{eq:collision-gk} and its smoothed versions converge
absolutely, with exponential tails uniform in $R$ and $\delta$. The
matrix $\Gamma_R$ is continuous and positive semidefinite on $I$, and
\begin{equation}\label{eq:collision-covariance-limit}
 \Gamma_R=\lim_{m\to\infty}\frac1m
   \int S_{m,R}h_R\otimes S_{m,R}h_R\dd\nu,
 \qquad
 \|\Gamma_{R,\delta}-\Gamma_R\|\le C\delta(1+|\log\delta|).
\end{equation}
The first limit is uniform in $R$.
\end{theorem}
\begin{proof}
Lemma~\ref{lem:physical-bv}, smoothing, and
Proposition~\ref{prop:bv-correlation} give the uniform exponential tails.
For each fixed $k$, the integrand in \eqref{eq:collision-gk} converges
almost everywhere as $R\to R_0$. Indeed, exclude the finitely many
collision singularity preimages and section boundaries encountered up
to time $k$ at $R_0$. On the remaining full-measure set the physical
roots and finite itinerary persist and the moving rectangles have
stable membership. All observables are uniformly bounded. Dominated
convergence proves continuity at each fixed lag; the uniform tail then
proves continuity of the infinite series. This argument establishes
long-time uniformity through the correlation estimate, not through a
fixed-record continuity constant.

Expansion of the square gives the finite covariance with lag weights
$1-k/m$. Absolute summability, uniformly in $R$, proves its limit;
indeed the error is $O(m^{-1})$ by the exponential bound. Positivity
follows from positivity of each finite covariance. At each lag the
change made by smoothing is at most $C\delta$, using invariance and the
$L^1$ approximation. Both the original and the smoothed lag correlations
are at most $Ce^{-\gamma k}$. Summing the minimum of these two bounds
proves the final estimate.
\end{proof}

\begin{remark}\label{rem:collision-gk-scope}
The sum in \eqref{eq:collision-gk} is for the bounded, compensated
collision observable $h_R$ under $T_R$. It is not the Green--Kubo sum
of the unbounded observable $G_R-\bar G_R$ under $F_R^*$. Its relation
to the Gaussian law of the latter will be proved by stopping the actual
collision sums. No bounded-variation estimate for a many-return coarea
density, and no second derivative of that density, follows from the
single-collision variation lemma.
\end{remark}

\section{A uniform Gaussian limit for the actual return record}
\label{sec:stopped-gaussian}
We first obtain a quantitative characteristic-function estimate on the
collision clock. Smooth approximation is used only inside the proof;
the limiting statement concerns the original record and the original
section probability.

\subsection{Smooth perturbation with a shrinking scale}
\begin{lemma}[A quantitative smooth collision expansion]
\label{lem:smooth-collision-expansion}
On each of the collision spaces of
Lemma~\ref{lem:collision-spectral-input}, set
\[
 \mathcal L_{R,\delta,z}
       =\mathcal L_R M_{\exp(i z\cdot h_{R,\delta})},
       \qquad z\in\mathbb C^4.
\]
There are $a,C>0$ and $\rho<1$, independent of $R$ and
$0<\delta<1/2$, such that for $|z|<a\delta^2$ this family has a simple
analytic eigenvalue $\lambda_{R,\delta}(z)$ near one and an analytic
rank-one projection $\Pi_{R,\delta}(z)$. Its complementary powers are
bounded by $C\rho^m$, and, on a smaller such ball,
\begin{equation}\label{eq:smooth-cubic-expansion}
 \begin{split}
 \log\lambda_{R,\delta}(z)
       &=-\tfrac12z^{\mathsf T}\Gamma_{R,\delta}z
                  +O(\delta^{-6}|z|^3),\\
 \|\Pi_{R,\delta}(z)-\Pi_R\|
       &\le C\delta^{-2}|z| .
 \end{split}
\end{equation}
All constants are uniform. The logarithm is the branch vanishing at
$z=0$.
\end{lemma}
\begin{proof}
The smooth multiplier estimate and the smoothing bounds give
$\|\mathcal L_{R,\delta,z}-\mathcal L_R\|
\le C|z|\delta^{-2}$ when $|z|\delta^{-2}$ is small.
The uniform resolvent contours around one and around the complementary
spectrum in the proof of Proposition~\ref{prop:soft-killing} therefore
apply on $|z|<a\delta^2$. The resolvent Neumann series gives the analytic
projection, its difference estimate, and complementary powers on a
fixed circle of radius $\rho<1$. Decrease $a$ so that the eigenvalue
stays in a fixed disk about one not containing zero.

The eigenvalue derivative at zero is zero since $\int h_{R,\delta}
\dd\nu=0$. We identify its second derivative, rather than assuming a
variance formula. Fix a real direction $v$, put $u=v\cdot h_{R,\delta}$,
and normalize a right eigenvector $e(t)$ by $\ell_j(e(t))=1$.
Differentiating the eigenvector equation at zero yields
\[
 e'(0)=i\sum_{k\ge1}\mathcal L_R^k(u\nu),
\]
where the series converges in $\mathcal B_j$ by the spectral splitting.
A second differentiation followed by the mass functional gives
\[
 \lambda''(0)=-\int u^2\dd\nu
       -2\sum_{k\ge1}\int u(u\circ T_R^k)\dd\nu
       =-v^{\mathsf T}\Gamma_{R,\delta}v.
\]
Polarization identifies the full Hessian. Since the first derivative
vanishes, the Hessian of the logarithm is the same. Cauchy's estimates
for the bounded analytic logarithm on a ball of radius $a\delta^2$
bound its third derivatives by $C\delta^{-6}$ on a smaller ball.
Taylor's formula proves \eqref{eq:smooth-cubic-expansion}.
\end{proof}

\begin{proposition}[Uniform collision characteristic estimate]
\label{prop:collision-gaussian}
For every $K<\infty$ there is $C_K$ such that, for $m\ge2$ and for
$\sigma_R=1$ or $\sigma_R=s_R$,
\begin{equation}\label{eq:collision-gaussian}
 \sup_{R\in I,\ |t|\le K}
 \left|\int \sigma_R e^{it\cdot S_{m,R}h_R/\sqrt m}\dd\nu
       -e^{-t^{\mathsf T}\Gamma_Rt/2}\right|
       \le C_K m^{-1/26}\sqrt{\log(2+m)}.
\end{equation}
\end{proposition}
\begin{proof}
Smooth both $h_R$ and the initial density at scale $\delta$. The
smoothed density $\sigma_{R,\delta}=\mathsf S_\delta\sigma_R$ is
nonnegative, has integral one, and satisfies
$\|\sigma_{R,\delta}\nu\|_{\mathcal B_j}\le C\delta^{-2}$.
For $z=t/\sqrt m$, the transfer pairing and the spectral decomposition
of Lemma~\ref{lem:smooth-collision-expansion} give
\[
 \int\sigma_{R,\delta}e^{it\cdot S_{m,R}h_{R,\delta}/\sqrt m}\dd\nu
 =\lambda_{R,\delta}(z)^m
       \ell_j\bigl(\Pi_{R,\delta}(z)\sigma_{R,\delta}\nu\bigr)
       +O(\delta^{-2}\rho^m).
\]
At $z=0$ the amplitude is one; its difference from one is
$O_K(\delta^{-4}m^{-1/2})$. The logarithmic expansion gives, whenever
$K/\sqrt m<a\delta^2$ and $\delta^{-6}m^{-1/2}$ is sufficiently small,
\begin{equation}\label{eq:smoothed-characteristic-error}
 \left|\int\sigma_{R,\delta}
 e^{it\cdot S_{m,R}h_{R,\delta}/\sqrt m}\dd\nu
       -e^{-t^{\mathsf T}\Gamma_{R,\delta}t/2}\right|
 \le C_K\bigl(\delta^{-6}m^{-1/2}
          +\delta^{-4}m^{-1/2}+\delta^{-2}\rho^m\bigr).
\end{equation}
Here positive semidefiniteness of $\Gamma_{R,\delta}$ bounds the
Gaussian factor by one. The estimate for the exponential of the
Taylor remainder follows, for example, from
$|e^w-1|\le |w|e^{|w|}$.

The change of initial density costs at most $C\delta$ in the original
pairing, because the exponential has modulus one. The centered
residual $r_{R,\delta}=h_R-h_{R,\delta}$ has uniformly bounded
supremum and variation norms and $L^1$ norm $O(\delta)$. Therefore
Proposition~\ref{prop:bv-correlation}, component by component, gives
\[
 \int\sigma_{R,\delta}
       |S_{m,R}r_{R,\delta}|^2\dd\nu
       \le C m\delta(1+|\log\delta|),
\]
since the initial density is bounded independently of $\delta$.
The inequality $|e^{ix}-e^{iy}|\le|x-y|$ and Cauchy--Schwarz bound
the change of observable by
$C_K\sqrt{\delta(1+|\log\delta|)}$. Finally,
Theorem~\ref{thm:collision-covariance} bounds the change of the Gaussian
factor by $C_K\delta(1+|\log\delta|)$.

Take $\delta=\tfrac14 m^{-1/13}$. The smoothing error is
$O(m^{-1/26}\sqrt{\log(2+m)})$ and the cubic spectral error is
$O(m^{-1/26})$. The two smallness requirements hold for all sufficiently
large $m$, depending only on $K$. The amplitude and complementary errors
are smaller. Increasing $C_K$ handles the remaining finitely many $m$.
This proves the estimate uniformly over the finite cover of collision
spaces. It does not use an induced spectral gap.
\end{proof}

\subsection{Stopping at the genuine return times}
\begin{lemma}[Exact compensation and a return-clock window]
\label{lem:stopped-compensation}
For $\nu_R^*$-almost every initial state and every $n\ge1$,
\begin{equation}\label{eq:exact-stopped-compensation}
 J_{n,R}-n\bar G_R=S_{N_{n,R},R}h_R.
\end{equation}
There is $C$ such that for $n\ge2$ and $2\le a\le n$,
\begin{equation}\label{eq:return-clock-window}
 \nu_R^*\{|N_{n,R}-n/c_*|>a\}\le Cn/a^2.
\end{equation}
\end{lemma}
\begin{proof}
A trajectory starting in $Y_R^*$ has exactly $n$ visits to this section
among collision times $0,\ldots,N_{n,R}-1$. These are its initial visit
and its first $n-1$ returns. Also
$J_{n,R}=S_{N_{n,R},R}f_R$. Substitution in
\eqref{eq:bounded-compensation} proves the identity with no remainder.

Let $H_{L,R}=\sum_{j=1}^L\eta_R\circ T_R^j$. Its stationary mean is
$c_*L$. Proposition~\ref{prop:bv-correlation} applied to
$\eta_R-c_*$ gives
\[
 \int|H_{L,R}-c_*L|^2\dd\nu\le CL.
\]
The same upper bound, with a changed constant, holds under $\nu_R^*$
because its density relative to $\nu$ is at most $c_*^{-1}$.
The exact equivalences
$N_{n,R}>L\iff H_{L,R}<n$ and
$N_{n,R}\le L\iff H_{L,R}\ge n$ reduce both tails to this estimate.
Choose integers within distance one of $n/c_*+a$ and $n/c_*-a$.
The corresponding deviations from $c_*L$ are at least $c_*a-O(1)$,
and $L\le Cn$. Chebyshev's inequality proves the assertion when $a$
exceeds a fixed constant. For the bounded remaining range, enlarge
$C$ so that the right side is at least one.
\end{proof}

\begin{lemma}[A deterministic-window maximal estimate]
\label{lem:dyadic-collision-maximal}
For $L\ge1$, $r\ge0$, and $R\in I$,
\begin{equation}\label{eq:dyadic-collision-maximal}
 \int\max_{0\le j\le L}
 |S_{r+j,R}h_R-S_{r,R}h_R|^2\dd\nu
       \le CL[\log(2+L)]^2.
\end{equation}
The same estimate holds under $\nu_R^*$ and for backward sums on a
deterministic interval contained in the nonnegative collision times.
\end{lemma}
\begin{proof}
Absolute summability of the correlations gives second moment at most
$C\ell$ for the sum over every interval of $\ell$ collisions, by
stationarity. Enlarge $L$ to a power of two smaller than $2L$.
Every prefix of that interval is a disjoint union of at most one dyadic
interval at each scale. Cauchy--Schwarz bounds the square of its sum by
the number of scales times the sum of the squares of these interval
sums. Bound the latter by the sum over all dyadic intervals at all
scales, independently of the prefix. At each scale, the expectations
sum to at most $CL$, since the intervals partition the full interval.
There are $O(\log(2+L))$ scales, proving the bound. The proof applies
to each component and to intervals in reverse order. A bounded initial
density gives the assertion under $\nu_R^*$ as well.
\end{proof}

\begin{theorem}[Parameter-uniform Gaussian limit of the physical record]
\label{thm:actual-record-gaussian}
Define $D_R=c_*^{-1}\Gamma_R$. It is a continuous positive semidefinite
matrix on $I$. For every $K<\infty$ there is $C_K$ such that
\begin{equation}\label{eq:actual-record-gaussian}
 \sup_{R\in I,\ |t|\le K}
 \left|\int_{Y_R^*}
 e^{it\cdot(J_{n,R}-n\bar G_R)/\sqrt n}\dd\nu_R^*
          -e^{-t^{\mathsf T}D_Rt/2}\right|
       \le C_K n^{-1/26}\sqrt{\log(2+n)}
       \qquad(n\ge2).
\end{equation}
In particular the original, unsmoothed return record has a centered
Gaussian weak limit with covariance $D_R$. The estimate is uniform in
the physical parameter, including along convergent parameter sequences.
\end{theorem}
\begin{proof}
Put $m_n=\lfloor n/c_*\rfloor$ and
$a_n=\lceil n^{3/5}\rceil$. The clock-window lemma, with a harmless
change for rounding, bounds the probability of
$|N_{n,R}-m_n|>a_n$ by $Cn/a_n^2$. On the complementary event, the
difference $S_{N_{n,R},R}h_R-S_{m_n,R}h_R$ is bounded by the larger
of the forward and backward maxima on intervals of length $a_n$ about
$m_n$. The intervals have nonnegative endpoints since $c_*<1$.
The preceding lemma and Cauchy--Schwarz consequently give
\begin{equation}\label{eq:stopped-characteristic-error}
 \begin{split}
 &\left|\int e^{it\cdot S_{N_{n,R},R}h_R/\sqrt n}\dd\nu_R^*
       -\int e^{it\cdot S_{m_n,R}h_R/\sqrt n}\dd\nu_R^*\right|\\
 &\hspace{20mm}\le C_K\left(\frac{n}{a_n^2}
                  +\sqrt{\frac{a_n}{n}}\log(2+a_n)\right)
       \le C_K n^{-1/5}\log(2+n).
 \end{split}
\end{equation}
The exceptional event costs at most twice its probability. No estimate
of an unbounded sum on that event is used.

Apply Proposition~\ref{prop:collision-gaussian} at time $m_n$ and
frequency $t\sqrt{m_n/n}$ with initial density $s_R$. This frequency
stays in a fixed compact set. Since $m_n/n=c_*^{-1}+O(n^{-1})$ and
$\Gamma_R$ is uniformly bounded, its Gaussian differs from the one in
\eqref{eq:actual-record-gaussian} by $O_K(n^{-1})$. The identity
\eqref{eq:exact-stopped-compensation} completes the estimate; the error
in \eqref{eq:stopped-characteristic-error} is smaller than the stated
one. The continuity of $D_R$ follows from
Theorem~\ref{thm:collision-covariance}. The characteristic-function
continuity theorem yields the weak limit, allowing a singular Gaussian.
Along $R_n\to R_0$, the same bound and continuity give the Gaussian
with covariance $D_{R_0}$.
\end{proof}

\subsection{The Gaussian kernel and measurable cohomology}
\begin{theorem}[The kernel is an actual $L^2$ coboundary]
\label{thm:gaussian-kernel-coboundary}
For $R\in I$ and $v\in\R^4$, the following three statements are
equivalent:
\begin{enumerate}
\item $v^{\mathsf T}D_Rv=0$;
\item there is $b\in L^2(M,\nu)$ such that
 $v\cdot h_R=b-b\circ T_R$ almost everywhere;
\item there is $b_*\in L^2(Y_R^*,\nu_R^*)$ such that
 $v\cdot(G_R-\bar G_R)=b_*-b_*\circ F_R^*$ almost everywhere.
\end{enumerate}
If the first condition holds, one may choose $b$ with
$\|b\|_2\le C|v|$, uniformly in $R$, and take $b_*=b|_{Y_R^*}$.
\end{theorem}
\begin{proof}
Put $u=v\cdot h_R$ and $c_k=\int u(u\circ T_R^k)\dd\nu$.
By \eqref{eq:bv-correlation}, $|c_k|\le C|v|^2e^{-\gamma k}$.
If $v^{\mathsf T}D_Rv=0$, then $c_0+2\sum_{k\ge1}c_k=0$.
The exact variance identity gives
\[
 \int|S_{m,R}u|^2\dd\nu
       =-2\sum_{k=1}^{m-1}k c_k-2m\sum_{k\ge m}c_k
       \le C|v|^2 \qquad(m\ge1).
\]
Let $Uu=u\circ T_R$, a unitary operator on $L^2(\nu)$, and form
$b_M=M^{-1}\sum_{m=1}^M S_{m,R}u$. This is a bounded sequence in
$L^2$, with norm at most $C|v|$. Moreover
\[
 (I-U)b_M=u-\frac1M\sum_{m=1}^M U^m u\longrightarrow u
       \quad\hbox{in }L^2.
\]
Here the mean ergodic theorem applies; the collision map is ergodic
and $u$ has mean zero. A weakly convergent subsequence of $b_M$ has
limit $b$, and bounded linearity of $I-U$ shows $(I-U)b=u$.
This proves the second assertion and its norm estimate.

Discard a countable union of collision translates of the null set
where this identity or a representative fails. On the remaining
invariant full-measure set, telescope until the actual first return.
The compensation identity for one block gives
\[
 v\cdot(G_R-\bar G_R)(x)=b(x)-b(F_R^*x).
\]
The restriction belongs to $L^2(\nu_R^*)$ and has norm at most
$c_*^{-1/2}\|b\|_2$. Thus the second assertion implies the third.
Conversely, the third assertion telescopes for the stationary induced
map, so the $L^2(\nu_R^*)$ norm of
$v\cdot(J_{n,R}-n\bar G_R)/\sqrt n$ is at most
$2\|b_*\|_2/\sqrt n$. Its characteristic function tends to one.
Theorem~\ref{thm:actual-record-gaussian} identifies the same limit as
$\exp(-t^2v^{\mathsf T}D_Rv/2)$ for every real $t$, proving the first
assertion. This also completes the equivalence without presupposing
summability of induced correlations.
\end{proof}

\begin{remark}[What the low-frequency theorem supplies]
\label{rem:gaussian-versus-raw}
The matrix $D_R$ is now established by an absolutely convergent
collision Green--Kubo formula and is the covariance of the Gaussian
law of the true induced record. The zero-Gaussian-variance condition
has been identified with a genuine induced $L^2$ coboundary. A
measurable representative of its transfer function has not thereby
been made continuous at a selected periodic orbit. Thus the periodic
rank calculation is not applied to this representative. The separate
phase argument in Theorem~\ref{thm:joint-uniform-nondegeneracy} proves
uniform positive definiteness without periodic evaluation. Neither
convergence of the induced Green--Kubo series nor convergence of its
normalized second moments follows from this stopping argument alone;
the actual second moments and the Cesaro formula are established in
Theorem~\ref{thm:marked-quadratic-moments} and
Corollary~\ref{cor:induced-cesaro-covariance}.

The estimate \eqref{eq:actual-record-gaussian} controls fixed compact
sets of rescaled central frequencies. It does not give the integrated
complementary-frequency bound in the raw local inversion theorem, an
induced eigenvalue expansion, or the complete critical/singular branch
sum. In particular it is not a raw density local limit, a weighted
local-density estimate, or a functional limit under a changed exact
conditioning event. These distinctions retain the original density
problem while supplying its actual Gaussian weak-limit input.
\end{remark}

\section{Functional Gaussian limits and the physical clock}
\label{sec:functional-gaussian}
A central-frequency estimate alone does not give tightness of the record
process. We prove it by smoothing on a slower scale, using fourth moments
for the smooth sums and a maximal second-moment estimate for the residual.
The estimates remain uniform over the moving physical family.

\subsection{Fourth moments and removal of smoothing}
\begin{lemma}[A fourth-moment bound at a specified smoothing scale]
\label{lem:smooth-fourth-moment}
For $m\ge1$, $R\in I$, and $0<\delta<1/2$,
\begin{equation}\label{eq:smooth-fourth-moment}
 \int |S_{m,R}h_{R,\delta}|^4\dd\nu
       \le C\bigl(m^2+(m+1)\delta^{-8}\bigr).
\end{equation}
For $r_{R,\delta}=h_R-h_{R,\delta}$ and $L\ge1$,
\begin{equation}\label{eq:residual-process-maximal}
 \int\max_{0\le j\le L}|S_{j,R}r_{R,\delta}|^2\dd\nu
 \le CL\delta(1+|\log\delta|)[\log(2+L)]^2.
\end{equation}
Both bounds hold, with changed constants, after multiplication by any
nonnegative initial density bounded by a fixed constant.
\end{lemma}
\begin{proof}
For the first assertion consider one real component $u$ of
$h_{R,\delta}$ and the one-variable family in
Lemma~\ref{lem:smooth-collision-expansion}. Write
$\psi(z)=\log\lambda(z)$ and $A(z)=\ell_j(\Pi(z)\nu)$. The spectral
pairing for the stationary characteristic function is
\[
 \int e^{izS_{m,R}u}\dd\nu=e^{m\psi(z)}A(z)+E_m(z).
\]
On a complex disk of radius $a\delta^2$, with $a$ decreased once,
$|E_m(z)|\le C\rho^m$, $|A(z)|\le C$, and $|\psi(z)|\le C$.
Cauchy's inequalities therefore give
$|A^{(j)}(0)|\le C_j\delta^{-2j}$ for $j\le4$,
$|\psi^{(j)}(0)|\le C_j\delta^{-2j}$ for $j=3,4$, and
$|E_m^{(4)}(0)|\le C\delta^{-8}\rho^m$.
In addition $A(0)=1$, $\psi'(0)=0$, and
$|\psi''(0)|\le C$, by the uniformly summable smoothed covariance.
The fourth derivative of the principal term at zero is exactly
\[
 A^{(4)}(0)+6m\psi''(0)A''(0)+4m\psi'''(0)A'(0)
       +m\psi^{(4)}(0)+3m^2\psi''(0)^2.
\]
It is bounded by $C(m^2+(m+1)\delta^{-8})$. Differentiating under the
finite-time integral is legitimate because $u$ is bounded. The fourth
derivative is its fourth moment. Summing the component estimates proves
\eqref{eq:smooth-fourth-moment}.

For the second assertion, \eqref{eq:small-bv-variance} bounds the second
moment of every interval sum of length $\ell$ by
$C\ell\delta(1+|\log\delta|)$. Apply the dyadic proof of
Lemma~\ref{lem:dyadic-collision-maximal} with this constant. The final
assertion follows by domination of the initial density; it does not
assert stationarity under that density.
\end{proof}

For $A_0<\infty$, let $X_{n,R}$ be the polygonal interpolation on
$[0,A_0]$ with values
\[
 X_{n,R}(j/n)=n^{-1/2}S_{j,R}h_R.
\]
Let $\mathcal W_C$ denote Brownian motion starting at zero with covariance
$\mathbb E[\mathcal W_C(s)\mathcal W_C(t)^{\mathsf T}]=(s\wedge t)C$;
$C$ is allowed to be singular. We use the bounded-Lipschitz distance on
probability laws on continuous path space, with its supremum norm.

\begin{theorem}[Uniform collision functional central limit]
\label{thm:collision-functional}
Fix $A_0,B_0<\infty$. Uniformly for $R\in I$ and all nonnegative
probability densities $\sigma_R$ satisfying
$\|\sigma_R\|_\infty+\|\sigma_R\|_{\mathrm{BV}}\le B_0$,
\begin{equation}\label{eq:collision-functional}
 d_{\rm BL}\bigl(\operatorname{Law}_{\sigma_R\nu}(X_{n,R}),
                     \operatorname{Law}(\mathcal W_{\Gamma_R})\bigr)
       \longrightarrow0.
\end{equation}
In particular this holds under $\nu$ and under $\nu_R^*$.
\end{theorem}
\begin{proof}
Set $\delta_n=\tfrac14 n^{-1/32}$ and let $X_{n,R}^{\delta_n}$ be the
polygonal process with $h_R$ replaced by $h_{R,\delta_n}$.
Lemma~\ref{lem:smooth-fourth-moment} implies
\[
 \sup_{R,\sigma_R}\int\|X_{n,R}-X_{n,R}^{\delta_n}\|_\infty^2
             \sigma_R\dd\nu
       \le C_{A_0,B_0}\delta_n[\log(2+n)]^3\longrightarrow0.
\]
Integer parts and interpolation change only the constant. Thus it
suffices to establish the limit for the smoothed process. This scale
is slower than the one optimizing the one-time estimate.

We first prove finite-dimensional convergence. Fix
$0=t_0<t_1<\cdots<t_q\le A_0$ and vectors $v_1,\ldots,v_q$.
Put $m_j=\lfloor nt_j\rfloor-\lfloor nt_{j-1}\rfloor$ and
$z_j=v_j/\sqrt n$. Smooth the initial density at the same scale;
its $L^1$ error is $O(\delta_n)$ and its distribution norm is
$O(\delta_n^{-2})$. The characteristic function of the consecutive
increments, before the harmless polygonal endpoint corrections, is
\[
 \ell_j\bigl(\mathcal L_{R,\delta_n,z_q}^{m_q}\cdots
       \mathcal L_{R,\delta_n,z_1}^{m_1}
                   (\mathsf S_{\delta_n}\sigma_R)\nu\bigr).
\]
The product is chronological, with the earliest segment on the right.
For each fixed $q$, substitute the spectral decompositions of
Lemma~\ref{lem:smooth-collision-expansion}. Every term with a
complementary power tends to zero uniformly, since $m_j$ is a positive
multiple of $n$ and its remaining factors are polynomially bounded in
$\delta_n^{-1}$. The product of principal amplitudes tends to one:
$\Pi(0)=\nu\ell_j$, and its error is
$O_q(\delta_n^{-4}n^{-1/2})$. The cubic error in the logarithm is
$O_q(\delta_n^{-6}n^{-1/2})=o(1)$. The resulting characteristic function
is therefore, uniformly in $R$ and in the initial densities,
\[
 \exp\left(-\frac12\sum_{j=1}^q(t_j-t_{j-1})
            v_j^{\mathsf T}\Gamma_Rv_j\right)+o(1).
\]
Here \eqref{eq:collision-covariance-limit} removes smoothing from the
covariance. This is the joint law of the independent Brownian
increments, not merely a collection of marginal Gaussian laws.

For tightness put $L_n=\lceil\delta_n^{-8}\rceil$ and interpolate the
smoothed sums only at collision indices $0,L_n,2L_n,\ldots$.
Call this coarse polygonal process $Z_{n,R}$. An interval of $m\ge L_n$
collisions has fourth moment at most $Cm^2$, by
\eqref{eq:smooth-fourth-moment} and stationarity of $\nu$. Domination
by $\sigma_R\le B_0$ gives the same bound from every deterministic
starting index under the initial law. It follows at coarse grid points
that the fourth moment of an increment is at most $C|t-s|^2$.
For arbitrary points, split an increment into its two fractional end
segments and the complete grid interval. Within one grid interval,
linear interpolation gives a factor $(|t-s|n/L_n)^4$ multiplying a
fourth moment $O((L_n/n)^2)$, which is bounded by $C|t-s|^2$.
The same estimate follows across adjacent intervals and then across
any number of intervals, with a fixed constant. Hence
\[
 \mathbb E_{\sigma_R\nu}|Z_{n,R}(t)-Z_{n,R}(s)|^4
       \le C_{A_0,B_0}|t-s|^2.
\]
The usual dyadic proof of the continuity and tightness criterion applies
uniformly to these continuous processes. Extend the last grid interval
past $A_0$ and restrict back to $[0,A_0]$ to avoid an endpoint convention.
Finally, boundedness of $h_{R,\delta_n}$ gives deterministically
\[
 \|Z_{n,R}-X_{n,R}^{\delta_n}\|_\infty
       \le C L_n/\sqrt n=O(n^{-1/4})\longrightarrow0.
\]
This proves tightness of the original processes as well.

To obtain the stated uniform distance, suppose it failed and select a
violating sequence of radii and initial densities. By compactness a
subsequence of the radii converges to $R_0$. The uniform tightness and
finite-dimensional limits identify every subsequential path law as
$\mathcal W_{\Gamma_{R_0}}$. Continuity of $\Gamma_R$, including at
singular matrices, gives continuity of the Brownian law; for example
couple by $\Gamma_R^{1/2}$ times one standard Brownian motion. This
contradicts the violating bounded-Lipschitz distance and proves
\eqref{eq:collision-functional}.
\end{proof}

\subsection{The induced record process}
Let $Y_{n,R}$ be the polygonal interpolation on $[0,1]$ of
$n^{-1/2}(J_{j,R}-j\bar G_R)$ at times $j/n$.

\begin{theorem}[Functional limit of the actual return process]
\label{thm:actual-return-functional}
Uniformly for $R\in I$,
\begin{equation}\label{eq:actual-return-functional}
 d_{\rm BL}\bigl(\operatorname{Law}_{\nu_R^*}(Y_{n,R}),
                         \operatorname{Law}(\mathcal W_{D_R})\bigr)
       \longrightarrow0,
       \qquad D_R=c_*^{-1}\Gamma_R.
\end{equation}
\end{theorem}
\begin{proof}
The count coordinate of Theorem~\ref{thm:uniform-return-fluid} gives
\[
 \sup_{0\le s\le1}|N_{\lfloor ns\rfloor,R}/n-s/c_*|
                  \longrightarrow0
\]
in probability, uniformly in $R$. Choose $A_0>c_*^{-1}+1$ and apply
Theorem~\ref{thm:collision-functional} on $[0,A_0]$ under $\nu_R^*$.
Tightness with continuous limits gives a uniform probabilistic modulus
of continuity for $X_{n,R}$. Consequently replacing its random argument
$N_{\lfloor ns\rfloor,R}/n$ by $s/c_*$ changes it by $o_P(1)$ in
supremum norm. The random argument remains in $[0,A_0]$ with probability
tending to one uniformly; independence from $X_{n,R}$ is not required.
The exact identity \eqref{eq:exact-stopped-compensation} identifies this
composition with the step version of the centered induced process.

For its polygonal version, stationarity under $F_R^*$ and the true
one-block tail give, for every $\epsilon>0$,
\[
 \nu_R^*\left\{\max_{j\le n}r_R^*\circ(F_R^*)^j>
                  \epsilon\sqrt n\right\}
       \le C(n+1)e^{-c\epsilon\sqrt n}\longrightarrow0.
\]
Each centered return increment has norm at most $C r_R^*$, by
boundedness of $h_R$ and exact compensation. Thus interpolation changes
the process by $o_P(1)$ uniformly. Brownian time change $s\mapsto s/c_*$
has covariance $c_*^{-1}\Gamma_R$, which proves the assertion by the
same compactness argument as above.
\end{proof}

\subsection{An unconditioned physical-time consequence}
Recall the actual stationary physical collision count $C_R(t)$ and
lattice displacement $W_R(t)$ of
Theorem~\ref{thm:uniform-physical-clock}. Define the linear map
\[
 B_R(x_1,x_2,x_3,x_4)=(x_1,x_2,x_3-x_4/\bar\tau_R),\qquad
 \mathcal V_R=\bar\tau_R^{-1}B_R\Gamma_R B_R^{\mathsf T}.
\]
\begin{corollary}[Physical displacement and collision fluctuations]
\label{cor:physical-functional-gaussian}
Under the stationary laws $\mu_R$, the processes
\begin{equation}\label{eq:physical-functional-gaussian}
 t^{-1/2}\bigl(W_R(ts),C_R(ts)-ts/\bar\tau_R\bigr),\quad 0\le s\le1,
\end{equation}
converge uniformly in $R$ to Brownian motion with covariance $\mathcal V_R$.
The convergence may be taken in the Skorokhod space, or for continuous
interpolations whose uniform difference from the displayed processes
is $O(t^{-1/2})$. The matrices $\mathcal V_R$ are continuous and positive
semidefinite. No conditioning event is changed or imposed in this result.
\end{corollary}
\begin{proof}
In the physical collision suspension, the marginal of the last outgoing
collision state $x$ is $\sigma_R\nu$, where
$\sigma_R=\tau_R/\bar\tau_R$. Its supremum and variation norms are
uniformly bounded by Lemma~\ref{lem:physical-bv} and
$\min_I\bar\tau_R>0$. The elapsed time since that collision is bounded
by the uniform single-flight horizon. Thus
Theorem~\ref{thm:collision-functional} applies with this initial density.

The section compensation disappears under $B_R$:
\[
 B_Rh_R=(\kappa_{{\rm coll},1},\kappa_{{\rm coll},2},
                    1-\tau_R/\bar\tau_R),
 \qquad B_R\bar G_R=0.
\]
The collision clock satisfies
$\sup_{s\le1}|C_R(ts)/t-s/\bar\tau_R|=o_P(1)$ uniformly, by
Theorem~\ref{thm:uniform-physical-clock}. Apply the continuous-limit
random-time argument to $B_RX_{\lfloor t\rfloor,R}$ on a fixed interval
larger than $1/\min_I\bar\tau_R+1$. Rounding $t$ affects neither the
normalization nor the limit.

The initial and final partial physical flights have uniformly bounded
durations. The collision sum of the lattice labels differs from
$W_R(ts)$ by a uniformly bounded cell-position term. Also the sum of
completed flight times differs from $ts$ by a uniformly bounded amount.
Consequently the stopped sum of $B_Rh_R$ differs from the numerator in
\eqref{eq:physical-functional-gaussian} by a uniformly bounded vector,
simultaneously for all $s$. The jumps of the lattice and collision counts
are uniformly bounded; interpolation between collision times, together
with the same cell-position bound, has error $O(t^{-1/2})$. The limiting
Brownian time change is $s/\bar\tau_R$, giving $\mathcal V_R$. Continuity follows
from that of $\Gamma_R$ and the explicit strictly positive mean roof.
\end{proof}

The functional results supply the unconditioned process input for the
original clock analysis. They neither evaluate a measurable transfer
function at the selected periods nor control a shrinking conditioning
event. Uniform positive definiteness is supplied separately by
Corollary~\ref{cor:actual-uniform-ellipticity}; the weighted raw density
estimates remain distinct from these path-space conclusions.

\section{One-step phase coercivity and the operator criterion}
\label{sec:operator-bridge}
The positive-measure estimate in Section~\ref{sec:periodic-coercivity}
uses an entire periodic tube. There is a stronger estimate if one first
locates a large one-step defect on the period, and only then thickens
that single state. It avoids the exponentially shrinking tube of the
whole period. We retain the constants $\rho_0,L,B_0$ of
Lemma~\ref{lem:uniform-period-tubes}, and set $r_0=\rho_0/(2L)$.

\begin{theorem}[One-step thickening of the periodic defect]
\label{thm:one-step-coercivity}
Fix $0<\alpha\le1$ and $1\le p<\infty$. Let $q$ be circle-valued and
measurable, with an $\alpha$-H\"older representative of constant $H\ge1$
on the $\rho_0$ balls around the states of the triple selected in
Theorem~\ref{thm:uniform-phase-separation}. For $b=|\beta|\ge1$ put
\[
 d_b=\frac1{4000000\pi M(b)},\qquad
 H_b=H(1+L^\alpha),\qquad
 r_b=\min\{r_0,d_b/(4bB_0),(d_b/(4H_b))^{1/\alpha}\}.
\]
For the actual one-step defect $\mathcal D_q$ defined in
Theorem~\ref{thm:positive-measure-coercivity},
\begin{equation}\label{eq:one-step-coercivity}
 \|\mathcal D_q\|_{L^p(\nu_R^*)}
       \ge\frac{d_b}{2}\left(\frac{r_b^2}{4c_*}\right)^{1/p}.
\end{equation}
If $H\le H_0b^\zeta$ with fixed $H_0\ge1$ and $\zeta\ge0$, this gives
\begin{equation}\label{eq:sharper-phase-power}
 \|\mathcal D_q\|_{L^p(\nu_R^*)}
 \ge C\,b^{-\lambda}(1+\log b)^{-\eta},\qquad
 \lambda=\frac2p\max\{1,\zeta/\alpha\},\quad
 \eta=1+\frac2{\alpha p}.
\end{equation}
The constants are uniform in the radius and all torus frequencies.
\end{theorem}
\begin{proof}
Telescoping the three exact periods as in
Corollary~\ref{cor:logarithmic-coercivity} shows that at some state $x_0$
of the selected triple the one-step defect has modulus at least $d_b$.
This is initially a statement about the specified continuous
representatives at finitely many states: the sum of the three induced
periods is at most $2M(b)$, while their product discrepancy is at least
$2/(4000000\pi)$.

On $B(x_0,r_0)$ the return has one fixed lattice label and one fixed
physical count, and $F_R^*$ maps that ball into the $\rho_0$ ball about
the successor of $x_0$. For $x$ in that ball,
\[
 |\mathcal D_q(x)-\mathcal D_q(x_0)|
 \le H_b|x-x_0|^\alpha+bB_0|x-x_0|.
\]
Each term on the right is at most $d_b/4$ on the ball of radius
$r_b$. On that ball the defect
is consequently at least $d_b/2$. Its actual invariant measure is
$r_b^2/(4c_*)$. Integration proves \eqref{eq:one-step-coercivity}.
Changes of representatives on null sets do not alter the integral.

Finally $M(b)\le3+\log b/\log(100/49)$ and
$H_b\le C b^\zeta$. The minimum defining $r_b$ is bounded
below by a positive constant times
$b^{-\max\{1,\zeta/\alpha\}}(1+\log b)^{-1/\alpha}$.
Substitution proves \eqref{eq:sharper-phase-power}. Unlike the
whole-period tube estimate, its exponent contains no $\log L$ term.
The regularity premise on $q$ remains essential.
\end{proof}

The next theorem specifies exactly what is required to turn this
physical coercivity into an operator resolvent. The hypothesis is a
reconstruction of a phase from an approximate spectral vector; it is
not silently identified with an eigenvalue equation.

\begin{theorem}[A quantitative phase-reconstruction criterion]
\label{thm:phase-resolvent-criterion}
For each $R,u,s,\beta$, let $\mathcal L_{R,u,s,\beta}$ be a bounded
operator on a Banach space $\mathcal B_R$. Fix the constants
$\alpha,p,H_0,\zeta$ of Theorem~\ref{thm:one-step-coercivity}, and let
$A\ge1$, $v\ge0$, $\theta>0$. Suppose that for every $b=|\beta|\ge1$,
$|z|=1$, and $h\in\mathcal B_R$ with $\|h\|=1$ and
$e=\|(z-\mathcal L_{R,u,s,\beta})h\|\le1$, one can construct a
circle-valued $q$ with the local H\"older bound $H_0b^\zeta$ and
\begin{equation}\label{eq:phase-reconstruction}
 \|\mathcal D_q\|_{L^p(\nu_R^*)}\le A b^v e^\theta,
 \qquad e^{ic}=z.
\end{equation}
Suppose also that $z-\mathcal L_{R,u,s,\beta}$ is Fredholm of index
zero. Then it is invertible and
\begin{equation}\label{eq:phase-resolvent}
 \|(z-\mathcal L_{R,u,s,\beta})^{-1}\|
 \le C b^{(\lambda+v)/\theta}
                (1+\log b)^{\eta/\theta},
\end{equation}
with $\lambda,\eta$ as in \eqref{eq:sharper-phase-power}.
\end{theorem}
\begin{proof}
Combining \eqref{eq:sharper-phase-power} and
\eqref{eq:phase-reconstruction} yields
\[
 e\ge (C_0/A)^{1/\theta}
       b^{-(\lambda+v)/\theta}(1+\log b)^{-\eta/\theta}
\]
when $e\le1$. Decreasing the positive constant to at most one makes
the same lower bound valid when $e>1$. Homogeneity gives a uniform
lower bound for $z-\mathcal L$ on every vector. It is therefore
injective and has closed range. Fredholm index zero makes its cokernel
zero as well, so it is surjective. The lower bound gives
\eqref{eq:phase-resolvent}. The Fredholm assumption cannot be omitted:
an injective operator bounded below need not be onto.
\end{proof}

For the unbounded induced twists, the phase reconstruction
\eqref{eq:phase-reconstruction}, including nonvanishing modulus and
quantitative regularity, is an unverified premise, not a result of
Lemma~\ref{lem:collision-spectral-input}. The latter treats a smooth
collision-space weight near zero. It provides the return-tail theorem
but is not a construction of $\mathcal B_R$ or of the raw high-frequency
operator. The criterion separates the necessary analytical construction
from the now explicit phase lower bound.

\paragraph{Scope of the retained edge calculation.}
The next section uses the original section $Y_R$ in \eqref{eq:Y}, not
$Y_R^*$. Its pointwise physical critical-word calculations are retained,
but an induced count or normalization is not transferred between the
sections without a new return computation.
\section{The exact raw density at a minimal flight length}
Let $L_N(x)=S_N\tau_R(x)$ on the collision section, and put
\[
 g=1-2R,\qquad \zeta_R=\operatorname{arcosh}(1+g/R)>0.
\]
Near the orbit alternating normally between centers $0$ and $a$, choose the initial arc coordinate $y$ in the upward tangent direction and the outgoing tangential velocity $p$. Both vanish on the orbit. Use the same upward signed arc coordinate $y_j$ at its successive endpoints.

\begin{lemma}[Reduced action Hessian]\label{lem:hessian}
The second variation of the length with all contact coordinates free is
\begin{equation}\label{eq:quadratic}
 \frac12\sum_{j=0}^{N-1}
 \left\{\frac{y_j^2+y_{j+1}^2}{R}
             +\frac{(y_{j+1}-y_j)^2}{g}\right\}.
\end{equation}
After solving the internal reflection equations, the endpoint Hessian is
\begin{equation}\label{eq:schur}
 S_N=\begin{pmatrix}A_N&-B_N\\-B_N&A_N\end{pmatrix},\quad
 A_N=\frac{\sinh\zeta_R}{g}\coth(N\zeta_R),\quad
 B_N=\frac{\sinh\zeta_R}{g\sinh(N\zeta_R)}.
\end{equation}
In the physical initial coordinates $(y,p)$, $L_N$ has a nondegenerate minimum $Ng$ and Hessian $H_N$ satisfying
\begin{equation}\label{eq:detH}
             \sqrt{\det H_N}=\sinh(N\zeta_R).
\end{equation}
\end{lemma}
\begin{proof}
At successive facing circles the longitudinal endpoint corrections are $y_j^2/(2R)$ and $y_{j+1}^2/(2R)$. Expanding their distance gives \eqref{eq:quadratic}. Its full Hessian is positive definite, since it is a sum of positive square terms including endpoint squares. The interior stationarity equations are
\[
 y_{j+1}-2(1+g/R)y_j+y_{j-1}=0.
\]
Their solution with given endpoints is
\[
 y_j=\frac{\sinh((N-j)\zeta_R)}{\sinh(N\zeta_R)}y_0
       +\frac{\sinh(j\zeta_R)}{\sinh(N\zeta_R)}y_N.
\]
Substitution into the endpoint gradients yields \eqref{eq:schur}. The nonlinear interior equations are uniquely solvable near zero by their positive definite Hessian. Their solution is the physical reflected path by the first variation and the strict incidence at the normal orbit.

The endpoint first variation gives $p=-\partial_{y_0}L_N$ at the linearized level, so $p=-A_Ny_0+B_Ny_N$. The linear map from $(y_0,p)$ to $(y_0,y_N)$ has determinant $1/B_N$. Thus
\[
 \det H_N=\frac{A_N^2-B_N^2}{B_N^2}
                         =\sinh^2(N\zeta_R).
\]
Positivity is preserved under this nonsingular coordinate change, proving the minimum assertion.
\end{proof}

\begin{theorem}[A physical edge coefficient]\label{thm:edge}
Let $\chi$ be a smooth nonnegative section cutoff supported in the regular alternating itinerary, equal to one near its normal initial point. The pushforward of $\chi\nu$ by $L_N$ has, near $Ng$, the form
\begin{equation}\label{eq:edge}
 f_{N,R}(Ng+s)=\one_{\{s\ge0\}}a_{N,R}(s),\qquad
 a_{N,R}(0)=J_{N,R}:=\frac1{2R\sinh(N\zeta_R)}.
\end{equation}
The function $a_{N,R}$ is smooth on a right neighborhood of zero. For each fixed $N$, the construction and smooth bounds may be made uniform on $I$. No assertion of uniform derivative bounds as $N\to\infty$ is made.
\end{theorem}
\begin{proof}
The section density in initial arc and tangential velocity coordinates is $(4\pi R)^{-1}\dd y\dd p$. The Morse lemma applied to Lemma~\ref{lem:hessian} gives $L_N=Ng+|z|^2/2$. The transformed density $w(z)$ is smooth with
\[
 w(0)=\frac1{4\pi R\sqrt{\det H_N}}.
\]
In polar coordinates its pushforward density is
$\one_{\{s\ge0\}}\int_0^{2\pi}w(\sqrt{2s}\,n_\phi)\dd\phi$ near zero. The angular integral of each odd Taylor term vanishes; Taylor's theorem at arbitrary order gives smoothness in $s\ge0$. Its right value is $2\pi w(0)$, which is \eqref{eq:edge}. Compactness in $R$ and the strictly positive determinant give the fixed-$N$ uniform statement.
\end{proof}

\begin{proposition}[Regular critical words and a uniform edge bound]
\label{prop:general-edge}
Fix a regular physical word with $N$ flights between circular obstacles. Every critical point of its total flight length in the two initial section coordinates is normal at both endpoints. There is at most one such point for the fixed sequence of lattice centers. It is a nondegenerate minimum. If $S=\left(\begin{smallmatrix}A&C\\C&D\end{smallmatrix}\right)$ is its reduced endpoint action Hessian, its section-density jump is
\begin{equation}\label{eq:general-edge}
 J=\frac{|C|}{2R\sqrt{AD-C^2}}.
\end{equation}
There is a fixed explicitly computable constant $C_*$ such that, for every regular word and $N\ge2$,
\begin{equation}\label{eq:edge-decay}
                  0<J\le C_*(47/53)^{N-2}.
\end{equation}
No lower incidence margin for the interior collisions is needed for this coefficient bound. This is an individual-word estimate, not a bound for the sum over words.
\end{proposition}
\begin{proof}
Let $q_j(y_j)$ be arc coordinates and $v_j$ the unit velocity of segment $j$. Its second length variation is
\[
 \frac{c_j}{R}y_j^2+\frac{c_{j+1}}{R}y_{j+1}^2+
 \frac{\{v_j^\perp\cdot(t_{j+1}y_{j+1}-t_jy_j)\}^2}{\ell_j},
\]
where $c_j,c_{j+1}>0$ are its outgoing and incoming incidence cosines. At an interior reflected contact the two cosines agree. Summing gives a positive definite tridiagonal matrix with nonzero off-diagonal entries. Its interior Schur complement $S$ is positive definite and $C\ne0$: eliminating the successive interior variables expresses $C$ as the nonzero product of neighboring entries divided by positive pivots.

The endpoint first variation of the reduced action is $-p_0\dd y_0+p_N\dd y_N$. Moreover $\partial_{p_0}y_N=-1/C\ne0$. Thus the derivative of the physical length in $p_0$ vanishes exactly when $p_N=0$, and its other derivative then vanishes exactly when $p_0=0$. At such a critical point the coordinate change gives $\det H=(AD-C^2)/C^2>0$. The same Morse calculation as in Theorem~\ref{thm:edge} proves \eqref{eq:general-edge}.

There cannot be two such regular points for one word. Regard the polygonal length as a convex function of all contact positions in the product of the closed disks. At a normal endpoint its gradient is $-n$, and at an interior reflected contact it is $-2c_j n_j$. These are strictly inward normals. For a different point $x_j$ of the same strictly convex disk, $n_j\cdot(x_j-q_j)<0$. The convex gradient inequality therefore makes every different tuple have strictly larger action. Hence the normal-to-normal critical tuple is the unique global minimizer of that convex problem. Occlusion by a disk not in the word is excluded separately by physical admissibility, not by this minimization argument.

To prove the uniform bound, set $z_j=c_jy_j$ and choose tangent orientations successively so that the off-diagonal entries become negative. At the normal endpoints $c_0=c_N=1$. Write $t_j=1/\ell_j\le t_*=50/3$. The interior matrix has off-diagonals $-t_j$ and diagonals
\[
         d_j=t_{j-1}+t_j+2/(Rc_j)\ge t_{j-1}+t_j+v_*,
             \qquad v_*=200/47.
\]
Its factorization $D_0(I-K)$ has $K\ge0$ with row sums at most
$q=2t_*/(2t_*+v_*)=47/53$. In the Neumann series, an entry connecting interior indices $1$ and $N-1$ vanishes before power $N-2$. Therefore
\[
 |C|\le t_*^2\frac{q^{N-2}}{v_*(1-q)}.
\]
The reduced quadratic form dominates the two endpoint curvature terms, so $S\ge (100/47)I$. Inserting these estimates in \eqref{eq:general-edge}, and using $R\ge9/20$, gives \eqref{eq:edge-decay} with
\[
 C_* = \frac{47}{90}\frac{t_*^2}{v_*(1-q)}.
\]
Arbitrarily small positive interior incidence only increases the diagonal potential, so it does not weaken the estimate.
\end{proof}

The proposition classifies regular critical contributions and bounds each jump uniformly. Singular itinerary boundaries and the number and arrangement of critical values remain to be controlled. Multiplying \eqref{eq:edge-decay} by an uncontrolled word count would not prove a summable global correction.

\begin{corollary}[Raw finite-record absolute continuity]\label{cor:raw-density}
For every $R\in I$ and $n\ge1$, the actual induced record
$(S_n\kappa_R,S_n r_R,S_n\tau_R^Y)$ under $\nu_R^Y$ has a density with respect to counting measure on $\Z^3$ times Lebesgue measure. For fixed total physical count $N$ and displacement $k$, it has the coarea representation, for almost every $t$,
\begin{equation}\label{eq:raw-coarea}
 p_{n,R}(k,N,t)=\frac1{4\pi\nu(Y_R)}
 \sum_{\mathbf d}\int_{D_{n,\mathbf d}\cap\{L_N=t\}}
              \frac{\dd\mathcal H^1}{|\nabla_{\alpha,p}L_N|}.
\end{equation}
Here $D_{n,\mathbf d}$ is the actual first-return event for the ordered center word $\mathbf d$, total count $N$ and displacement $k$. Bounded measurable insertions have the same formula with the insertion in the numerator. Neither a uniform density bound nor a long-time local limit is asserted by this statement.
\end{corollary}
\begin{proof}
Partition the recurrent regular initial states by total collision count and ordered center word. Uniform finite horizon makes the available lattice steps finite; for fixed $N$ there are finitely many words. On every regular branch the physical length is smooth, and Proposition~\ref{prop:general-edge} makes its critical set empty or a singleton. Away from that null set, the coarea formula applied on a countable exhaustion where $|\nabla L_N|$ is bounded below proves absolute continuity and \eqref{eq:raw-coarea}. Restriction to the measurable first-return event does not alter absolute continuity of the initial measure. The singular initial states have zero flux measure, and summing over the countable possible $N$ gives the full law. The same exhaustion with a bounded insertion proves its signed version.
\end{proof}

For the induced system \eqref{eq:Y}, fix $n$ and restrict its joint law to $\kappa=0$, $r=2n$. Every physical flight is at least $g$. Equality of total length with $2ng$ requires every flight to have length $g$, and hence a nearest-neighbor normal alternating orbit. Of the six possible normal starting states only $(0,0)$ lies in $Y_R$. A neighborhood of this point makes exactly $n$ returns of physical length two each. Compactness, or directly the equality cases of the distance bound, excludes other itineraries from a sufficiently small right neighborhood of $2ng$. Consequently this lattice component of the induced law has right density
\begin{equation}\label{eq:inducededge}
                \frac{J_{2n,R}}{\nu(Y_R)}>0
\end{equation}
at its left endpoint, and zero density on its left.

\begin{corollary}[The boundary term must be retained]\label{cor:nonLone}
For the specified inducing set, the raw joint characteristic function $\Phi_{n,R}$ is not in $L^1(\T^3\times\R)$, for any $n\ge1$. In particular a global integrable bound for this raw function cannot be justified merely by making the normal itinerary exponentially rare.
\end{corollary}
\begin{proof}
If it were integrable, taking its $(0,0,2n)$ lattice Fourier coefficient would give an integrable function of the roof frequency. Its inverse would be a continuous density of that component. A continuous representative cannot agree almost everywhere with zero to the left of $2ng$ and with a function having the positive right limit \eqref{eq:inducededge}. This contradicts Theorem~\ref{thm:edge}.
\end{proof}

The central local-limit problem is not decided by this corollary. The endpoint may be far outside the central window, and $J_{2n,R}$ decays exponentially. What changes is the global inversion method, not the raw observable or the intended central theorem.

\begin{proposition}[Exact edge subtraction]\label{prop:subtract}
Choose a smooth roof cutoff supported in the right neighborhood of Theorem~\ref{thm:edge}, equal to one near its endpoint, and denote the resulting subdensity by $f$. With $J=J_{N,R}$, one has the exact identity
\begin{equation}\label{eq:subtract}
 f(Ng+s)=J e^{-s}\one_{\{s\ge0\}}+q(s),\qquad
 \widehat f(b)=e^{ibNg}\left(\frac{J}{1-ib}+\widehat q(b)\right),
\end{equation}
where $q\in L^1(\R)$, its distributional second derivative is a finite signed measure, and
$|\widehat q(b)|\le C_{N,R}(1+b^2)^{-1}$.
\end{proposition}
\begin{proof}
Writing $f(Ng+s)=\one_{s\ge0}\eta(s)a(s)$, the right value of $\eta a-Je^{-s}$ at zero is zero. Thus $q$ is continuous across zero, piecewise smooth, and exponentially decaying at infinity. Its first derivative may jump at zero, so its second distributional derivative has one finite atom and an integrable density. Distributional integration by parts gives $b^2|\widehat q(b)|\le\|D^2q\|_{\TV}$; the $L^1$ norm gives the bound near zero. Direct integration of the exponential proves \eqref{eq:subtract}.
\end{proof}
This is subtraction and exact addition of an explicit density, not convolution with noise or removal of a physical event. It proves an integrable remainder for the isolated edge. Bounds for all remaining branches and edges are still needed for a long-time result.


\section{A raw residual-inversion theorem}

\paragraph{Frequency convention.}
The rescaled frequency is $v=\sqrt n\,\omega$. When the proved growing
central band is used below, its radii are $2n^{1/200}$ in $v$ and
$2n^{-99/200}$ in $\omega$. General cutoffs in the following statements
retain their stated hypotheses; a Gaussian-tail estimate is not a bound
for the complementary transform of the raw measure.

Let $\eta_{n,R}$ be finite measures on $\Z^3\times\R$ with characteristic functions $\Phi_{n,R}$. Frequencies are $\omega=(u,s,b)\in[-\pi,\pi]^3\times\R$, and our transform convention is $e^{i\omega\cdot z}$. Suppose there are explicit signed edge measures $E_{n,R}$ with densities $e_{n,R}$, and put $H_{n,R}=\Phi_{n,R}-\widehat E_{n,R}$. The following hypotheses are an interface, not consequences of the mechanical constructions above.

\begin{theorem}[Residual criterion for the mixed density LLT]\label{thm:LLT}
Assume the following uniformly for $R\in I$.

\smallskip\noindent
(i) There are $c,C,\delta>0$, $\rho<1$, means $m_R\in\R^4$, and symmetric matrices $cI\le\Sigma_R\le CI$ such that, for $|\omega|\le\delta$,
\[
 \Phi_{n,R}(\omega)=e^{inm_R\cdot\omega}
 \{A_R(\omega)\lambda_R(\omega)^n+O(\rho^n)\},
\]
where $A_R(0)=1$, $|A_R(\omega)-1|\le C|\omega|$,
$\log\lambda_R(\omega)=-\tfrac12\omega^T\Sigma_R\omega+O(|\omega|^3)$, and $|\lambda_R(\omega)|\le e^{-c|\omega|^2}$.

\smallskip\noindent
(ii) $H_{n,R}\in L^1(\T^3\times\R)$ and
\begin{equation}\label{eq:tailcriterion}
 \int_{|\omega|>n^{-2/5}}|H_{n,R}(\omega)|\dd\omega=o(n^{-2}).
\end{equation}

\smallskip\noindent
(iii) $\|E_{n,R}\|_{\TV}=o(n^{-2})$ and, for each fixed $K$,
\[
 n^2\mathop{\rm ess\,sup}_{|(z-nm_R)/\sqrt n|\le K}|e_{n,R}(z)|\longrightarrow0.
\]
Then $\eta_{n,R}$ has a density with respect to counting measure times Lebesgue measure, and on each such central window
\begin{equation}\label{eq:LLT}
 n^2 p_{n,R}(k,m,t)=
 \frac{\exp(-\xi^T\Sigma_R^{-1}\xi/2)}{(2\pi)^2\sqrt{\det\Sigma_R}}+o(1),
 \quad \xi=\frac{(k_1,k_2,m,t)-nm_R}{\sqrt n},
\end{equation}
with uniform essential-supremum error. No smoothing of $\eta_{n,R}$ is used.
\end{theorem}
\begin{proof}
Fourier inversion of $H_{n,R}$ gives a continuous density of the signed measure $\eta_{n,R}-E_{n,R}$. To justify this without assuming a density, take each lattice Fourier coefficient of $H$. It is an integrable roof transform of the corresponding signed component. Convolve that component with a Gaussian approximate identity, use integrability to pass to the inverse-transform limit, and test against compactly supported continuous functions. Uniqueness of finite Fourier transforms identifies each component, hence the full measure. This is a proof of inversion, not a change to the datum. Adding the known density of $E_{n,R}$ gives $p_{n,R}$.

On $|\omega|\le n^{-2/5}$, replacing $H$ by $\Phi$ costs at most a fixed volume times $\|E\|_{\TV}=o(n^{-2})$. Inserting (i), let $v=\sqrt n\,\omega$. The cubic error is uniformly $O(n^{-1/5})$ on this region; the bound on $\lambda$ supplies an integrable Gaussian dominator. Uniform Taylor bounds, truncation first to $|v|\le M$, and then $M\to\infty$, give the inverse Gaussian with error $o(n^{-2})$. The inverse normalization is $(2\pi)^{-4}$ and the Gaussian integral is $(2\pi)^2/\sqrt{\det\Sigma_R}$, giving the constant in \eqref{eq:LLT}. Hypothesis (ii) bounds the remaining frequencies, and (iii) bounds the added density in the central window. All estimates are uniform in $\xi$ in that window.
\end{proof}

For weighted inserted measures, the same proof applies with $A_R(0)=a_R$ and uniformly bounded amplitudes, provided the edge and tail hypotheses are checked for those measures as well. It gives $a_R$ times the Gaussian. It does not assert uniformity for arbitrary $n$-dependent path functionals without such estimates.

\begin{lemma}[An explicit frequency splice]\label{lem:splice}
Suppose a residual transform has a compact minor-arc bound $Ce^{-cn}$ and, for $1\le|b|\le e^{\kappa n}$, a bound $C(1+|b|)^A e^{-cn}$. Beyond $e^{\kappa n}$ suppose it is bounded by
\[
 C_M\rho^n(1+|b|)^{-2}+C_Me^{c_M n}(1+|b|)^{-M}.
\]
If one fixed integer $M>1$ and one $\kappa>0$ satisfy
\begin{equation}\label{eq:splice}
       \frac{c_M}{M-1}<\kappa<\frac{c}{A+1},
\end{equation}
these regions contribute exponentially small $L^1$ mass. Together with the Gaussian major-arc bound between $n^{-2/5}$ and $\delta$, this proves \eqref{eq:tailcriterion}.
\end{lemma}
\begin{proof}
The middle integral is bounded by $C\exp(-(c-(A+1)\kappa)n)$. The two outer integrals are bounded by $C_M\exp((\log\rho-\kappa)n)$ and $C_M\exp(-( (M-1)\kappa-c_M)n)/(M-1)$. The torus has finite volume. The major-arc Gaussian tail is bounded by a polynomial factor times $e^{-c n^{1/5}}$, and the compact remainder is exponential. Strict positivity of the three exponent margins proves the claim.
\end{proof}
The inequality \eqref{eq:splice} must be verified with the actual derivative-growth constants. The phrases ``take $\kappa$ small'' and ``then take $M$ large'' do not imply it: for example, $c_M=2(M-1)$ and $c/(A+1)=1$ leave no possible $\kappa$. This checks a proposed estimate, not the truth or falsity of a central LLT. The residual formulation also requires an independently proved bound for the sum of all extracted edge terms; Proposition~\ref{prop:subtract} treats one physical family only.


\section{Local inversion with nonnegligible remote edge mass}

\paragraph{Frequency convention.}
The rescaled frequency is $v=\sqrt n\,\omega$. When the proved growing
central band is used below, its radii are $2n^{1/200}$ in $v$ and
$2n^{-99/200}$ in $\omega$. General cutoffs in the following statements
retain their stated hypotheses; a Gaussian-tail estimate is not a bound
for the complementary transform of the raw measure.

\label{sec:localized-inversion}
Theorem~\ref{thm:LLT} gives a sufficient global smallness condition on
extracted edges. A central local limit needs less: the extracted measure
must have negligible density and negligible low-frequency convolution
on the central window. Its total variation need not tend to zero.
The following refinement retains the raw datum and the density norm.

Fix an even real function $\chi\in C_c^\infty(\R^4)$ equal to one on
the unit ball and supported in the ball of radius two. Let $a_n\to\infty$
with $a_n=o(n^{1/6})$, put $B_n=a_n/\sqrt n$, and set
$\chi_n(\omega)=\chi(\omega/B_n)$. For all sufficiently large $n$ its
support fits in the fundamental torus interval in the first three
coordinates. With the transform convention already fixed, its inverse
kernel on $\Z^3\times\R$ is
\begin{equation}\label{eq:localized-kernel}
 K_n(k,t)=B_n^4\check\chi(B_nk,B_nt),\qquad
 |K_n(k,t)|\le C_M B_n^4(1+B_n|(k,t)|)^{-M}.
\end{equation}
The check denotes the inverse transform on $\R^4$, including its
$(2\pi)^{-4}$ factor. Convolution here and below uses counting measure
in the discrete coordinates and Lebesgue measure in the roof coordinate.
For fixed $K$, let $W_{n,R,K}=\{z:|z-nm_R|\le K\sqrt n\}$.

\begin{theorem}[A local edge criterion]\label{thm:local-edge-LLT}
Assume the small-frequency hypothesis (i) of Theorem~\ref{thm:LLT}.
Let $E_{n,R}$ be finite signed measures with densities $e_{n,R}$ and
suppose $H_{n,R}=\Phi_{n,R}-\widehat E_{n,R}$ belongs to
$L^1(\T^3\times\R)$. Assume, uniformly in $R$, that
\begin{equation}\label{eq:local-tail}
 \int |1-\chi_n(\omega)|\,|H_{n,R}(\omega)|\dd\omega=o(n^{-2}).
\end{equation}
For every fixed central window suppose also that
\begin{equation}\label{eq:local-edge-errors}
 n^2\mathop{\rm ess\,sup}_{W_{n,R,K}}|e_{n,R}|\longrightarrow0,
 \qquad
 n^2\sup_{W_{n,R,K}}|K_n*E_{n,R}|\longrightarrow0.
\end{equation}
Then the raw density conclusion \eqref{eq:LLT} holds on that window.
There is no small-total-variation hypothesis on $E_{n,R}$.
\end{theorem}
\begin{proof}
As in the proof of Theorem~\ref{thm:LLT}, integrability of $H$ first
identifies a continuous density for $\eta-E$, and adding $e$ gives
an actual mixed density for $\eta$. Splitting this identity with
$\chi_n$ gives the exact equality almost everywhere
\begin{equation}\label{eq:exact-local-decomposition}
 p_{n,R}=K_n*\eta_{n,R}
       +(e_{n,R}-K_n*E_{n,R})
       +\mathcal F^{-1}[(1-\chi_n)H_{n,R}].
\end{equation}
All terms are defined: $E$ is finite, $K_n$ is bounded and integrable,
and the last inverse transform is absolutely convergent.
The last term is at most $(2\pi)^{-4}$ times the integral in
\eqref{eq:local-tail}. The middle term is controlled by
\eqref{eq:local-edge-errors}.

For the first term, its Fourier integral has the factor $\chi_n$.
Insert the assumed small-frequency expansion and use
$v=\sqrt n\,\omega$. The support has $|v|\le2a_n$ and
$nO(|\omega|^3)=O(a_n^3/\sqrt n)=o(1)$. The uniform quadratic
bound on $\lambda_R$ gives a Gaussian dominator, while the amplitude
converges to one. The exponentially small spectral remainder, after
integration and multiplication by $n^2$, is bounded by
$Ca_n^4\rho^n$ and tends to zero. Truncate first at a fixed $|v|$,
then let that truncation grow. This proves uniformly in $R$ and the
central variable that $n^2K_n*\eta_{n,R}$ converges to the Gaussian
in \eqref{eq:LLT}. Combining the three terms proves the assertion.

The kernel is only an analytical device in the exact decomposition
\eqref{eq:exact-local-decomposition}. The two correction terms are
retained and estimated; the result is not a local limit for a
smoothed replacement of the observed law.
\end{proof}

\begin{corollary}[Remote mass need not be small]\label{cor:remote-edges}
Suppose $\|E_{n,R}\|_{\TV}\le C$ and, for the window under
consideration, the support of $E_{n,R}$ is at distance at least
$\varepsilon\sqrt n$ from $W_{n,R,K}$, for a fixed
$\varepsilon>0$. Then \eqref{eq:local-edge-errors} holds.
In particular an extracted density of bounded, nonvanishing mass
can satisfy the local edge hypotheses.
\end{corollary}
\begin{proof}
The density vanishes almost everywhere in the central window. For
$z$ in that window, \eqref{eq:localized-kernel} gives
\[
 n^2|(K_n*E_{n,R})(z)|
 \le CC_M a_n^4(1+\varepsilon a_n)^{-M}\longrightarrow0
\]
by taking $M>4$. The Schwartz bound is available for every fixed $M$
because $\chi$ is smooth. These constants are kernel constants, not
the derivative-growth constants of the billiard in Lemma~\ref{lem:splice}.
\end{proof}

The same conclusion holds for a sum of remote extracted terms whose
total variations are summable with bounded total, and a remaining
nearby extracted density satisfying \eqref{eq:local-edge-errors}.
This follows by linearity and the triangle inequality. Thus individual
edge coefficients do not all have to be $o(n^{-2})$ merely because
one proves a central limit. Nevertheless, one must still construct
the full extraction and prove \eqref{eq:local-tail}. Neither this
corollary nor the single-word estimate in Proposition~\ref{prop:general-edge}
controls the sum of all central critical words or singular boundaries.
A restriction of a density to a sharp window can itself create new
jump terms; it is not automatically an integrable residual.

For weighted measures the same proof gives the corresponding amplitude
times the Gaussian, provided the small-frequency expansion, residual
estimate and local edge errors are verified for that weighted class.
This is the precise additional input needed before applying
Proposition~\ref{prop:conditioning} to a proposed family of insertions.

\section{An integrated growing central band for the physical record}

\paragraph{Frequency convention.}
The rescaled frequency is $v=\sqrt n\,\omega$. When the proved growing
central band is used below, its radii are $2n^{1/200}$ in $v$ and
$2n^{-99/200}$ in $\omega$. General cutoffs in the following statements
retain their stated hypotheses; a Gaussian-tail estimate is not a bound
for the complementary transform of the raw measure.

\label{sec:growing-major-arc}
The estimate on compact rescaled frequency sets can be strengthened to
an integrated estimate on an expanding band. The distinction matters
for local inversion: convergence after integration controls the inverse
transform uniformly in the observation variable. We prove this estimate
for the actual return record, including bounded-variation initial
insertions. No induced spectral realization is used.

An admissible initial insertion is a complex-valued function $a_R$ on
$M$, vanishing almost everywhere outside $Y_R^*$, such that
\[
 \|a_R\|_\infty+\|a_R\|_{\mathrm{BV}}\le B.
\]
Its integral is denoted by $\alpha_R$. The original probability has
$a_R=s_R$ and $\alpha_R=1$. Write
\begin{equation}\label{eq:inserted-central-characteristic}
 C^a_{n,R}(v)=\int_{Y_R^*}a_R(x)
 \exp\left(\frac{i v\cdot(J_{n,R}(x)-n\bar G_R)}{\sqrt n}\right)
 \dd\nu(x),\qquad v\in\R^4.
\end{equation}
Here $v$ is the rescaled frequency; the physical Fourier frequency is
$\omega=v/\sqrt n$. Constants below depend on $B$ but not on the
insertion, the radius, or $n$. No continuity of the insertion in $R$
is required.

\subsection{A frequency-dependent error bound}
\begin{lemma}[Collision estimate with explicit frequency dependence]
\label{lem:explicit-frequency-error}
Put $\ell_\delta=1+|\log\delta|$. There are $c,C>0$ and $\rho<1$
such that, for $m\ge2$, $0<\delta<1/2$, and real $t$ satisfying
\[
 |t|\le c\sqrt m\,\delta^2,\qquad
 \delta^{-6}|t|^3m^{-1/2}\le c,
\]
every complex function $a$ with
$\|a\|_\infty+\|a\|_{\mathrm{BV}}\le B$ satisfies
\begin{align}
 &\left|\int a\exp(it\cdot S_{m,R}h_R/\sqrt m)\dd\nu
       -\left(\int a\dd\nu\right)e^{-t^{\mathsf T}\Gamma_Rt/2}\right|
       \notag\\[-2mm]
 &\quad\le C\left\{
 \delta+|t|\sqrt{\delta\ell_\delta}+|t|^2\delta\ell_\delta
 +\frac{|t|\delta^{-4}+|t|^3\delta^{-6}}{\sqrt m}
 +\delta^{-2}\rho^m\right\}.\label{eq:explicit-frequency-error}
\end{align}
\end{lemma}
\begin{proof}
Set $a_\delta=\mathsf S_\delta a$ and $z=t/\sqrt m$.
The operator in Lemma~\ref{lem:smooth-collision-expansion} gives
\[
 \int a_\delta e^{it\cdot S_{m,R}h_{R,\delta}/\sqrt m}\dd\nu
 =\lambda_{R,\delta}(z)^m
   \ell_j\bigl(\Pi_{R,\delta}(z)a_\delta\nu\bigr)
       +O(\delta^{-2}\rho^m).
\]
The embedding bound in Lemma~\ref{lem:density-norm-details} below
makes the initial-vector factor explicit. The principal amplitude
at zero is $\int a\dd\nu$, and its change is at most
$C|t|\delta^{-4}/\sqrt m$. Moreover
\[
 m\log\lambda_{R,\delta}(t/\sqrt m)
 =-\tfrac12t^{\mathsf T}\Gamma_{R,\delta}t
       +O(\delta^{-6}|t|^3/\sqrt m).
\]
The stated smallness conditions put $z$ in the analytic ball and bound
the last remainder. Positive semidefiniteness of
$\Gamma_{R,\delta}$ therefore bounds the principal power by a fixed
constant. The inequality $|e^w-1|\le |w|e^{|w|}$ proves the two
spectral error terms in \eqref{eq:explicit-frequency-error}.

Mean-preserving smoothing changes the initial insertion by $O(\delta)$
in $L^1$. If $r=h_R-h_{R,\delta}$, then
\[
 \int |a_\delta|\,|S_{m,R}r|^2\dd\nu
       \le C m\delta\ell_\delta,
\]
by \eqref{eq:small-bv-variance} and $|a_\delta|\le B$.
Cauchy--Schwarz bounds the characteristic error from removing this
smoothing by $C|t|\sqrt{\delta\ell_\delta}$. This argument uses
absolute values of the insertion and so applies to signed and complex
insertions as well. Finally, the segment joining the two covariance
matrices consists of positive semidefinite matrices. Differentiation
along that segment and \eqref{eq:collision-covariance-limit} give
\[
 \left|e^{-t^{\mathsf T}\Gamma_{R,\delta}t/2}
             -e^{-t^{\mathsf T}\Gamma_Rt/2}\right|
       \le C|t|^2\delta\ell_\delta.
\]
Combining these estimates proves the lemma.
\end{proof}

\begin{proposition}[Stopped estimate on finite frequency balls]
\label{prop:stopped-band-error}
For $U\ge1$, let
\[
 \mathcal E_n(U)=\sup_{R,a_R}
   \int_{|v|\le2U}
    |C^a_{n,R}(v)-\alpha_Re^{-v^{\mathsf T}D_Rv/2}|\dd v.
\]
If $U\le c\sqrt n\,\delta^2$ and
$\delta^{-6}U^3/\sqrt n\le c$, with $c$ sufficiently small, then
\begin{align}
 \mathcal E_n(U)\le C\bigl\{
 &\delta U^4+\sqrt{\delta\ell_\delta}\,U^5
       +\delta\ell_\delta U^6
       +n^{-1/2}\delta^{-4}U^5+n^{-1/2}\delta^{-6}U^7
       \notag\\
 &+\delta^{-2}\rho_0^nU^4
       +n^{-1/5}U^4+n^{-1/5}\log(2+n)U^5+n^{-1}U^6
                          \bigr\},\label{eq:integrated-band-error}
\end{align}
where $\rho_0<1$. The supremum is over the admissible insertions.
\end{proposition}
\begin{proof}
Put $m_n=\lfloor n/c_*\rfloor$ and $b_n=\lceil n^{3/5}\rceil$.
The exact compensation identity \eqref{eq:exact-stopped-compensation}
permits replacement of the centered record by $S_{N_{n,R},R}h_R$.
The exceptional set $|N_{n,R}-m_n|>b_n$ has $\nu_R^*$-probability
$O(n/b_n^2)$ by Lemma~\ref{lem:stopped-compensation}; rounding changes
only the constant. Its contribution with insertion $a_R$ is bounded
by the same order, since $|a_R|\nu\le Bc_*\nu_R^*$ on the section.
On its complement use the forward and backward deterministic-window
maxima from Lemma~\ref{lem:dyadic-collision-maximal}. The elementary
exponential difference bound and Cauchy--Schwarz give
\begin{align}
 &\left|C^a_{n,R}(v)
       -\int a_R e^{iv\cdot S_{m_n,R}h_R/\sqrt n}\dd\nu\right|
       \notag\\
 &\hspace{12mm}\le C\left\{n^{-1/5}
                  +|v|n^{-1/5}\log(2+n)\right\}.
                 \label{eq:weighted-stopping-error}
\end{align}
No second moment of an unbounded stopped sum on the exceptional set
is used. In particular the estimate is valid for complex insertions.

Apply Lemma~\ref{lem:explicit-frequency-error} at $m_n$ with
$t=v\sqrt{m_n/n}$. These frequencies satisfy its restrictions after
changing the uniform constant $c$. The covariance of the deterministic
sum is $(m_n/n)\Gamma_R$, which differs from $D_R=c_*^{-1}\Gamma_R$
by $O(n^{-1})$. Its Gaussian characteristic function consequently
changes by at most $C|v|^2/n$. Integrating these bounds and
\eqref{eq:weighted-stopping-error} uses only
$\int_{|v|\le2U}|v|^j\dd v=C_jU^{4+j}$, for $j\ge0$.
This proves \eqref{eq:integrated-band-error}. Finitely many small values
of $n$ can be absorbed by increasing the constants wherever the stated
smallness conditions hold.
\end{proof}

\subsection{A polynomially expanding major arc}
\begin{theorem}[Integrated major arc of the actual return record]
\label{thm:growing-major-arc}
Let
\[
 0<\varepsilon<\frac1{134},\qquad
 10\varepsilon<\theta<\frac{1/2-7\varepsilon}{6},\qquad
 r=\min\left\{\frac\theta2-5\varepsilon,
              \frac12-6\theta-7\varepsilon,
              \frac15-5\varepsilon\right\}.
\]
Then $r>0$ and
\begin{equation}\label{eq:general-growing-arc}
 \sup_{R,a_R}\int_{|v|\le2n^\varepsilon}
 |C^a_{n,R}(v)-\alpha_Re^{-v^{\mathsf T}D_Rv/2}|\dd v
       \le C n^{-r}\log(2+n).
\end{equation}
In particular, with $\varepsilon=1/200$ and $\theta=1/14$,
\begin{equation}\label{eq:explicit-growing-arc}
 \sup_{R,a_R}\int_{|v|\le2n^{1/200}}
 |C^a_{n,R}(v)-\alpha_Re^{-v^{\mathsf T}D_Rv/2}|\dd v
       \le C n^{-3/280}\sqrt{\log(2+n)}.
\end{equation}
These estimates hold for $n\ge2$. They require only the positive
semidefinite covariance already established, not positive definiteness.
\end{theorem}
\begin{proof}
Take $U=n^\varepsilon$ and $\delta=\tfrac14n^{-\theta}$ in
Proposition~\ref{prop:stopped-band-error}. Its analytic restrictions hold
for all sufficiently large $n$: in particular
$3\varepsilon+6\theta<1/2$ follows from
$7\varepsilon+6\theta<1/2$, and
$\varepsilon+2\theta<1/2$ follows as well. Constants in the radius
of the analytic ball do not affect these strict inequalities.

The three potentially slow terms in \eqref{eq:integrated-band-error}
are, respectively,
\[
 n^{-(\theta/2-5\varepsilon)}\sqrt{\log(2+n)},\qquad
 n^{-(1/2-6\theta-7\varepsilon)},\qquad
 n^{-(1/5-5\varepsilon)}\log(2+n).
\]
All other polynomial terms decay at least as fast as one of these;
this follows by subtracting their exponents, using $\theta>10\varepsilon$.
For the rounding term also note that
$1-6\varepsilon>1/5-5\varepsilon$ in the indicated range.
The exponentially decaying term is smaller than every fixed negative
power. The interval for $\theta$ is nonempty precisely when
$67\varepsilon<1/2$, proving \eqref{eq:general-growing-arc}.

At the displayed explicit choice the three margins are
\[
 \frac3{280},\qquad \frac{51}{1400},\qquad \frac7{40}.
\]
Only the first is minimal, and that term carries a square root of the
logarithm. The strict excess of the other margins absorbs their
logarithmic factors. Increasing the constant handles the finite initial
range, since both characteristic functions are bounded uniformly by
a constant depending on $B$. This proves \eqref{eq:explicit-growing-arc}.
\end{proof}

\subsection{Insertion into exact raw inversion}
Define the finite complex measure $\mu^a_{n,R}$ by pushing forward
$a_R\nu$ under $J_{n,R}$, and write $\Phi^a_{n,R}$ for its transform.
Take $a_n=n^{1/200}$ in the cutoff and kernel of
\eqref{eq:localized-kernel}, and use $m_R=\bar G_R$ in the central
windows. The preceding theorem supplies a proved replacement for the
spectral small-frequency hypothesis in the following inversion route.

\begin{corollary}[Raw inversion from the proved central band]
\label{cor:raw-from-growing-arc}
Use the uniform bound $D_R\ge d_0 I_4$ of
Theorem~\ref{thm:joint-uniform-nondegeneracy}.
For an admissible family of insertions, let $E^a_{n,R}$ be finite complex
measures with mixed densities $e^a_{n,R}$ and put
$H^a_{n,R}=\Phi^a_{n,R}-\widehat E^a_{n,R}$. Assume that
$H^a_{n,R}\in L^1(\T^3\times\R)$ and, uniformly in the family,
\begin{equation}\label{eq:new-complementary-tail}
 \int |1-\chi_n(\omega)|\,|H^a_{n,R}(\omega)|\dd\omega=o(n^{-2}).
\end{equation}
Assume also the two local edge conditions
\eqref{eq:local-edge-errors} with $e,E$ replaced by $e^a,E^a$.
Then $\mu^a_{n,R}$ has a mixed density $p^a_{n,R}$, and for each fixed
$K$,
\begin{align}
 \mathop{\rm ess\,sup}_{|z-n\bar G_R|\le K\sqrt n}
 \left|n^2p^a_{n,R}(z)
 -\frac{\alpha_R}{(2\pi)^2\sqrt{\det D_R}}
   \exp\left(-\tfrac12\xi^{\mathsf T}D_R^{-1}\xi\right)\right|
 &\longrightarrow0,
 \quad \xi=\frac{z-n\bar G_R}{\sqrt n},
 \label{eq:inserted-raw-conclusion}
\end{align}
uniformly in $R$ and the admissible family. No fixed-neighborhood
spectral expansion of an induced operator is assumed in this corollary.
\end{corollary}
\begin{proof}
Fourier inversion for $H^a$ identifies a continuous density for
$\mu^a-E^a$. Add $e^a$ and use the exact decomposition
\eqref{eq:exact-local-decomposition}, which is linear and remains valid
for complex measures. After $v=\sqrt n\,\omega$, the difference between
$n^2K_n*\mu^a$ and the inverse Gaussian integral with cutoff
$\chi(v/n^{1/200})$ is bounded, independently of $z$, by
$(2\pi)^{-4}\|\chi\|_\infty$ times the integral in
\eqref{eq:explicit-growing-arc}. Positive definiteness bounds the
removed Gaussian tail by
$C\int_{|v|>n^{1/200}}e^{-d_0|v|^2/2}\dd v$, which tends to zero
uniformly. This proves the main term in
\eqref{eq:inserted-raw-conclusion}. The other two terms of the exact
decomposition vanish on the stated scale by the assumed complementary
tail and local edge conditions. These steps prove the raw density
conclusion, not a conclusion only for a convolved observation.
\end{proof}

\begin{remark}[Location and scope of the remaining estimates]
\label{rem:growing-arc-boundary}
The physical radius of the proved central band is
$n^{-1/2+1/200}$. It is smaller than $n^{-2/5}$. Thus
\eqref{eq:new-complementary-tail} is not a consequence of the outer-tail
hypothesis in Theorem~\ref{thm:LLT}: the intervening annulus also needs
an estimate. The present theorem removes the need for an assumed induced
spectral expansion on the central band only. It does not assert a
resolvent estimate on its complement or the full physical branch
decomposition. Uniform positive definiteness is proved independently
in Theorem~\ref{thm:joint-uniform-nondegeneracy}. Both inversion routes are
retained, with their distinct frequency boundaries.

The insertion class is an initial-state class with a uniform norm on
the collision cylinder. Terminal insertions, multiple-time insertions,
and exact-event indicators need not belong to it. In particular
\eqref{eq:explicit-growing-arc} does not justify a change of exact
conditioning event or provide the weighted raw tail required for such
a change. The count truncation and the initial-coordinate variation
bounds continue to concern different objects from the derivative sums
of the raw residual densities.
\end{remark}

\section{Central integrals marked at an arbitrary actual return}
\label{sec:marked-return-band}
A weight at a terminal or intermediate return cannot simply be smoothed
as a function of the initial state: composition with a long first-return
iterate need not preserve a uniform variation bound. Instead we change
the origin of the orbit to the marked return. The record then has a
backward and a forward part, while the weight is a function of the new
origin. Both parts can be estimated with the same forward collision
operator. This argument does not require invariance of the section under
time reversal, or mixing of the induced map.

All measures and normalizations in this section are as follows. The
measure $\nu$ is the probability on the full collision cylinder;
$\nu(Y_R^*)=c_*$, and $\nu_R^*=c_*^{-1}\nu|_{Y_R^*}$. For a complex
function $a$ that is zero almost everywhere outside $Y_R^*$ put
\begin{equation}\label{eq:marked-norms}
 M_a=\|a\|_\infty,\qquad V_a=\|a\|_{\mathrm{BV}},\qquad
 \alpha_a=\int_M a\dd\nu.
\end{equation}
For complex functions the variation norm is the sum of the norms of the
real and imaginary parts. It includes their $L^1$ norms and is taken
in the original collision coordinates, including the section boundary.
No upper bound for $V_a$ independent of $n$ is imposed below.

For $0\le k\le n$ define the finite complex measure and its centered,
rescaled transform by
\begin{align}
 \mu^{[k],a}_{n,R}(B)
   &=\int_{Y_R^*}a((F_R^*)^kx)\one_{\{J_{n,R}(x)\in B\}}\dd\nu(x),
       \label{eq:marked-measure}\\
 C^{[k],a}_{n,R}(v)
   &=\int_{Y_R^*}a((F_R^*)^kx)
          e^{iv\cdot(J_{n,R}(x)-n\bar G_R)/\sqrt n}\dd\nu(x).
       \label{eq:marked-transform}
\end{align}
Its mass is $\alpha_a$, which may be complex. The endpoints $k=0,n$
are the initial and terminal insertions. Taking $a=s_R$ gives the
original probability for every $k$. The physical frequency is always
$\omega=v/\sqrt n$.

\subsection{Changing the origin to the mark}
Let $N^-_{j,R}(y)$ be the collision distance to the $j$th preceding
section visit, so that
\[
 (F_R^*)^{-j}y=T_R^{-N^-_{j,R}(y)}y,\qquad N^-_{0,R}=0.
\]
The forward distance is $N^+_{j,R}=N_{j,R}$, with $N^+_{0,R}=0$.
These distances are finite almost everywhere, by recurrence and
invertibility. For nonnegative integers $r,s$ write
\begin{equation}\label{eq:two-sided-collision-sum}
 \mathsf H_{r,s,R}(y)=\sum_{j=-r}^{s-1}h_R(T_R^jy),
\end{equation}
where a sum over an empty interval is zero.

\begin{lemma}[Exact compensation about a marked return]
\label{lem:marked-compensation}
For almost every $y\in Y_R^*$,
\begin{equation}\label{eq:marked-compensation}
 J_{n,R}((F_R^*)^{-k}y)-n\bar G_R
   =\mathsf H_{N^-_{k,R}(y),N^+_{n-k,R}(y),R}(y).
\end{equation}
Consequently
\begin{equation}\label{eq:marked-recentering}
 C^{[k],a}_{n,R}(v)=\int_{Y_R^*}a(y)
 e^{iv\cdot\mathsf H_{N^-_{k,R}(y),N^+_{n-k,R}(y),R}(y)/\sqrt n}
 \dd\nu(y).
\end{equation}
\end{lemma}
\begin{proof}
The collision interval from $(F_R^*)^{-k}y$ up to, but excluding,
$(F_R^*)^{n-k}y$ has left endpoint $-N^-_{k,R}(y)$ and right endpoint
$N^+_{n-k,R}(y)$ in coordinates based at $y$. The sum of $f_R$ on this
interval is precisely the original four-coordinate record. Its negative
part contains exactly $k$ section visits, including the left endpoint
and excluding zero. Its nonnegative part contains exactly $n-k$ visits,
including zero if that part is nonempty and excluding the right endpoint.
Thus the sum of $\eta_R$ over the entire interval is $n$. Substitute
$h_R=f_R-\bar G_R\eta_R$. This proves the first identity, including
$k=0,n$. The actual first-return map preserves $\nu|_{Y_R^*}$, since it
preserves its normalization $\nu_R^*$. The change of variable
$y=(F_R^*)^kx$ gives the second identity. Only additivity and invariance
are used; the two parts are not asserted independent.
\end{proof}

\begin{lemma}[Forward and backward return-clock windows]
\label{lem:two-sided-clock-window}
For $j\ge0$, $b\ge2$, and $m_j=\lfloor j/c_*\rfloor$, uniformly in $R$,
\begin{equation}\label{eq:two-sided-clock-window}
 \nu_R^*\{|N^\pm_{j,R}-m_j|>b\}\le C\frac{j+b}{b^2}.
\end{equation}
Moreover the collision maximal estimate of
Lemma~\ref{lem:dyadic-collision-maximal} holds for an interval beginning
at any integer collision time, positive or negative, and in either order.
\end{lemma}
\begin{proof}
The case $j=0$ is immediate. For either sign let
$H_L^\pm=\sum_{i=1}^L\eta_R\circ T_R^{\pm i}$. Its mean under $\nu$
is $c_*L$. The variance is at most $CL$ by
Proposition~\ref{prop:bv-correlation}. For the negative sign the scalar
lag correlation is the same as for the positive sign by invariance.
Domination $\nu_R^*\le c_*^{-1}\nu$ gives
\[
 \int|H_L^\pm-c_*L|^2\dd\nu_R^*\le CL.
\]
The equivalences $N_j^\pm>L\iff H_L^\pm<j$ and
$N_j^\pm\le L\iff H_L^\pm\ge j$ hold for the actual visits. Take
integer thresholds within a fixed distance of $m_j+b$ and $m_j-b$.
When the latter threshold is negative, its lower-tail event is empty.
Otherwise the required deviation from $c_*L$ is at least $c_*b-O(1)$,
while $L\le C(j+b)$. Chebyshev's inequality proves the estimate for
$b$ above a fixed constant. For the remaining bounded range enlarge $C$
so that its right-hand side is at least one. In particular no assumption
$b\le j$ is needed.

For the maximal assertion, translate any deterministic integer interval
to a nonnegative one using invariance of $\nu$. Every subinterval still
has second moment bounded by its length, from the absolutely summable
collision correlations. The same dyadic prefix proof applies in the
reverse order. Domination gives its version under $\nu_R^*$ as well.
The assertion concerns fixed intervals, not a conditional distribution
of the past given the future.
\end{proof}

\begin{proposition}[Uniform stopping error about the mark]
\label{prop:marked-stopping}
For all $0\le k\le n$, $n\ge2$, and $v\in\R^4$,
\begin{align}
 &\left|C^{[k],a}_{n,R}(v)
  -\int_M a(y)e^{iv\cdot\mathsf H_{m_k,m_{n-k},R}(y)/\sqrt n}\dd\nu(y)
       \right|\notag\\
 &\hspace{12mm}\le C M_a\left(n^{-1/5}
             +|v|n^{-1/5}\log(2+n)\right).
       \label{eq:marked-stopping-error}
\end{align}
\end{proposition}
\begin{proof}
Use \eqref{eq:marked-recentering} and let
$b_n=\lceil n^{3/5}\rceil$. By Lemma~\ref{lem:two-sided-clock-window},
the union of the two events on which a clock differs from its
corresponding $m_j$ by more than $b_n$ has probability at most
$Cn^{-1/5}$ under $\nu_R^*$. Its contribution to the difference of
characteristic functions is at most $C M_an^{-1/5}$, since
$|a|\nu\le M_ac_*\nu_R^*$ on the section.

On the complement, the difference of the two collision sums is bounded
by the sum of two maxima on deterministic windows of length at most
$2b_n+2$ around their respective endpoints. The preceding maximal
estimate bounds its squared integral under $\nu$ by
$C b_n[\log(2+b_n)]^2$. It remains valid when a window crosses collision
time zero. Cauchy--Schwarz and $|e^{ix}-e^{iy}|\le|x-y|$ give
$C M_a|v|\sqrt{b_n/n}\log(2+b_n)$, which is the second term in
\eqref{eq:marked-stopping-error}. Nothing is estimated in an unbounded
stopped sum on the exceptional event. Finitely many initial values of
$n$ are included by increasing the constant.
\end{proof}

\subsection{A collision multiplier between two spectral powers}
Set $a_\delta=\mathsf S_\delta a$ and $h_{R,\delta}=\mathsf S_\delta h_R$.
The former need not be supported in the section. This is harmless:
\begin{equation}\label{eq:marked-smoothing-norms}
 \|a-a_\delta\|_1\le C\delta V_a,\qquad
 \|a_\delta\|_\infty\le M_a,\qquad
 \|M_{a_\delta}\|_{\mathcal B_j\to\mathcal B_j}
       \le C M_a\delta^{-2}.
\end{equation}
The last bound is Lemma~\ref{lem:density-norm-details}. In particular
its constant depends on the supremum norm, not on $V_a$.

\begin{lemma}[A marked deterministic collision integral]
\label{lem:marked-spectral}
Fix $c_0,C_0>0$. Suppose $c_0n\le r+s\le C_0n$ with $r,s\ge0$,
and let $q\ge1$ be an integer satisfying $2q\le r+s$. There are
$c,C>0$ and $\rho<1$ such that, whenever
\[
 |v|/\sqrt n\le c\delta^2,\qquad
 \delta^{-6}|v|^3/\sqrt n\le c,
\]
the following bound holds, with $\ell_\delta=1+|\log\delta|$:
\begin{align}
 &\left|\int_M a\,e^{iv\cdot\mathsf H_{r,s,R}/\sqrt n}\dd\nu
      -\alpha_a e^{-\frac{r+s}{2n}v^{\mathsf T}\Gamma_Rv}\right|
       \notag\\
 &\quad\le C\delta V_a+C M_a\left\{
 |v|\sqrt{\delta\ell_\delta}+|v|^2\delta\ell_\delta
 +\frac{|v|\delta^{-4}+|v|^3\delta^{-6}}{\sqrt n}
 +\delta^{-2}\rho^q+\frac{|v|q}{\sqrt n}
 +\frac{|v|^2q}{n}\right\}.
       \label{eq:marked-spectral-bound}
\end{align}
All constants are independent of the location of the mark and of $a$.
\end{lemma}
\begin{proof}
First replace $a$ by $a_\delta$. The cost is at most $C\delta V_a$.
Let $u_\delta=h_R-h_{R,\delta}$. Stationarity and
\eqref{eq:small-bv-variance} give
\[
 \int\left|\sum_{i=-r}^{s-1}u_\delta(T_R^iy)\right|^2\dd\nu(y)
       \le C(r+s)\delta\ell_\delta.
\]
The bounded weight $|a_\delta|\le M_a$ and Cauchy--Schwarz therefore
bound the cost of replacing $h_R$ by $h_{R,\delta}$ by
$C M_a|v|\sqrt{\delta\ell_\delta}$.

Put $z=v/\sqrt n$ and $L_z=\mathcal L_{R,\delta,z}$, using one of
the fixed local collision spaces. Its mass functional will be denoted
by $\ell$ to avoid confusion with block indices. Invariance and the
transfer convention give the exact chronological pairing
\begin{equation}\label{eq:marked-chronological-pairing}
 \int_M a_\delta(y)
 e^{iz\cdot\sum_{i=-r}^{s-1}h_{R,\delta}(T_R^iy)}\dd\nu(y)
       =\ell\bigl(L_z^s M_{a_\delta} L_z^r\nu\bigr).
\end{equation}
Indeed, starting at $x=T_R^{-r}y$, the rightmost power records collisions
$0,\ldots,r-1$, the multiplier is evaluated at $T_R^rx=y$, and the
leftmost power records collisions $r,\ldots,r+s-1$. Thus all weights
are evaluated at the prescribed states.

Suppose first that $r,s\ge q$. Write
$L_z^m=\lambda(z)^m\Pi(z)+E_m(z)$ with
$\|E_m(z)\|\le C\rho^m$. The hypotheses and
\eqref{eq:smooth-cubic-expansion} bound each real-frequency principal
power of length at most $r+s$ by a uniform constant. Expanding
\eqref{eq:marked-chronological-pairing}, every complementary term is
bounded by $C M_a\delta^{-2}(\rho^r+\rho^s)$; the term with both
complementary powers is included in this bound. For the principal
amplitude use $\Pi(0)=\nu\ell$ and \eqref{eq:marked-smoothing-norms}:
\begin{align*}
 \ell(\Pi(0)M_{a_\delta}\Pi(0)\nu)&=\alpha_a,\\
 |\ell(\Pi(z)M_{a_\delta}\Pi(z)\nu)-\alpha_a|
       &\le C M_a\delta^{-4}|z|.
\end{align*}
The cubic remainder accumulated over $r+s$ collisions is at most
$C\delta^{-6}|v|^3/\sqrt n$. Since the quadratic form is nonnegative,
split the principal comparison as
$\lambda(z)^{r+s}(A(z)-\alpha_a)
 +\alpha_a(\lambda(z)^{r+s}-e^{-(r+s)z^{\mathsf T}\Gamma_{R,\delta}z/2})$,
where $A(z)$ is its principal amplitude. The first term uses the
amplitude bound above; the second uses $|\alpha_a|\le M_a$ and
$|e^w-1|\le |w|e^{|w|}$. Thus the cubic error is multiplied only
by $C M_a$, not by the multiplier norm. This proves the required spectral bound in the
case of two long blocks.

If $r<q$, delete the first $r$ terms from the real exponent on the left
of \eqref{eq:marked-chronological-pairing}. The cost is at most
$C M_a|v|q/\sqrt n$, by boundedness of $h_{R,\delta}$. The remaining
pairing is $\ell(L_z^s(a_\delta\nu))$. Its principal amplitude differs
from $\alpha_a$ by at most $C M_a\delta^{-4}|z|$, by the smooth-density
embedding, and its complementary term is bounded by
$C M_a\delta^{-2}\rho^s$. Here $s\ge q$. If $s<q$, the remaining
pairing is instead $\ell(M_{a_\delta}L_z^r\nu)$, with exactly the same
bounds. At most one of $r,s$ is less than $q$. Restoring the removed
Gaussian time coefficient changes its value by at most
$C M_a|v|^2q/n$. This handles zero blocks as well; no spectral expansion
of an identity operator is used.

Finally replace $\Gamma_{R,\delta}$ by $\Gamma_R$. The segment joining
these matrices is positive semidefinite, and
\eqref{eq:collision-covariance-limit} bounds the change by
$C M_a|v|^2\delta\ell_\delta$. All the displayed errors combine to
\eqref{eq:marked-spectral-bound}. Only forward collision operators
occurred in the proof.
\end{proof}

\subsection{The growing band and norm growth}
\begin{theorem}[Integrated central band with an arbitrary return mark]
\label{thm:marked-return-band}
There is a constant $C$ such that for every $n\ge2$, $R\in I$,
$0\le k\le n$, and complex $a$ as in \eqref{eq:marked-norms},
\begin{align}
 &\int_{|v|\le2n^{1/200}}
   \left|C^{[k],a}_{n,R}(v)-\alpha_a e^{-v^{\mathsf T}D_Rv/2}\right|\dd v
      \notag\\
 &\hspace{15mm}\le C\left\{
 M_an^{-3/280}\sqrt{\log(2+n)}+V_an^{-9/175}\right\}.
      \label{eq:marked-return-band}
\end{align}
In particular the error tends to zero uniformly for bounded $M_a$ and
$V_a=O(n^\kappa)$ with any $\kappa<9/175$.
\end{theorem}
\begin{proof}
Take $r=m_k$ and $s=m_{n-k}$, so
\[
 (r+s)/n=c_*^{-1}+O(n^{-1}),
\]
with an error independent of $k$. They satisfy the length bounds in
Lemma~\ref{lem:marked-spectral}. Choose
\begin{equation}\label{eq:marked-scales}
 U_n=n^{1/200},\qquad \delta_n=\tfrac14n^{-1/14},\qquad
 q_n=\left\lceil\left(1+\frac4{|\log\rho|}\right)\log(2+n)\right\rceil.
\end{equation}
Then $\rho^{q_n}\le(2+n)^{-4}$ and $2q_n\le r+s$ for sufficiently
large $n$, uniformly in the mark. On $|v|\le2U_n$, both analytic
restrictions hold for large $n$: their positive margins are
$1/2-2/14-1/200$ and $1/2-6/14-3/200$.

Apply the marked spectral estimate and
Proposition~\ref{prop:marked-stopping}. Replacing $(r+s)\Gamma_R/n$
by $D_R$ costs at most $C M_a|v|^2/n$. Integration uses
$\int_{|v|\le2U}|v|^j\dd v=C_jU^{4+j}$. The contribution involving
$V_a$ is only
\[
 C V_a\delta_nU_n^4=C'V_an^{-9/175}.
\]
For clarity, the negative-power margins of the remaining polynomial
terms are displayed below; logarithmic factors are recorded separately.
\begin{center}
\begin{tabular}{lll}
Source & Margin & Logarithmic factor\\ \hline
Observable smoothing & $3/280$ & $\sqrt{\log(2+n)}$\\
Covariance smoothing & $29/700$ & $\log(2+n)$\\
Principal amplitude & $53/280$ & $1$\\
Cubic remainder & $51/1400$ & $1$\\
Clock exception & $9/50$ & $1$\\
Clock window & $7/40$ & $\log(2+n)$\\
Short collision block & $19/40$ & $\log(2+n)$\\
Time coefficient and rounding & $97/100$ & $\log(2+n)$
\end{tabular}
\end{center}
The complementary-power contribution is at most
$C M_an^{2/14+4/200}(2+n)^{-4}$ and is smaller than the displayed
polynomial errors. The first margin in the table is strictly smaller
than all the others, so it absorbs their logarithmic factors.
This proves \eqref{eq:marked-return-band} for large $n$. For the
remaining finite range, both transforms have modulus at most $M_a$,
and the frequency ball has finite volume; enlargement of $C$ proves
the stated estimate for all $n\ge2$. The last assertion follows by
substitution, without changing the smoothing scale with the insertion.
\end{proof}

\begin{corollary}[Conditioning on a state at a marked return]
\label{cor:marked-state-conditioning}
Let $E_{n,R}\subset Y_R^*$ have positive measure and a zero-extended
indicator in $\mathrm{BV}(M)$. Set
$V_{n,R}=\|\one_{E_{n,R}}\|_{\mathrm{BV}}$. Under the section probability,
condition on the exact event
\[
 A_{n,k,R}=\{x:(F_R^*)^kx\in E_{n,R}\}
\]
Uniformly for $0\le k\le n$,
\begin{align}
 &\int_{|v|\le2n^{1/200}}
 \left|\mathbb E_{\nu_R^*}\left[
 e^{iv\cdot(J_{n,R}-n\bar G_R)/\sqrt n}\mid A_{n,k,R}\right]
          -e^{-v^{\mathsf T}D_Rv/2}\right|\dd v\notag\\
 &\quad\le\frac{C}{\nu(E_{n,R})}
  \left(n^{-3/280}\sqrt{\log(2+n)}+V_{n,R}n^{-9/175}\right).
       \label{eq:marked-conditional-band}
\end{align}
For example this tends to zero if, uniformly in $R$,
$\nu_R^*(E_{n,R})\ge c n^{-\beta}$ and $V_{n,R}\le Cn^\kappa$,
where $\beta<3/280$ and $\beta+\kappa<9/175$.
\end{corollary}
\begin{proof}
Invariance gives $\nu_R^*(A_{n,k,R})=\nu(E_{n,R})/c_*$.
The conditional characteristic function is exactly
$C^{[k],\one_{E_{n,R}}}_{n,R}/\nu(E_{n,R})$. Apply
Theorem~\ref{thm:marked-return-band} with $M_a=1$. The two strict
exponent inequalities absorb the logarithm. This argument uses the
same event throughout and makes no comparison with another event.
\end{proof}

\paragraph{Relation to the raw density theorem.}
These estimates control the central integral of the actual complex
measures \eqref{eq:marked-measure}. The exact local decomposition
\eqref{eq:exact-local-decomposition} also applies to these measures by
linearity. More explicitly, use $D_R\ge d_0 I_4$ from
Theorem~\ref{thm:joint-uniform-nondegeneracy}; suppose they admit
extracted measures $E^{[k],a}_{n,R}$ with mixed densities; and suppose
the residual transform $H^{[k],a}_{n,R}$ is in
$L^1(\T^3\times\R)$, satisfies the complementary-tail condition
\eqref{eq:new-complementary-tail}, and satisfies the local edge
conditions \eqref{eq:local-edge-errors}, all uniformly in the chosen
marks and insertions. If the right-hand side of
\eqref{eq:marked-return-band} tends to zero uniformly, the proof of
Corollary~\ref{cor:raw-from-growing-arc} gives its raw density conclusion
for these marked measures. Thus its central hypothesis is now proved
also for terminal and single intermediate state insertions. The other
hypotheses have not been proved by the recentering argument.

The band still has physical radius $n^{-99/200}$, not $n^{-2/5}$.
Multiple separated return-state factors are not one insertion at a
single mark, and composing them with a long induced iterate does not
give the variation norm used here. Likewise an exact final lattice or
flight-time constraint is a condition on the record, not generally a
collision-state set $E_{n,R}$ of the class in
Corollary~\ref{cor:marked-state-conditioning}. Its relative error bound
is consequently not a replacement for the weighted raw local limits
needed for those constraints. Nor does a two-sided orbit identity
provide a pointwise representative of the zero-variance coboundary.
Covariance nondegeneracy is now proved in
Theorem~\ref{thm:joint-uniform-nondegeneracy}. The complete complementary
integral and full critical/singular residual estimates retain their
separate roles.

\section{Unsmoothed moments and the covariance of actual returns}
\label{sec:unsmoothed-moments}
The covariance of a weak Gaussian limit need not be the limit of the
normalized second moments. We establish that identification for the
actual return record, including a bounded-variation insertion at any
return. The argument also gives fourth moments for the unsmoothed
collision sums and improves the two-sided stopping estimate. It uses
only the collision splitting, not decay of correlations for $F_R^*$.

Write $\|u\|_{\mathcal V}=\|u\|_\infty+\|u\|_{\mathrm{BV}}$.
All expectations without a subscript in this section are with respect
to the collision probability $\nu$. Constants are uniform in $R\in I$.
The smoothing operators preserve means and have the bounds in
Lemma~\ref{lem:bv-smoothing}.

\subsection{Decoupling at a deterministic collision gap}
\begin{lemma}[A finite-product decoupling estimate]
\label{lem:bv-finite-product}
For $2\le q\le4$, integers $t_1\le\cdots\le t_q$, and
$1\le p<q$, bounded-variation functions $u_1,\ldots,u_q$ satisfy
\begin{align}
 &\left|\mathbb E\prod_{j=1}^q u_j\circ T_R^{t_j}
 -\left(\mathbb E\prod_{j=1}^p u_j\circ T_R^{t_j}\right)
  \left(\mathbb E\prod_{j=p+1}^q u_j\circ T_R^{t_j}\right)\right|
 \notag\\
 &\hspace{18mm}\le C\exp\{-\gamma(t_{p+1}-t_p)\}
                         \prod_{j=1}^q\|u_j\|_{\mathcal V}.
 \label{eq:bv-finite-product}
\end{align}
The times may be negative or repeated. The constants $C,\gamma>0$
can be chosen for all the indicated values of $q$ and $p$.
\end{lemma}
\begin{proof}
By multilinearity assume $\|u_j\|_{\mathcal V}\le1$; a zero norm
makes the assertion immediate. Replace every $u_j$ by
$u_{j,\epsilon}=\mathsf S_\epsilon u_j$. Telescoping the products,
invariance of $\nu$, and the uniform supremum bounds show that the
change in the left-hand expression is at most $C_q\epsilon$,
independently of all the times. In particular, no variation norm of
$u_j\circ T_R^{t_j}$ is estimated.

Translate $t_1$ to zero. Put $d_j=t_{j+1}-t_j$ and use the transfer
convention of Lemma~\ref{lem:collision-spectral-input}. The smoothed
joint expectation is
\[
 \ell\bigl(M_{u_{q,\epsilon}}\mathcal L_R^{d_{q-1}}
 M_{u_{q-1,\epsilon}}\cdots
 \mathcal L_R^{d_1}M_{u_{1,\epsilon}}\nu\bigr).
\]
Replace the power across the chosen gap by $\Pi_R=\nu\ell$.
The resulting expression is exactly the product of the two smoothed
expectations in \eqref{eq:bv-finite-product}. The difference is bounded
by $C_q\epsilon^{-2q}\vartheta^{d_p}$. Indeed each smooth multiplier
costs at most $C\epsilon^{-2}$, every other power is uniformly bounded,
and the complementary power across the gap costs $C\vartheta^{d_p}$.
For a gap of length zero use
$\mathcal L_R^0-\Pi_R=I-\Pi_R$, with its uniform bound.

Choose $\epsilon=\tfrac14\vartheta^{d_p/(2q+1)}$. Smoothing and
operator errors are at most $C_q\vartheta^{d_p/(2q+1)}$. The finite
parameter cover permits $\gamma=-\log\vartheta/9$ after changing $C$.
Invariance justifies the initial translation also for negative times.
\end{proof}

\begin{proposition}[Fourth moments without smoothing]
\label{prop:unsmoothed-fourth}
Let $u$ be real, centered, and $\|u\|_{\mathcal V}\le B$.
For every deterministic integer interval of length $m\ge1$,
\begin{equation}\label{eq:unsmoothed-fourth}
 \mathbb E\left|\sum_{j=r}^{r+m-1}u\circ T_R^j\right|^4\le C_Bm^2,
 \qquad
 \mathbb E\max_{0\le l\le m}
 \left|\sum_{j=r}^{r+l-1}u\circ T_R^j\right|^4\le C_Bm^2.
\end{equation}
Both estimates hold in the reverse order of the interval and under
any nonnegative initial density bounded by a fixed constant. They hold
componentwise, and hence in Euclidean norm, for $h_R$.
\end{proposition}
\begin{proof}
Put $c_u(a)=\mathbb E[u(u\circ T_R^a)]$. For $a,b,c\ge0$, apply
Lemma~\ref{lem:bv-finite-product} to four copies of $u$ at times
$0,a,a+b,a+b+c$. Splitting across the middle gap bounds
\[
 \left|\mathbb E[u(u\circ T_R^a)(u\circ T_R^{a+b})
                   (u\circ T_R^{a+b+c})]-c_u(a)c_u(c)\right|
       \le C_B e^{-\gamma b}.
\]
Splitting off the first or last centered factor bounds the fourfold
expectation by $C_Be^{-\gamma a}$ and $C_Be^{-\gamma c}$,
respectively. Also $|c_u(a)c_u(c)|\le C_Be^{-\gamma(a+c)}$.
Taking the minimum of these three bounds gives
\begin{align}\label{eq:four-product-gap}
 &\left|\mathbb E[u(u\circ T_R^a)(u\circ T_R^{a+b})
                   (u\circ T_R^{a+b+c})]-c_u(a)c_u(c)\right|
 \le C_B e^{-\gamma\max\{a,b,c\}}.
\end{align}
This includes repeated times.

Expand the fourth power and order its four indices; their multiplicity
is at most $24$. For each initial index the sum of
$e^{-\gamma(a+c)}$ over the two outer gaps is bounded, while the middle
gap has at most $m$ choices. This contributes $O(m^2)$. The sum over
all three nonnegative gaps of $e^{-\gamma\max\{a,b,c\}}$ is finite,
so the remaining contribution is $O(m)$. This proves the first bound.

For the second enlarge $m$ to a power of two $M<2m$. Every prefix is
a disjoint union of at most one dyadic interval of each length.
Consequently its absolute sum is bounded by the sum, over scales, of
the maximum absolute interval sum at that scale. At interval length
$l$, the fourth power of this maximum has expectation at most
$C_B(M/l)l^2=C_BMl$. Minkowski's inequality therefore gives an $L^4$
norm at most
\[
 C_B\sum_{j=0}^{\log_2 M}(M^2 2^{-j})^{1/4}
       \le C_B M^{1/2}.
\]
There is no logarithmic loss. Stationarity translates deterministic
intervals; reversal gives suffixes of the same interval. Domination
by a bounded initial density proves the remaining assertions.
\end{proof}

\subsection{Fourth-moment windows and two-sided stopping}
Retain $N^\pm_{j,R}$, $m_j=\lfloor j/c_*\rfloor$ and
$\mathsf H_{r,s,R}$ from Section~\ref{sec:marked-return-band}. Define
\begin{equation}\label{eq:moment-stopped-pair}
 Z_{n,k,R}=\mathsf H_{N^-_{k,R},N^+_{n-k,R},R},\qquad
 W_{n,k,R}=\mathsf H_{m_k,m_{n-k},R}.
\end{equation}
Thus $Z_{n,k,R}(y)=J_{n,R}((F_R^*)^{-k}y)-n\bar G_R$ almost
everywhere. Zero-length backward and forward intervals are empty.

\begin{lemma}[Fourth-moment return-clock window]
\label{lem:fourth-clock-window}
For $j\ge0$, $b\ge2$, and either sign,
\begin{equation}\label{eq:fourth-clock-window}
 \nu_R^*\{|N^\pm_{j,R}-m_j|>b\}
       \le C\frac{(j+b)^2}{b^4}.
\end{equation}
\end{lemma}
\begin{proof}
The case $j=0$ is immediate. The centered visit count
$\sum_{i=1}^L(\eta_R-c_*)\circ T_R^{\pm i}$ has fourth moment at
most $CL^2$ under $\nu_R^*$, by
Proposition~\ref{prop:unsmoothed-fourth}. The exact visit equivalences
in Lemma~\ref{lem:two-sided-clock-window} turn the upper and lower
return-clock events into deviations of these counts. At integer
thresholds within a bounded distance of $m_j\pm b$, the required
deviation is at least $c_*b-O(1)$ and $L\le C(j+b)$. Markov's
inequality proves \eqref{eq:fourth-clock-window}. A negative lower
threshold contributes no event. Increase $C$ for bounded $b$.
\end{proof}

\begin{theorem}[Moment control of the actual marked stopping]
\label{thm:marked-L2-stopping}
Uniformly for $n\ge2$, $0\le k\le n$, and $R\in I$,
\begin{align}
 \int (|Z_{n,k,R}|^4+|W_{n,k,R}|^4)\dd\nu_R^*&\le Cn^2,
       \label{eq:actual-marked-fourth}\\
 \left\|\frac{Z_{n,k,R}-W_{n,k,R}}{\sqrt n}
                  \right\|_{L^2(\nu_R^*)}&\le Cn^{-1/6}.
       \label{eq:actual-marked-L2}
\end{align}
For every real $v\in\R^4$ and insertion $a$ of
\eqref{eq:marked-norms}, the stopping comparison also satisfies
\begin{equation}\label{eq:improved-marked-stopping}
 \left|C^{[k],a}_{n,R}(v)
       -\int_M a e^{iv\cdot W_{n,k,R}/\sqrt n}\dd\nu\right|
       \le C M_a(1+|v|)n^{-2/9}.
\end{equation}
\end{theorem}
\begin{proof}
The deterministic interval has length at most $Cn$, so its fourth
moment follows from Proposition~\ref{prop:unsmoothed-fourth} and
$\nu_R^*\le c_*^{-1}\nu$. For the random interval put
$Q=N^-_{k,R}+N^+_{n-k,R}$. Invariance of the actual return map gives
$N^-_{j,R}\overset{\rm law}=N^+_{j,R}$ under $\nu_R^*$: use
$N^-_{j,R}((F_R^*)^jx)=N^+_{j,R}(x)$. Hence
Theorem~\ref{thm:cumulative-return-tail} and a union bound imply
\[
 \nu_R^*(Q>L)\le C e^{an-cL/2}.
\]
No independence of the two clocks is used. Fix $A>2a/c+2$.
On $Q\le An$, $|Z_{n,k,R}|$ is bounded by the sum of the backward
and forward collision maxima up to $\lceil An\rceil$. Their fourth
moments are $O(n^2)$. On $Q>An$, boundedness of $h_R$ gives
$|Z_{n,k,R}|\le C Q$. Integrating the displayed exponential tail
bounds that fourth-moment contribution by
$C(1+n)^4e^{-c' n}$ for some $c'>0$. This proves
\eqref{eq:actual-marked-fourth}.

Let $2\le b\le n$ and let $\mathcal E_b$ be the union of the two
clock-window exceptions. Lemma~\ref{lem:fourth-clock-window} gives
$\nu_R^*(\mathcal E_b)\le Cn^2/b^4$. Off this event, the difference
$Z_{n,k,R}-W_{n,k,R}$ is bounded by maxima on two deterministic
collision windows of length $O(b)$, including windows that cross zero.
Their second moments are $O(b)$ by the fourth-moment maximal estimate.
On the exceptional event, Cauchy--Schwarz and
\eqref{eq:actual-marked-fourth} give
\[
 \int |Z_{n,k,R}-W_{n,k,R}|^2\one_{\mathcal E_b}\dd\nu_R^*
       \le Cn\,\nu_R^*(\mathcal E_b)^{1/2}\le Cn^2/b^2.
\]
Thus the normalized squared $L^2$ error is at most
$C(b/n+n/b^2)$. Taking $b=\lceil n^{2/3}\rceil$ proves
\eqref{eq:actual-marked-L2}.

For the characteristic comparison, the exceptional event costs only
$CM_a n^2/b^4$. On its complement, the same deterministic windows and
$|e^{ix}-e^{iy}|\le|x-y|$ give $CM_a|v|\sqrt{b/n}$.
Choose instead $b=\lceil n^{5/9}\rceil$. Both powers are
$n^{-2/9}$. The recentering identity
\eqref{eq:marked-recentering} and $|a|\nu\le c_*M_a\nu_R^*$
complete the proof. All rounding exceptions are absorbed in $C$.
\end{proof}

\subsection{Quadratic insertions and the induced covariance}
The following estimate supplies the moment identification that is not
asserted by the characteristic-function proof in
Remark~\ref{rem:gaussian-versus-raw}.

\begin{lemma}[A centered insertion in a quadratic collision sum]
\label{lem:quadratic-collision-insertion}
Let $a$ be a complex bounded-variation function on $M$, and let
$\alpha=\int a\dd\nu$. For all $r,s\ge0$, with $m=r+s$,
\begin{align}
 \left|\int a\mathsf H_{r,s,R}\dd\nu\right|
       &\le C(\|a\|_\infty+\|a\|_{\mathrm{BV}}),
       \label{eq:linear-insertion-moment}\\
 \left\|\int a\mathsf H_{r,s,R}\otimes\mathsf H_{r,s,R}\dd\nu
                          -\alpha m\Gamma_R\right\|
       &\le C(\|a\|_\infty+\|a\|_{\mathrm{BV}}).
       \label{eq:quadratic-insertion-moment}
\end{align}
The constants do not depend on the location of the origin in the interval.
\end{lemma}
\begin{proof}
Center $a$ by subtracting $\alpha$. Each component of $h_R$ is
centered. The two-factor version of
Lemma~\ref{lem:bv-finite-product} bounds
$\int(a-\alpha)h_R\circ T_R^i\dd\nu$ by
$C\|a\|_{\mathcal V}e^{-\gamma|i|}$. Summation proves the first
assertion.

For components $u,v$ of $h_R$, sort the three times $0,i,j$ and
denote their two consecutive gaps by $d_1,d_2$. All three factors
$a-\alpha,u,v$ are centered. Splitting off the first and the last
factor in \eqref{eq:bv-finite-product} gives
\[
 \left|\int(a-\alpha)(u\circ T_R^i)(v\circ T_R^j)\dd\nu\right|
       \le C\|a\|_{\mathcal V}e^{-\gamma\max\{d_1,d_2\}}.
\]
Because $\max(d_1,d_2)\ge\max(|i|,|j|)/2$, the right side is
summable over all $(i,j)\in\Z^2$. The centered insertion therefore
changes the quadratic moment of any such interval by at most
$C\|a\|_{\mathcal V}$. Finally the exponentially summable collision
covariances give
\[
 \left\|\int\mathsf H_{r,s,R}\otimes\mathsf H_{r,s,R}\dd\nu
                                  -m\Gamma_R\right\|\le C.
\]
This is the finite covariance expansion with the lag weights $m-l$;
the sum of $l$ times the absolute lag covariances is finite. Multiplying
by $\alpha$ proves the second assertion. Real and imaginary parts
handle complex $a$ without any positivity assumption.
\end{proof}

\begin{theorem}[Moments at an arbitrary actual return mark]
\label{thm:marked-quadratic-moments}
Put $U_{n,R}=J_{n,R}-n\bar G_R$. For $a$ as in
\eqref{eq:marked-norms}, uniformly for $n\ge2$ and $0\le k\le n$,
\begin{align}
 \left|\int_{Y_R^*}a((F_R^*)^kx)\frac{U_{n,R}(x)}{\sqrt n}\dd\nu(x)\right|
 &\le C\left(M_an^{-1/6}+(M_a+V_a)n^{-1/2}\right),
       \label{eq:marked-first-moment}\\
 \left\|\int_{Y_R^*}a((F_R^*)^kx)
       \frac{U_{n,R}(x)\otimes U_{n,R}(x)}n\dd\nu(x)-\alpha_a D_R\right\|
 &\le C\left(M_an^{-1/6}+(M_a+V_a)n^{-1}\right).
       \label{eq:marked-second-moment}
\end{align}
These are moments of the finite complex measure of
\eqref{eq:marked-measure}, using $\nu|_{Y_R^*}$, not its normalization.
In particular
\begin{equation}\label{eq:actual-covariance-rate}
 \sup_{R\in I}\left\|\frac1n\int U_{n,R}\otimes U_{n,R}\dd\nu_R^*
                                               -D_R\right\|
       \le C n^{-1/6}.
\end{equation}
\end{theorem}
\begin{proof}
Recenter at $y=(F_R^*)^kx$. The weighted moments then involve
$a(y)Z_{n,k,R}(y)$ and $a(y)Z_{n,k,R}(y)\otimes Z_{n,k,R}(y)$.
Theorem~\ref{thm:marked-L2-stopping} changes the normalized first
moment by at most $CM_an^{-1/6}$. For the second use
\[
 \|z\otimes z-w\otimes w\|\le |z-w|(|z|+|w|)
\]
and Cauchy--Schwarz. Equations~\eqref{eq:actual-marked-fourth} and
\eqref{eq:actual-marked-L2} show that its normalized change is also
at most $CM_an^{-1/6}$. The factor $c_*$ in passing from
$\nu|_{Y_R^*}$ to $\nu_R^*$ is fixed.

Apply Lemma~\ref{lem:quadratic-collision-insertion} to
$W_{n,k,R}=\mathsf H_{m_k,m_{n-k},R}$. Its length satisfies
$m_k+m_{n-k}=n/c_*+O(1)$ uniformly in $k$. Since
$D_R=\Gamma_R/c_*$, division by $\sqrt n$ or $n$ gives
\eqref{eq:marked-first-moment} and \eqref{eq:marked-second-moment}.
For \eqref{eq:actual-covariance-rate} take $a=s_R=c_*^{-1}\eta_R$,
whose mass is one and whose two norms are uniformly bounded. Invariance
of $F_R^*$ and the Kac mean give $\int U_{n,R}\dd\nu_R^*=0$.
Thus the last display is convergence of the actual covariance matrices.
\end{proof}

\begin{corollary}[The induced Green--Kubo formula with Cesaro weights]
\label{cor:induced-cesaro-covariance}
Let $g_R=G_R-\bar G_R$ and
$C^*_{j,R}=\int g_R\otimes(g_R\circ(F_R^*)^j)\dd\nu_R^*$.
Then
\begin{equation}\label{eq:induced-cesaro-covariance}
 \left\|C^*_{0,R}+\sum_{j=1}^{n-1}\left(1-\frac jn\right)
          (C^*_{j,R}+(C^*_{j,R})^{\mathsf T})-D_R\right\|
       \le C n^{-1/6}.
\end{equation}
\end{corollary}
\begin{proof}
The exponential one-return moment makes each term finite. Expand the
second moment of $\sum_{j=0}^{n-1}g_R\circ(F_R^*)^j$ and use stationarity
under $F_R^*$. The left matrix before subtracting $D_R$ is exactly
$n^{-1}\int U_{n,R}\otimes U_{n,R}\dd\nu_R^*$. Apply
\eqref{eq:actual-covariance-rate}.
\end{proof}

\begin{corollary}[Moments under the unchanged marked-state event]
\label{cor:marked-conditional-moments}
Use the event $A_{n,k,R}$ and variation $V_{n,R}$ of
Corollary~\ref{cor:marked-state-conditioning}, and put
$p_{n,R}=\nu_R^*(E_{n,R})>0$. The conditional first moment of
$U_{n,R}/\sqrt n$ is bounded by
\[
 \frac C{p_{n,R}}\left(n^{-1/6}+(1+V_{n,R})n^{-1/2}\right).
\]
Its conditional second moment about the original center $n\bar G_R$
differs from $D_R$ by at most
\[
 \frac C{p_{n,R}}\left(n^{-1/6}+(1+V_{n,R})n^{-1}\right).
\]
Consequently, if $p_{n,R}\ge c n^{-\beta}$ and
$V_{n,R}\le Cn^\kappa$, the latter difference tends to zero when
$\beta<1/6$ and $\beta+\kappa<1$. The conditional mean tends to zero,
and the conditional covariance tends to $D_R$, under the stronger
conditions $\beta<1/6$ and $\beta+\kappa<1/2$.
\end{corollary}
\begin{proof}
Take $a=c_*^{-1}\one_{E_{n,R}}$ in
Theorem~\ref{thm:marked-quadratic-moments}. Its mass is $p_{n,R}$.
Divide the two estimates by this unchanged probability. Subtracting
the tensor square of the conditional mean identifies the conditional
covariance. The strict exponent inequalities prove convergence.
\end{proof}

\paragraph{Role in raw local inversion.}
Theorems~\ref{thm:marked-L2-stopping} and
\ref{thm:marked-quadratic-moments} concern the unsmoothed physical
record, not a replacement process. They identify its covariance by
actual second moments, establish uniform integrability of its squared
normalized size, and strengthen the stopping part of Theorem C.
In particular all the rare marked-state events in
Corollary~\ref{cor:marked-state-conditioning} also satisfy the new
conditional moment conclusions. Cesaro convergence in
\eqref{eq:induced-cesaro-covariance} does not assert absolute summability
of the induced correlations.

Positive definiteness is a distinct assertion. The estimates above give
no positive lower bound for the smallest eigenvalue of $D_R$ and no
periodic-evaluable representative of its zero-variance coboundary.
Theorem~\ref{thm:joint-uniform-nondegeneracy} supplies positive
definiteness by a separate measurable phase argument, not from these
moment bounds. Likewise \eqref{eq:improved-marked-stopping} is not a complementary
frequency estimate. The physical central radius remains
$2n^{-99/200}$, or $2n^{1/200}$ after rescaling. The annulus beyond it,
the complete raw critical and singular branch sum, and the weighted
tails for an exact physical observation retain their roles in
Theorem~\ref{thm:LLT}. No conditioning event has been replaced.

\section{Nondegeneracy of the collision-count component}
\label{sec:count-nondegeneracy}
\begin{corollary}[The count component is uniformly nondegenerate]
\label{cor:positive-count-variance}
Let $e_3=(0,0,1,0)$, the collision-count direction. There is a constant
$d_N>0$ such that
\begin{equation}\label{eq:positive-count-variance}
 e_3^{\mathsf T}D_Re_3\ge d_N\qquad(R\in I).
\end{equation}
In particular the actual normalized variance of $N_{n,R}$ is bounded
below by $d_N/2$ for all sufficiently large $n$, uniformly in $R$.
\end{corollary}
\begin{proof}
First, $T_R$ is mixing for the collision probability. The splitting in
Lemma~\ref{lem:collision-spectral-input} proves mixing for smooth
densities and smooth tests. Approximation in $L^2(\nu)$ and invariance
extend it to arbitrary $L^2$ functions: the approximation error in a
correlation is bounded, uniformly in its lag, by Cauchy--Schwarz.
Explicitly, if $f_\epsilon,g_\epsilon$ approximate $f,g$ in $L^2(\nu)$,
then, uniformly in $m$,
\[
 \left|\int\overline f(g\circ T_R^m)\dd\nu
       -\int\overline{f_\epsilon}(g_\epsilon\circ T_R^m)\dd\nu\right|
 \le\|f-f_\epsilon\|_2\|g\|_2
      +\|f_\epsilon\|_2\|g-g_\epsilon\|_2.
\]
The product of the two means converges under the same approximation.
This implication is used for each fixed $R$, without a rate for the
nonsmooth functions below.

Suppose $e_3^{\mathsf T}D_Re_3=0$. By
Theorem~\ref{thm:gaussian-kernel-coboundary}, there is a real
$b\in L^2(\nu)$ with
\[
 1-c_*^{-1}\eta_R=b-b\circ T_R
 \quad\hbox{almost everywhere}.
\]
Set $q=\exp(2\pi i c_*b)$, so $|q|=1$. Since
$2\pi c_*=91/5000$ and $\eta_R$ is integer-valued,
\[
 q\circ T_R=e^{-2\pi i c_*}q.
\]
Here $0<c_*<1$, so the multiplier $\lambda=e^{-2\pi i c_*}$ is
not one. Invariance gives $\int q\dd\nu=\lambda\int q\dd\nu$,
hence $\int q\dd\nu=0$. Mixing then gives
$\int\overline q(q\circ T_R^m)\dd\nu\to0$. But this integral is
exactly $\lambda^m$, whose modulus is one. This contradiction proves
strict positivity at every $R$. The continuity of $D_R$ and compactness
of $I$ give the uniform constant in \eqref{eq:positive-count-variance}.
Finally \eqref{eq:actual-covariance-rate} identifies the actual count
variance and gives the asserted lower bound for large $n$.

The argument uses the exact integer-valued section indicator and the
nonintegral section mass. It neither chooses a pointwise representative
at a periodic orbit nor, by itself, proves positive definiteness in
every joint direction. The latter conclusion is proved separately in
Theorem~\ref{thm:joint-uniform-nondegeneracy}.
\end{proof}


\section{Measurable collision phases and symplectic area}
\label{sec:measurable-phase-rigidity}
The measurable coboundary in Theorem~\ref{thm:gaussian-kernel-coboundary}
need not have a value at a prescribed periodic orbit. We shall instead
use conditional pairs on a positive-measure hyperbolic product set.
The discontinuous section contribution will be removed algebraically
before this argument is applied. Throughout this section $R$ is fixed;
none of the local product constants is asserted uniform in $R$.

Write $\kappa_R$ for the one-collision lattice label, in the basis
$a=(1,0)$, $b=(1/2,\sqrt3/2)$. Thus
$f_R=(\kappa_{R,1},\kappa_{R,2},1,\tau_R)$. All identities between
measurable functions are understood almost everywhere with respect to
$\nu$. Let $\mathbb S^1=\{z\in\mathbb C:|z|=1\}$.

\subsection{A qualitative local product input}
\begin{lemma}[Regular product sets for the collision map]
\label{lem:regular-collision-product}
There is a compact positive-$\nu$-measure set $\mathcal P$ in a regular
collision coordinate rectangle, away from grazing, with the following
properties. It has product coordinates
$\Psi:K^u\times K^s\longrightarrow\mathcal P$. Write a product point as $\Psi(u,s)$. Fixing $u$ places points on one
local homogeneous stable curve; fixing $s$ places them on one local
homogeneous unstable curve. Thus $K^u$ labels stable leaves and $K^s$
labels unstable leaves. These curves are $C^1$, have transverse tangent cones,
and each stable curve meets each unstable curve exactly once.

The probability $m=\nu(\mathcal P)^{-1}\nu|_{\mathcal P}$ is equivalent,
in these coordinates, to a product $m^u\otimes m^s$ of nonatomic
transversal probabilities. A stable arc joining two points on the same
curve remains a connected regular arc under every forward collision
iterate, and its length tends to zero. The corresponding statement
holds for unstable arcs under backward iterates. In particular the
successive lattice labels agree along the respective arcs.

For every sufficiently small $\varepsilon>0$, the product coordinates
can be restricted to positive-measure subsets so that every four-corner
stable--unstable quadrilateral lies in a coordinate box of diameter
less than $\varepsilon$. Almost every such quadrilateral with distinct
coordinates is a simple, nondegenerate Jordan domain.
\end{lemma}
\begin{proof}
We use the qualitative billiard product construction in
\cite[Sections 8.2--8.3]{Young}. Section 8.2B gives transverse stable
and unstable cones. Homogeneous local curves and their regular iterates
are described in Section 8.2D and constructed in Section 8.3,
Sublemma 1. The product set chosen there has positive invariant area;
this is stated at the end of its construction, before the verification
of (P3)--(P5). The absolute-continuity formula (P5)(b), verified there,
has positive finite Jacobian. Apply the same local statement in the
reverse time direction. Since the invariant collision measure is
smooth and positive away from grazing, disintegration and the two
holonomy formulas identify its null sets on the product rectangle with
those of the product of the transversal measures. These measures are
nonatomic, being equivalent to restricted arc length on transversals.
Restriction to compact positive-measure subsets gives the stated form.
No assertion that $\mathcal P$ contains an open set is needed.

For completeness, contraction here means contraction of the actual
curves, not just their endpoints in a symbolic metric. The billiard
estimates of \cite[Section 8.2C]{Young} give exponential contraction of
the stable arcs in its adapted metric and
$\operatorname{length}(\gamma)\le C\sqrt{\operatorname{length}_{\rm ad}(\gamma)}$
for a stable or unstable curve. Thus, for fixed $R$ and the chosen product set, the ordinary lengths
of these iterated arcs are bounded by $C_{\mathcal P}\theta_{\mathcal P}^n$
for some $C_{\mathcal P}<\infty$ and $0<\theta_{\mathcal P}<1$.
In particular they tend to zero. No uniform choice in $R$ is made. Passing from $(\alpha,\varphi)$ to $(\alpha,p)$ uses only
$p=\sin\varphi$, whose derivative is bounded; no inverse change at
grazing is made. Homogeneous local curves do not cross the relevant
finite-iterate collision singularities. The label $\kappa_R$ is
locally constant on each such physical branch, which proves the label
assertion.

The hypotheses of that construction hold here: the disks are disjoint
with separation at least $3/50$, their boundaries are smooth with
strictly positive curvature, the free domain is connected, and the
horizon is finite. These are properties of the physical Euclidean
table. The same local construction works on its flat triangular torus;
alternatively the rectangular two-scatterer cover in
Lemma~\ref{lem:collision-spectral-input} gives the local product set in
a covering chart. There is no affine change of the reflection law.

Finally choose a product support point and restrict each transversal
to a sufficiently short positive-measure interval around it. Continuity
of the local product map and the transverse graph cones put the
connecting arcs into a small coordinate box. Unique intersections of
the graph curves make the four distinct corners bound a simple Jordan
domain. Nonatomicity excludes repeated transversal coordinates almost
everywhere. This also proves the last assertion for product sets with
gaps.
\end{proof}


\paragraph{Conditional pairs in product coordinates.}
To fix the measure conventions, write $\mu^u=m^u$, $\mu^s=m^s$,
$\mu=\mu^u\otimes\mu^s$ and
\[
 \dd m_{\mathcal P}=\rho(u,s)\dd\mu^u(u)\dd\mu^s(s),\qquad
 m_{\mathcal P}=\nu(\mathcal P)^{-1}\nu|_{\mathcal P}.
\]
Here $0<\rho<\infty$ almost everywhere and $\int\rho\dd\mu=1$.
Put $\rho_u(u)=\int\rho(u,s)\dd\mu^s(s)$ and
$\rho_s(s)=\int\rho(u,s)\dd\mu^u(u)$. The stable and unstable pair
probabilities are explicitly
\begin{align}\label{eq:explicit-conditional-pair-densities}
 \dd\mathfrak m_s(u,s_0,s_1)
 &=\frac{\rho(u,s_0)\rho(u,s_1)}{\rho_u(u)}
           \dd\mu^u(u)\dd\mu^s(s_0)\dd\mu^s(s_1),\\
 \dd\mathfrak m_u(u_0,u_1,s)
 &=\frac{\rho(u_0,s)\rho(u_1,s)}{\rho_s(s)}
           \dd\mu^u(u_0)\dd\mu^u(u_1)\dd\mu^s(s).\notag
\end{align}
The denominators are positive and finite almost everywhere. Integration
of either unused coordinate shows that both marginals are
$m_{\mathcal P}$. The displayed positive finite densities also give
exactly the triple-product measure equivalences used below, without
a boundedness assertion for $\rho$ or its reciprocal.

\subsection{The physical action one-form}
On the collision cylinder define
\begin{equation}\label{eq:collision-action-form}
 \vartheta_R=Rp\,\dd\alpha,\qquad
 \dd\vartheta_R=R\,\dd p\wedge\dd\alpha.
\end{equation}
The one-form is globally defined on the cylinder and bounded in both
the flat and the original collision coordinates.

\begin{lemma}[Exact action on a regular branch]
\label{lem:exact-collision-action}
On every regular collision branch,
\begin{equation}\label{eq:exact-collision-action}
 \dd\tau_R=T_R^*\vartheta_R-\vartheta_R.
\end{equation}
Consequently, if an arc $\gamma$ from $x$ to $y$ stays on regular
branches through time $n$, then
\begin{equation}\label{eq:action-telescope-arc}
 S_{n,R}\tau_R(y)-S_{n,R}\tau_R(x)
       =\int_{T_R^n\gamma}\vartheta_R-\int_\gamma\vartheta_R.
\end{equation}
For backward regular arcs the corresponding identity is
\begin{equation}\label{eq:backward-action-arc}
 S_{n,R}\tau_R(T_R^{-n}y)-S_{n,R}\tau_R(T_R^{-n}x)
       =\int_\gamma\vartheta_R-\int_{T_R^{-n}\gamma}\vartheta_R.
\end{equation}
\end{lemma}
\begin{proof}
Lift the flight to its actual initial disk and the next disk with fixed
label $d$. Its endpoints are $q=Rn_\alpha$ and
$q'=d+Rn_{\alpha'}$. If $v$ is its unit velocity, differentiation of
its Euclidean length gives
$\dd\tau_R=v\cdot\dd q'-v\cdot\dd q$.
The first tangential component is $p$, and specular reflection preserves
the tangential component at arrival, which is $p'$ in the outgoing
coordinates of $T_Rx$. Hence
$\dd\tau_R=Rp'\dd\alpha'-Rp\dd\alpha$. The lattice label is constant
on the branch and contributes no derivative. This proves
\eqref{eq:exact-collision-action}. Summing its pullbacks and integrating
over the arc proves \eqref{eq:action-telescope-arc}; applying that
identity to $T_R^{-n}\gamma$ proves \eqref{eq:backward-action-arc}.
\end{proof}

\begin{proposition}[A measurable physical phase has zero roof frequency]
\label{prop:measurable-roof-rigidity}
Suppose $q:M\to\mathbb S^1$ is measurable and, for constants
$u\in\R^2$ and $s,t\in\R$,
\begin{equation}\label{eq:physical-measurable-phase}
 q\circ T_R=\exp\{i(u\cdot\kappa_R+s+t\tau_R)\}\,q.
\end{equation}
Then $t=0$.
\end{proposition}
\begin{proof}
Discard an invariant null set so that every finite iterate of the phase
identity holds at the remaining points. Use the product set and the
probability $m$ of Lemma~\ref{lem:regular-collision-product}.
Disintegrate $m$ over its stable curves and, conditionally on the curve,
draw two independent points $x,y$. Denote this pair probability by
$\mathfrak m_s$. Both of its marginals are exactly $m$; in particular
they are bounded by $\nu(\mathcal P)^{-1}\nu$.

The following argument justifies endpoint agreement for a merely
measurable $q$. Let $q_\epsilon$ be a smooth approximation in $L^1(\nu)$,
with supremum norm at most one. Invariance of $\nu$ and the marginal
bound give, for every $n$,
\begin{align}\label{eq:conditional-pair-approximation}
 &\int |q(T_R^ny)-q(T_R^nx)|\dd\mathfrak m_s(x,y)\notag\\
 &\quad\le \frac{2}{\nu(\mathcal P)}\|q-q_\epsilon\|_{L^1(\nu)}
       +\int|q_\epsilon(T_R^ny)-q_\epsilon(T_R^nx)|\dd\mathfrak m_s(x,y).
\end{align}
For fixed $\epsilon$, the last term tends to zero by contraction and
bounded convergence. Then let $\epsilon$ tend to zero. Thus the left
side tends to zero. This uses bounded marginals, not invariance of the
pair probability or a claim about every stable curve.

Let $\gamma_s(x,y)$ be the stable arc oriented from $x$ to $y$.
Its forward labels agree at all times. Iteration of
\eqref{eq:physical-measurable-phase} and
\eqref{eq:action-telescope-arc} gives
\[
 \frac{q(T_R^ny)}{q(T_R^nx)}
 =\exp\!\left(it\left[\int_{T_R^n\gamma_s(x,y)}\vartheta_R
                         -\int_{\gamma_s(x,y)}\vartheta_R\right]\right)
                         \frac{q(y)}{q(x)}.
\]
The first arc integral tends to zero, because $\vartheta_R$ is bounded
and the arc length tends to zero. Since $|q|=1$, the left side minus one
has $L^1(\mathfrak m_s)$ norm tending to zero by
\eqref{eq:conditional-pair-approximation}. The right side has the
pointwise limit just displayed. Fatou's lemma therefore gives
\begin{equation}\label{eq:measurable-stable-holonomy}
 \frac{q(y)}{q(x)}
       =\exp\!\left(it\int_{\gamma_s(x,y)}\vartheta_R\right)
       \quad\text{for }\mathfrak m_s\text{-almost every }(x,y).
\end{equation}
Disintegrating instead over unstable curves and using backward iterates
in \eqref{eq:conditional-pair-approximation} gives, by
\eqref{eq:backward-action-arc},
\begin{equation}\label{eq:measurable-unstable-holonomy}
 \frac{q(y)}{q(x)}
       =\exp\!\left(it\int_{\gamma_u(x,y)}\vartheta_R\right)
       \quad\text{for almost every unstable pair}.
\end{equation}
The sign is the same in the two formulas: the backward action in
\eqref{eq:backward-action-arc} has limiting value
$\int_{\gamma_u(x,y)}\vartheta_R$.

We now use the measure-class part of the local product input. The stable
and unstable pair probabilities are equivalent, in product coordinates,
to the corresponding three-coordinate product measures. Fubini's theorem
therefore permits \eqref{eq:measurable-stable-holonomy} and
\eqref{eq:measurable-unstable-holonomy} to be multiplied around almost
every four-corner product quadrilateral $Q$. The oriented corner order is
\begin{gather*}
 x_{00}=\Psi(u_0,s_0)\ \longrightarrow\
 x_{01}=\Psi(u_0,s_1)\ \longrightarrow\ x_{11}=\Psi(u_1,s_1),\\
 x_{11}\ \longrightarrow\ x_{10}=\Psi(u_1,s_0)
            \ \longrightarrow\ x_{00}.
\end{gather*}
The edges are stable, unstable, stable with reversed $s$-orientation,
and unstable with reversed $u$-orientation. This specifies the boundary
orientation used in Stokes' formula; its sign need not be fixed when
only the absolute area is used. Their left sides telescope, so
\begin{equation}\label{eq:measurable-area-obstruction}
 \exp\!\left(it\oint_{\partial Q}\vartheta_R\right)=1.
\end{equation}
This assertion is on a positive-measure family of quadrilaterals; it does
not assign values to $q$ on a preselected closed orbit.

Suppose $t\ne0$. Restrict to a positive-measure product subrectangle
small enough that every such quadrilateral has coordinate area less
than $2\pi/(R|t|)$. Almost every quadrilateral is a nondegenerate Jordan
domain. Stokes' formula and \eqref{eq:collision-action-form} give
\[
 0<\left|t\oint_{\partial Q}\vartheta_R\right|
       =|t|R\operatorname{area}_{\alpha,p}(Q)<2\pi,
\]
contradicting \eqref{eq:measurable-area-obstruction}. Hence $t=0$.
\end{proof}

\subsection{Rotations and the constant collision phase}
Let $S_j(\alpha,p)=(\alpha+j\pi/3,p)$, $j\in\Z$. In lattice
coordinates rotation by $\pi/3$ is
\begin{equation}\label{eq:triangular-rotation-matrix}
 A=\begin{pmatrix}0&-1\\1&1\end{pmatrix},\qquad
 A^3=-I_2,\qquad I_2+A^2+A^4=0.
\end{equation}
The actual physical symmetries give
\begin{equation}\label{eq:collision-rotation-identities}
 T_RS_j=S_jT_R,\qquad
 \kappa_R\circ S_j=A^j\kappa_R,\qquad
 \tau_R\circ S_j=\tau_R,\qquad (S_j)_*\nu=\nu.
\end{equation}
They follow by rotating both endpoints, the velocity and the selected
next disk. The first admissible physical root remains the first root.
No invariance of $Y_R^*$ under these rotations is asserted.

\begin{theorem}[Rigidity of the roof and constant parts of a phase]
\label{thm:physical-phase-rigidity}
If a measurable circle-valued $q$ satisfies
\eqref{eq:physical-measurable-phase}, then
\begin{equation}\label{eq:physical-phase-conclusion}
 t=0,\qquad s\in2\pi\Z.
\end{equation}
\end{theorem}
\begin{proof}
The proposition gives $t=0$. Define
\[
 Q_2=q\,(q\circ S_3),\qquad
 Q_3=q\,(q\circ S_2)(q\circ S_4).
\]
By \eqref{eq:triangular-rotation-matrix} and
\eqref{eq:collision-rotation-identities}, the lattice phases cancel
exactly and
\[
 Q_2\circ T_R=e^{2is}Q_2,\qquad Q_3\circ T_R=e^{3is}Q_3.
\]
Both functions have modulus one. The mixing collision map has no
nontrivial circle eigenvalue, by the $L^2$ argument in the proof of
Corollary~\ref{cor:positive-count-variance}. Therefore
$e^{2is}=e^{3is}=1$, and their quotient gives $e^{is}=1$.
\end{proof}

The theorem is qualitative. It excludes exact measurable physical
phases with nonzero roof frequency or a nonintegral constant phase.
It is not a norm estimate for approximate spectral vectors, and supplies
no decay rate for a high-frequency collision or induced resolvent.

\section{Uniform nondegeneracy of the joint return record}
\label{sec:joint-nondegeneracy}
We now combine measurable phase rigidity with the exact integer-valued
section compensation. The proof treats every joint direction of the
same four-coordinate record. It does not pass from the scalar count
variance to a determinant by an inequality.

\subsection{A finite-cover obstruction for displacement}
\begin{lemma}[No real spatial collision coboundary]
\label{lem:no-spatial-coboundary}
If $w\in\R^2$ and $b\in L^2(\nu;\R)$ satisfy
\begin{equation}\label{eq:spatial-coboundary}
 w\cdot\kappa_R=b-b\circ T_R,
\end{equation}
then $w=0$.
\end{lemma}
\begin{proof}
Regard $w$ as a row vector and suppose $w\ne0$.
Composition with $S_1$, using \eqref{eq:collision-rotation-identities},
gives $(wA)\cdot\kappa_R=b\circ S_1-(b\circ S_1)\circ T_R$.
The matrix with rows $w=(w_1,w_2)$ and
$wA=(w_2,-w_1+w_2)$ has determinant
\begin{equation}\label{eq:spatial-rotation-determinant}
 -(w_1^2-w_1w_2+w_2^2),
 \qquad w_1^2-w_1w_2+w_2^2\ge\tfrac12|w|^2>0.
\end{equation}
Solving these two real identities gives a vector-valued
$H\in L^2(\nu;\R^2)$ with $\kappa_R=H-H\circ T_R$.

Consider the physical billiard on the flat torus $\R^2/(2\Lambda)$.
There are four equal scatterers, labeled by
$\ell\in\Lambda/(2\Lambda)\simeq(\Z/2\Z)^2$. In collision
coordinates its map and invariant probability are exactly
\begin{equation}\label{eq:four-sheet-collision}
 \widetilde T_R(x,\ell)
       =(T_Rx,\ell+\kappa_R(x)\bmod2),\qquad
 \widetilde\nu=\nu\otimes\operatorname{unif}((\Z/2\Z)^2).
\end{equation}
This map is ergodic. The planar lift has precisely the original disk
array, so a free flight in the cover has the original uniform length
bound. The enlarged torus contains finitely many pairwise disjoint
embedded disks; their complement is connected. Thus the enlarged table
has smooth strictly convex scatterers, connected free space and finite
horizon; the finite-horizon dispersing-billiard theorem
\cite[Section 8.1, Theorem 6]{Young} applies. To work only with a
rectangular fundamental domain, take the physical cover with periods
$(2,0)$ and $(0,2\sqrt3)$. It has eight scatterers and satisfies the
same hypotheses. Its quotient by the deck group is precisely the
four-scatterer table above, so its ergodicity implies that of
\eqref{eq:four-sheet-collision}. Neither argument changes Euclidean
length or specular reflection.

However,
\[
 F(x,\ell)=\exp\{i\pi(H_1(x)+\ell_1)\}
\]
is invariant under \eqref{eq:four-sheet-collision}: the integer
increment $\kappa_{R,1}$ cancels $H_1\circ T_R-H_1$, and reduction
of $\ell_1$ modulo two does not alter the exponential. The function
has modulus one and integral zero by averaging over $\ell_1$.
It cannot be almost everywhere constant, contradicting ergodicity.
This proves the lemma. No value of $b$ or $H$ at a periodic point is
used.
\end{proof}

\subsection{Removing the section contribution}
\begin{theorem}[Uniform positive definiteness]
\label{thm:joint-uniform-nondegeneracy}
There exist constants $0<d_0\le D_0<\infty$ such that
\begin{equation}\label{eq:joint-uniform-nondegeneracy}
 d_0I_4\le D_R\le D_0I_4\qquad(R\in I).
\end{equation}
Equivalently the actual $L^2$ coboundaries in
Theorem~\ref{thm:gaussian-kernel-coboundary} have only the zero joint
direction. The discontinuity of $\eta_R=\one_{Y_R^*}$ is allowed.
\end{theorem}
\begin{proof}
Fix $R$ and suppose $v^{\mathsf T}D_Rv=0$, where $v=(w,a,b)$ with
$w\in\R^2$. The collision-clock coboundary characterization gives a
real $B\in L^2(\nu)$ such that
\begin{equation}\label{eq:joint-zero-coboundary}
 w\cdot\kappa_R+a+b\tau_R-\zeta\eta_R=B-B\circ T_R,
 \qquad \zeta=\frac{a+b\bar\tau_R}{c_*}.
\end{equation}
All terms are the physical single-collision quantities. We separate
whether the coefficient of the section indicator vanishes.

If $\zeta\ne0$, set $q=\exp(2\pi i B/\zeta)$. Since $\eta_R$ is
integer-valued, exponentiation removes it exactly and yields
\[
 q\circ T_R
 =\exp\!\left(-\frac{2\pi i}{\zeta}
          (w\cdot\kappa_R+a+b\tau_R)\right)q.
\]
Theorem~\ref{thm:physical-phase-rigidity} implies $b=0$ and
$-2\pi a/\zeta\in2\pi\Z$. But now
$\zeta=a/c_*\ne0$, so $a/\zeta=c_*$. Since $0<c_*<1$, this is
impossible. Notice that the rotations are applied only after the
section indicator has disappeared. Their use requires no symmetry of
the actual return section.

If $\zeta=0$, set $q=\exp(iB)$ in
\eqref{eq:joint-zero-coboundary}. Proposition~\ref{prop:measurable-roof-rigidity}
implies $b=0$. The definition of $\zeta$ then gives $a=0$.
The remaining identity is the real spatial coboundary
$w\cdot\kappa_R=B-B\circ T_R$, and
Lemma~\ref{lem:no-spatial-coboundary} gives $w=0$.

Thus $v=0$ in all cases. The symmetric matrix $D_R$ is positive definite
for each $R$. Theorem~\ref{thm:collision-covariance} gives its continuity
on the compact interval $I$. The smallest eigenvalue is consequently
bounded below by a positive number, and the largest is bounded above.
This proves \eqref{eq:joint-uniform-nondegeneracy}. Uniformity is obtained
at this last step, not by assuming uniform constants for measurable
phase regularity or for the product set.
\end{proof}

\begin{remark}[The role of the section]
The algebra in \eqref{eq:joint-zero-coboundary} uses only that the
section indicator is integer-valued and that its mean lies strictly
between zero and one. The two rotation products force an integral
constant phase without an arithmetic assumption on this mean. The
specific section geometry is used earlier to establish the bounded
collision covariance and its $L^2$ kernel characterization. Thus the
present proof neither evaluates a discontinuous section weight on
stable arcs nor replaces it by a smoother section.
\end{remark}

\begin{corollary}[Actual covariances and the physical-time projection]
\label{cor:actual-uniform-ellipticity}
For all sufficiently large $n$, uniformly in $R$,
\begin{equation}\label{eq:finite-n-joint-ellipticity}
 \frac1n\int U_{n,R}\otimes U_{n,R}\dd\nu_R^*
       \ge\frac{d_0}{2}I_4.
\end{equation}
The physical-time fluctuation covariance $\mathcal V_R$ in
Corollary~\ref{cor:physical-functional-gaussian} is also uniformly
positive definite on its three-dimensional space.
\end{corollary}
\begin{proof}
Equation~\eqref{eq:actual-covariance-rate} and
\eqref{eq:joint-uniform-nondegeneracy} give
\eqref{eq:finite-n-joint-ellipticity}. For the physical covariance recall
$\Gamma_R=c_*D_R$ and
$\mathcal V_R=\bar\tau_R^{-1}B_R\Gamma_RB_R^{\mathsf T}$, where
$B_R(x_1,x_2,x_3,x_4)=(x_1,x_2,x_3-x_4/\bar\tau_R)$.
For $z\in\R^3$, $|B_R^{\mathsf T}z|\ge|z|$. Hence
$z^{\mathsf T}\mathcal V_Rz\ge
 c_*d_0(\max_I\bar\tau_R)^{-1}|z|^2$. All constants are positive.
\end{proof}

\subsection{The Gaussian part of exact local inversion}
\begin{corollary}[Uniform Gaussian inversion and the remaining terms]
\label{cor:elliptic-central-inversion}
The centered Gaussian density associated with the actual return record is
\begin{equation}\label{eq:uniform-joint-gaussian-density}
 g_R(\xi)=\frac{1}{(2\pi)^2\sqrt{\det D_R}}
              \exp\!\left(-\tfrac12\xi^{\mathsf T}D_R^{-1}\xi\right).
\end{equation}
It is defined for every $R\in I$. There are $C,c>0$ such that
\begin{equation}\label{eq:elliptic-gaussian-tail}
 \sup_R\int_{|v|>n^{1/200}}
          e^{-v^{\mathsf T}D_Rv/2}\dd v
       \le C\exp(-c n^{1/100}).
\end{equation}
In Corollary~\ref{cor:raw-from-growing-arc}, uniform positive definiteness
is therefore a proved conclusion for this physical family, not an
additional hypothesis. Under its stated residual integrability,
complementary-transform estimate and local-edge conditions, the raw
mixed-density conclusion follows with the density
\eqref{eq:uniform-joint-gaussian-density}.
\end{corollary}
\begin{proof}
Positive definiteness gives the Gaussian inversion formula. In four
dimensions, integration in polar coordinates gives, for $A>0$,
\[
 \int_{|v|>A}e^{-d_0|v|^2/2}\dd v
   =2\pi^2\left(\frac{A^2}{d_0}+\frac{2}{d_0^2}\right)
                         e^{-d_0A^2/2}.
\]
Absorb the polynomial into a weaker exponential and put $A=n^{1/200}$.
The last assertion follows from the exact decomposition and the proved
central comparison in Corollary~\ref{cor:raw-from-growing-arc}.
\end{proof}

The two Fourier tails in this last argument concern different objects.
Equation~\eqref{eq:elliptic-gaussian-tail} estimates the limiting Gaussian.
It does not estimate the complementary transform of the physical raw
record or of its edge-subtracted residual. The proved central band is
still $|v|\le2n^{1/200}$, equivalently
$|\omega|\le2n^{-99/200}$. The annulus beyond it, the full critical and
singular edge extraction, and the weighted tails and relative event
comparison for exact physical conditioning remain the analytical tasks
in the raw mixed-density theorem. The covariance obstruction has been
removed without changing that theorem's record or target measure.

\section{Complete physical phase arithmetic}
\label{sec:complete-phase-arithmetic}
The roof and constant parts of a physical collision phase have already
been determined in Theorem~\ref{thm:physical-phase-rigidity}. We now
determine its two spatial parts. The argument uses an atom of the phase
on a positive-measure product set, not a value at a periodic point.
It applies to arbitrary real spatial frequencies, including irrational
ones.

Write
\[
 \Omega=\T^2\times\T\times\R,\qquad
 \phi_{R,\omega}=u\cdot\kappa_R+s+t\tau_R,
 \quad \omega=(u,s,t).
\]
The three torus coordinates have period $2\pi$. The origin of $\Omega$
therefore means $u\in2\pi\Z^2$, $s\in2\pi\Z$, $t=0$.

\begin{lemma}[Separation of a proper integer subgroup]
\label{lem:integer-subgroup-separation}
If $L$ is a proper subgroup of $\Z^d$, then there are a prime $p$ and a
nonzero $a\in(\Z/p\Z)^d$ with $a\cdot L=0$ modulo $p$.
\end{lemma}
\begin{proof}
An integer subgroup has an integer basis. If its rank is less than $d$,
choose a nonzero integer row annihilating it and a prime not dividing
all coordinates of that row. If it has full rank, its index is the
absolute determinant of a basis matrix. This integer exceeds one.
Modulo a prime divisor of the index, the basis matrix is singular and
has a nonzero left null vector. These two constructions prove the
claim without a rationality assumption on the character that might
have defined $L$.
\end{proof}

\begin{theorem}[Complete measurable physical phase rigidity]
\label{thm:complete-physical-phase}
For every $R\in I$, a measurable $q:M\to\mathbb S^1$ satisfies
\begin{equation}\label{eq:complete-physical-phase}
 q\circ T_R=e^{i\phi_{R,\omega}}q
\end{equation}
if and only if $\omega=0$ in $\Omega$ and $q$ is constant almost
everywhere.
\end{theorem}
\begin{proof}
Theorem~\ref{thm:physical-phase-rigidity} gives $t=0$ and
$s\in2\pi\Z$. Both holonomy formulas
\eqref{eq:measurable-stable-holonomy} and
\eqref{eq:measurable-unstable-holonomy} consequently have right side
one. In the product coordinates of
Lemma~\ref{lem:regular-collision-product}, the first says that $q$
is almost everywhere a function only of the unstable transversal
coordinate, and the second that it is almost everywhere a function
only of the stable transversal coordinate. Product measure equivalence
and Fubini imply that $q=c$ almost everywhere on $\mathcal P$, for
one $c\in\mathbb S^1$.

The saturation $\bigcup_{j\ge0}T_R^j\mathcal P$ has full measure by
ergodicity and recurrence. Iterating \eqref{eq:complete-physical-phase}
from $\mathcal P$ therefore shows that
\begin{equation}\label{eq:countable-phase-range}
 q(x)/c\in\{e^{iu\cdot k}:k\in\Z^2\}
       \quad\hbox{for almost every }x.
\end{equation}
Let $L=\{k\in\Z^2:u\cdot k\in2\pi\Z\}$. The map
$[k]\mapsto e^{iu\cdot k}$ identifies the countable group $\Z^2/L$
with the displayed range. Its inverse on that countable range is
measurable. Thus \eqref{eq:countable-phase-range} defines a measurable
$h:M\to\Z^2/L$ with
\begin{equation}\label{eq:countable-phase-cocycle}
 h(T_Rx)=h(x)+[\kappa_R(x)].
\end{equation}
This is a countable-valued lift, not a real logarithm with a chosen
branch. No integrability of the lift is required.

If $L\ne\Z^2$, apply Lemma~\ref{lem:integer-subgroup-separation}.
The resulting nonzero row $a$ defines a homomorphism
$\chi:\Z^2/L\to\Z/p\Z$. On the actual physical billiard modulo
$p\Lambda$ the collision map is
\[
 \widetilde T_R(x,\ell)
 =(T_Rx,\ell+\kappa_R(x)\bmod p),\qquad
 \ell\in(\Z/p\Z)^2.
\]
It preserves $\nu$ times the uniform sheet probability and is ergodic.
Indeed its lift to the plane has exactly the original disk array, so
its free-flight bound is unchanged. Its finitely many disjoint circular
scatterers have positive curvature, and the complement of these disks
in the torus is connected. Thus the finite-horizon billiard theorem
\cite[Section 8.1]{Young} applies. Equivalently use the rectangular
Euclidean cover with periods $(p,0)$ and $(0,p\sqrt3)$, with $2p^2$
scatterers, and then its quotient. No affine deformation or independent
sheet dynamics is introduced.

Equation~\eqref{eq:countable-phase-cocycle} makes
\[
 \exp\!\left(\frac{2\pi i}{p}
             [a\cdot\ell-\chi(h(x))]\right)
\]
invariant under this collision map. The brackets in this formula are
interpreted modulo $p$. The function has modulus one and mean zero
upon averaging over $\ell$, since $a\ne0$. This contradicts ergodicity.
Hence $L=\Z^2$, or $u\in2\pi\Z^2$. The original phase equation now
says that $q$ is invariant, so it is constant. The converse is immediate.
\end{proof}

\subsection{Finite positive-measure return witnesses}
This consequence will replace a nonquantitative appeal to absence of
exact eigenfunctions in the next section. The auxiliary set below is
not the return section $Y_R^*$ of the physical record.

\begin{lemma}[A finite generating family of product-set returns]
\label{lem:finite-return-witnesses}
Fix $R$ and a positive-measure regular product set $\mathcal P$.
Let $r_{\mathcal P}$ be its actual first-return time under $T_R$.
There are finitely many pairs $(n_i,k_i)\in\Z_{>0}\times\Z^2$ such that
\begin{equation}\label{eq:positive-return-witness}
 E_i=\{x\in\mathcal P:r_{\mathcal P}(x)=n_i,
                  S_{n_i,R}\kappa_R(x)=k_i\},\qquad \nu(E_i)>0,
\end{equation}
and the vectors $g_i=(k_{i,1},k_{i,2},n_i)$ generate $\Z^3$ over $\Z$.
In particular, integers $b_{ji}$ can be chosen with
$e_j=\sum_i b_{ji}g_i$ for $1\le j\le3$.
\end{lemma}
\begin{proof}
Let $L$ be the subgroup generated by all positive-measure cells in
\eqref{eq:positive-return-witness}. The cells form a countable partition
of $\mathcal P$ modulo a null set. If $L$ were proper,
Lemma~\ref{lem:integer-subgroup-separation} would give a nontrivial
finite character $\theta=(u,s)\in\T^3$ annihilating every cell record.
For $\phi=u\cdot\kappa_R+s$ this means
$e^{iS_{r_{\mathcal P},R}\phi}=1$ on $\mathcal P$ almost everywhere.

The first-return tower over $\mathcal P$ represents almost every point
of the invertible ergodic collision system uniquely as $T_R^j y$,
where $y\in\mathcal P$ and $0\le j<r_{\mathcal P}(y)$.
Set $q(T_R^j y)=e^{iS_{j,R}\phi(y)}$ on these measurable tower levels.
The preceding return identity verifies the phase equation also at
the top of each level. This is an exact measurable physical phase at
$(u,s,0)\ne0$, contradicting
Theorem~\ref{thm:complete-physical-phase}. Thus $L=\Z^3$.
Writing the three standard basis vectors as integer combinations uses
only finitely many cell records, which gives the asserted family.
The witnesses have positive measure; they are not isolated orbits.
\end{proof}

\section{Quantitative defects of regular collision phases}
\label{sec:quantitative-phase-defects}
Absence of exact phases does not itself control approximate vectors.
We give a quantitative version on the space of bounded functions of
bounded variation. The phase is reconstructed from a complex vector;
no pointwise lower bound for its modulus is assumed. The logarithmic
bounds in this section are at a fixed radius and on a fixed compact
frequency set. A separate compactness argument gives a parameter-uniform
bound when the regularity budget is fixed.

For $\omega\in\Omega$ put
\begin{equation}\label{eq:collision-phase-defects}
 \begin{split}
 d_{R,\omega}(q)&=\|q\circ T_R-e^{i\phi_{R,\omega}}q\|_{L^1(\nu)},
                      \qquad (|q|=1),\\
 \mathcal E_{R,\omega}(f)&=\|f\circ T_R-e^{i\phi_{R,\omega}}f\|_{L^2(\nu)}.
 \end{split}
\end{equation}
The variation is always in the original collision coordinates. In
particular it is neither an anisotropic distribution norm nor a norm
of a many-return pushforward density.

\subsection{Finite-time measurable holonomy with a defect}
Use the product coordinates and the pair probabilities
\eqref{eq:explicit-conditional-pair-densities}. For a stable or unstable
oriented arc from $x$ to $y$ define
\[
 H_{s/u}(x,y)=
 \left|\frac{q(y)}{q(x)}-
            \exp\!\left(it\int_{\gamma_{s/u}(x,y)}\vartheta_R\right)\right|.
\]
These are bounded measurable functions of pairs. Null sets are taken
with respect to the indicated pair probabilities.

\begin{lemma}[A finite-time holonomy defect]
\label{lem:finite-holonomy-defect}
Fix a regular product set $\mathcal P$ at radius $R$. There are
$C_{\mathcal P}<\infty$ and $0<\theta_{\mathcal P}<1$ such that,
for $m\ge1$, $0<\epsilon<1/2$, and $q:M\to\mathbb S^1$ with
$\|q\|_{\mathrm{BV}}\le H$, $H\ge2$,
\begin{equation}\label{eq:finite-holonomy-defect}
 \int H_{s/u}\dd\mathfrak m_{s/u}
 \le C_{\mathcal P}\left[
 m d_{R,\omega}(q)+H\epsilon+
                  (\epsilon^{-1}+|t|)\theta_{\mathcal P}^{m}\right].
\end{equation}
The same expression controls both signs of time. The constants do not
depend on $u,s,t,H,m,\epsilon$.
\end{lemma}
\begin{proof}
Iterating the approximate phase equation along one orbit introduces
an error at most the sum of its $m$ one-step absolute defects. The
stable-pair version therefore has errors at its two endpoints. Both
pair marginals equal $m_{\mathcal P}=\nu(\mathcal P)^{-1}\nu|_{\mathcal P}$.
Invariance of $\nu$ bounds the integrated sum of endpoint defects by
$2m\nu(\mathcal P)^{-1}d_{R,\omega}(q)$.

Let $q_\epsilon=\mathsf S_\epsilon q$. Positivity of smoothing gives
$|q_\epsilon|\le1$, and
$\|q-q_\epsilon\|_1\le C H\epsilon$,
$\|q_\epsilon\|_{C^1(\alpha,p)}\le C\epsilon^{-1}$.
The two endpoint approximation errors are bounded at every iterate
by $2\nu(\mathcal P)^{-1}\|q-q_\epsilon\|_1$.
The actual stable arc lengths are at most
$C_{\mathcal P}\theta_{\mathcal P}^m$ after $m$ forward iterates,
by the contraction estimate in the proof of
Lemma~\ref{lem:regular-collision-product}. Thus the smoothed endpoint
difference costs $C_{\mathcal P}\epsilon^{-1}\theta_{\mathcal P}^m$.
The terminal action integral in \eqref{eq:action-telescope-arc} costs
at most $C_{\mathcal P}|t|\theta_{\mathcal P}^m$.
The lattice and constant phases cancel along the stable pair exactly.
These estimates prove \eqref{eq:finite-holonomy-defect} for stable arcs.

For an unstable pair telescope backwards. Its defect is the sum of the
one-step defects at $T_R^{-j}x,T_R^{-j}y$, $1\le j\le m$; the same
marginal domination applies by invertibility and invariance. Use
\eqref{eq:backward-action-arc} and backward contraction. The holonomy
sign is again positive, as in
\eqref{eq:measurable-unstable-holonomy}. This proves the other case.
No invariant pair measure and no independent orbit assumption is used.
\end{proof}

We make explicit how a measure-class assertion can be used
quantitatively without replacing it by an unproved bounded-density
assertion. Write $\mu=\mu^u\otimes\mu^s$ and
$\dd m_{\mathcal P}=\rho\dd\mu$ in the product coordinates.
Let $\lambda_s=\mu^u\otimes\mu^s\otimes\mu^s$ and
$\lambda_u=\mu^u\otimes\mu^u\otimes\mu^s$ be the reference triple
probabilities, and set
\[
 w_s=\frac{\dd\lambda_s}{\dd\mathfrak m_s},\qquad
 w_u=\frac{\dd\lambda_u}{\dd\mathfrak m_u}.
\]
These derivatives are positive and finite almost everywhere. For
$0\le F\le2$ and either pair type,
\begin{equation}\label{eq:pair-density-truncation}
 \int F\dd\lambda_{s/u}
 \le K\int F\dd\mathfrak m_{s/u}
            +2\lambda_{s/u}\{w_{s/u}>K\}.
\end{equation}
The last probability tends to zero as $K\to\infty$. Its value is part
of the constant below; no uniform pointwise bound on $\rho$ is needed.

\begin{lemma}[An almost constant phase on the product set]
\label{lem:almost-constant-product-phase}
For every $\eta>0$ there is a finite $C_{\mathcal P,\eta}$ such that
one can choose $c\in\mathbb S^1$ with
\begin{align}\label{eq:almost-constant-product-phase}
 \int_{\mathcal P}|q-c|\dd\nu
 \le\eta+C_{\mathcal P,\eta}\bigl[
 |t|+m d_{R,\omega}(q)+H\epsilon
              +(\epsilon^{-1}+|t|)\theta_{\mathcal P}^m\bigr].
\end{align}
This holds for the same $q,H,m,\epsilon$ as in the preceding lemma.
\end{lemma}
\begin{proof}
The action integrals on the original arcs are uniformly bounded on
$\mathcal P$. Hence the stable and unstable pair averages of
$|q(y)-q(x)|$ are bounded by the right side of
\eqref{eq:finite-holonomy-defect} plus $C_{\mathcal P}|t|$.

Here are the density choices needed to pass to two arbitrary product
points. Choose $J$ so that
$\int_{\{\rho>J\}}\rho\dd\mu$ is less than a prescribed small number.
Choose $K$ so that both tail probabilities in
\eqref{eq:pair-density-truncation} are less than that number divided
by $J$. Joining $(u,s)$ to $(u,s')$ and then to $(u',s')$, Fubini
bounds the $\mu\otimes\mu$ average of $|q(x)-q(y)|$ by the sum of
the two reference triple averages. Thus there is a value $c=q(y)$
for which the $\mu$ average of $|q-c|$ is at most that sum. Finally
\[
 \int|q-c|\dd m_{\mathcal P}
 \le J\int|q-c|\dd\mu+2\int_{\{\rho>J\}}\rho\dd\mu.
\]
Take the prescribed small numbers sufficiently small in terms of
$\eta$ and $\nu(\mathcal P)$, and apply
\eqref{eq:pair-density-truncation}. The remaining multiplier is a
finite function of $J,K,C_{\mathcal P},\nu(\mathcal P)$.
This proves \eqref{eq:almost-constant-product-phase}.
\end{proof}

\subsection{A quantitative roof obstruction}
\begin{proposition}[A logarithmic defect bound on a roof band]
\label{prop:logarithmic-roof-defect}
Fix $R$ and $0<t_0\le T<\infty$. There is $c_{R,t_0,T}>0$ such that
\begin{equation}\label{eq:logarithmic-roof-defect}
 d_{R,(u,s,t)}(q)\ge\frac{c_{R,t_0,T}}{1+\log H}
\end{equation}
whenever $t_0\le|t|\le T$, $|q|=1$, $H\ge2$, and
$\|q\|_{\mathrm{BV}}\le H$. The estimate is independent of $u,s$.
\end{proposition}
\begin{proof}
Restrict a regular product set so that every oriented four-corner
quadrilateral has action $\mathcal A=\oint\vartheta_R$ satisfying
$|\mathcal A|\le\pi/T$. By Stokes' formula and nonatomicity,
$|\mathcal A|>0$ for almost every quadruple under
$\lambda_4=(\mu^u)^2\otimes(\mu^s)^2$.
Choose $a_0>0$ such that
\[
 p_0:=\lambda_4\{|\mathcal A|\ge a_0\}>0.
\]
For this set and the entire roof band,
$|e^{it\mathcal A}-1|\ge2t_0a_0/\pi$.

Each of the four edge marginals of $\lambda_4$ is one of the reference
triple measures. Choose $K$ so that the sum of their probabilities
of $w_{s/u}>K$ is less than $p_0/2$. Intersect the set
$\{|\mathcal A|\ge a_0\}$ with all four complementary conditions;
call the resulting set $\mathcal Q$. It has measure at least $p_0/2$,
and each of its unnormalized edge marginals is bounded by
$K\mathfrak m_s$ or $K\mathfrak m_u$.

Multiply the four oriented approximate holonomies. The product of
$q$-ratios is exactly one, and the error in a product of unit complex
numbers is at most the sum of its four errors. Integration over
$\mathcal Q$ and Lemma~\ref{lem:finite-holonomy-defect} give
\begin{equation}\label{eq:quantitative-area-defect}
 \frac{p_0t_0a_0}{\pi}
 \le4K C_{\mathcal P}\bigl[
 m d_{R,\omega}(q)+H\epsilon
                  +(\epsilon^{-1}+T)\theta_{\mathcal P}^m\bigr].
\end{equation}
All constants have now been selected from the product set, its density
tails and its action distribution, independently of $q$ and $H$.
Choose $\epsilon=c_1/H$, with a fixed sufficiently small $c_1>0$.
Then choose $m=\lceil c_2+c_3\log H\rceil$ so that the final term in
brackets is also less than one quarter of the required lower bound.
The positive contraction exponent makes this possible. Equation
\eqref{eq:quantitative-area-defect} leaves $m d_{R,\omega}(q)\ge c_4>0$.
This proves \eqref{eq:logarithmic-roof-defect}.
\end{proof}

\subsection{All compact nonzero physical frequencies}
\begin{theorem}[A logarithmic phase defect on a compact band]
\label{thm:compact-log-phase-defect}
For a fixed $R\in I$ and a compact set
$\mathcal K\subset\Omega\setminus\{0\}$ there is $c_{R,\mathcal K}>0$
such that, for every $\omega\in\mathcal K$,
\begin{equation}\label{eq:compact-log-phase-defect}
 |q|=1,\quad \|q\|_{\mathrm{BV}}\le H,\quad H\ge2
 \quad\Longrightarrow\quad
 d_{R,\omega}(q)\ge\frac{c_{R,\mathcal K}}{1+\log H}.
\end{equation}
The set includes compact nonzero spatial frequencies with $s=t=0$.
\end{theorem}
\begin{proof}
Set
\[
 \rho(u,s)=|e^{iu_1}-1|+|e^{iu_2}-1|+|e^{is}-1|,
 \qquad \rho_*:=\min_{\mathcal K}(|t|+\rho(u,s))>0.
\]
Let $T=1+\max_{\mathcal K}|t|$. Fix a product set $\mathcal P$ and
the finite witnesses $(n_i,k_i),E_i,b_{ji}$ in
Lemma~\ref{lem:finite-return-witnesses}. Write $p_i=\nu(E_i)>0$ and
$d=d_{R,\omega}(q)$. For a constant $c\in\mathbb S^1$, iteration
of the approximate phase equation on $E_i$ gives
\begin{align}\label{eq:witness-phase-defect}
 p_i|e^{i(u\cdot k_i+s n_i)}-1|
 \le 2\int_{\mathcal P}|q-c|\dd\nu+n_i d
                     +p_i|t|n_i\|\tau_R\|_\infty.
\end{align}
Indeed $T_R^{n_i}E_i\subset\mathcal P$, its contribution to the
endpoint error is bounded by the integral over $\mathcal P$ by
invariance, and the $n_i$ one-step defects integrate to at most $n_i d$.
The roof sum is bounded by $n_i\|\tau_R\|_\infty$.

The integer identities $e_j=\sum_i b_{ji}g_i$ and
$|z^b-1|\le |b||z-1|$ for $|z|=1$, $b\in\Z$, show that finite
constants $A,B,C$ satisfy
\begin{equation}\label{eq:witness-resonance-control}
 \rho(u,s)\le A\int_{\mathcal P}|q-c|\dd\nu+B d+C|t|.
\end{equation}
For example one may take
$A=2\sum_{j,i}|b_{ji}|/p_i$,
$B=\sum_{j,i}|b_{ji}|n_i/p_i$ and
$C=\|\tau_R\|_\infty\sum_{j,i}|b_{ji}|n_i$.
These constants include the actual positive masses of the witnesses.

Choose $\eta>0$ with $A\eta<\rho_*/16$, and use
Lemma~\ref{lem:almost-constant-product-phase} with this $\eta$.
Choose $t_0\in(0,\rho_*/4)$ so small that all terms proportional
to $|t|$ in \eqref{eq:witness-resonance-control}, after substitution
of that lemma, are less than $\rho_*/16$ when $|t|\le t_0$.
Choose $\epsilon=c_5/H$, then $m=\lceil c_6+c_7\log H\rceil$,
so that the two remaining approximation terms contribute less than
$\rho_*/8$. For $|t|\le t_0$, the left side is at least
$3\rho_*/4$. The inequalities thus force $m d\ge c_8>0$.
For $|t|\ge t_0$ use Proposition~\ref{prop:logarithmic-roof-defect}
on $[t_0,T]$. Taking the smaller of the two constants proves the
assertion. All selections precede the choice of $q,H$.
\end{proof}

\subsection{Reconstructing a phase without a lower-modulus assumption}
\begin{lemma}[A bounded-variation circle truncation]
\label{lem:bv-circle-truncation}
For every complex $f\in\mathrm{BV}(M)$ there is a measurable
$q:M\to\mathbb S^1$ such that
\begin{equation}\label{eq:bv-circle-truncation}
 \|q\|_{\mathrm{BV}}\le C(1+\|f\|_{\mathrm{BV}}),\qquad
 \|q-f\|_2\le3\||f|-1\|_2
\end{equation}
whenever the latter norm is finite.
\end{lemma}
\begin{proof}
Put $a=|f|$. The BV chain inequality gives $\TV(a)\le C\TV(f)$.
The coarea formula gives a level $\tau\in[1/4,1/2]$ for which
$E=\{a>\tau\}$ has perimeter at most $4\TV(a)$.
These chain, product and coarea statements are for first distributional
derivatives in the collision coordinates; see \cite[Chapter 3]{AFP}.
Let $N_\tau(z)=z/\max\{|z|,\tau\}$ and define
$q=N_\tau(f)\one_E+\one_{M\setminus E}$.
The map $N_\tau$ is uniformly Lipschitz for the selected levels and
has modulus at most one. The BV product inequality, including the
jump on the finite-perimeter boundary of $E$, proves the first bound.
On $E$, $|q-f|=|1-a|$. On its complement $a\le1/2$ and
$|q-f|\le1+a\le3(1-a)$. This proves the second bound.
In particular no value of $f/|f|$ is used at a zero of $f$.
\end{proof}

\begin{theorem}[A logarithmic obstruction for complex approximate vectors]
\label{thm:compact-vector-defect}
Fix $R$ and a compact $\mathcal K\subset\Omega\setminus\{0\}$.
There is $c'_{R,\mathcal K}>0$ such that
\begin{equation}\label{eq:compact-vector-defect}
 \|f\|_2=1,\quad \|f\|_\infty+\|f\|_{\mathrm{BV}}\le H,\quad H\ge2
 \quad\Longrightarrow\quad
 \inf_{\omega\in\mathcal K}\mathcal E_{R,\omega}(f)
             \ge\frac{c'_{R,\mathcal K}}{(1+\log H)^2}.
\end{equation}
The bound applies to complex $f$, with no restriction on its zero set.
\end{theorem}
\begin{proof}
Write $\varepsilon=\mathcal E_{R,\omega}(f)$, $a=|f|$,
$\bar a=\int a\dd\nu$ and $x=\|a-\bar a\|_2$.
The reverse triangle inequality gives
$\|a\circ T_R-a\|_2\le\varepsilon$, hence
$\|a\circ T_R^m-a\|_2\le m\varepsilon$.
The BV correlation bound of Proposition~\ref{prop:bv-correlation},
rescaled by homogeneity, gives constants $C_0,\gamma>0$ with
\begin{equation}\label{eq:approximate-modulus-variance}
 x^2\le C_0 H^2e^{-\gamma m}+x m\varepsilon.
\end{equation}
Indeed add and subtract the lag-$m$ covariance of the centered $a$ and
use Cauchy--Schwarz on the remaining term. The constants in this
collision estimate can be taken uniform in $R$.

Given $0<\eta<1$, choose
$m\le C[1+\log H+|\log\eta|]$ with
$C_0H^2e^{-\gamma m}\le\eta^2$.
Equation~\eqref{eq:approximate-modulus-variance} implies
$x\le\eta+m\varepsilon$. Since $\|a\|_2=1$ and $\bar a\ge0$,
\[
 \|a-1\|_2^2=2(1-\bar a)
       \le2(1-\bar a^2)=2x^2.
\]
Lemma~\ref{lem:bv-circle-truncation} constructs $q$ with
$\|q\|_{\mathrm{BV}}\le C(1+H)$ and
$\|q-f\|_2\le3\sqrt2(\eta+m\varepsilon)$. Invariance and the
unit modulus of the multiplier yield
\begin{equation}\label{eq:phase-reconstruction-defect}
 d_{R,\omega}(q)
 \le\mathcal E_{R,\omega}(q)
 \le\varepsilon+6\sqrt2(\eta+m\varepsilon).
\end{equation}
This is an explicit reconstruction and defect comparison.

Let $L=1+\log H$. Theorem~\ref{thm:compact-log-phase-defect}, with
the bound $C(1+H)$ for $q$, makes the left side at least $c_9/L$.
Choose $\eta=c_{10}/L$ so that $6\sqrt2\eta<c_9/(2L)$.
The selected $m$ is at most $C_{R,\mathcal K}L$, since
$\log L\le L$. Equation~\eqref{eq:phase-reconstruction-defect}
then gives $\varepsilon\ge c'_{R,\mathcal K}/L^2$.
No lower-modulus estimate for $f$ was assumed; it was replaced by
\eqref{eq:approximate-modulus-variance} and the controlled truncation.
\end{proof}

\begin{corollary}[Polynomial regularity budgets on a fixed band]
\label{cor:polynomial-bv-vector-budget}
For fixed $R,\mathcal K,H_0,\zeta$ with $H_0\ge2$, $\zeta\ge0$,
normalized vectors with
$\|f_N\|_\infty+\|f_N\|_{\mathrm{BV}}\le H_0N^\zeta$
satisfy $\inf_{\omega\in\mathcal K}\mathcal E_{R,\omega}(f_N)
\ge c/(1+\log(2+N))^2$. In particular no such vectors have defect
$O(N^{-a})$ for any fixed $a>0$.
\end{corollary}
\begin{proof}
Substitute the regularity budget into
\eqref{eq:compact-vector-defect}. A positive power of $N$ dominates
$(1+\log(2+N))^2$.
\end{proof}

\subsection{Uniform compact separation across the physical family}
\begin{theorem}[Parameter-uniform separation at bounded regularity]
\label{thm:uniform-compact-bv-separation}
For every $H<\infty$ and compact
$\mathcal K\subset\Omega\setminus\{0\}$ there is
$c_{H,\mathcal K}>0$ such that
\begin{equation}\label{eq:uniform-compact-bv-separation}
 \inf_{\substack{R\in I,\ \omega\in\mathcal K,\ \|f\|_2=1\\
                \|f\|_\infty+\|f\|_{\mathrm{BV}}\le H}}
          \mathcal E_{R,\omega}(f)\ge c_{H,\mathcal K}.
\end{equation}
An empty admissible class causes no difficulty. The theorem asserts
no uniform logarithmic dependence of this constant on $H$.
\end{theorem}
\begin{proof}
Suppose a sequence made the infimum zero. Pass to subsequences with
$R_j\to R$, $\omega_j\to\omega\in\mathcal K$ and $f_j\to f$ in
$L^1(\nu)$. The last subsequence exists by BV compactness. In the present
coordinates it also follows directly from
Lemma~\ref{lem:bv-smoothing}: the approximation is uniform in $L^1$,
and at every fixed smoothing scale the bounded smoothed functions
form a precompact family on the reflected compact cylinder.
The supremum bound upgrades this convergence to $L^2$, preserves
$\|f\|_2=1$, and gives $|f|\le H$ almost everywhere.

Outside the finite one-collision singularity set at $R$,
$T_{R_j}x\to T_Rx$, $\kappa_{R_j}(x)\to\kappa_R(x)$ and
$\tau_{R_j}(x)\to\tau_R(x)$, by persistence of the regular physical
root. Thus $e^{i\phi_{R_j,\omega_j}}\to e^{i\phi_{R,\omega}}$
almost everywhere and in every finite $L^p$ norm. For a smooth $g$,
invariance of the common probability gives
\begin{align*}
 \|f_j\circ T_{R_j}-f\circ T_R\|_2
 \le\|f_j-f\|_2+2\|f-g\|_2
                      +\|g\circ T_{R_j}-g\circ T_R\|_2.
\end{align*}
Take $j\to\infty$, then $g\to f$ in $L^2$. Dominated convergence
controls the last term. The multiplier term also converges in $L^2$.
The limiting defect is therefore zero.

The identity for $f$ makes $|f|$ invariant. Ergodicity and $\|f\|_2=1$
give $|f|=1$ almost everywhere. Theorem~\ref{thm:complete-physical-phase}
now forces $\omega=0$, contradicting $\mathcal K\subset\Omega\setminus\{0\}$.
This proves the positive infimum, without claiming continuity of
product sets or their density-tail constants in $R$.
\end{proof}

\paragraph{Relation to the complementary integral.}
Theorems~\ref{thm:compact-vector-defect} and
\ref{thm:uniform-compact-bv-separation} are estimates on actual
collision functions, with their stated norms. They cannot be applied
to an anisotropic distribution until a map to such functions, its
normalization, regularity and defect comparison have been proved.
They also do not assert that these BV classes are invariant under
the twisted collision operator. Absence of exact phases alone has not
been substituted for either assertion.

The logarithmic estimate allows a growing regularity budget on a fixed
compact nonzero band at each fixed $R$. The parameter-uniform theorem
instead fixes that budget. No estimate of $c_{R,\mathcal K}$ is supplied
when the roof band grows or $\mathcal K$ approaches the origin.
Consequently the physical annulus beginning at
$2n^{-99/200}$, equivalently the rescaled boundary $2n^{1/200}$,
is not declared controlled by these theorems. The full complementary
integral, the complete critical and singular raw residual sum and the
weighted exact-event comparison retain their distinct roles in
Theorem~\ref{thm:LLT}. The present estimates give a quantitative
phase and function-vector step toward that same inversion problem.

\section{Uniform defect bounds at frequencies approaching zero}
\label{sec:near-origin-defects}
The compact bands of Section~\ref{sec:quantitative-phase-defects}
exclude the origin by a fixed amount. We now allow the physical
frequency to approach zero and the variation budget to grow. The
observable is the bounded compensation $h_R$, including its original
section discontinuity. Uniform ellipticity, rather than compact-band
phase rigidity, supplies the obstruction.

For $z\in\R^4$, $\sigma\in\R$, and a collision function $f$, put
\begin{equation}\label{eq:compensated-defect-definition}
 \mathcal D_{R,z,\sigma}f
       =f\circ T_R-e^{i(\sigma+z\cdot h_R)}f.
\end{equation}
The constant phase $\sigma$ is arbitrary. All norms in this section
are with respect to $\nu$ on the collision cylinder. Constants are
uniform in $R$; none depends on $\sigma$.

\subsection{A twisted middle block with untwisted ends}
\begin{lemma}[Damping on the diffusive collision scale]
\label{lem:diffusive-middle-damping}
There exist $r_0>0$, $A>1$, and $C<\infty$ with the following
properties. If $0<r=|z|\le r_0$, set
\[
 \delta=\tfrac14r^{1/12},\qquad m=\lceil A/r^2\rceil,\qquad
 P_{R,z}=\mathcal L_R M_{\exp(-iz\cdot h_{R,\delta})}.
\]
On every relevant local collision space,
\begin{equation}\label{eq:diffusive-middle-damping}
 \sup_{j\ge0}\|P_{R,z}^{j}\|\le C,
       \qquad |\ell(P_{R,z}^{m}\nu)|\le\tfrac18.
\end{equation}
Moreover
\begin{equation}\label{eq:middle-unsmoothing-error}
 r\|S_{m,R}(h_R-h_{R,\delta})\|_2
       \le C r^{1/24}\sqrt{1+|\log r|}.
\end{equation}
\end{lemma}
\begin{proof}
By Theorem~\ref{thm:joint-uniform-nondegeneracy},
$\Gamma_R=c_*D_R\ge c_*d_0I_4$. The covariance smoothing estimate
makes $\Gamma_{R,\delta}\ge(c_*d_0/2)I_4$ for small $r$.
The spectral expansion of Lemma~\ref{lem:smooth-collision-expansion}
applies since $r\delta^{-2}=O(r^{5/6})$. Its cubic error is
$O(r^3\delta^{-6})=O(r^{5/2})$. Reducing $r_0$ gives
\[
 \mathop{\rm Re}\log\lambda_{R,\delta}(-z)\le-c_1r^2
\]
for some uniform $c_1>0$. The projections and complementary powers
are uniformly bounded, proving the power bound. The stationary
principal amplitude is $1+O(r\delta^{-2})$. Choose $A$ large enough
that $2e^{-c_1A}<1/16$, and then reduce $r_0$ so that the
complementary term at $m$ is less than $1/16$. This proves the
second bound. Finally Proposition~\ref{prop:bv-correlation} gives
$\|S_m(h-h_\delta)\|_2^2\le Cm\delta(1+|\log\delta|)$.
Since $mr^2\le A+1$, it proves
\eqref{eq:middle-unsmoothing-error}. No multiplier bound for the
unsmoothed $h_R$ is asserted.
\end{proof}

\begin{theorem}[Quadratic coercivity for collision phases]
\label{thm:near-origin-circle-defect}
There exist $a_0,c_0,r_0>0$ such that, if $H\ge2$,
$0<r=|z|\le r_0$, and $r(1+\log H)\le a_0$, then
\begin{equation}\label{eq:near-origin-circle-defect}
 |q|=1,\quad \|q\|_{\mathrm{BV}}\le H
 \quad\Longrightarrow\quad
 \inf_{R\in I,\ \sigma\in\R}
        \|\mathcal D_{R,z,\sigma}q\|_1\ge c_0r^2.
\end{equation}
Here the infimum means the same uniform bound for every admissible
$q$ at each radius; $q$ itself may depend on $R,z,\sigma$.
\end{theorem}
\begin{proof}
Fix the radius and write $d=\|\mathcal D_{R,z,\sigma}q\|_1$.
Let $q_\epsilon=\mathsf S_\epsilon q$, where
$\epsilon=b_0/H$ and $b_0>0$ is a sufficiently small fixed number.
Then $|q_\epsilon|\le1$, $\|q-q_\epsilon\|_1\le Cb_0$, and
$\|q_\epsilon\|_{C^2}\le C\epsilon^{-2}$.
Choose integers
\[
 l=\lceil b_1+b_2\log H\rceil,\qquad N=m+2l,
\]
where $m$ is as in Lemma~\ref{lem:diffusive-middle-damping}.
The constants $b_1,b_2$ will be fixed before $H,r,q$ are chosen.

Iteration of the defect and invariance give
\begin{equation}\label{eq:near-origin-lower-pairing}
 \left|\int\overline q\,(q\circ T_R^N)
              e^{-iz\cdot S_{N,R}h_R}\dd\nu\right|\ge1-Nd.
\end{equation}
Indeed the corresponding unit-modulus integrand differs in $L^1$
from $e^{iN\sigma}$ by at most $Nd$.
Remove the phase from the first and last $l$ collisions. Since $h_R$
is bounded, this changes the pairing by at most $2lr\|h_R\|_\infty$.
Replacing its two endpoint factors by $q_\epsilon$ costs at most
$2\|q-q_\epsilon\|_1$. Smoothing only the middle collision observable
costs at most \eqref{eq:middle-unsmoothing-error}, by stationarity
and the unit supremum bounds on the endpoint factors.

The resulting pairing is exactly
\begin{equation}\label{eq:untwisted-end-word}
 \ell\left(M_{q_\epsilon}\mathcal L_R^l
                 P_{R,z}^m\mathcal L_R^l(\overline{q_\epsilon}\nu)\right).
\end{equation}
Replace both untwisted end powers by $\Pi_R=\nu\ell$.
The total error is bounded by
$C\epsilon^{-4}\vartheta^l$: both smooth endpoint factors cost
$C\epsilon^{-2}$, all other untwisted powers are bounded, and the
middle power is bounded by \eqref{eq:diffusive-middle-damping}.
The remaining expression is
\[
 \left|\int q_\epsilon\dd\nu\right|^2
                          \ell(P_{R,z}^m\nu),
\]
whose modulus is at most $1/8$.

Choose $b_0$ to make the endpoint smoothing error less than $1/16$;
then choose $b_1,b_2$ to make $C\epsilon^{-4}\vartheta^l<1/16$
for every $H\ge2$. Reducing $r_0$ makes the middle unsmoothing error
less than $1/16$. Finally choose $a_0$ so that
$r(1+\log H)\le a_0$ makes $2lr\|h_R\|_\infty<1/16$.
Thus the modulus in \eqref{eq:near-origin-lower-pairing} is at most
$1/2$. Since $N\le C/r^2$ under the same constraint, $Nd\ge1/2$
proves the assertion. Notice that the end blocks are untwisted only
inside this comparison; the original record is unchanged.
\end{proof}

\subsection{Complex vectors and their modulus}
\begin{theorem}[A uniform small-frequency function bound]
\label{thm:near-origin-vector-defect}
There exist $a_1,c_1,r_1>0$ such that, with
\[
 \Lambda(H,r)=1+\log(H/r),\qquad H\ge2,\quad 0<r\le r_1,
\]
$r\Lambda(H,r)\le a_1$ implies
\begin{equation}\label{eq:near-origin-vector-defect}
 \|f\|_2=1,\quad\|f\|_\infty+\|f\|_{\mathrm{BV}}\le H
 \quad\Longrightarrow\quad
 \|\mathcal D_{R,z,\sigma}f\|_2
                    \ge\frac{c_1r^2}{\Lambda(H,r)}
       \quad (|z|=r).
\end{equation}
The assertion is uniform in $R$ and $\sigma$, and allows zeros of $f$.
\end{theorem}
\begin{proof}
Put $e=\|\mathcal D_{R,z,\sigma}f\|_2$, $a=|f|$ and
$x=\|a-\int a\dd\nu\|_2$. The modulus estimate
\eqref{eq:approximate-modulus-variance} uses only the unit modulus of
the twist and the collision BV correlation estimate. It therefore
applies here with constants uniform in $R$:
\[
 x^2\le C H^2e^{-\gamma j}+xje.
\]
Given $\eta=b r^2$ with fixed small $b>0$, choose
$j\le C\Lambda(H,r)$ so that the first term is at most $\eta^2$.
Then $x\le\eta+je$. The circle truncation of
Lemma~\ref{lem:bv-circle-truncation} gives a phase $q$ with
\[
 \|q\|_{\mathrm{BV}}\le C(1+H),\qquad
 \|q-f\|_2\le3\sqrt2(\eta+je).
\]
Consequently
\[
 \|\mathcal D_{R,z,\sigma}q\|_1
       \le e+6\sqrt2(\eta+je).
\]
After decreasing $a_1,r_1$, the hypotheses of
Theorem~\ref{thm:near-origin-circle-defect} hold for the budget
$C(1+H)$. Choose $b$ so that the $\eta$ term is less than
$c_0r^2/2$. The remaining inequality gives
$e\ge c r^2/(1+j)\ge c_1r^2/\Lambda(H,r)$.
This proves the claim without dividing by $f$ on its zero set.
\end{proof}

\paragraph{The frequency being estimated.}
The twist in \eqref{eq:compensated-defect-definition} contains
$z\cdot h_R$, not just the physical one-collision displacement and
flight time. Its sum over a genuine return is exactly
$z\cdot(G_R-\bar G_R)$. This exact relation is used next to obtain
bounds for the actual induced record. The preceding theorems bound
defects of functions; neither is an integrated Fourier-tail estimate.

\section{Regularizing the actual tower and excluding return defects}
\label{sec:actual-return-defects}
A function of an actual return state lifts to the collision tower with
an exactly localized defect. That lift need not have bounded variation.
We construct a regular approximation with a controlled cost before
using Section~\ref{sec:near-origin-defects}. The construction concerns
functions and first distributional derivatives, not anisotropic
vectors or inverse-coarea densities.

\subsection{Finite collision records in the initial coordinates}
\begin{lemma}[A quantitative finite-record variation bound]
\label{lem:finite-record-variation}
There is $C\ge1$ such that, for every $L\ge1$ and $R\in I$, the
following bounded functions have variation norm at most
\begin{equation}\label{eq:finite-record-variation-budget}
 B_L=\exp\{C L\log(2+L)\}:
\end{equation}
all components of
$\chi\circ T_R^j$, $h_R\circ T_R^j$, and $\eta_R\circ T_R^j$,
$|j|\le L$, where $\chi(\alpha,p)=(\cos\alpha,\sin\alpha,p)$.
Their supremum norms are bounded independently of $L$. Null singular
sets are assigned any fixed bounded semialgebraic convention.
Moreover, if $g$ is smooth on the reflected collision cylinder, then
\begin{equation}\label{eq:finite-record-smooth-composition}
 \|g\circ T_R^j\|_{\mathrm{BV}}
              \le C\|g\|_{C^1(\alpha,p)}B_L\qquad(|j|\le L).
\end{equation}
The constant in $B_L$ is independent of the coefficients specifying
the moving section rectangles.
\end{lemma}
\begin{proof}
We supply the finite-complexity argument, since a one-collision BV
bound does not alone imply this lemma. A fixed finite set of disk
labels contains every next collision, by the uniform horizon. Use
normal vectors $n_j\in\mathbb S^1$, outgoing velocities
$v_j\in\mathbb S^1$, and flight lengths $\tau_j$ as auxiliary
variables. For a fixed label word $d_j$, each step is given by
\begin{equation}\label{eq:polynomial-flight-graph}
 Rn_j+\tau_jv_j-d_j=Rn_{j+1},\qquad
 v_{j+1}=v_j-2(v_j\cdot n_{j+1})n_{j+1},\qquad
 p_j=v_j\cdot Jn_j,
\end{equation}
where $J$ rotates vectors by $\pi/2$. Require $\tau_j>0$,
$v_j\cdot n_j>0$, and $v_j\cdot n_{j+1}<0$.
Require also that the open flight not enter any other candidate disk.
These last conditions are finite Boolean combinations of polynomial
sign conditions of bounded degree: for each center $d$, minimize
$|Rn_j+t v_j-d|^2$ over $0\le t\le\tau_j$.
The minimizing value occurs at an endpoint or at
$t=(d-Rn_j)\cdot v_j$; splitting into these three cases gives the
required inequalities without a quantified time variable. Endpoints
and tangencies are included in the declared singular convention.
Thus the graph records the actual first collision, not merely some
root of the flight equation.

On each bounded rational angular chart, the initial normal and outgoing
velocity also have such a graph, using a positive square-root variable
for $\sqrt{1-p^2}$. A word of length at most $2L+2$ therefore uses
at most $C_2L$ variables and $C_2L$ polynomials, of degree bounded by
a fixed $d$. All these variables can be confined to a bounded ball.
The same representation gives negative iterates by fixing the terminal
rather than the initial state; the regular billiard map is invertible.
Every lifted fiber over a regular initial state is unique.

Section membership is equally finite. Each angular interval has length
less than $\pi$ and is described by two oriented half-plane tests on
$n_j$, together with the two inequalities for $p_j$. The moving angular
centers enter only as real coefficients. Taking the finite union of
rectangles gives $\eta_R$; adjoining its two possible values also gives
the graph of $h_R=f_R-\bar G_R\eta_R$. All degree and variable counts
remain as above. There are at most $C_3^L$ choices of label word and
of these finite membership alternatives.

Here is the topological bound we use. The finite sign-condition bound
recorded in \cite[Section 3.2, equation (3.3)]{BPR}, with the zeroth
Betti number and ambient dimension $k$, bounds the total number of
connected components by
\[
 d(2d-1)^{k-1}\sum_{i=0}^{k}\binom{s}{i}4^i
                       \le (C(1+s)d)^{Ck}.
\]
The bound applies uniformly in the polynomial coefficients and to any
Boolean union of the sign conditions. We use this finite inequality,
not an asymptotic formula whose dimension is held fixed.
Fix the second collision coordinate and a real level $a$. Add the
condition that a chosen scalar output $u$ exceed $a$. The lifted
semialgebraic set has the preceding complexity. Its projection onto
the remaining coordinate has at most as many connected components:
every connected component projects to an interval or a point. Summing
over the at most $C_3^L$ alternatives bounds the number of components
of the slice superlevel set by $\exp\{C L\log(2+L)\}$.
The same reasoning applies with the two input coordinates interchanged.

If $|u|\le M_0$, the perimeter of a one-dimensional superlevel set
with at most $b$ components is at most $2b$ in the interior. The
one-dimensional coarea formula therefore bounds the variation of
that slice by $4M_0b$. Integrating both slice bounds gives the two
first distributional derivative bounds, including jumps. The fixed
angular chart cover has bounded derivatives and inverse derivatives;
its seams add only uniformly bounded boundary terms. This proves
\eqref{eq:finite-record-variation-budget}, after increasing $C$.
No inverse derivative near grazing has been bounded pointwise.

Finally extend $g$ from the embedded cylinder
$\chi(M)\subset\R^3$ to a neighborhood, with Lipschitz norm at most
$C\|g\|_{C^1(\alpha,p)}$, by a fixed tubular chart and cutoff. The
BV chain inequality applied to the three components of
$\chi\circ T_R^j$ proves \eqref{eq:finite-record-smooth-composition}.
This does not assert a BV composition bound for an arbitrary nonsmooth
$g$; only the smoothed functions below require this estimate.
\end{proof}

\subsection{Exact lifting and its regular approximation}
Write $Y=Y_R^*$, $F=F_R^*$, $r_R^*$ for its actual return time and
$g_R=G_R-\bar G_R$. For almost every collision state let its age since
the most recent section visit be $j$. The corresponding disjoint levels
are
\begin{equation}\label{eq:actual-age-levels}
 D_{j,R}=T_R^j\{y\in Y:r_R^*(y)>j\},\qquad j\ge0.
\end{equation}
They partition a full-measure set, by recurrence and invertibility.
For a measurable section function $f$, define its exact phase lift by
\begin{equation}\label{eq:exact-phase-tower-lift}
 (\mathsf L_{R,z}f)(T_R^jy)
       =e^{iz\cdot S_{j,R}h_R(y)}f(y),\qquad 0\le j<r_R^*(y).
\end{equation}
This is a function lift, not an induced transfer operator.

\begin{lemma}[The defect is supported at the true tower tops]
\label{lem:exact-tower-defect}
For $1\le p<\infty$, whenever the terms are finite,
\begin{align}
 \|\mathsf L_{R,z}f\|_{L^p(\nu)}^p
       &=c_*\int_Y r_R^*|f|^p\dd\nu_R^*,
                      \label{eq:tower-lift-norm}\\
 \|\mathcal D_{R,z,0}(\mathsf L_{R,z}f)\|_{L^p(\nu)}^p
       &=c_*\int_Y|f\circ F-e^{iz\cdot g_R}f|^p\dd\nu_R^*.
                      \label{eq:tower-exact-defect}
\end{align}
The analogous supremum bound is
$\|\mathsf L_{R,z}f\|_\infty\le\|f\|_\infty$.
\end{lemma}
\begin{proof}
Invariance and the disjoint levels give, for every nonnegative $b$,
\[
 \int_M b\dd\nu
   =c_*\int_Y\sum_{j=0}^{r_R^*-1}b(T_R^jy)\dd\nu_R^*(y).
\]
Apply this first to the modulus of the lift. At every nontop level
its next value is exactly $e^{iz\cdot h_R}$ times its current value.
At the last level the defect is
\[
 f(Fy)-e^{iz\cdot S_{r_R^*,R}h_R(y)}f(y)
               =f(Fy)-e^{iz\cdot g_R(y)}f(y),
\]
using the one-return compensation identity. This proves
\eqref{eq:tower-exact-defect}. Only one top occurs per return; its
measure is integrated against $c_*\nu_R^*$, not against a
return-time-biased probability.
\end{proof}

For a section function $f$ let $f^0$ be its zero extension to $M$,
$M_f=\|f\|_\infty$ and $V_f=\|f^0\|_{\mathrm{BV}}$.
Set $f_\epsilon=\mathsf S_\epsilon f^0$ and define
\begin{equation}\label{eq:finite-age-approximation}
 Q_{L,\epsilon}(x)=\sum_{j=0}^{L-1}\one_{D_{j,R}}(x)
 e^{iz\cdot\sum_{i=1}^j h_R(T_R^{-i}x)}
                              f_\epsilon(T_R^{-j}x).
\end{equation}
The sum for $j=0$ is empty. The approximant is zero on ages at least
$L$; the original lift is not truncated in any theorem statement.

\begin{proposition}[A uniform cost for approximating the lift]
\label{prop:tower-bv-approximation}
For $L\ge2$ and $0<\epsilon<1/2$, uniformly in $R,z,f$,
\begin{align}
 \|Q_{L,\epsilon}\|_\infty&\le M_f,
                                      \label{eq:tower-sup-bound}\\
 \|\mathsf L_{R,z}f-Q_{L,\epsilon}\|_1
        &\le C(L\epsilon V_f+M_fe^{-cL}),
                                      \label{eq:tower-L1-approximation}\\
 \|\mathsf L_{R,z}f-Q_{L,\epsilon}\|_2^2
        &\le C(M_fL\epsilon V_f+M_f^2e^{-cL}),
                                      \label{eq:tower-L2-approximation}\\
 \|Q_{L,\epsilon}\|_{\mathrm{BV}}
        &\le M_f\exp\{C L\log(2+L)\}(1+\epsilon^{-1}+|z|).
                                      \label{eq:tower-BV-budget}
\end{align}
\end{proposition}
\begin{proof}
The level indicators are disjoint, proving the supremum bound.
On the first $L$ levels, invariance bounds the total $L^1$ change
of the initial function by
$L\|f^0-f_\epsilon\|_1\le CL\epsilon V_f$.
The remaining age probability is exactly
\[
 \nu\{\text{age}\ge L\}
       =c_*\int_Y(r_R^*-L)_+\dd\nu_R^*\le Ce^{-cL},
\]
by Theorem~\ref{thm:exponential-return}. This proves the $L^1$
bound. Both functions have supremum at most $M_f$; multiplying the
$L^1$ bound by $2M_f$ gives the $L^2$ bound.

For variation, the exact level indicator is
\[
 \one_{D_{j,R}}=(\eta_R\circ T_R^{-j})
                    \prod_{i=0}^{j-1}(1-\eta_R\circ T_R^{-i}).
\]
The finite-record lemma and the BV product inequality bound its
variation by $C(j+1)B_L$. The same lemma bounds the phase sum
variation by $CjB_L$. The exponential is a Lipschitz function of this
real vector sum, with Lipschitz constant at most $|z|$.
Moreover $\|f_\epsilon\|_{C^1}\le C M_f\epsilon^{-1}$, so
\eqref{eq:finite-record-smooth-composition} applies to its pullback.
Apply the BV product inequality to each summand and then sum over
$j<L$. Absorb the resulting fixed powers of $L$ into $B_L$ by
increasing its constant. This proves \eqref{eq:tower-BV-budget}.
Neither the levels nor their images have been declared smooth; their
indicator jumps are included in the BV estimates.
\end{proof}

\subsection{Defects for the actual unbounded return record}
For $H\ge2$, $0<r<1/2$ define
\begin{equation}\label{eq:return-log-budget}
 \Lambda=1+\log(H/r),\qquad \mathcal K(H,r)=\Lambda\log(2+\Lambda).
\end{equation}
The letter $\mathcal K$ here denotes a scalar budget, not a compact
frequency set as in Section~\ref{sec:quantitative-phase-defects}.

\begin{theorem}[Uniform return-phase coercivity near the origin]
\label{thm:actual-return-circle-defect}
There exist $a_2,c_2,r_2>0$ such that, for $H\ge2$,
$0<r=|z|\le r_2$, and $r\mathcal K(H,r)\le a_2$,
\begin{equation}\label{eq:actual-return-circle-defect}
 |q|=1\text{ on }Y_R^*,\quad\|q^0\|_{\mathrm{BV}}\le H
 \quad\Longrightarrow\quad
 \|q\circ F_R^*-e^{iz\cdot(G_R-\bar G_R)}q\|_{L^1(\nu_R^*)}
                         \ge c_2r^2.
\end{equation}
All constants are independent of $R$, and no induced mixing hypothesis
is used.
\end{theorem}
\begin{proof}
The exact lift $Q=\mathsf L_{R,z}q$ has modulus one. Given
$\eta=b r^2$ with $b>0$ fixed and small, choose
$L\le C(1+|\log\eta|)$ and
$\epsilon=\eta/(C L H)$ so that
\eqref{eq:tower-L1-approximation} gives
$\|Q-Q_{L,\epsilon}\|_1\le\eta$.
The circle truncation in Lemma~\ref{lem:bv-circle-truncation}, applied
to $Q_{L,\epsilon}$, gives a phase $\widetilde q$ satisfying
\[
 \|\widetilde q-Q\|_1\le4\eta,\qquad
 1+\log(2+\|\widetilde q\|_{\mathrm{BV}})
                                     \le C\mathcal K(H,r).
\]
For the first bound, use the pointwise inequality from that lemma,
$|\widetilde q-Q_{L,\epsilon}|\le3\big||Q_{L,\epsilon}|-1\big|
\le3|Q_{L,\epsilon}-Q|$. The second follows from
\eqref{eq:tower-BV-budget}; the constants include the chosen $b$.

After decreasing $a_2,r_2$, Theorem~\ref{thm:near-origin-circle-defect}
applies to $\widetilde q$. Changing a function by $4\eta$ in $L^1$
changes its collision defect by at most $8\eta$.
Equation~\eqref{eq:tower-exact-defect} therefore gives
\[
 c_0r^2\le c_*\|q\circ F_R^*-e^{iz\cdot g_R}q\|_1+8br^2.
\]
Choose $b<c_0/16$. This proves the assertion with a uniform $c_2>0$.
In particular the regularity of the exact lift was not assumed.
\end{proof}

\begin{theorem}[Uniform coercivity for actual return functions]
\label{thm:actual-return-vector-defect}
There exist $a_3,c_3,r_3>0$ such that
$0<r=|z|\le r_3$, $H\ge2$, and $r\mathcal K(H,r)\le a_3$ imply
\begin{align}
 &\|f\|_{L^2(\nu_R^*)}=1,\quad M_f+V_f\le H
 \quad\Longrightarrow\notag\\
 &\qquad
 \|f\circ F_R^*-e^{iz\cdot(G_R-\bar G_R)}f\|_{L^2(\nu_R^*)}
                     \ge\frac{c_3r^2}{\mathcal K(H,r)}.
                         \label{eq:actual-return-vector-defect}
\end{align}
The section function may be complex and may vanish. The conclusion
is uniform in the radius and in the direction of $z$.
\end{theorem}
\begin{proof}
Let $Q=\mathsf L_{R,z}f$ and $S=\|Q\|_2$. By
\eqref{eq:tower-lift-norm}, $\sqrt{c_*}\le S\le H$; the upper bound
also follows from $\|Q\|_\infty\le H$ and $\nu(M)=1$.
If $e$ denotes the defect on the left of
\eqref{eq:actual-return-vector-defect}, then
\begin{equation}\label{eq:normalized-tower-defect}
 \|\mathcal D_{R,z,0}(Q/S)\|_2=\sqrt{c_*}\,e/S\le e.
\end{equation}

Choose an accuracy
$\eta=b r^2/\mathcal K(H,r)$, where $0<b<1$ will be fixed.
The bounds in Proposition~\ref{prop:tower-bv-approximation} allow
$L\le C[\Lambda+|\log b|]$ and
$\epsilon\ge c\eta^2/(H^2L)$ with
$\|Q-Q_{L,\epsilon}\|_2\le\sqrt{c_*}\eta$.
For example choose $L$ to make the exponential term at most
$c_*\eta^2/2$, then choose $\epsilon$ equal to a sufficiently small
fixed multiple of $\eta^2/(H^2L)$. Both choices are uniform in $R$.
Normalize $Q_{L,\epsilon}/S$ to an $L^2$ unit function $v$.
For $\eta<1/4$ this is well defined and
\[
 \|v-Q/S\|_2\le2\eta,\qquad
 1+\log(2+\|v\|_\infty+\|v\|_{\mathrm{BV}})+|\log r|
       \le C(1+|\log b|)^2\mathcal K(H,r).
\]
The last bound is \eqref{eq:tower-BV-budget}, using the lower bound
on $S$, and $\log\mathcal K(H,r)\le C\Lambda$.
The displayed dependence on $b$ prevents a circular choice of accuracy.

Choose $b$ so small that
$4b<c_1/[2C(1+|\log b|)^2]$; such a choice exists because
$b(1+|\log b|)^2\to0$. Fix this $b$ once and for all. Decrease
$a_3,r_3$ so that Theorem~\ref{thm:near-origin-vector-defect} applies
to $v$ under the preceding budget. That theorem and
\eqref{eq:normalized-tower-defect} give
\[
 e+4\eta\ge\|\mathcal D_{R,z,0}v\|_2
   \ge\frac{c_1r^2}{C(1+|\log b|)^2\mathcal K(H,r)}.
\]
The selected accuracy absorbs $4\eta$ and proves
\eqref{eq:actual-return-vector-defect}. No lower bound for $|f|$
and no mixing rate for $F_R^*$ enter the proof.
\end{proof}

\begin{corollary}[The intervening annulus with polynomial budgets]
\label{cor:actual-annulus-defect}
Fix $H_0\ge2$ and $\zeta\ge0$. For all sufficiently large $n$,
uniformly in $R$ and in physical frequencies
\begin{equation}\label{eq:actual-annulus-frequencies}
 2n^{-99/200}\le|z|\le n^{-2/5},\qquad
        2n^{1/200}\le|\sqrt n\,z|\le n^{1/10},
\end{equation}
unit section phases with zero-extension variation at most $H_0n^\zeta$
have defect at least $c|z|^2$. Normalized complex section functions
with $M_f+V_f\le H_0n^\zeta$ have defect at least
\begin{equation}\label{eq:actual-annulus-vector-rate}
 \frac{c|z|^2}{\log(2+n)\log\log(e^e+n)}.
\end{equation}
In particular this last lower bound is at least
$c n^{-99/100}/[\log(2+n)\log\log(e^e+n)]$ on the whole annulus.
\end{corollary}
\begin{proof}
On this annulus $\Lambda(H_0n^\zeta,|z|)\le C\log(2+n)$, so
$\mathcal K\le C\log(2+n)\log\log(e^e+n)$.
The largest possible value of $|z|\mathcal K$ is bounded by
$C n^{-2/5}\log(2+n)\log\log(e^e+n)$, which tends to zero.
Thus both preceding theorems apply. Squaring the lower frequency
endpoint gives the final power. The interval is nonempty for all
sufficiently large $n$, since its endpoint ratio is
$n^{19/200}/2\to\infty$.
\end{proof}

\paragraph{What the annulus estimate supplies.}
The norm in \eqref{eq:actual-annulus-vector-rate} is a one-step
cohomological defect for the \emph{actual unbounded return record},
with the normalization and first-variation budget stated above.
The lifted defect and the regularization error have both been proved,
rather than assumed to come from an induced spectral construction.
No conclusion is made here for arbitrary eigenvalues of the induced
operator: the multiplier in the theorem is precisely
$e^{iz\cdot(G_R-\bar G_R)}$.

This estimate is a quantitative input on the intervening annulus, not
the complementary Fourier integral itself. To deduce decay of powers
from it on a distribution space still requires a construction of that
space, a function reconstruction with the displayed budgets, and a
bound relating the operator defect to the present physical $L^2$
defect. Nor does the finite-record BV lemma estimate any second
derivative of a raw pushforward density. The growing roof region,
the complete critical/singular extraction, and the weighted exact-event
comparison remain distinct requirements of Theorem~\ref{thm:LLT}.
All conditioning indicators in the inherited central and moment
theorems are unchanged.

\section{Conditioning, clocks, and geometric prediction}
\begin{proposition}[Central density conditioning]\label{prop:conditioning}
Suppose an unweighted and an $f$-weighted joint density satisfy
\[
 n^2p_n(z)=G_n(z)+r_n(z),\qquad
 n^2p_n^f(z)=a_fG_n(z)+r_n^f(z)
\]
on a common central window, with $G_n\ge g_0>0$ and $|r_n|,|r_n^f|\le\varepsilon_n<g_0$. Then their density ratio differs from $a_f$ by at most
\begin{equation}\label{eq:ratio}
          \frac{(1+|a_f|)\varepsilon_n}{g_0-\varepsilon_n}.
\end{equation}
The same bound holds after conditioning on fixed lattice coordinates and any positive-length roof interval contained in that window, including intervals whose lengths shrink with $n$.
\end{proposition}
\begin{proof}
Subtract $a_f$ times the first identity from the second and divide by $G_n+r_n\ge g_0-\varepsilon_n$. For an interval, integrate both identities first; the errors are at most its length times $\varepsilon_n$ and the denominator is at least its length times $g_0-\varepsilon_n$. The length cancels.
\end{proof}
At an exact roof value this is a density-defined version of a regular conditional expectation for almost every value with positive density. It is not conditioning on a positive-probability singleton. It supplies no estimate outside the common central window and no functional path convergence by itself. These are the precise limitations of this interface to stopped-path large deviations.

\begin{proposition}[A conditional physical-clock transfer]\label{prop:clock}
Let $(X_j,\tau_j)\in\R^d\times(0,\infty)$ be stationary dynamical observations with means $(a,\mu)$, $\mu>0$, and a joint functional CLT with covariance $\Sigma$. Assume the uniform law of large numbers for their clock and that initial and terminal residual records on bounded rescaled time intervals are $o_P(\sqrt t)$. Then the reward accumulated up to physical time $t$, centered by $(a/\mu)t$, has limiting covariance per unit time
\begin{equation}\label{eq:clock}
 \frac1\mu Q\Sigma Q^T,\qquad Q=(I_d,-a/\mu).
\end{equation}
For $X=(\kappa_1,\kappa_2,r)$ this transfers induced displacement and collision counts, not the number of experimental bits.
\end{proposition}
\begin{proof}
The clock law gives its inverse index $N(t)/t\to1/\mu$. The centered functional CLT evaluated at this inverse time and the identity $S_{N(t)}\tau=t+o_P(\sqrt t)$ give
\[
 S_{N(t)}X-(a/\mu)t
 = (S_{N(t)}X-N(t)a)-(a/\mu)(S_{N(t)}\tau-N(t)\mu)+o_P(\sqrt t).
\]
The continuous limit permits the random-time substitution. Projection by $Q$ and evaluation at time $1/\mu$ prove \eqref{eq:clock}. The assumed residual estimate includes unfinished induced returns, which are not bounded merely because individual flights are bounded.
\end{proof}

\begin{lemma}[Covariance continuity from a finite-time comparison]\label{lem:covariance}
Suppose a centered correlation matrix satisfies $\|C_j(R)\|\le Ce^{-cj}$ and
$\|C_j(R)-C_j(R')\|\le C\delta^\alpha e^{a j}$, where $\delta=|R-R'|$, $a,c,\alpha>0$. Then the absolutely convergent Green--Kubo covariance satisfies
\[
 \|\Sigma_R-\Sigma_{R'}\|\le C'\delta^{\alpha c/(a+c)}.
\]
If $a=0$, the bound is instead $C'\delta^\alpha(1+|\log\delta|)$.
\end{lemma}
\begin{proof}
Split the defining correlation series at $J=\lfloor\alpha\log(1/\delta)/(a+c)\rfloor$. For $a>0$, the finite sum is bounded by $C\delta^\alpha e^{aJ}$ and the remaining sum by $Ce^{-cJ}$. Both have the stated order. For $a=0$, sum the constant comparison bound over $J+1$ terms and use the exponential tail. Symmetrizing the correlations only changes the constant.
\end{proof}
The hypotheses of this lemma must concern the actual induced observables and their varying domains. Spectral continuity for a different observable norm is not a substitute for this comparison.

\begin{proposition}[Geometric error and the central prediction scale]\label{prop:prediction}
Assume \eqref{eq:LLT} with error $\epsilon_n$, uniformly elliptic covariances, and estimates $(\widehat m,\widehat\Sigma)$ satisfying $|\widehat m-m_R|+\|\widehat\Sigma-\Sigma_R\|\le\omega(\nu)$. If $\sqrt n\,\omega(\nu)\le1$, replacing the Gaussian parameters in \eqref{eq:LLT} changes its right side on a fixed central window by at most $C(1+\sqrt n)\omega(\nu)$. The total prediction error is bounded by this quantity plus $\epsilon_n$.
\end{proposition}
\begin{proof}
The argument of the Gaussian moves by $\sqrt n(\widehat m-m_R)$. On a slightly enlarged compact window the Gaussian and its derivatives in its argument and covariance are bounded uniformly under ellipticity. The mean value theorem gives the bound. Adding the LLT error finishes the proof.
\end{proof}

There is also a direct equilibrium conclusion that needs no local-limit hypothesis. The mean collision rate is
\[
 \lambda(R)=\frac{2R}{\sqrt3/2-\pi R^2}.
\]
Its derivative is $2/A_R+4\pi R^2/A_R^2$, where $A_R=\sqrt3/2-\pi R^2$. On $I$, $A_R>17/100$, using $\sqrt3>173/100$ and $\pi<22/7$, and direct substitution gives $0<\lambda'(R)<110$. Therefore a radius estimate with error at most $\nu$ gives collision-rate error at most $110\nu$. On the circular submodel the average support function identifies the radius, independently of translation. The finite geometric conclusion of \cite{A2G} therefore transfers to this rate with its geometric observation cost, without recovering the entire unknown launch law.


\section{A logarithmic unfinished-block bound under conditioning}
\label{sec:conditioned-clock}
The stationary law of an induced block is length-biased. We use this law
for the initial block and the exact visit intensity for subsequent blocks.
These two ingredients avoid applying a launch-law union bound to a
stationary initial state.

Write $\overline\tau_R^*=\bar\tau_R/c_*$ and represent the flow as the
suspension over $(Y_R^*,F_R^*,\nu_R^*)$ with roof $\tau_R^*$. Let $t_j$
be its two-sided visit times and $x_j$ their outgoing section states.
The minimum time between consecutive visits is at least $3/50$.

\begin{lemma}[Marked visit intensity]
\label{lem:marked-visit-intensity}
For every nonnegative measurable $f$ on $Y_R^*$ and every $t\ge0$,
\begin{equation}\label{eq:marked-visit-intensity}
 \mathbb E_{\mu_R}\sum_{t_j\in(0,t]} f(x_j)
       =\frac{t}{\overline\tau_R^*}\int f\dd\nu_R^*.
\end{equation}
\end{lemma}
\begin{proof}
First take bounded $f$. Stationarity makes the locally finite measure
$A\mapsto\mathbb E_{\mu_R}\sum_{t_j\in A}f(x_j)$ translation-invariant
on $\R$, hence equal to $i_f$ times Lebesgue measure. Local finiteness
follows from the positive lower roof bound. Choose $0<\delta<3/50$.
At a stationary time, being in the first $\delta$ units of an induced
roof is the same event as a visit in the preceding interval of length
$\delta$, apart from null endpoints. There is at most one such visit,
and its mark is the current base state. Integration in the suspension
therefore gives
\[
 \delta i_f=\frac1{\overline\tau_R^*}
       \int_{Y_R^*}\int_0^\delta f(x)\dd s\dd\nu_R^*(x)
       =\frac{\delta}{\overline\tau_R^*}\int f\dd\nu_R^*.
\]
This determines $i_f$ and proves the formula for bounded $f$.
Monotone convergence extends it to arbitrary nonnegative marks, including
the case of an infinite integral. No renewal independence is involved.
\end{proof}

Let $\mathcal M_R(t)$ be the maximum of $Q_R$ over all complete induced
blocks whose time intervals meet $[0,t]$, including the block containing
time zero. Endpoint conventions at a visit are immaterial for $\mu_R$.
The marked size of any unfinished piece is bounded by a fixed multiple
of $1+\mathcal M_R(t)$.

\begin{theorem}[Uniform exponential maximum and conditioned residual]
\label{thm:conditioned-maximal-block}
There are constants $a_2,C_2>0$ such that, for every $R\in I$, $t\ge0$
and $q\ge0$,
\begin{equation}\label{eq:exponential-maximal-block}
 \mu_R\{\mathcal M_R(t)>q\}\le C_2(1+t)e^{-a_2q}.
\end{equation}
Consequently, for any measurable event $B_{R,t}$ of probability at least
$d_t>0$,
\begin{equation}\label{eq:conditioned-maximal-block}
 \mu_R\{\mathcal M_R(t)>q\mid B_{R,t}\}
       \le\min\{1,C_2(1+t)e^{-a_2q}/d_t\}.
\end{equation}
In particular, if $d_t\ge d(1+t)^{-A}$ with $d>0$ and $A\ge0$, then
$\mathcal M_R(t)=O_P(\log(2+t))$ under these conditional laws, uniformly
in $R$ and in the events. The maximum marked unfinished return during
$[0,t]$ is then $o_P(\sqrt t)$ under the same laws. More generally the
last assertion holds if $\log((1+t)/d_t)=o(\sqrt t)$.
\end{theorem}
\begin{proof}
The initial containing block has law
$\tau_R^*\dd\nu_R^*/\overline\tau_R^*$. If the maximum exceeds $q$,
either that block has $Q_R>q$ or a visit in $(0,t]$ starts a block with
$Q_R>q$. The union bound and Lemma~\ref{lem:marked-visit-intensity} give
\begin{equation}\label{eq:exact-maximal-tail-budget}
 \mu_R\{\mathcal M_R(t)>q\}
 \le\frac1{\overline\tau_R^*}
       \left(\int\tau_R^*\one_{\{Q_R>q\}}\dd\nu_R^*
                       +t\,\nu_R^*(Q_R>q)\right).
\end{equation}
The denominator is bounded below on $I$. Since $\tau_R^*\le Q_R$,
the exponential moment in \eqref{eq:exponential-return} bounds the first
integral by $C e^{-a_1q/2}$ and the probability by $C e^{-a_1q}$.
For example, $Q\le (2/a_1)e^{a_1Q/2}$ and
$e^{a_1Q/2}\one_{Q>q}\le e^{a_1Q}e^{-a_1q/2}$ suffice for the
first bound. This proves \eqref{eq:exponential-maximal-block}.

Dividing the unconditional exceptional probability by $\mu_R(B_{R,t})$
proves \eqref{eq:conditioned-maximal-block}. With a polynomial lower
bound on $d_t$, setting $q=K\log(2+t)$ and taking
$a_2K>A+1$ makes the right side tend to zero. With the more general
hypothesis, set $q=\epsilon\sqrt t$ for any $\epsilon>0$; the logarithm
of the displayed upper bound tends to minus infinity. The deterministic
bound of each marked piece by $C(1+Q_R)$ proves the assertions for all
unfinished pieces at once. The event may depend on the whole orbit and
on an estimator from that orbit; no independence is assumed.
\end{proof}

For a stationary starting point $(x,s)$ put
$V_R(v)=\max\{j:T_{j,R}(x)\le s+v\}$. Define the two path records
\[
 Z_R(v)=(V_R(v),C_R(v),W_R(v)),\qquad
 Z_R^\circ(v)=(V_R(v),N_{V_R(v),R}(x),K_{V_R(v),R}(x)).
\]
The second record counts full blocks from the beginning of the initial
containing roof, not from the observation time. It is used solely as an
exact intermediate record for comparison.

\begin{corollary}[A conditional path-comparison budget]
\label{cor:conditional-path-budget}
There is $C_3$ such that, pathwise outside the collision singularities,
\begin{equation}\label{eq:deterministic-prefix-bound}
 \sup_{0\le v\le t}|Z_R(v)-Z_R^\circ(v)|
                         \le C_3(1+\mathcal M_R(t)).
\end{equation}
Let $\widehat Z_R$ and $\widehat Z_R^\circ$ be the corresponding paths
on $[0,1]$, divided by $\sqrt t$ and with any identical deterministic
centering. For $|F|\le1$ Lipschitz with constant $L_F$ for the uniform
path metric and $\mu_R(B_{R,t})\ge d_t>0$,
\begin{align}\label{eq:conditional-path-budget}
 \left|\mathbb E_{\mu_R}[F(\widehat Z_R)\mid B_{R,t}]
       -\mathbb E_{\mu_R}[F(\widehat Z_R^\circ)\mid B_{R,t}]\right|
 &\le\frac{C_3L_F(1+q)}{\sqrt t}
        +\frac{2C_2(1+t)e^{-a_2q}}{d_t}.
\end{align}
For $t\ge2$ choose
\begin{equation}\label{eq:conditional-log-budget}
 q=a_2^{-1}\log\left(\frac{\max\{1,2C_2\}(1+t)\sqrt t}{d_t}\right).
\end{equation}
The second term is at most $t^{-1/2}$; for polynomially bounded
conditioning probabilities the total error is
$O((1+L_F)\log t/\sqrt t)$.
\end{corollary}
\begin{proof}
The collision counts differ by the initial prefix and the final
incomplete block. Each is bounded by the full block's collision count.
The same holds for the lattice displacement after multiplying by the
uniform bound on a single-flight lattice step and adding a fixed
fundamental-cell boundary constant. The visit counts agree exactly by
construction. These estimates are simultaneous for all $v\in[0,t]$,
and prove \eqref{eq:deterministic-prefix-bound}. On
$\{\mathcal M_R(t)\le q\}$ the Lipschitz bound gives the first term in
\eqref{eq:conditional-path-budget}. On its complement the difference of
two values of $F$ is at most two. Apply
\eqref{eq:conditioned-maximal-block}. The stated choice of $q$ and
substitution prove the final assertions.
\end{proof}

\begin{remark}[Which conditional interface has been closed]
\label{rem:conditional-interface-scope}
The square-root unfinished-block requirement is now a theorem for the
actual mechanical family, even after any conditioning of polynomially
small positive probability and over the whole observation interval. A
functional central limit theorem for completed records is still a
separate premise in Proposition~\ref{prop:clock}. For a local limit, the
conditioning event on the two sides of \eqref{eq:conditional-path-budget}
is the same event. Replacing an exact lattice observation by the
completed-block observation changes that event and needs a further
boundary estimate. If two events $B,B'$ have probabilities at least $d$,
then for $|f|\le1$ direct subtraction of the ratios gives
\[
 |\mathbb E[f\mid B]-\mathbb E[f\mid B']|
             \le2\mu_R(B\mathbin\triangle B')/d.
\]
Thus an $O(\log t)$ error in an integer record is not discarded when
conditioning on one lattice value. The weighted raw-density estimates
of Proposition~\ref{prop:conditioning} remain the requisite local input.
A zero-probability singleton in the continuous coordinate is not an event
to which \eqref{eq:conditioned-maximal-block} applies.
\end{remark}

\section{Induced covariance and raw-summation criteria}
\label{sec:closure-contracts}
We record the exact analytical implications needed in the low-frequency
and outer-frequency parts of the original density theorem. These
criteria distinguish proved finite-record facts from the two infinite
sums that still require estimates. The collision matrix $\Gamma_R$ and
the actual Gaussian covariance $D_R$ have already been established in
Theorems~\ref{thm:collision-covariance} and
\ref{thm:actual-record-gaussian}. The following induced Green--Kubo
criterion is stronger: neither its infinite sum nor its periodic-evaluable
representative is presumed by those theorems. The direct argument of
Theorem~\ref{thm:joint-uniform-nondegeneracy} now proves positive
definiteness of $D_R$ without either of these additional premises;
the criterion below is retained for the stronger, absolutely convergent
induced Green--Kubo representation.

\begin{proposition}[A covariance closure criterion]
\label{prop:covariance-closure}
Put $g_R=G_R-\bar G_R$. Suppose that
\begin{equation}\label{eq:covariance-tail-contract}
 \lim_{m\to\infty}\sup_R\sum_{k>m}
       \left\|\int g_R\otimes(g_R\circ(F_R^*)^k)\dd\nu_R^*\right\|=0.
\end{equation}
Suppose also that zero asymptotic variance of $v\cdot g_R$ implies a
coboundary whose transfer function has a representative for which the
coboundary identity telescopes on the five real-rank periods of
Corollary~\ref{cor:real-period-rank}. Then the Green--Kubo matrices
\[
 \Sigma_R=\int g_R\otimes g_R\dd\nu_R^*
  +\sum_{k\ge1}\int\bigl(g_R\otimes(g_R\circ F^k)
                    +(g_R\circ F^k)\otimes g_R\bigr)\dd\nu_R^*
\]
are continuous and satisfy $cI_4\le\Sigma_R\le CI_4$ on $I$ for some
$c,C>0$.
\end{proposition}
\begin{proof}
Corollary~\ref{cor:all-moment-continuity} makes every fixed-lag matrix
continuous, including the zero-lag term. The uniformly summable tail in
\eqref{eq:covariance-tail-contract} therefore gives continuity of the
series. The usual expansion of the covariance of $\sum_{j<n}g_R\circ F^j$
is a finite identity; dividing it by $n$ and using absolute summability
identifies its limit with $\Sigma_R$. Thus $\Sigma_R$ is positive
semidefinite and its directional quadratic form is the asymptotic
variance.

If $v^T\Sigma_Rv=0$, the stipulated coboundary identity gives on each of
the five selected periods
\[
 v_1K_1+v_2K_2+v_3N+v_4T-h(v\cdot\bar G_R)=0.
\]
The augmented vector $(v_1,v_2,v_3,-v\cdot\bar G_R,v_4)$ annihilates
their five-dimensional record matrix, whose determinant is nonzero by
Corollary~\ref{cor:real-period-rank}. Hence $v=0$. Each matrix is
positive definite. Continuity and compactness of $I$ give uniform lower
and upper eigenvalue bounds. The periodic representative premise is
used explicitly here; a merely measurable identity is insufficient.
\end{proof}

\begin{lemma}[An absolute raw-residual summation criterion]
\label{lem:raw-residual-sum}
Let $\mathcal W$ be countable. For $w\in\mathcal W$ assign a lattice
label $k_w\in\Z^3$ and a function $q_w\in L^1(\R)$ whose second
distributional derivative is a finite signed measure. Suppose
\[
 A_0=\sum_w\|q_w\|_1<\infty,\qquad
 A_2=\sum_w\|D^2q_w\|_{\TV}<\infty.
\]
Then
$H(u,b)=\sum_w e^{iu\cdot k_w}\widehat q_w(b)$ belongs to
$L^1(\T^3\times\R)$ and, for $B\ge1$,
\begin{align}\label{eq:absolute-residual-sum}
 \int_{\T^3\times\R}|H|&\le2(2\pi)^3(A_0+A_2),\\
 \int_{\T^3\times\{|b|>B\}}|H|&\le2(2\pi)^3A_2/B.\nonumber
\end{align}
The same estimates hold uniformly for a parameter family whenever the
corresponding sums are uniformly bounded.
\end{lemma}
\begin{proof}
Distributional differentiation and the stated Fourier convention give
$-b^2\widehat q_w(b)=\widehat{D^2q_w}(b)$. Thus
\[
 |\widehat q_w(b)|\le\min\{\|q_w\|_1,
                    |b|^{-2}\|D^2q_w\|_{\TV}\}\quad(b\ne0).
\]
Integrate the first estimate on $|b|\le1$ and the second on $|b|>1$,
and then sum by Tonelli. The torus has volume $(2\pi)^3$.
The series itself converges uniformly as a transform of the absolutely
summable $L^1$ densities. Integrating only $|b|>B$ gives the second
bound. All lattice multiplicities are included in the two sums; no
bound on the number of words has been suppressed.
\end{proof}

To use this lemma, the residuals must be those of a decomposition of the
actual full measure, including singular itinerary boundaries. A jump left
in $q_w$ contributes a derivative of a point mass to $D^2q_w$ and does
not satisfy the hypothesis. The individual jump estimates of
Proposition~\ref{prop:general-edge} do not prove the sum $A_2<\infty$.
Nor does the scalar tail \eqref{eq:exponential-return} bound inverse
Jacobian factors or the variation of a density near a critical value.
If $A_2$ grows with $n$, its actual growth must be used in the frequency
splice rather than treated as a fixed constant.

There is nevertheless a useful exact truncation budget. For an insertion
$|w|\le M_0$, deleting the physical words with $N_{n,R}>L$ changes its
record measure, and its raw transform in uniform norm, by at most
\begin{equation}\label{eq:physical-count-truncation}
 M_0C_1e^{-a_1L/n}.
\end{equation}
Indeed total variation contracts under pushforward, and
\eqref{eq:block-exponential-moment} bounds the discarded mass.
On a frequency band of volume $V$ the integrated error is at most
$VM_0C_1e^{-a_1L/n}$. Thus the sufficient physical-count budget for a
prescribed integrated error $\varepsilon$ is
\[
 L\ge\frac{n}{a_1}\max\{0,\log(VM_0C_1/\varepsilon)\}.
\]
This controls a finite-band error, not a pointwise density error or an
infinite frequency band. It is compatible with retaining the original
raw datum and with the separate all-branch derivative contract above.

\section{The joint local-limit conclusion and its realization}
\label{sec:final-realization}
The endpoint of this article remains the parameter-uniform raw mixed
density local limit \eqref{eq:LLT} for the physical record
$(K_{n,R},N_{n,R},T_{n,R})$, with the weighted positive-denominator
interface of Proposition~\ref{prop:conditioning}. Neither the return
section nor the recorded observables are replaced by independent flights,
a smoothed observation, or an unrelated geometric problem.

The change from $Y_R$ to $Y_R^*$ is explicit, and the periodic records
are those of its actual first return. The joint periodic annihilator is
trivial. The adjacent-increment comparison yields a uniform positive
period discrepancy with logarithmic collision length, and
Theorem~\ref{thm:one-step-coercivity} gives its positive-measure
consequence for controlled regular phases. The phase-reconstruction and
Fredholm hypotheses of Theorem~\ref{thm:phase-resolvent-criterion} state
precisely how these estimates would give an induced resolvent bound.
Their verification on the actual induced strong/weak spaces is still
required. No division by an uncontrolled spectral modulus is implicit.

The collision-space spectral input has a different, now completed use:
Theorem~\ref{thm:exponential-return} proves uniform exponential moments
of the genuine return block. It strengthens finite-record stability to
all finite moments. The common-chart construction in
Proposition~\ref{prop:finite-common-patches} gives an explicit
initial-state remainder for the moving-domain density comparison.
The bounded compensation of Section~\ref{sec:collision-covariance}
then proves a uniformly summable collision Green--Kubo formula. The
exact stopped identity gives the actual Gaussian covariance
$D_R=c_*^{-1}\Gamma_R$, its continuity, and its $L^2$ cohomological
kernel. These conclusions no longer depend on an assumed induced
spectral expansion. The stronger induced-correlation criterion of
Proposition~\ref{prop:covariance-closure} remains available, but is not
used to obtain these Gaussian laws. The actual normalized second
moments now converge by Theorem~\ref{thm:marked-quadratic-moments},
and Theorem~\ref{thm:joint-uniform-nondegeneracy} proves uniform positive
definiteness. Its measurable phase argument removes the section term
before using symplectic area and finite physical covers; periodic
evaluation of the transfer function is not required.

At the physical-clock level, the first-order functional law and its
geometric rate comparison remain valid. The stronger maximum estimate
in Theorem~\ref{thm:conditioned-maximal-block} now controls all unfinished
returns on a growing interval, at logarithmic scale, even after
conditioning of polynomially small positive probability. Corollary
\ref{cor:conditional-path-budget} supplies the corresponding explicit
square-root path error. Thus the unfinished-block premise no longer
needs to be assumed in that range of conditioning probabilities.
Theorem~\ref{thm:actual-return-functional} now supplies the unconditioned
functional Gaussian theorem for the completed physical return record;
Corollary~\ref{cor:physical-functional-gaussian} gives unconditioned
physical-time fluctuations as well. Neither an unconditioned functional
limit nor a path comparison under one event justifies changing an exact
lattice conditioning event. That boundary comparison is a separate
weighted local-density problem, as recorded in
Remark~\ref{rem:conditional-interface-scope}.

For the raw inversion itself, the exact edge subtraction and the local
edge criterion retain every correction term. The all-branch residual
must include central critical words and singular itinerary boundaries.
Lemma~\ref{lem:raw-residual-sum} specifies an absolute derivative-norm
criterion for that sum; no uncontrolled word multiplicity is hidden in
it. The required uniform sum has not been established here. Its growth
constants, together with the actual induced operator estimates, must
satisfy the strict splice inequalities \eqref{eq:splice} or directly
prove \eqref{eq:local-tail}. The weighted insertion class needs the same
verification. A positive conditioning denominator then follows from the
raw LLT on central windows by Proposition~\ref{prop:conditioning}.

The resulting dependency chain is exact: the geometry and periodic
coercivity, the genuine exponential return tail, the common finite-record
charts, and the conditioned unfinished-block estimate are proved;
the continuous collision Green--Kubo formula, actual Gaussian laws,
functional limits, and measurable zero-variance cohomology are also
proved. The marked moment laws and full joint nondegeneracy are proved
as well. Quantitative induced phase reconstruction or a direct complete
complementary-frequency estimate, together with the global raw residual
and weighted bounds, remains necessary to instantiate the full density
theorem. No complete unconditional raw
LLT or full A3 density input is asserted before those premises are
verified. The collision spectral theory is used at its stated smooth-multiplier
scope for killing and for the compensated Gaussian limit. No spectral
theorem about the unbounded induced observables is imported.

\subsection*{Acknowledgment and reproducibility}
AI-assisted analysis contributed to this manuscript. The supplementary
source audit fixes the mathematical baseline, the referee report and
all imported source files; the response describes the revisions to each
numbered request. Finite computations check algebra and explicit orbit
margins and do not replace the analytical proofs. Independent human
specialist review is not claimed by this manuscript.

\appendix
\section{Uniform variation and the collision-space embeddings}
\label{app:bv-details}
We give the geometric and norm details used in
Section~\ref{sec:collision-covariance}. Variation is always taken in
initial collision coordinates. The statements below do not concern
variation of a density after a many-collision pushforward.

\begin{lemma}[A finite set of possible next centers]
\label{lem:finite-next-centers}
There is a finite subset $\mathcal D$ of the triangular lattice such
that every regular next collision, for every $R\in I$, is with a disk
centered at a member of $\mathcal D$. The cardinality and the set may
be chosen independently of the initial collision state.
\end{lemma}
\begin{proof}
Let $H$ be a uniform upper bound for the free flight, as supplied by the
finite-horizon geometry in Lemma~\ref{lem:collision-spectral-input}.
A flight from $q=Rn_\alpha$ to the disk centered at $d$ ends at
$q'=d+Rn'$ and has length at most $H$. Hence
$|d|\le |q|+|q'-q|+|q'-d|\le H+2R_{\max}$.
The intersection of the lattice with this fixed ball is finite. For
example the open disks of radius $1/2$ about its lattice points are
disjoint, and lie in a ball of radius $H+2R_{\max}+1/2$; area comparison
gives a bound of $4(H+2R_{\max}+1/2)^2$ for their number. Removing the
initial disk or inadmissible centers only decreases this bound.
\end{proof}

\begin{proposition}[Uniform semialgebraic slicing]
\label{prop:semialgebraic-slicing}+Let $\Lambda$ be a semialgebraic parameter set and
$Q=(a,b)\times(c,d)$ a bounded rectangle. Suppose
$u:\Lambda\times Q\to\R$ is semialgebraic and $|u|\le M$.
There is an integer $N$, depending on the family but not on
$\lambda\in\Lambda$, such that $u_\lambda\in\mathrm{BV}(Q)$ and
\begin{equation}\label{eq:slicing-variation-bound}
 |D_1u_\lambda|(Q)\le 4NM(d-c),\qquad
 |D_2u_\lambda|(Q)\le 4NM(b-a).
\end{equation}
The zero extension outside $Q$ has bounded variation as well, with an
additional bound $2M[(b-a)+(d-c)]$ for the sum of its two variations.
The assertions depend only on the almost-everywhere class of the
functions, not on chosen values on lower-dimensional sets.
\end{proposition}
\begin{proof}
Apply differentiable cell decomposition to the graph of $u$, ordering
$(\lambda,y)$ before $x$, and refine the resulting cells by the sign
of the derivative in $x$ wherever that derivative is defined. Uniform
finiteness and one-dimensional monotonicity give a single bound $N$
for the number of monotone intervals and exceptional points in every
fiber. These are the parameterized cell and monotonicity theorems of
\cite[Sections 2.1--2.2 and 6.2]{Coste}. On a monotone interval the
variation is at most $2M$. Jumps between intervals are also at most
$2M$. Isolated assigned values are immaterial for distributional
variation. Thus the variation of the almost-everywhere one-dimensional
fiber is bounded by $4NM$, after enlarging $N$ if necessary.

For $\phi\in C_c^1(Q)$, Fubini's theorem and the one-dimensional
integration-by-parts inequality now give
\[
 \left|\int_Q u_\lambda\,\partial_x\phi\dd x\dd y\right|
 \le 4NM(d-c)\|\phi\|_\infty.
\]
This bounded functional defines a finite distributional derivative
measure, proving the first estimate. Reorder the parameters as
$(\lambda,x)$ before $y$ and repeat the argument for the second.
Taking the larger of the two finite fiber bounds gives the stated
single $N$. Extending each fiber by zero adds at most two endpoint
jumps, each of size at most $M$. The same Fubini argument proves the
claimed zero-extension bound. In particular no bound on a derivative
in the parameter is needed.
\end{proof}

\begin{proposition}[Angular chart gluing for the physical record]
\label{prop:physical-chart-gluing}
The one-collision record $f_R$ has a uniformly bounded supremum norm
and flat-coordinate variation on $M$. Adding any fixed finite union
of moving rectangular section indicators with uniformly bounded
perimeters preserves these bounds. This proves the variation assertions
for $h_R$ and $s_R$ in Lemma~\ref{lem:physical-bv}.
\end{proposition}
\begin{proof}
Choose a fixed finite angular atlas on which
$t=\tan((\alpha-\alpha_0)/2)$ lies in a bounded interval, and choose a
smooth partition of unity with supports strictly inside these charts.
The chart maps and their inverses have bounded first derivatives.
The components of $n_\alpha$ and $t_\alpha$ are rational in $t$, and
$\sqrt{1-p^2}n_\alpha+pt_\alpha$ is semialgebraic in $(t,p)$.
Lemma~\ref{lem:finite-next-centers} reduces the next collision to
finitely many candidate roots. The conditions for a root to be real,
positive, and the first admissible entry, and all comparisons between
two such roots, are semialgebraic. Therefore its minimum and the finite
lattice label define bounded semialgebraic families in $(R,t,p)$.
Ties and singular trajectories can be assigned a fixed bounded value.
They do not change the almost-everywhere record.

Proposition~\ref{prop:semialgebraic-slicing} gives a uniform variation
bound on each chart. Multiplication by a fixed smooth partition
function bounds the variation by its supremum times the original
variation plus its derivative supremum times the original $L^1$ norm.
A bounded smooth change of the angular variable preserves a variation
bound with a constant depending only on the chart and its inverse.
For completeness, this follows for smooth functions by the chain rule
and change of variables, and for a bounded-variation function by
convolution on a slightly smaller chart followed by lower semicontinuity
of distributional variation. The cutoff has support in that smaller
chart. There is therefore no unaccounted jump on the chart boundary.
The sum of the finitely many localized functions gives the asserted
bound on the cylinder. At $p=\pm1$, the fiber proof already includes
finite one-sided traces; no inverse transformation from $p$ to the
angle of incidence is used.

A rectangle indicator has variation equal to its coordinate perimeter,
up to the fixed normalization of $\nu$. The variation of the indicator
of a finite union is bounded by the sum of these perimeters, including
rectangles represented across the angular seam. Its moving center does
not change that estimate. The bounded coefficients in
\eqref{eq:compensation-normalizations} finish the assertion for $h_R$
and $s_R$.
\end{proof}

\begin{lemma}[Smooth densities and test pairings in the local spaces]
\label{lem:density-norm-details}
On each local space $(\mathcal B_j,\mathcal B_{w,j})$ used in
Lemma~\ref{lem:collision-spectral-input}, smooth functions $a,g,b$
satisfy, uniformly for the radii in that neighborhood,
\begin{align}
 \|a\nu\|_{\mathcal B_j}+|a\nu|_{w,j}
       &\le C\|a\|_{C^1},\label{eq:smooth-density-embedding}\\
 \|M_gV\|_{\mathcal B_j}&\le C\|g\|_{C^2}\|V\|_{\mathcal B_j},
 &|\ell_j(M_bV)|&\le C\|b\|_{C^2}\|V\|_{\mathcal B_j}.
 \label{eq:smooth-multiplier-pairing}
\end{align}
The function norms are in the collision coordinates used to define
the local spaces. In particular, for the mean-preserving smoothing,
\begin{equation}\label{eq:smoothed-density-vector}
 \| (\mathsf S_\delta a)\nu\|_{\mathcal B_j}
       \le C\delta^{-2}\|a\|_\infty .
\end{equation}
Here $a$ is a density relative to $\nu$, not relative to unweighted
area in the incidence-angle coordinates.
\end{lemma}
\begin{proof}
We record the norm identification explicitly. In
\cite[Sections 3.2--3.3, (3.11)--(3.13)]{DZ}, a smooth function is
identified with its density times the invariant reference measure
$\mu=c\cos\varphi\dd r\dd\varphi$. Under the local boundary
reparametrization $r=R_0\alpha$, this probability is precisely $\nu$.
Thus the smooth function representing $a\nu$ is $a$; an additional
factor of $\cos\varphi$ is not inserted into that function.

For the weak norm, each stable-curve integral is bounded by
$C\|a\|_\infty$ because the stable curves have uniformly bounded
length and the test functions have bounded supremum. For the strong
stable norm, tests on a curve $W$ have supremum at most $|W|^{-\alpha}$.
The integral is therefore at most
$\|a\|_\infty |W|^{1-\alpha}$, uniformly bounded since $\alpha<1$.
For the strong unstable norm, take two matched graph curves at distance
at most $\epsilon$ in the source metric. Their unmatched parameter
intervals have length $O(\epsilon)$. On their common interval the graph
coordinates and arclength factors differ by $O(\epsilon)$, and the
matched tests differ in supremum by at most $\epsilon$.
The difference of the two integrals is consequently bounded by
$C\|a\|_{C^1}\epsilon$. Dividing by $\epsilon^\beta$ is harmless,
since the source exponent satisfies $\beta<1$. These three estimates
prove \eqref{eq:smooth-density-embedding}, with constants uniform in
the finite local cover and its fixed charts.

For a multiplier, stable and weak tests are simply multiplied by the
restriction of $g$ to their curves. In the matched-curve term, the two
restrictions differ in $C^q$ by at most
$C\|g\|_{C^2}\epsilon^{1-q}$, by interpolation of their supremum
and Lipschitz bounds. The admissible source exponents have
$\beta<1-q$. Splitting into matched tests and this difference gives
the first bound in \eqref{eq:smooth-multiplier-pairing}; it is the
multiplier estimate already used in
Lemma~\ref{lem:collision-spectral-input}. Applying the bounded mass
functional proves the second. Finally
\eqref{eq:bv-smoothing-bounds} bounds the smoothed $C^2$ norm by
$C\delta^{-2}\|a\|_\infty$. Composition with $p=\sin\varphi$
uses only bounded derivatives of this forward map. This proves
\eqref{eq:smoothed-density-vector}.
\end{proof}

The estimates above concern smooth functions on the local collision
spaces. They do not make the hard section projection a bounded
multiplier on those spaces, nor do they realize an unbounded induced
twist. The use of that projection in the compensation argument remains
entirely on the level of ordinary functions and smooth approximations.

\section{Chronological spectral products and coarse interpolation}
\label{app:functional-details}
This appendix supplies explicit product and interpolation estimates for
Section~\ref{sec:functional-gaussian}. It also separates the treatment
of short increments from the long-block spectral argument.

\begin{lemma}[Products of differently twisted collision powers]
\label{lem:chronological-product}
Fix an integer $q\ge1$. Use one local collision space, a common smoothing
scale $\delta$, and real vectors $z_1,\ldots,z_q$ with
$|z_j|\le c\delta^2$. Let $m_j\ge1$ and assume
\[
 \delta^{-6}\sum_{j=1}^q m_j|z_j|^3\le c.
\]
Write $L_j=\mathcal L_{R,\delta,z_j}$ and let $a$ be a complex initial
insertion with bounded supremum and variation norms. Then
\begin{align}
 &\left|\ell_j\bigl(L_q^{m_q}\cdots L_1^{m_1}
                   (\mathsf S_\delta a)\nu\bigr)
       -\left(\int a\dd\nu\right)
       \exp\left(-\tfrac12\sum_{j=1}^q
                      m_j z_j^{\mathsf T}\Gamma_{R,\delta}z_j\right)
        \right|\notag\\
 &\quad\le C_q\left\{
       \delta^{-4}\sum_{j=1}^q|z_j|
       +\delta^{-6}\sum_{j=1}^qm_j|z_j|^3
       +\delta^{-2}\sum_{j=1}^q\rho^{m_j}\right\}.
       \label{eq:chronological-product-bound}
\end{align}
The index on $\ell_j$ names the chosen local space and is fixed in the
formula. The powers are in physical chronological order: the earliest
increment acts first, on the right.
\end{lemma}
\begin{proof}
Write $L_j^{m_j}=\lambda_j^{m_j}\Pi_j+E_j$.
Lemma~\ref{lem:smooth-collision-expansion} gives
$\|\Pi_j\|\le C$, $\|E_j\|\le C\rho^{m_j}$ and
$\|\Pi_j-\Pi_R\|\le C\delta^{-2}|z_j|$.
The real-frequency quadratic terms in $m_j\log\lambda_j$ are
nonpositive. The assumption on the sum of the cubic remainders bounds
each principal scalar power and every product of these powers by a
fixed constant. Expand the product into its $2^q$ terms. Every term
with a complementary factor is bounded, after evaluation, by a
constant times $\delta^{-2}\rho^{m_j}$ for at least one of its
complementary indices. The factor $\delta^{-2}$ comes only from the
initial vector, through \eqref{eq:smoothed-density-vector}; the other
projectors and powers have uniform bounds. Summing the finitely many
terms gives the last term of \eqref{eq:chronological-product-bound}.

For the all-principal term, telescope the product of projectors about
$\Pi_R$, retaining their order. Since $\Pi_R^2=\Pi_R=\nu\ell_j$,
\[
 \|\Pi_q\cdots\Pi_1-\Pi_R\|
       \le C_q\delta^{-2}\sum_j|z_j|.
\]
Evaluation on the initial vector gives the first error in
\eqref{eq:chronological-product-bound}. Its zero-frequency amplitude
is $\int a\dd\nu$. The logarithmic expansion and
$|e^w-1|\le |w|e^{|w|}$ give the second error. The principal scalar
powers remain bounded during this last comparison. This proves the
claim without introducing a factor $\delta^{-2q}$.
\end{proof}

\begin{corollary}[Finite-dimensional scales, zero blocks, and short blocks]
\label{cor:finite-dimensional-details}
For fixed distinct times $0=t_0<t_1<\cdots<t_q$ and fixed vectors
$v_j$, put
$m_j=\lfloor nt_j\rfloor-\lfloor nt_{j-1}\rfloor$,
$z_j=v_j/\sqrt n$, and $\delta_n=\tfrac14n^{-1/32}$.
The spectral product error is bounded by
\[
 C_q\left\{\delta_n^{-4}n^{-1/2}
             +\delta_n^{-6}n^{-1/2}
             +\delta_n^{-2}\sum_j\rho^{m_j}\right\}=o(1).
\]
Repeated times give zero blocks, which may be omitted exactly. A fixed
number of additional deterministic increments of lengths $m_j=o(n)$
may instead be removed from a real characteristic exponent at a cost
$C\sum_j |v_j|\sqrt{m_j/n}$ for the original collision observable.
The estimates are uniform for initial probability densities with uniformly
bounded supremum and variation norms. The short-increment estimate uses
only the supremum bound.
\end{corollary}
\begin{proof}
Distinct fixed times make all $m_j$ positive multiples of $n$ for
large $n$. Substitute the scales in
\eqref{eq:chronological-product-bound}. Smoothing the initial density
costs $O(\delta_n)$; smoothing the covariance costs
$O(\delta_n\ell_{\delta_n})$ by
\eqref{eq:collision-covariance-limit}. The residual-path bound
\eqref{eq:residual-process-maximal} removes smoothing from the original
finite-dimensional process. Integer parts in the covariance time
coefficients cost $O(n^{-1})$. This gives precisely the independent
Gaussian increments in Theorem~\ref{thm:collision-functional}.

When an interval has length zero its operator is the identity, not a
principal spectral projector, so it is omitted before this expansion.
For a deterministic short interval, uniform collision correlation
summability bounds the second moment of its original centered sum by
$Cm_j$, from every deterministic starting index under a bounded
initial density. The difference of the two characteristic exponents,
with the other increments left at their original times, is therefore
bounded in $L^1$ by $C|v_j|\sqrt{m_j/n}$. The elementary exponential
difference bound proves the last assertion. This argument does not
use a complementary spectral power for a short block, and does not
move the other increments in physical time.
\end{proof}

\begin{lemma}[The complex bound before the fourth derivative]
\label{lem:complex-fourth-details}
For one real component $u$ of $h_{R,\delta}$, the stationary pairing
has, on a complex disk $|z|<c\delta^2$, a decomposition
\[
 \int e^{izS_{m,R}u}\dd\nu=e^{m\psi(z)}A(z)+E_m(z)
\]
with $|A(z)|+|\psi(z)|\le C$ and $|E_m(z)|\le C\rho^m$, uniformly
on a slightly smaller such disk. Consequently
$|E_m^{(4)}(0)|\le C\delta^{-8}\rho^m$.
\end{lemma}
\begin{proof}
The two fixed resolvent contours in
Lemma~\ref{lem:smooth-collision-expansion} give uniformly bounded
spectral projections and complementary powers throughout the complex
ball, not just on its real diameter. Both the vector $\nu$ and the
mass functional have bounded norms. Evaluating the complementary
power between them gives the displayed complex bound for $E_m$.
The eigenvalue lies in a fixed disk about one not meeting zero, so its
chosen logarithm and the principal amplitude are uniformly bounded.
Cauchy's estimate on a circle of radius proportional to $\delta^2$
now yields the fourth-derivative estimate. There is no need to bound
the principal power uniformly for complex $z$ and arbitrary $m$.
Its derivative is calculated by the product rule, as in
Lemma~\ref{lem:smooth-fourth-moment}.
\end{proof}

\begin{proposition}[Coarse polygonal interpolation]
\label{prop:coarse-interpolation}
Let $(u_k)$ be bounded real-vector random variables, $|u_k|\le M$.
Let $L\ge1$ be an integer and suppose, under the probability in question,
that every deterministic interval of $m\ge L$ indices has
\[
 \mathbb E\left|\sum_{k=r}^{r+m-1}u_k\right|^4\le C m^2.
\]
At times $jL/n$, interpolate the sums normalized by $\sqrt n$ to obtain
$Z_n$. Let $X_n$ be the polygonal interpolation using every index.
Then, with constants independent of $L,n$,
\begin{equation}\label{eq:coarse-interpolation-bound}
 \mathbb E|Z_n(t)-Z_n(s)|^4\le C'|t-s|^2,
 \qquad \|Z_n-X_n\|_\infty\le 2ML/\sqrt n.
\end{equation}
A final incomplete coarse interval is defined by extending one full
block beyond the observation horizon and restricting the interpolation.
\end{proposition}
\begin{proof}
Write $\Delta=L/n$. For a segment of duration $h\le\Delta$ inside
one block, linear interpolation multiplies the full-block increment by
$h/\Delta$. Its fourth moment is at most
$C\Delta^2(h/\Delta)^4\le Ch^2$. For an interval spanning two adjacent
blocks, apply this estimate separately to the durations $h_1,h_2$ and
use the $L^4$ triangle inequality. Since
$(\sqrt{h_1}+\sqrt{h_2})^4\le4(h_1+h_2)^2$, the same form holds.

For a general interval, split it into its first fractional block, its
complete-block middle interval, and its last fractional block. The
middle interval, when nonempty, contains at least $L$ indices, so the
hypothesis gives the required fourth moment for its total duration.
The $L^4$ triangle inequality and
$(\sqrt{h_1}+\sqrt{h_2}+\sqrt{h_3})^4
 \le9(h_1+h_2+h_3)^2$
prove the first bound with one fixed constant. Empty pieces are simply
omitted. Within one coarse block both interpolations remain within
$ML/\sqrt n$ of the sum at the left endpoint, proving the deterministic
second bound. The endpoint convention uses a full block, so these same
estimates apply on a fixed observation interval without a new constant.
\end{proof}

Applying this proposition with $u_k=h_{R,\delta_n}\circ T_R^k$ and
$L_n=\lceil\delta_n^{-8}\rceil$ gives the interpolation estimates in
Theorem~\ref{thm:collision-functional}. Indeed
\eqref{eq:smooth-fourth-moment} is bounded by $Cm^2$ for $m\ge L_n$.
Stationarity gives the interval estimates under $\nu$, and bounded
initial densities give them by domination. The fourth-moment constants
are uniform. Together with the residual maximal estimate this proves
tightness of the original process, not just of its coarse approximation.

\begin{lemma}[Continuity of the covariance square root]
\label{lem:covariance-square-root}
On every bounded set of positive semidefinite matrices, the map
$C\mapsto C^{1/2}$ is continuous in operator norm. Consequently
$R\mapsto\operatorname{Law}(\mathcal W_{\Gamma_R})$ is continuous on
continuous path space, including at singular matrices.
\end{lemma}
\begin{proof}
Uniformly approximate the scalar square-root function on a bounded
interval containing all the spectra by polynomials. The spectral theorem
bounds the matrix approximation by the same uniform scalar error, and
each polynomial in the matrix is norm-continuous. This proves the first
claim. On one probability space take a standard four-dimensional
Brownian motion $W$ and set $\mathcal W_C=C^{1/2}W$. On any compact time
interval its sample paths are bounded almost surely. Norm continuity
of the square root therefore gives almost-sure uniform convergence of
the coupled paths and hence convergence of their laws.
\end{proof}

\clearpage
\section{Locations of the collision norm definitions}
\label{app:collision-norm-map}
The following map fixes the source conventions in
Lemma~\ref{lem:density-norm-details}. Locations refer to the 45-page
preprint of \cite{DZ}; the last column records the deduction used here,
not an additional theorem attributed to that source.

\begin{center}
\renewcommand{\arraystretch}{1.2}
\begin{tabular}{>{\raggedright\arraybackslash}p{0.24\textwidth}>{\raggedright\arraybackslash}p{0.24\textwidth}>{\raggedright\arraybackslash}p{0.40\textwidth}}
Object & Source location & Use in this article\\ \hline
Relative density & Section 3.2; Section 3.3, footnote 7 & A smooth $a$
represents $a\nu$, without another cosine factor.\\
Weak norm & Equation (3.11) & Bounded stable-curve integrals.\\
Stable norm & Equation (3.12) & The test normalization gives
$\|a\|_\infty |W|^{1-\alpha}$.\\
Unstable norm & Equation (3.13) & Matched curves give
$C\|a\|_{C^1}\epsilon^{1-\beta}$.\\
Exponent restrictions & Section 3.3, before (3.12) &
$\beta\le p-q$, $p\le1/3$, hence $\beta<1-q$.\\
Mass functional & Lemma 3.4, with test $1$ &
$|\ell(V)|\le C|V|_w$.\\
Complementary powers & Theorem 2.1(3) &
Local uniform spectral power bound.
\end{tabular}
\end{center}

Smooth multiplication is proved in
Lemma~\ref{lem:collision-spectral-input} from these definitions.
The changing boundary length is handled by Remark 2.7(b) of the source.
The table concerns collision-space distributions. It neither embeds the
hard induced twist nor estimates second derivatives of raw densities.

\clearpage
\begin{thebibliography}{99}
\bibitem{A2G} Q. Qi, \emph{Scalar collision laws and recognition of periodic dispersing billiards}, A2 v43, primary article and technical supplement, theta-theory repository, 2026.
\bibitem{SV} D. Sz\'asz and T. Varj\'u, \emph{Local limit theorem for the Lorentz process and its recurrence in the plane}, Ergodic Theory Dynam. Systems \textbf{24} (2004), 257--278.
\bibitem{DN} D. Dolgopyat and P. N\'andori, \emph{On mixing and the local central limit theorem for hyperbolic flows}, Ergodic Theory Dynam. Systems \textbf{40} (2020), 142--174; arXiv:1710.08568v1, Definition 2.2 and Proposition 6.3.
\bibitem{DZ} M. F. Demers and H.-K. Zhang, \emph{A functional analytic approach to perturbations of the Lorentz gas}, Comm. Math. Phys. \textbf{324} (2013), 767--830; arXiv:1210.1261.
\bibitem{BDL} V. Baladi, M. F. Demers and C. Liverani, \emph{Exponential decay of correlations for finite horizon Sinai billiard flows}, Invent. Math. \textbf{211} (2018), 39--177; arXiv:1506.02836v3, including the correction note.
\bibitem{Young} L.-S. Young, \emph{Statistical properties of dynamical systems with some hyperbolicity}, Ann. of Math. (2) \textbf{147} (1998), 585--650; author preprint ESI 445 (1997), Section 8, Theorem 6.
\bibitem{BHN} H. Bruin, M. Holland and M. Nicol, \emph{Liv\v{s}ic regularity for Markov systems}, Ergodic Theory Dynam. Systems \textbf{25} (2005), 1739--1765; doi:10.1017/S0143385705000179; arXiv:math/0503690.
\bibitem{DZ11} M. F. Demers and H.-K. Zhang, \emph{Spectral analysis of the transfer operator for the Lorentz gas}, J. Mod. Dyn. \textbf{5} (2011), 665--709; doi:10.3934/jmd.2011.5.665.
\bibitem{Coste} M. Coste, \emph{An introduction to o-minimal geometry}, lecture notes, November 1999; Dottorato di Ricerca in Matematica, Universit\`a di Pisa, Istituti Editoriali e Poligrafici Internazionali, 2000.
\bibitem{GouezelRenewal} S. Gou\"ezel, \emph{Sharp polynomial estimates for the decay of correlations}, Israel J. Math. \textbf{139} (2004), 29--65; arXiv:math/0202147.
\bibitem{MN} I. Melbourne and M. Nicol, \emph{Almost sure invariance principle for nonuniformly hyperbolic systems}, Comm. Math. Phys. \textbf{260} (2005), 131--146; arXiv:math/0503693.
\bibitem{DPZ} M. F. Demers, F. P\`ene and H.-K. Zhang, \emph{Local limit theorem for randomly deforming billiards}, preprint, 2019, arXiv:1902.06850v1.
\bibitem{AFP}
L.~Ambrosio, N.~Fusco and D.~Pallara,
\emph{Functions of Bounded Variation and Free Discontinuity Problems},
Oxford Mathematical Monographs, Clarendon Press, Oxford, 2000.
\bibitem{BPR}
S.~Basu, R.~Pollack and M.-F.~Roy,
\emph{An asymptotically tight bound on the number of semi-algebraically
connected components of realizable sign conditions},
Combinatorica \textbf{29} (2009), 523--546;
arXiv:math/0603256v3, Section 3.2, equation (3.3).
\end{thebibliography}

\end{document}
